\documentclass[11pt,letterpaper,oneside,openany]{book}
\usepackage{fontspec}
\setmainfont{Latin Modern Roman}[SmallCapsFont={Latin Modern Roman Caps}]
\setsansfont{Latin Modern Sans}
\usepackage[provide=*,spanish]{babel}
\usepackage{microtype}
\usepackage[margin=26mm,headheight=22pt,headsep=8mm]{geometry}
\usepackage{amsmath,amssymb,amsthm,mathtools,mathrsfs}
\usepackage{enumitem,longtable,booktabs,array}
\usepackage{fancyhdr,emptypage}
\usepackage{imakeidx}
\usepackage[hidelinks,unicode,bookmarksnumbered]{hyperref}
\usepackage{bookmark}
\hypersetup{pdftitle={Espacios métricos, normados y de Hilbert},pdfauthor={Eder Said Martinez Ramirez},pdfsubject={Análisis. Volumen I. Teoría y ejercicios resueltos},pdfkeywords={espacios métricos, Banach, Hilbert, análisis}}
\makeindex[title={Índice alfabético},intoc,options=-s indice.ist]
\setcounter{tocdepth}{1}
\setcounter{secnumdepth}{2}
\setlength{\parindent}{1.2em}
\setlength{\parskip}{3pt plus 1pt minus 1pt}
\emergencystretch=5em
\allowdisplaybreaks[1]
\setlist[enumerate,1]{label=(\roman*),leftmargin=*,itemsep=2pt}
\setlist[itemize]{leftmargin=*,itemsep=2pt}
\pagestyle{fancy}
\fancyhf{}
\fancyhead[L]{\small\nouppercase{\leftmark}}
\fancyfoot[C]{\thepage}
\renewcommand{\headrulewidth}{0.3pt}
\renewcommand{\chaptermark}[1]{\markboth{\thechapter.\ #1}{}}
\fancypagestyle{plain}{\fancyhf{}\fancyfoot[C]{\thepage}\renewcommand{\headrulewidth}{0pt}}
\theoremstyle{plain}
\newtheorem{teorema}{Teorema}[chapter]
\newtheorem{proposicion}[teorema]{Proposición}
\newtheorem{lema}[teorema]{Lema}
\newtheorem{corolario}[teorema]{Corolario}
\theoremstyle{definition}
\newtheorem{definicion}[teorema]{Definición}
\newtheorem{ejemplo}[teorema]{Ejemplo}
\newtheorem{ejercicio}{Ejercicio}[chapter]
\theoremstyle{remark}
\newtheorem{observacion}[teorema]{Observación}
\newenvironment{solucion}{\begin{proof}[Solución]}{\end{proof}}
\newenvironment{estrategia}{\par\noindent\textbf{Estrategia.}\ }{\par}
\newcommand{\R}{\mathbb R}
\newcommand{\C}{\mathbb C}
\newcommand{\N}{\mathbb N}
\newcommand{\Q}{\mathbb Q}
\newcommand{\Z}{\mathbb Z}
\newcommand{\F}{\mathbb F}
\newcommand{\eps}{\varepsilon}
\newcommand{\norm}[1]{\left\lVert#1\right\rVert}
\newcommand{\abs}[1]{\left\lvert#1\right\rvert}
\newcommand{\cl}[1]{\overline{#1}}
\newcommand{\interior}[1]{#1^{\circ}}
\newcommand{\ip}[2]{\left\langle#1,#2\right\rangle}
\newcommand{\Bcal}{\mathcal B}
\newcommand{\Lcal}{\mathcal L}
\newcommand{\Kcal}{\mathcal K}
\DeclareMathOperator{\dist}{dist}
\DeclareMathOperator{\diam}{diam}
\DeclareMathOperator{\spanop}{span}
\DeclareMathOperator{\supp}{supp}
\DeclareMathOperator{\id}{id}
% Identidad gráfica: fuentes y criterios en editorial/identidad-esfm-ipn.md.
\usepackage{xcolor,graphicx,tikz,titlesec,needspace}
\usepackage[titles]{tocloft}
\usepackage[skins,breakable]{tcolorbox}
\definecolor{ipnGuinda}{HTML}{6C1D45}
\definecolor{ipnGris}{HTML}{63666A}
% Acentos muestreados del escudo original publicado por la ESFM.
\definecolor{esfmOcre}{HTML}{836F3D}
\definecolor{tinta}{HTML}{232225}
\color{tinta}
\hypersetup{colorlinks=true,linkcolor=ipnGuinda,citecolor=ipnGuinda,urlcolor=ipnGuinda}
\renewcommand{\contentsname}{Índice general}
\renewcommand{\proofname}{Demostración}
\renewcommand{\qedsymbol}{\textcolor{ipnGuinda}{\ensuremath{\square}}}

\fancyhf{}
\fancyhead[L]{\fontsize{8}{10}\selectfont\color{ipnGris}\nouppercase{\leftmark}}
\fancyhead[R]{\fontsize{8}{10}\selectfont\color{ipnGuinda}ESFM\enspace/\enspace IPN}
\fancyfoot[L]{\fontsize{7.5}{9}\selectfont\color{ipnGris}Eder Said Martinez Ramirez}
\fancyfoot[R]{\fontsize{8}{10}\selectfont\hyperlink{indice-general}{\thepage}}
\renewcommand{\headrulewidth}{0.4pt}
\renewcommand{\headrule}{\hbox to\headwidth{\color{ipnGuinda}\leaders\hrule height\headrulewidth\hfill}}
\fancypagestyle{plain}{%
  \fancyhf{}
  \fancyfoot[L]{\fontsize{7.5}{9}\selectfont\color{ipnGris}Análisis\enspace/\enspace Volumen I}
  \fancyfoot[R]{\fontsize{8}{10}\selectfont\hyperlink{indice-general}{\thepage}}
  \renewcommand{\headrulewidth}{0pt}}

\titleformat{\chapter}[display]
  {\normalfont\bfseries\color{ipnGuinda}\raggedright}
  {\parbox[b]{.72\textwidth}{\fontsize{10}{13}\selectfont\mdseries
     Análisis\enspace/\enspace Volumen I\\[3pt]
     {\color{ipnGris}Capítulo \thechapter}}
   \hfill{\setlength{\fboxsep}{8pt}\colorbox{ipnGuinda}{%
     \makebox[19mm]{\color{white}\fontsize{38}{42}\selectfont\thechapter}}}}
  {14pt}{\fontsize{25}{30}\selectfont}
  [\vspace{9pt}{\color{esfmOcre}\titlerule[1pt]}]
\titleformat{name=\chapter,numberless}[display]
  {\normalfont\bfseries\color{ipnGuinda}\raggedright}
  {}{0pt}{\fontsize{25}{30}\selectfont}
  [\vspace{9pt}{\color{esfmOcre}\titlerule[1pt]}]
\titlespacing*{\chapter}{0pt}{-5pt}{24pt}
\titleformat{\section}{\normalfont\large\bfseries\color{ipnGuinda}}{\thesection}{.7em}{}
\titleformat{\subsection}{\normalfont\normalsize\bfseries\color{tinta}}{\thesubsection}{.7em}{}
\titlespacing*{\section}{0pt}{18pt plus 3pt minus 2pt}{8pt}
\titlespacing*{\subsection}{0pt}{12pt plus 2pt minus 2pt}{5pt}
% La reserva evita separar el título de ejercicios de su primer recuadro.
\NewCommandCopy{\seccionOriginal}{\section}
\RenewDocumentCommand{\section}{s o m}{%
  \IfBooleanTF{#1}{\seccionOriginal*{#3}}{%
    \ifstrequal{#3}{Ejercicios resueltos}{\Needspace{8\baselineskip}}{}%
    \IfNoValueTF{#2}{\seccionOriginal{#3}}{\seccionOriginal[#2]{#3}}}}

\renewcommand{\cftchapfont}{\bfseries\color{ipnGuinda}}
\renewcommand{\cftchappagefont}{\bfseries\color{ipnGuinda}}
\renewcommand{\cftsecfont}{\small}
\renewcommand{\cftsecpagefont}{\small}
\setlength{\cftbeforechapskip}{7pt}
\setlength{\cftbeforesecskip}{1pt}
\setlength{\cftchapnumwidth}{2.4em}
\setlength{\cftsecindent}{2.4em}
\setlength{\cftsecnumwidth}{3em}
\renewcommand{\cftchapleader}{\cftdotfill{\cftdotsep}}

% El recuadro crece hacia el margen para conservar la anchura de las fórmulas.
\tcolorboxenvironment{ejercicio}{enhanced,breakable,lines before break=4,sharp corners,
  colback=esfmOcre!4!white,colframe=ipnGuinda!45!white,boxrule=.45pt,
  borderline west={2pt}{0pt}{esfmOcre},
  grow to left by=7pt,grow to right by=7pt,
  boxsep=0pt,left=6.55pt,right=6.55pt,top=5pt,bottom=5pt,
  before skip=10pt,after skip=8pt}

\begin{document}
\renewcommand{\bibname}{Referencias}
\frontmatter
\hypersetup{pageanchor=false}
\begin{titlepage}
\thispagestyle{empty}
\setlength{\parindent}{0pt}
\begin{tikzpicture}[remember picture,overlay]
  \fill[ipnGuinda] (current page.south west) rectangle ++(\paperwidth,5mm);
  \fill[ipnGris!16] (current page.south east) -- ++(-55mm,0) -- ++(55mm,55mm) -- cycle;
  \fill[esfmOcre] (current page.south east) -- ++(-37mm,0) -- ++(37mm,37mm) -- cycle;
\end{tikzpicture}
\noindent
\begin{minipage}[c]{.72\textwidth}
{\color{ipnGuinda}\fontsize{15}{20}\selectfont\bfseries Instituto Politécnico Nacional\par}
\vspace{3mm}
{\fontsize{12}{17}\selectfont Escuela Superior de Física\par y Matemáticas\par}
\vspace{3mm}
{\small Sección de Estudios de Posgrado\par e Investigación\par}
\end{minipage}\hfill
\begin{minipage}[c]{.22\textwidth}
\raggedleft\includegraphics[height=39mm]{assets/escudo-esfm.png}
\end{minipage}
\vspace{13mm}
\par\noindent{\color{ipnGuinda}\rule{\textwidth}{.7pt}}\par
\vspace{12mm}
{\color{ipnGuinda}\fontsize{15}{18}\selectfont Análisis\hfill Volumen I\par}
\vspace{7mm}
{\color{ipnGuinda}\fontsize{31}{38}\selectfont\bfseries
Espacios métricos,\par normados y de Hilbert\par}
\vspace{8mm}
{\fontsize{13}{18}\selectfont Teoría desarrollada\par y ejercicios resueltos\par}
\vspace{14mm}
{\color{esfmOcre}\rule{17mm}{2pt}\par}
\vspace{5mm}
{\fontsize{15}{20}\selectfont\bfseries Eder Said Martinez Ramirez\par}
\vspace{4mm}
{\fontsize{10.5}{15}\selectfont Maestría en Ciencias Fisicomatemáticas\par
Línea de investigación: Matemáticas Puras\par}
\vfill
{\small Edición ampliada de las lecturas de\par Análisis Matemático 1\par}
\vspace{5mm}
{\small Ciudad de México\enspace/\enspace 2026\par}
\vspace{7mm}
\end{titlepage}

\hypersetup{pageanchor=true}
\chapter*{Presentación}
\addcontentsline{toc}{chapter}{Presentación}
Este volumen desarrolla el primer curso de una colección de cuatro libros de análisis. Su hilo conductor es el paso de la distancia a la estructura lineal y, finalmente, a la geometría del producto interno. La intención es que un estudiante de la mitad de la licenciatura en Matemáticas pueda seguir las pruebas y entender por qué se eligen sus pasos, sin rebajar la precisión de los enunciados.

El punto de partida son las 32 lecturas de \emph{Análisis Matemático 1, Licenciatura en Física y Matemáticas}, de Miguel Ángel Valencia Bucio \cite{valencia}. Las lecturas se integran en seis partes y dieciséis capítulos temáticos. Se conserva la identificación de los ejercicios originales y una tabla final permite localizar el contenido de cada lectura. La autoría de esas lecturas se reconoce expresamente: la organización editorial, las explicaciones ampliadas, las soluciones desarrolladas y las rectificaciones forman esta edición de Eder Said Martinez Ramirez.

Un resultado falso no puede completarse añadiendo pasos a una prueba. Cuando un enunciado necesita una hipótesis o admite un contraejemplo, el texto explica el problema y formula una versión válida. Se distinguen asimismo las erratas matemáticas de las pérdidas de barras, índices o signos ocasionadas por la extracción del PDF. Las soluciones incluyen el enunciado que se utiliza y señalan las correcciones pertinentes.

\section*{Cómo estudiar este volumen}
Conviene conocer operaciones con conjuntos, cuantificadores, inducción, sucesiones y series reales, cálculo de una variable y álgebra lineal elemental. No se presupone un curso previo de topología ni de análisis funcional. La completitud de los números reales y los resultados básicos del cálculo se usan como conocimientos previos. El lema de Zorn se enuncia cuando interviene como principio de existencia; las construcciones que dependen de él se separan de las construcciones explícitas.

Antes de leer una demostración, identifique sus hipótesis y escriba lo que se quiere producir: un radio, una subsucesión, una cota o un operador. Lea después la estrategia y trate de reconocer dónde se usa cada hipótesis. En los ejercicios, conviene intentar una solución antes de consultar el desarrollo. Un contraejemplo merece el mismo cuidado que una prueba: hay que verificar tanto las hipótesis que satisface como la conclusión que contradice.

La jerarquía de la colección es volumen, parte, capítulo y sección. Cada volumen corresponde a un curso; sus partes reúnen etapas conceptuales y sus capítulos desarrollan unidades matemáticas. En este volumen, las seis partes recorren espacios métricos y continuidad, completitud, compacidad y estructura topológica, espacios normados y operadores, espacios de funciones y geometría de Hilbert. Los capítulos se numeran consecutivamente dentro del volumen.

Los ejercicios se distribuyen cerca de las herramientas que utilizan. Su número dentro del volumen permite citarlos sin ambigüedad; la indicación entre paréntesis conserva la lectura y el número de origen. Las referencias enlazadas llevan al resultado o ejercicio correspondiente. La tabla final de correspondencia distingue el orden temático del libro del orden de las lecturas originales.

\section*{Convenciones y fuentes}
Salvo indicación contraria, $\N=\{1,2,\ldots\}$ y $\F$ es $\R$ o $\C$. El producto interno es lineal en su primer argumento. Las bolas abiertas y cerradas se distinguen del cierre de una bola abierta: en un espacio métrico arbitrario esos dos últimos conjuntos pueden ser diferentes. El dual $V^*$ contiene solamente los funcionales lineales continuos.

Rudin \cite{rudin1976} ofrece un desarrollo de topología métrica, sucesiones de funciones y aproximación; Munkres \cite{munkres2000} amplía el contexto topológico; Kreyszig \cite{kreyszig1978} y Conway \cite{conway1990} permiten continuar hacia espacios normados, Hilbert y operadores. La bibliografía final identifica obras reales y ediciones comprobadas. Las referencias externas son orientación y contraste; las demostraciones desarrolladas aquí deben poder seguirse a partir de los resultados previos indicados, sin buscar una página no especificada en otro libro.
\hypertarget{indice-general}{}
\tableofcontents
\mainmatter
\parte{Espacios métricos y continuidad}{%
La distancia permite expresar con precisión qué significa aproximar un punto, converger a un límite y conservar la proximidad mediante una función. Esta parte construye ese lenguaje a partir de ejemplos: normas, bolas, conjuntos abiertos y cerrados, sucesiones y subespacios.

Después se estudian la continuidad, la densidad y las condiciones de numerabilidad, hasta distinguir el control local de una función del control uniforme de todas las distancias. Se requieren operaciones con conjuntos, cuantificadores y sucesiones reales; los conceptos topológicos se introducen desde sus definiciones.}
\chapter{Distancias, normas y topología métrica}\label{cap:nuevo01}\label{chap:01}\label{chap:02}\label{chap:03}\label{chap:04}

Este capítulo construye el lenguaje métrico y topológico que se utilizará en todo el volumen. Partiremos de distancias y normas sobre vectores, sucesiones y funciones; pasaremos de las bolas a las vecindades, el interior y la clausura; y terminaremos describiendo la convergencia y el cambio de espacio ambiente. El recorrido integra las lecturas 1--4 de las notas de partida \cite{valencia}, conservando sus desarrollos y ejercicios con una organización continua. Se requieren las propiedades del valor absoluto, el lenguaje de conjuntos y cuantificadores, las series de términos no negativos y la integral de Riemann de funciones continuas. La referencia a integración en varias variables se precisa al introducirla; el caso de un intervalo permite seguir las pruebas sin ese prerrequisito adicional.

\section{Qué significa medir la proximidad}\label{lectura:01}
Medir la proximidad de números mediante $|x-y|$ sugiere una definición que también pueda aplicarse a vectores, sucesiones, funciones y operadores. Comenzamos por la distancia; después incorporaremos su compatibilidad con la estructura vectorial mediante una norma. Las desigualdades de Hölder y Minkowski justificarán los ejemplos, en lugar de dejar su verificación como una afirmación pendiente.

\begin{definicion}\label{ch01:metrica}\index{metrica@métrica}
Una \emph{métrica} en un conjunto $X$ es una función $d:X\times X\to[0,\infty)$ tal que, para $x,y,z\in X$,
\[
d(x,y)=0\Longleftrightarrow x=y,\qquad d(x,y)=d(y,x),\qquad
d(x,y)\le d(x,z)+d(z,y).
\]
El par $(X,d)$ es un espacio métrico. La última condición es la desigualdad triangular.
\end{definicion}
La desigualdad triangular expresa que pasar por un punto intermedio no acorta la distancia directa. No presupone que existan segmentos o caminos en $X$. El conjunto vacío admite formalmente una métrica; salvo indicación expresa trabajaremos con espacios no vacíos.

\begin{ejemplo}[Métrica discreta]\index{metrica@métrica!discreta}
En cualquier conjunto, $d(x,y)=1$ si $x\ne y$ y $d(x,x)=0$ define una métrica. Para comprobar la desigualdad triangular, si $x=y$ el lado izquierdo es cero. Si $x\ne y$, no pueden cumplirse simultáneamente $x=z$ y $z=y$; al menos uno de los dos sumandos del lado derecho es $1$, de modo que su suma es al menos $d(x,y)=1$.
\end{ejemplo}
En $\R$ y $\C$, $d(x,y)=|x-y|$ es una métrica: la separación y la simetría se deducen del valor absoluto, y $|x-y|=|(x-z)+(z-y)|\le|x-z|+|z-y|$.

\begin{definicion}\label{ch01:norma}\index{norma}
Sea $V$ un espacio vectorial sobre $\F=\R$ o $\C$. Una norma es una función $\norm{\cdot}:V\to[0,\infty)$ tal que
\[
\norm{v}=0\Longleftrightarrow v=0,\qquad
\norm{\alpha v}=|\alpha|\norm{v},\qquad
\norm{u+v}\le\norm{u}+\norm{v}.
\]
El par $(V,\norm{\cdot})$ es un espacio normado. Una \emph{seminorma} satisface las mismas condiciones salvo que puede anularse en vectores no nulos.
\end{definicion}
\begin{proposicion}\label{ch01:inducida}
Toda norma induce la métrica $d(u,v)=\norm{u-v}$.
\end{proposicion}
\begin{proof}
La norma es no negativa y $\norm{u-v}=0$ equivale a $u-v=0$. La homogeneidad con escalar $-1$ da $\norm{u-v}=\norm{v-u}$. Finalmente,
\[
\norm{u-v}=\norm{(u-w)+(w-v)}\le\norm{u-w}+\norm{w-v}.
\]
Estas son exactamente las tres propiedades de una métrica.
\end{proof}
La métrica inducida es invariante por traslación: $d(u+a,v+a)=d(u,v)$, y satisface $d(\alpha u,\alpha v)=|\alpha|d(u,v)$. Una métrica arbitraria no tiene por qué poseer estas propiedades.

\section{Las desigualdades de Hölder y Minkowski}
En los ejemplos con varias variables, un compacto $K\subset\R^m$ se llama \emph{Jordan-medible} si su función indicadora $\mathbf1_K$ es integrable en el sentido de Riemann en un rectángulo que lo contiene. La integral $\int_K h$ se entiende como la integral en ese rectángulo de la función igual a $h$ en $K$ y a cero fuera de él. Para $h$ continua en $K$, esta extensión es integrable: sus posibles discontinuidades exteriores a la continuidad relativa están en la frontera de $K$, que tiene contenido de Jordan cero. El lector que solo haya estudiado cálculo de una variable puede tomar $K=[a,b]$, $a<b$, en todas las pruebas integrales de este capítulo; los argumentos de Hölder y Minkowski son idénticos y no requieren todavía teoría de la medida.
Para justificar las normas $p$ no basta con afirmar que la desigualdad triangular es conocida. La herramienta decisiva es Hölder. Para $1<p<\infty$, el exponente conjugado es $q=p/(p-1)$; así $1/p+1/q=1$. Para $p=1$ se toma $q=\infty$, y para $p=\infty$ se toma $q=1$.

\begin{lema}[Desigualdad de Young]\label{ch01:young}\index{Young, desigualdad de}
Si $a,b\ge0$ y $1<p<\infty$, entonces $ab\le a^p/p+b^q/q$.
\end{lema}
\begin{proof}
Fijado $b$, consideremos $h(t)=t^p/p-bt$ para $t\ge0$. Si $b=0$, la afirmación es inmediata. Si $b>0$, $h'(t)=t^{p-1}-b$ cambia de negativo a positivo en $t=b^{1/(p-1)}$. Ese punto es el mínimo global; su valor es $-b^q/q$. Por tanto $a^p/p-ab\ge-b^q/q$.
\end{proof}

\begin{teorema}[Hölder]\label{ch01:holder}\index{Holder@Hölder, desigualdad de}
Para familias finitas de escalares,
\[
\sum_{k=1}^N|a_kb_k|\le
\left(\sum_{k=1}^N|a_k|^p\right)^{1/p}
\left(\sum_{k=1}^N|b_k|^q\right)^{1/q},\qquad 1<p<\infty.
\]
La misma desigualdad vale para series si ambos factores del lado derecho son finitos. En los extremos, se reemplaza el factor de exponente infinito por el supremo. Para funciones continuas en un compacto Jordan-medible $K\subset\R^m$, vale también
\[
\int_K|fg|\le \left(\int_K|f|^p\right)^{1/p}
\left(\int_K|g|^q\right)^{1/q}.
\]
\end{teorema}
\begin{proof}
Para sumas finitas, escribamos $A=(\sum|a_k|^p)^{1/p}$ y $B=(\sum|b_k|^q)^{1/q}$. Si alguno es cero, todos sus términos son cero y la desigualdad queda probada. Si $A,B>0$, Young aplicada a $|a_k|/A$ y $|b_k|/B$ da
\[
\frac{|a_kb_k|}{AB}\le\frac{|a_k|^p}{pA^p}+\frac{|b_k|^q}{qB^q}.
\]
Al sumar se obtiene $\sum|a_kb_k|/(AB)\le1/p+1/q=1$. Para series, cada suma parcial está acotada por el producto de las normas completas; las sumas parciales son crecientes, por lo que convergen y satisfacen la misma cota. En el caso $p=1$, basta sumar $|a_kb_k|\le|a_k|\sup_j|b_j|$; el caso $p=\infty$ intercambia los papeles.

Para integrales con ambos factores positivos se aplica Young a $|f|/A$ y $|g|/B$ punto a punto y se integra; la linealidad y la monotonía de la integral producen nuevamente $1/p+1/q$. Si $\int_K|f|^p=0$, Hölder no puede usarse todavía para concluir que el producto tiene integral cero. Para $\delta>0$, la desigualdad de Young aplicada a $|f|/\delta$ y $|g|$ da
\[
\int_K|fg|\le \frac{\delta^{1-p}}p\int_K|f|^p+
\frac\delta q\int_K|g|^q=\frac\delta q\int_K|g|^q.
\]
Como esto vale para todo $\delta>0$, la integral del producto es cero. El otro caso nulo se trata intercambiando factores. En los extremos se integra $|fg|\le |f|\norm{g}_\infty$, o su versión intercambiada.
\end{proof}

\begin{teorema}[Minkowski]\label{ch01:minkowski}\index{Minkowski, desigualdad de}
Para $1\le p<\infty$, las expresiones $p$ de sumas finitas, series e integrales satisfacen la desigualdad triangular.
\end{teorema}
\begin{proof}
Para $p=1$ se suma o integra $|a+b|\le|a|+|b|$. Supongamos $p>1$ y comencemos por una suma finita. Si $S=(\sum|a_k+b_k|^p)^{1/p}=0$, no hay nada que probar. En otro caso,
\begin{align*}
S^p&\le\sum|a_k|\,|a_k+b_k|^{p-1}+\sum|b_k|\,|a_k+b_k|^{p-1}\\
&\le\left[\left(\sum|a_k|^p\right)^{1/p}
+\left(\sum|b_k|^p\right)^{1/p}\right]
\left(\sum|a_k+b_k|^{(p-1)q}\right)^{1/q}.
\end{align*}
Como $(p-1)q=p$, el último factor es $S^{p-1}$. Dividir por ese número positivo prueba la desigualdad. Si $a,b$ tienen sumas $p$ finitas, aplicamos la desigualdad a sus primeras $N$ coordenadas. Las sumas parciales de $\sum|a_k+b_k|^p$ quedan acotadas por $(\norm a_p+\norm b_p)^p$. Su límite es finito y, al pasar al límite, se obtiene la desigualdad para series. Este argumento demuestra a la vez que la suma pertenece al espacio: no lo presupone.

Para funciones continuas en $K$, $|f+g|^p$ es integrable. Si su integral es positiva, se repite el cálculo anterior con integrales y Hölder; el factor final es $(\int|f+g|^p)^{1/q}$. Si la integral es cero, la desigualdad es automática. Para el supremo, $|f(x)+g(x)|\le\norm f_\infty+\norm g_\infty$ para cada $x$, y tomar el supremo concluye la prueba.
\end{proof}

\section{Los espacios modelo}
\begin{ejemplo}[Normas en dimensión finita]\index{norma!p en dimension finita@$p$ en dimensión finita}
En $\F^n$, para $n\ge1$,
\[
\norm x_1=\sum_{k=1}^n|x_k|,\qquad
\norm x_p=\left(\sum_{k=1}^n|x_k|^p\right)^{1/p},\qquad
\norm x_\infty=\max_{1\le k\le n}|x_k|.
\]
Todas son normas. La separación se sigue de que una suma de términos no negativos solo es cero si cada término lo es; para el máximo se usa el mismo razonamiento. La homogeneidad se obtiene sacando $|\alpha|^p$ de la suma, o $|\alpha|$ del máximo. La desigualdad triangular es Minkowski para $1<p<\infty$, la desigualdad del valor absoluto para $p=1$, y la estimación del supremo para $p=\infty$. La norma $2$ es euclidiana; la norma $1$ produce la distancia del taxista.
\end{ejemplo}

\begin{definicion}\index{ell p@\(\ell^p\)}\index{ell infinito@\(\ell^\infty\)}
Para $1\le p<\infty$, se define
\[
\ell^p=\left\{a=(a_k)_{k\ge1}:\sum_{k=1}^{\infty}|a_k|^p<\infty\right\},
\qquad \norm a_p=\left(\sum_{k=1}^{\infty}|a_k|^p\right)^{1/p}.
\]
Además, $\ell^\infty$ es el espacio de sucesiones acotadas, con $\norm a_\infty=\sup_k|a_k|$.
\end{definicion}
El cero pertenece a cada espacio; la multiplicación por escalares conserva la condición de pertenencia porque multiplica la serie por $|\alpha|^p$; y la suma la conserva por la prueba de Minkowski. En $\ell^\infty$, la cota $|a_k+b_k|\le\norm a_\infty+\norm b_\infty$ prueba lo mismo. Las verificaciones de separación y homogeneidad son las de dimensión finita, coordenada por coordenada. Así se justifica que estos objetos sean espacios vectoriales normados. Hölder dice en particular que si $a\in\ell^p$ y $b\in\ell^q$, entonces el producto coordenado $ab$ está en $\ell^1$.

\begin{ejemplo}[Funciones continuas]\index{C(K)@$C(K)$}\label{ch01:funciones}
En un compacto no vacío $K\subset\R^m$, la norma uniforme es
\[
\norm f_\infty=\max_{x\in K}|f(x)|,\qquad f\in C(K).
\]
El máximo existe por el teorema elemental de valores extremos para funciones continuas en compactos euclidianos. La separación se comprueba evaluando en cada punto; la homogeneidad y la desigualdad triangular se siguen de las propiedades del máximo.

Si $a<b$, también son normas en $C([a,b])$ las expresiones
\[
\norm f_p=\left(\int_a^b|f(t)|^p\,dt\right)^{1/p},\qquad 1\le p<\infty.
\]
Solo falta justificar la separación. Si $f(t_0)\ne0$, la continuidad proporciona un intervalo relativo de longitud positiva en el que $|f(t)|>|f(t_0)|/2$. La integral de $|f|^p$ es al menos esa longitud multiplicada por $(|f(t_0)|/2)^p$, y es positiva. La homogeneidad sigue de la integral y la triangular de Minkowski.
\end{ejemplo}
\begin{observacion}[Corrección de la hipótesis del dominio]\index{seminorma}
En un compacto Jordan-medible arbitrario, la expresión integral puede ser solo una seminorma. Si $K=\{0\}$ y $f(0)=1$, entonces $f\ne0$ pero $\int_K|f|^p=0$. Una condición suficiente para que sea norma es que toda vecindad relativa no vacía de $K$ tenga medida positiva. En ese caso la continuidad vuelve a producir un conjunto de medida positiva donde $|f|$ tiene una cota inferior positiva. Esta condición suele expresarse diciendo que la medida tiene \emph{soporte pleno} en $K$. No se necesita imponerla para la norma uniforme ni para las desigualdades de Hölder y Minkowski.
\end{observacion}

\section{Ejercicios sobre métricas, normas y ultramétricas}
\begin{ejercicio}[Lectura 1, ejercicio 1]\label{ej:lectura-1-1}
Sea $d:X\times X\to\R$ tal que $d(x,y)=0$ si y solo si $x=y$ y $d(x,y)\le d(x,z)+d(y,z)$ para todo $x,y,z$. Demostrar que $d$ es una métrica.
\end{ejercicio}


\begin{ejercicio}[Lectura 1, ejercicio 2]\label{ej:lectura-1-2}
Sean $d_1,\ldots,d_n$ métricas en $X$. Demostrar que (i) su suma y (ii) su máximo son métricas. (iii) Dar un ejemplo en el que $\min(d_1,d_2)$ no lo sea.
\end{ejercicio}


\begin{ejercicio}[Lectura 1, ejercicio 3]\label{ej:lectura-1-3}
Si $d$ es una métrica, demostrar que $d'(x,y)=d(x,y)/(1+d(x,y))$ también lo es.
\end{ejercicio}


\begin{ejercicio}[Lectura 1, ejercicio 4]\label{ej:lectura-1-4}
Sean $(d_n)$ métricas en $X$ y $a_n>0$ con $\sum_na_n<\infty$. Demostrar que
\[
D(x,y)=\sum_{n=1}^{\infty}a_n\frac{d_n(x,y)}{1+d_n(x,y)}
\]
es una métrica.
\end{ejercicio}


\begin{ejercicio}[Lectura 1, ejercicio 5]\label{ej:lectura-1-5}
Para métricas $d_1,\ldots,d_n$ y $p>1$, demostrar que
\[
d_p(x,y)=\left(\sum_{i=1}^n d_i(x,y)^p\right)^{1/p}
\]
es una métrica.
\end{ejercicio}


\begin{ejercicio}[Lectura 1, ejercicio 6]\label{ej:lectura-1-6}
Demostrar que $|d(x,z)-d(z,y)|\le d(x,y)$.
\end{ejercicio}


\begin{ejercicio}[Lectura 1, ejercicio 7]\label{ej:lectura-1-7}\index{ultrametrica@ultramétrica}
Una métrica es ultramétrica si $d(x,y)\le\max\{d(x,z),d(y,z)\}$. Demostrar que, cuando $d(x,z)\ne d(y,z)$, se tiene igualdad.
\end{ejercicio}


\section{De una distancia a una noción de entorno}\label{lectura:02}
La distancia ya permite comparar puntos. A partir de ella construiremos ahora los conjuntos que rodean a un punto y los que contienen sus puntos límite: esta es la información topológica que se extrae de una métrica. Las bolas conectan las estimaciones numéricas anteriores con las nociones de vecindad, interior y clausura.

\begin{definicion}\label{ch02:bolas}\index{bola!abierta}\index{bola!cerrada}\index{esfera}
En un espacio métrico $(X,d)$, para $r>0$, se definen
\[
B(x,r)=\{y:d(x,y)<r\},\quad
\overline B(x,r)=\{y:d(x,y)\le r\},\quad
S(x,r)=\{y:d(x,y)=r\}.
\]
Se llaman bola abierta, bola cerrada y esfera, respectivamente. El símbolo $\overline B(x,r)$ designa aquí la bola definida por $\le r$; todavía no afirma que sea la clausura de $B(x,r)$.
\end{definicion}
En la métrica discreta, $B(x,r)=\{x\}$ para $0<r\le1$ y $B(x,r)=X$ para $r>1$. La bola cerrada es $\{x\}$ para $0<r<1$ y $X$ para $r\ge1$. En $\R^n$ con la norma euclidiana,
\[
B(x,r)=\left\{y:\sum_{k=1}^n(x_k-y_k)^2<r^2\right\}.
\]
Se usa el cuadrado de la diferencia, no la diferencia de los cuadrados.

En $C(K)$ con la norma uniforme,
\[
B(f,r)=\{g\in C(K):\max_{t\in K}|g(t)-f(t)|<r\}.
\]
Como $K$ es compacto y $f,g$ continuas, esto equivale a $|g(t)-f(t)|<r$ para todo $t\in K$: el máximo se alcanza en algún punto y, en ese punto, la desigualdad también es estricta. En un dominio no compacto, las desigualdades estrictas punto a punto no garantizan un supremo estrictamente menor que $r$; por ejemplo, $t/(1+t)<1$ en $[0,\infty)$, pero su supremo es $1$.

\begin{definicion}\index{vecindad}\label{ch02:vecindad}
Un conjunto $V\subset X$ es una vecindad de $x$ si existe $r>0$ tal que $B(x,r)\subset V$. Toda vecindad contiene al punto, porque $d(x,x)=0<r$.
\end{definicion}
Una vecindad no tiene por qué ser una bola ni un conjunto abierto. En $\R$, $[-1,1]\cup\{4\}$ es una vecindad de $0$, pues contiene $B(0,1)$.

\begin{proposicion}\label{ch02:vecindades}
Si $U,V$ son vecindades de $x$, también lo es $U\cap V$. Todo conjunto que contiene una vecindad de $x$ es vecindad de $x$.
\end{proposicion}
\begin{proof}
Elijamos radios $r,s>0$ con $B(x,r)\subset U$ y $B(x,s)\subset V$. Entonces $B(x,\min\{r,s\})\subset U\cap V$. La segunda afirmación se deduce conservando la bola contenida en la vecindad original.
\end{proof}
Con el orden dado por inclusión, la intersección es la máxima cota inferior de dos vecindades: está contenida en ambas, y cualquier conjunto contenido en ambas está contenido en su intersección. La elección del radio mínimo explica por qué funciona una familia finita; para una familia infinita, el ínfimo de los radios puede ser cero.

\section{Interior y clausura}
\begin{definicion}\index{interior}\index{clausura}\label{ch02:interiorclausura}
El interior y la clausura de $A\subset X$ son
\begin{align*}
\interior A&=\{x:\text{existe }r>0\text{ con }B(x,r)\subset A\},\\
\cl A&=\{x:\text{para todo }r>0,\ B(x,r)\cap A\ne\varnothing\}.
\end{align*}
La segunda condición equivale a que toda vecindad de $x$ intersecte a $A$. Un conjunto es abierto si coincide con su interior y cerrado si coincide con su clausura.
\end{definicion}
La equivalencia sobre vecindades tiene dos direcciones: las bolas son vecindades; y, si toda bola intersecta a $A$, una vecindad cualquiera también lo hace porque contiene una bola. Este razonamiento permite reducir muchas pruebas a bolas, pero es necesario explicar por qué la reducción conserva la propiedad buscada.

\begin{proposicion}\label{ch02:operaciones}
Para $A,B\subset X$ se cumple:
\begin{enumerate}[label=(\roman*)]
\item $\interior A\subset A\subset\cl A$.
\item Si $A\subset B$, entonces $\interior A\subset\interior B$ y $\cl A\subset\cl B$.
\item $\interior{(\interior A)}=\interior A$ y $\cl{\cl A}=\cl A$.
\item $\interior A\cup\interior B\subset\interior{(A\cup B)}$ y $\interior{(A\cap B)}=\interior A\cap\interior B$.
\end{enumerate}
\end{proposicion}
\begin{proof}
\begin{enumerate}[label=(\roman*)]
\item Si una bola centrada en $x$ está contenida en $A$, entonces $x\in A$. Si $x\in A$, cada bola centrada en $x$ intersecta a $A$ en el propio $x$.
\item Una bola contenida en $A$ está contenida en $B$; una bola que intersecta a $A$ también intersecta a $B$. Estas dos observaciones dan las respectivas inclusiones.
\item Sea $x\in\interior A$ y elijamos $r>0$ con $B(x,r)\subset A$. Para $y\in B(x,r)$, el número $s=r-d(x,y)$ es positivo. Si $d(y,z)<s$, la desigualdad triangular da $d(x,z)<r$, por lo que $B(y,s)\subset A$. Así cada $y\in B(x,r)$ pertenece a $\interior A$, y $B(x,r)\subset\interior A$. Esto prueba $x\in\interior{(\interior A)}$; la otra inclusión es (i).

Ahora sea $x\in\cl{\cl A}$ y fijemos $r>0$. Existe $y\in\cl A$ con $d(x,y)<r/2$, y existe $z\in A$ con $d(y,z)<r/2$. Por tanto $d(x,z)<r$ y $B(x,r)$ intersecta a $A$. Como $r$ era arbitrario, $x\in\cl A$. La inclusión inversa se sigue de (i).
\item Si $x\in\interior A\cup\interior B$, alguna bola centrada en $x$ está contenida en uno de esos conjuntos y, por tanto, en su unión. Si $x$ pertenece a ambos interiores, la proposición~\ref{ch02:vecindades} da una bola contenida en $A\cap B$. Recíprocamente, una bola contenida en $A\cap B$ está contenida tanto en $A$ como en $B$.
\end{enumerate}
\end{proof}
La idempotencia significa que tomar interior o clausura una segunda vez no añade información. En particular, el interior siempre es abierto y la clausura siempre es cerrada. El interior describe puntos con margen para moverse dentro del conjunto; la clausura agrega todos los puntos que pueden aproximarse desde él.

\begin{ejemplo}
En $\R$, para $A=(0,1]\cup\{2\}$, se tiene $\interior A=(0,1)$ y $\cl A=[0,1]\cup\{2\}$. En cambio, para $A=\Q$, el interior es vacío y la clausura es $\R$. Un conjunto puede ser cerrado sin ser abierto, abierto sin ser cerrado, ambas cosas o ninguna; estas palabras no son negaciones lógicas entre sí.
\end{ejemplo}

\begin{definicion}[Diámetro]\index{diametro@diámetro}
El diámetro de un subconjunto no vacío $A$ de $(X,d)$ es
\[
 \diam A=\sup\{d(x,y):x,y\in A\}\in[0,\infty].
\]
Adoptamos $\diam\varnothing=0$. El valor infinito expresa que las distancias entre puntos de $A$ no tienen una cota finita. Si $A\subset B(a,r)$, la desigualdad triangular da $d(x,y)<2r$ para $x,y\in A$, y por tanto $\diam A\le2r$. Recíprocamente, si $A\ne\varnothing$ y $\diam A<\infty$, al fijar $a\in A$ obtenemos $A\subset B(a,\diam A+1)$. Así, tener diámetro finito equivale a estar contenido en una bola de radio finito.
\end{definicion}

\section{Ejercicios sobre bolas, interior y clausura}
\begin{ejercicio}[Lectura 2, ejercicio 1]\label{ej:lectura-2-1}
En $\R^{\N}$, sea $A=\{a:|a_n|<1/n\text{ para todo }n\}$. Decidir si $A$ es vecindad de $0$ para (i) la métrica discreta; (ii) la métrica $d(a,b)=1/\min\{n:a_n\ne b_n\}$ cuando $a\ne b$, y $d(a,a)=0$; (iii) la métrica
\[
d'(a,b)=\sum_{n=1}^{\infty}\frac{|a_n-b_n|}{2^n(1+|a_n-b_n|)}.
\]
\end{ejercicio}


\begin{ejercicio}[Lectura 2, ejercicio 2]\label{ej:lectura-2-2}
Demostrar que una intersección finita de vecindades de un mismo punto es una vecindad de él.
\end{ejercicio}


\begin{ejercicio}[Lectura 2, ejercicio 3]\label{ej:lectura-2-3}
Probar $\cl{A\cup B}=\cl A\cup\cl B$ y $\cl{A\cap B}\subset\cl A\cap\cl B$.
\end{ejercicio}


\begin{ejercicio}[Lectura 2, ejercicio 4]\label{ej:lectura-2-4}
Dar ejemplos en los que (i) $\interior{(A\cup B)}\not\subset\interior A\cup\interior B$ y (ii) $\cl A\cap\cl B\not\subset\cl{A\cap B}$.
\end{ejercicio}


\begin{ejercicio}[Lectura 2, ejercicio 5]\label{ej:lectura-2-5}
En un espacio ultramétrico demostrar: (i) si $y\in B(x,r)$, entonces $B(y,r)=B(x,r)$; (ii) las bolas abiertas y las bolas cerradas de radio positivo son conjuntos abiertos y cerrados; (iii) dos bolas que se intersectan están ordenadas por inclusión.
\end{ejercicio}


\begin{ejercicio}[Lectura 2, ejercicio 6]\label{ej:lectura-2-6}
En cualquier espacio métrico, demostrar que (i) las bolas abiertas son abiertas, (ii) las bolas cerradas son cerradas y (iii) las esferas son cerradas.
\end{ejercicio}


\begin{ejercicio}[Lectura 2, ejercicio 7]\label{ej:lectura-2-7}
Demostrar que (i) $A$ es abierto si y solo si es vecindad de todos sus puntos; (ii) $A$ es cerrado si y solo si todo punto cuyas vecindades intersectan a $A$ pertenece a $A$.
\end{ejercicio}


\section{Clausura de una bola y bola cerrada}
Siempre $\cl{B(x,r)}\subset\overline B(x,r)$, porque la bola cerrada es cerrada y contiene a la abierta. La igualdad puede fallar: en un espacio discreto con al menos dos puntos, $B(x,1)=\{x\}$ tiene clausura $\{x\}$, mientras que $\overline B(x,1)=X$. Por eso el tipo de desigualdad que define una bola debe distinguirse de la operación topológica de tomar clausura.

\section{Operaciones que preservan la topología}\label{lectura:03}
Una vez definidos el interior y la clausura, podemos reconocer conjuntos abiertos y cerrados. Estudiaremos las operaciones que preservan esas propiedades y cómo la clausura se detecta mediante sucesiones. La distinción decisiva será entre pertenecer simplemente a un conjunto y ser un punto al que se aproximan \emph{otros} puntos de él.

\begin{proposicion}\label{ch03:axiomas}\index{topologia@topología!axiomas}
En un espacio métrico $X$, una unión arbitraria de abiertos es abierta y una intersección finita de abiertos es abierta. Los conjuntos $\varnothing$ y $X$ son abiertos.
\end{proposicion}
\begin{proof}
Sea $G=\bigcup_{i\in I}G_i$, con cada $G_i$ abierto. Si $x\in G$, existe $i$ con $x\in G_i$. Alguna bola centrada en $x$ está contenida en $G_i$ y, por ello, en $G$. Si $I$ es vacío, la unión es vacía y no hay punto en el que verificar la condición.

Para $x\in\bigcap_{j=1}^mG_j$, elijamos $r_j>0$ con $B(x,r_j)\subset G_j$. La bola de radio $\min_jr_j$ está contenida en la intersección. Por último, $X$ contiene cualquier bola centrada en cualquiera de sus puntos; $\varnothing$ es abierto porque la afirmación «es vecindad de todos sus puntos» no tiene contraejemplos.
\end{proof}

\begin{teorema}\label{ch03:complemento}
Un conjunto $F\subset X$ es cerrado si y solo si $X\setminus F$ es abierto.
\end{teorema}
\begin{proof}
Si $F$ es cerrado y $x\notin F=\cl F$, alguna bola centrada en $x$ no intersecta a $F$. Esa bola está contenida en el complemento, que es por tanto vecindad de cada punto suyo. Recíprocamente, si el complemento es abierto y $x\notin F$, existe una bola centrada en $x$ contenida en él. Entonces $x\notin\cl F$. Hemos probado $\cl F\subset F$, y la inclusión opuesta vale siempre.
\end{proof}

\begin{corolario}\label{ch03:cerrados}
Las intersecciones arbitrarias de cerrados y las uniones finitas de cerrados son cerradas. Los conjuntos $\varnothing$ y $X$ son cerrados.
\end{corolario}
\begin{proof}
Las identidades de De Morgan dan
\[
X\setminus\bigcap_iF_i=\bigcup_i(X\setminus F_i),\qquad
X\setminus\bigcup_{j=1}^mF_j=\bigcap_{j=1}^m(X\setminus F_j).
\]
Los lados derechos son abiertos por la proposición~\ref{ch03:axiomas}. Sus complementos son cerrados por el teorema anterior. Los complementos de $\varnothing$ y $X$ son abiertos, lo que prueba la última afirmación. Se admite la intersección vacía, igual a $X$.
\end{proof}
Por ejemplo, $S(x,r)=\overline B(x,r)\cap(X\setminus B(x,r))$ es intersección de dos cerrados. Esta es una segunda prueba de que las esferas son cerradas.

\begin{teorema}\label{ch03:unionbolas}
Un conjunto es abierto si y solo si es unión de bolas abiertas.
\end{teorema}
\begin{proof}
Si $G$ es abierto, para cada $x\in G$ elijamos $r_x>0$ con $B(x,r_x)\subset G$. Todo $x\in G$ pertenece a su propia bola, así que $G=\bigcup_{x\in G}B(x,r_x)$. Si $G=\varnothing$, se toma la unión vacía. Recíprocamente, cada bola abierta es abierta por el ejercicio~\ref{ej:lectura-2-6}, y una unión de abiertos es abierta.
\end{proof}
Este teorema permite comprobar propiedades de todos los abiertos a partir de bolas, siempre que la propiedad sea compatible con uniones. Esa compatibilidad habrá que verificarla en cada aplicación.

\section{Adherencia, acumulación y aislamiento}
\begin{definicion}\index{punto!de adherencia}\index{punto!de acumulacion@de acumulación}\index{conjunto derivado}
Un punto interior de $A$ pertenece a $\interior A$; un punto de adherencia pertenece a $\cl A$. Un punto $x$ es de acumulación de $A$ si
\[
B(x,r)\cap(A\setminus\{x\})\ne\varnothing\quad\text{para todo }r>0.
\]
El conjunto de estos puntos se denota $A'$ y se llama conjunto derivado. Un punto de $A$ que no está en $A'$ es aislado en $A$.
\end{definicion}
Todo punto de acumulación es de adherencia, porque la condición de intersectar a $A\setminus\{x\}$ implica la de intersectar a $A$. Más precisamente,
\begin{equation}\label{ch03:clderivada}
\cl A=A\cup A'.
\end{equation}
En efecto, si $x\in\cl A\setminus A$, cada bola intersecta a $A$, y el punto encontrado es diferente de $x$ porque $x\notin A$. Esto da $x\in A'$. La inclusión contraria usa $A\subset\cl A$ y $A'\subset\cl A$.

Un conjunto finito no tiene puntos de acumulación. Para verlo en un punto arbitrario $x$, consideremos el conjunto finito de distancias $d(x,a)$ con $a\in A\setminus\{x\}$. Si ese conjunto es vacío, cualquier bola sirve. Si no lo es, su mínimo $m$ es positivo y $B(x,m)$ no contiene puntos de $A\setminus\{x\}$. Esta prueba cubre también los puntos exteriores a $A$ y los conjuntos con cero o un elemento.

\begin{lema}\label{ch03:infinitospuntos}
Si $x\in A'$, toda bola centrada en $x$ contiene infinitos puntos de $A$.
\end{lema}
\begin{proof}
Si $B(x,r)\cap(A\setminus\{x\})$ fuera finito, eligiendo un radio menor que $r$ y menor que todas las distancias positivas de $x$ a esos puntos obtendríamos una bola sin puntos de $A\setminus\{x\}$. Esto contradice la definición de acumulación.
\end{proof}

\section{Sucesiones como pruebas de clausura}
\begin{definicion}\index{sucesion@sucesión!convergente}
Una sucesión $(x_n)$ converge a $x$ si para cada $\eps>0$ existe $N\in\N$ tal que $d(x_n,x)<\eps$ para todo $n\ge N$. Se escribe $x_n\to x$.
\end{definicion}
El índice $N$ depende de la precisión $\eps$. No basta con encontrar un término cerca de $x$; todos los términos a partir de cierto momento deben estar cerca de él.

\begin{teorema}\label{ch03:secuencial}
Para $A\subset X$ y $x\in X$:
\begin{enumerate}[label=(\roman*)]
\item $x\in\cl A$ si y solo si existe una sucesión de puntos de $A$ que converge a $x$.
\item $x\in A'$ si y solo si existe una sucesión de puntos de $A$ distintos dos a dos que converge a $x$.
\end{enumerate}
\end{teorema}
\begin{proof}
\begin{enumerate}[label=(\roman*)]
\item Si $x\in\cl A$, para cada $n$ elegimos $x_n\in A\cap B(x,1/n)$. Dado $\eps>0$, tomamos $N$ con $1/N<\eps$; entonces $n\ge N$ implica $d(x_n,x)<1/n\le1/N<\eps$. Recíprocamente, si $x_n\in A$ y $x_n\to x$, cualquier bola $B(x,r)$ contiene todos los términos a partir de cierto índice; por ello intersecta a $A$.
\item Si $x\in A'$, el lema~\ref{ch03:infinitospuntos} permite elegir recursivamente
\[
x_n\in A\cap B(x,1/n)\setminus\{x,x_1,\ldots,x_{n-1}\}.
\]
Solo se excluye un conjunto finito de una intersección infinita, así que siempre hay una elección. Las distancias menores que $1/n$ prueban la convergencia. En sentido inverso, dada una bola $B(x,r)$, la convergencia pone en ella todos los términos suficientemente avanzados. Como son distintos dos a dos, a lo sumo uno es igual a $x$; existe, por tanto, un término de la bola diferente de $x$. Esto prueba acumulación.
\end{enumerate}
\end{proof}
\begin{corolario}\label{ch03:cerradosecuencial}
Un subconjunto $F$ de un espacio métrico es cerrado si y solo si contiene el límite de toda sucesión de puntos de $F$ que converja en $X$.
\end{corolario}
\begin{proof}
Si $F$ es cerrado, el teorema anterior coloca cualquier límite en $\cl F=F$. Si contiene todos esos límites, cada $x\in\cl F$ es límite de una sucesión en $F$ y pertenece a $F$; así $\cl F=F$.
\end{proof}

\section{Ejercicios sobre cerradura, acumulación y sucesiones}
\begin{ejercicio}[Lectura 3, ejercicio 1]\label{ej:lectura-3-1}
Demostrar que (i) $\interior A$ es la unión de todos los abiertos contenidos en $A$ y (ii) $\cl A$ es la intersección de todos los cerrados que contienen a $A$.
\end{ejercicio}


\begin{ejercicio}[Lectura 3, ejercicio 2]\label{ej:lectura-3-2}
Demostrar que (i) $c_0$, el conjunto de sucesiones que convergen a cero, es cerrado en $\ell^\infty$; (ii) fijado $t\in K$, el conjunto $\{f\in C(K):f(t)=0\}$ es cerrado en la norma uniforme.
\end{ejercicio}


\begin{ejercicio}[Lectura 3, ejercicio 3]\label{ej:lectura-3-3}\index{Hausdorff, propiedad de}
Para $x\ne y$, encontrar abiertos disjuntos $U,V$ con $x\in U$ y $y\in V$. Deducir la unicidad del límite de una sucesión.
\end{ejercicio}


\begin{ejercicio}[Lectura 3, ejercicio 4]\label{ej:lectura-3-4}
Demostrar que $\{x\}$ es la intersección de todas las vecindades cerradas de $x$, y deducir que todo conjunto finito es cerrado.
\end{ejercicio}


\begin{ejercicio}[Lectura 3, ejercicio 5]\label{ej:lectura-3-5}
En $\ell^1$, sea $e_k=(\delta_{n,k})_{n\ge1}$ y $A=\{e_k:k\in\N\}$. Demostrar que $A$ no tiene puntos de acumulación.
\end{ejercicio}


\begin{ejercicio}[Lectura 3, ejercicio 6]\label{ej:lectura-3-6}
Demostrar que $A$ es cerrado si y solo si $A'\subset A$.
\end{ejercicio}


\begin{ejercicio}[Lectura 3, ejercicio 7]\label{ej:lectura-3-7}\index{distancia!a un conjunto}
Para $A\ne\varnothing$, sea $d(x,A)=\inf\{d(x,a):a\in A\}$. Probar $x\in\cl A$ si y solo si $d(x,A)=0$.
\end{ejercicio}


\begin{ejercicio}[Lectura 3, ejercicio 8]\label{ej:lectura-3-8}
Para una familia no vacía $\mathcal S$ de subconjuntos de $X$, demostrar
\[
\left(\bigcap_{A\in\mathcal S}A\right)'\subset\bigcap_{A\in\mathcal S}A',\qquad
\bigcup_{A\in\mathcal S}A'\subset\left(\bigcup_{A\in\mathcal S}A\right)',
\]
y dar ejemplos de inclusiones estrictas.
\end{ejercicio}


\begin{ejercicio}[Lectura 3, ejercicio 9]\label{ej:lectura-3-9}
En un espacio ultramétrico demostrar: (i) las relaciones $d(x,y)<r$ y $d(x,y)\le r$, con $r>0$, son de equivalencia; (ii) las bolas abiertas de radio $r$ particionan $X$, y lo mismo las cerradas; (iii) si $0<s<r$, una bola $B(x,r)$ es particionada por las bolas abiertas de radio $s$ que la intersectan.
\end{ejercicio}


\section{El espacio ambiente forma parte de la afirmación}\label{lectura:04}
Las caracterizaciones anteriores cambian de interpretación cuando restringimos el espacio ambiente. Decir que un conjunto es abierto o cerrado requiere indicar dentro de qué espacio se trabaja: $[0,1]$ es cerrado en $\R$, pero abierto dentro de sí mismo. No hay contradicción, porque las bolas disponibles cambian al restringir el universo. Las propiedades de bolas, abiertos y clausuras ya demostradas permitirán describir con precisión ese cambio.

\begin{definicion}\index{subespacio!metrico@métrico}\label{ch04:subespacio}
Si $(X,d)$ es métrico y $Y\subset X$, la restricción $d|_{Y\times Y}$ es la métrica inducida en $Y$. El espacio resultante se llama subespacio métrico.
\end{definicion}
La restricción es no negativa, simétrica y separa puntos porque conserva los valores de $d$. Para $u,v,w\in Y$, la desigualdad triangular ya es válida en $X$ y, por tanto, en $Y$. Esto prueba todos los axiomas, sin ninguna hipótesis de que $Y$ sea abierto o cerrado.

Indicaremos el ambiente mediante subíndices cuando pueda haber confusión. Para $y\in Y$,
\begin{equation}\label{ch04:bolas}
B_Y(y,r)=Y\cap B_X(y,r).
\end{equation}
La identidad se deduce porque ambos miembros exigen simultáneamente pertenecer a $Y$ y tener distancia a $y$ menor que $r$. También vale con bolas cerradas y esferas.

\begin{proposicion}\label{ch04:vecindades}
Un conjunto $V\subset Y$ es vecindad de $y\in Y$ en $Y$ si y solo si existe una vecindad $W$ de $y$ en $X$ tal que $V=W\cap Y$.
\end{proposicion}
\begin{proof}
Si $B_Y(y,r)\subset V$, definamos $W=V\cup B_X(y,r)$. Es vecindad de $y$ en $X$ porque contiene la bola ambiental. Además,
\[
W\cap Y=(V\cap Y)\cup(B_X(y,r)\cap Y)
=V\cup B_Y(y,r)=V.
\]
La última igualdad usa $B_Y(y,r)\subset V$. Es una unión, no una intersección, en este cálculo distributivo. En sentido inverso, si $V=W\cap Y$ y $B_X(y,r)\subset W$, entonces $B_Y(y,r)=B_X(y,r)\cap Y\subset V$.
\end{proof}

\begin{teorema}\label{ch04:relativa}\index{topologia@topología!relativa}
Para $A\subset Y\subset X$:
\begin{enumerate}[label=(\roman*)]
\item $A$ es abierto en $Y$ si y solo si $A=Y\cap G$ para algún abierto $G$ de $X$.
\item $A$ es cerrado en $Y$ si y solo si $A=Y\cap F$ para algún cerrado $F$ de $X$.
\end{enumerate}
\end{teorema}
\begin{proof}
\begin{enumerate}[label=(\roman*)]
\item Si $A$ es abierto en $Y$, para cada $a\in A$ elegimos $r_a>0$ con $B_Y(a,r_a)\subset A$. Pongamos $G=\bigcup_{a\in A}B_X(a,r_a)$, abierto en $X$. Entonces
\[
Y\cap G=\bigcup_{a\in A}B_Y(a,r_a)=A,
\]
pues cada bola está en $A$ y cada $a$ pertenece a su bola. Si $A$ es vacío, se toma $G=\varnothing$. Recíprocamente, si $A=Y\cap G$ y $a\in A$, la apertura de $G$ proporciona una bola $B_X(a,r)\subset G$; su intersección con $Y$ está en $A$ y demuestra apertura relativa.
\item $A$ es cerrado en $Y$ si y solo si $Y\setminus A$ es abierto en $Y$. Por (i), esto equivale a $Y\setminus A=Y\cap G$ para algún abierto $G$ en $X$. Tomando complementos relativos, $A=Y\cap(X\setminus G)$, donde $X\setminus G$ es cerrado en $X$. En sentido inverso, si $A=Y\cap F$ con $F$ cerrado, $Y\setminus A=Y\cap(X\setminus F)$ es abierto en $Y$ por (i).
\end{enumerate}
\end{proof}

\begin{proposicion}\label{ch04:clausura}
Si $A\subset Y$, entonces $\cl A^{\,Y}=Y\cap\cl A^{\,X}$, donde los superíndices indican el espacio de clausura. En consecuencia, $A$ es cerrado en $Y$ si y solo si $A=Y\cap\cl A^{\,X}$.
\end{proposicion}
\begin{proof}
Para $y\in Y$, la condición de pertenecer a la clausura relativa es que $B_Y(y,r)\cap A\ne\varnothing$ para todo $r>0$. Como $A\subset Y$, esta intersección es exactamente $B_X(y,r)\cap A$. La condición equivale, por tanto, a $y\in\cl A^{\,X}$. Para ser cerrado, $A$ debe coincidir con su clausura relativa.
\end{proof}
La clausura ambiental puede añadir puntos que no están disponibles en el subespacio. En $Y=(0,1)$, el conjunto $A=Y$ es cerrado relativamente, pero su clausura en $\R$ es $[0,1]$.

\begin{observacion}[Convergencia relativa]\label{ch04:convergencia}
Una sucesión de puntos de $Y$ converge a $y\in Y$ en el subespacio si y solo si converge a $y$ en $X$: las distancias numéricas son idénticas. Sin embargo, puede converger en $X$ a un punto exterior a $Y$, en cuyo caso no tiene ese límite en el subespacio. Por ejemplo, $1/n$ converge a $0$ en $\R$, pero no converge en $(0,1]$.
\end{observacion}

\section{Ejercicios sobre topología relativa}
\begin{ejercicio}[Lectura 4, ejercicio 1]\label{ej:lectura-4-1}
Demostrar que los abiertos y cerrados relativos son intersecciones con abiertos y cerrados del espacio ambiente.
\end{ejercicio}


\begin{ejercicio}[Lectura 4, ejercicio 2]\label{ej:lectura-4-2}
Sea $A\subset Y\subset X$. Probar: (i) si $A$ es abierto en $Y$ y $Y$ abierto en $X$, entonces $A$ es abierto en $X$; (ii) si $A$ es cerrado en $Y$ y $Y$ cerrado en $X$, entonces $A$ es cerrado en $X$.
\end{ejercicio}


\begin{ejercicio}[Lectura 4, ejercicio 3]\label{ej:lectura-4-3}
Considerar $\Q\subset\R$ con la métrica usual. Para un intervalo abierto relativo no vacío y acotado $I=(a,b)\cap\Q$, con $a<b$, demostrar que $I$ es cerrado en $\Q$ si y solo si ambos extremos de su clausura en $\R$ son irracionales. Precisar el caso de intervalos no acotados.
\end{ejercicio}


\begin{ejercicio}[Lectura 4, ejercicio 4]\label{ej:lectura-4-4}
Para $A\subset Y\subset X$, demostrar que el derivado relativo $A'_Y$ está contenido en $A'_X$, y dar un ejemplo de inclusión estricta.
\end{ejercicio}


\begin{ejercicio}[Lectura 4, ejercicio 5: corrección del enunciado]\label{ej:lectura-4-5}
La lectura propone que la imagen de una subsucesión convergente de una sucesión acotada es cerrada en la imagen de la sucesión original. Examinar la afirmación y establecer una versión válida.
\end{ejercicio}


\begin{ejercicio}[Lectura 4, ejercicio 6]\label{ej:lectura-4-6}
¿Es cierto que $A\subset Y$ es abierto en $Y$ si y solo si $A=Y\cap\operatorname{int}_X A$?
\end{ejercicio}


\section{Distancias, límites y espacio ambiente}
Una norma genera una distancia, pero las distancias tienen un ámbito más amplio. Hölder controla productos y Minkowski controla sumas. Antes de llamar norma a una expresión, hay que verificar que esté definida, que su dominio sea un espacio vectorial y que únicamente se anule en el vector cero. Esta última comprobación es esencial cuando aparecen integrales.

En espacios métricos, las sucesiones detectan exactamente la clausura y la cerradura. El requisito de términos distintos detecta acumulación y elimina la información trivial de las sucesiones constantes. Estas distinciones serán decisivas al estudiar compacidad, donde las sucesiones dejarán de ser solo una herramienta de lenguaje y se convertirán en una fuente de existencia de límites.

La clausura relativa se obtiene tomando la clausura ambiental y eliminando los puntos exteriores al subespacio. Para el interior no basta esa misma receta: primero hay que permitir todos los puntos exteriores a $Y$. Registrar el ambiente en la notación evita confundir afirmaciones sobre límites, apertura y cerradura que pueden cambiar al pasar de un espacio a otro.

\chapter{Continuidad, densidad y uniformidad}\label{cap:nuevo02}\label{chap:05}\label{chap:06}\label{chap:07}

La topología métrica permite precisar qué significa que una función respete la proximidad. En este capítulo estudiaremos primero la continuidad local y sus diferentes criterios; después veremos cómo los conjuntos densos y las bases numerables organizan las aproximaciones; finalmente distinguiremos el control local del control uniforme, que preserva sucesiones de Cauchy y permite extender funciones a puntos límite. Se integran las lecturas 5--7 de \cite{valencia}. Usaremos las propiedades de distancias, clausuras y subespacios del capítulo~\ref{cap:nuevo01}; los resultados posteriores se apoyarán en esos criterios y en el control explícito de los cuantificadores.

\section{Controlar la imagen de puntos próximos}\label{lectura:05}
Para estudiar la continuidad hay que distinguir la \emph{imagen} $f(A)$ de la \emph{imagen inversa} $f^{-1}(B)$. La primera se aplica a subconjuntos del dominio y la segunda a subconjuntos del codominio; la notación de imagen inversa no exige que $f$ sea invertible. Con esta distinción, conectaremos las formulaciones mediante bolas y tolerancias, abiertos y sucesiones.

\begin{definicion}\label{ch05:continua}\index{continuidad!en un punto}
Sean $(X,d)$ y $(Y,\rho)$ espacios métricos. Una función $f:X\to Y$ es continua en $x\in X$ si para toda vecindad $V$ de $f(x)$, el conjunto $f^{-1}(V)$ es vecindad de $x$.
\end{definicion}

\begin{teorema}[Criterio de tolerancias]\label{ch05:epsdelta}
La función $f$ es continua en $x$ si y solo si
\[
\forall\eps>0\ \exists\delta>0\ \forall u\in X:
\quad d(x,u)<\delta\ \Longrightarrow\ \rho(f(x),f(u))<\eps.
\]
\end{teorema}
\begin{proof}
Si $f$ es continua en $x$, la imagen inversa de la vecindad $B_Y(f(x),\eps)$ contiene alguna bola $B_X(x,\delta)$. La inclusión de bolas es exactamente la implicación del enunciado. Recíprocamente, sea $V$ una vecindad de $f(x)$ y elijamos $\eps>0$ con $B_Y(f(x),\eps)\subset V$. La condición proporciona $\delta>0$ con $B_X(x,\delta)\subset f^{-1}(B_Y(f(x),\eps))\subset f^{-1}(V)$, así que esta última es vecindad de $x$.
\end{proof}
El orden de los cuantificadores explica el significado: primero se prescribe una precisión en las imágenes; después se elige una tolerancia en el dominio que sirve para \emph{todos} los puntos suficientemente próximos. La tolerancia puede depender del punto $x$; exigir una sola tolerancia para todos los puntos conduce más adelante a continuidad uniforme.

\begin{teorema}[Criterios globales]\label{ch05:criterios}\index{continuidad!criterios topologicos@criterios topológicos}
Para $f:X\to Y$, son equivalentes:
\begin{enumerate}[label=(\roman*)]
\item $f^{-1}(U)$ es abierto en $X$ para todo abierto $U$ en $Y$.
\item Para todo $A\subset X$, $f(\cl A)\subset\cl{f(A)}$.
\item $f^{-1}(F)$ es cerrado en $X$ para todo cerrado $F$ en $Y$.
\item $f$ es continua en cada punto de $X$.
\end{enumerate}
Cuando se cumplen, se dice que $f$ es continua en $X$.
\end{teorema}
\begin{proof}
\emph{(i)$\Rightarrow$(ii).} Si $x\in\cl A$ y $U$ es abierto con $f(x)\in U$, entonces $f^{-1}(U)$ es un abierto que contiene a $x$. Por ello intersecta a $A$, y la imagen de un punto de esa intersección pertenece a $U\cap f(A)$. Como toda bola alrededor de $f(x)$ es uno de esos abiertos, $f(x)\in\cl{f(A)}$.

\emph{(ii)$\Rightarrow$(iii).} Sea $F$ cerrado y $E=f^{-1}(F)$. La inclusión $f(E)\subset F$ implica $\cl{f(E)}\subset F$. Si $x\in\cl E$, por (ii), $f(x)\in f(\cl E)\subset\cl{f(E)}\subset F$, de modo que $x\in E$. Así $E$ es cerrado.

\emph{(iii)$\Rightarrow$(i).} Para $U$ abierto, $Y\setminus U$ es cerrado y
\[
X\setminus f^{-1}(U)=f^{-1}(Y\setminus U)
\]
es cerrado. Por tanto $f^{-1}(U)$ es abierto.

\emph{(i)$\Rightarrow$(iv).} Si $V$ es vecindad de $f(x)$, contiene una bola abierta $U$ centrada en $f(x)$. Por (i), $f^{-1}(U)$ es abierto y contiene a $x$; contiene entonces una bola alrededor de $x$. Como $f^{-1}(U)\subset f^{-1}(V)$, esta última es vecindad de $x$.

\emph{(iv)$\Rightarrow$(i).} Sea $x\in f^{-1}(U)$, con $U$ abierto. Este $U$ es vecindad de $f(x)$, así que $f^{-1}(U)$ es vecindad de $x$. Lo es para cada uno de sus puntos y, por tanto, es abierto.
\end{proof}
No se afirma en (ii) igualdad. Por ejemplo, $f(t)=\arctan t$ es continua en $\R$, pero para $A=\R$ su imagen es $(-\pi/2,\pi/2)$ y la clausura de esa imagen contiene además los dos extremos.

\section{Ejemplos y reglas de construcción}
Toda función cuyo dominio lleva la topología discreta es continua, porque cualquier imagen inversa es un subconjunto abierto del dominio. La identidad de un espacio en sí mismo es continua, ya que sus imágenes inversas no cambian los conjuntos.

Si $f:X\to Y$ y $g:Y\to Z$ son continuas, también lo es $g\circ f$: para un abierto $U\subset Z$,
\[
(g\circ f)^{-1}(U)=f^{-1}(g^{-1}(U)),
\]
y primero la continuidad de $g$ y luego la de $f$ garantizan apertura. Se está usando la preimagen de conjuntos, no una inversa funcional de $f$ o $g$.

Una función $L$-Lipschitz, es decir, una función con $\rho(f(x),f(y))\le Ld(x,y)$, es continua. Si $L>0$, basta tomar $\delta=\eps/L$; si $L=0$, todas las imágenes coinciden y cualquier $\delta>0$ sirve. Así, la función $x\mapsto d(x,A)$ para $A\ne\varnothing$ es continua por el ejercicio~\ref{ej:lectura-3-7}, y la norma $v\mapsto\norm v$ es continua porque
\[
|\norm u-\norm v|\le\norm{u-v}.
\]
Esta última desigualdad se prueba aplicando la triangular a $u=(u-v)+v$ y luego intercambiando $u,v$.

La inclusión $Y\hookrightarrow X$ de un subespacio es continua: la preimagen de un abierto ambiental es su intersección con $Y$, abierta relativamente. Una restricción $f|_Y$ de una función continua es continua por composición con esa inclusión. Además, si una función toma valores en un subespacio $Z\subset Y$, es continua como función a $Z$ si y solo si lo es como función a $Y$. En una dirección se compone con la inclusión; en la otra, para $U=Z\cap G$ abierto en $Z$, se tiene $f^{-1}(U)=f^{-1}(G)$ porque toda imagen pertenece a $Z$. Estos hechos incorporan el resultado sobre subespacios que aparece repetido en la lectura original.

\begin{teorema}[Criterio secuencial]\label{ch05:secuencial}\index{continuidad!criterio secuencial}
Una función entre espacios métricos es continua en $x$ si y solo si $x_n\to x$ implica $f(x_n)\to f(x)$.
\end{teorema}
\begin{proof}
Supongamos continuidad y fijemos $\eps>0$. Elegimos $\delta>0$ según el teorema~\ref{ch05:epsdelta}. Como $x_n\to x$, hay $N$ tal que $d(x_n,x)<\delta$ para $n\ge N$. Entonces $\rho(f(x_n),f(x))<\eps$ para esos índices.

Para el recíproco usamos la negación exacta de continuidad: existe $\eps_0>0$ tal que para \emph{todo} $\delta>0$ existe $u\in X$ con $d(u,x)<\delta$ y $\rho(f(u),f(x))\ge\eps_0$. Tomando $\delta=1/n$, elegimos $x_n$ con esas propiedades. Entonces $x_n\to x$, pero ninguna imagen entra en la bola $B(f(x),\eps_0)$. Esto contradice el criterio secuencial. La desigualdad de la negación es $\ge\eps_0$, y no es necesario reemplazarla por una estricta.
\end{proof}

\section{Funciones abiertas, cerradas y homeomorfismos}
\begin{definicion}\index{funcion@función!abierta}\index{funcion@función!cerrada}\index{homeomorfismo}
Una función $f:X\to Y$ es abierta si $f(U)$ es abierto para cada abierto $U\subset X$, y cerrada si $f(F)$ es cerrado para cada cerrado $F\subset X$. Es un homeomorfismo si es biyectiva, continua y su inversa $f^{-1}:Y\to X$ es continua. En este caso se escribe $X\cong Y$.
\end{definicion}
Estas condiciones sobre imágenes directas son diferentes de continuidad. La identidad de $\R$ con su métrica discreta hacia $\R$ usual es continua, pero no es abierta, pues la imagen de $\{0\}$ no es abierta, ni cerrada, pues la imagen del conjunto discreto cerrado $\Q$ no es cerrada en la recta usual. La identidad en la dirección contraria es abierta y cerrada, porque todo subconjunto del codominio discreto es abierto y cerrado; sin embargo, no es continua.

\begin{ejemplo}\label{ch05:homeorecta}
La función $h:\R\to(-1,1)$ dada por $h(t)=t/(1+|t|)$ es un homeomorfismo. Su valor tiene módulo menor que $1$. Para $s\in(-1,1)$, definamos $k(s)=s/(1-|s|)$. Como $|k(s)|=|s|/(1-|s|)$,
\[
h(k(s))=\frac{s/(1-|s|)}{1+|s|/(1-|s|)}=s.
\]
Del mismo modo, $1-|h(t)|=1/(1+|t|)$ implica $k(h(t))=t$. Por tanto son funciones inversas y ambas son continuas por las reglas elementales de cálculo: los denominadores son positivos en sus respectivos dominios. Esto demuestra biyectividad y continuidad de las dos direcciones, no solamente una semejanza geométrica.
\end{ejemplo}

\section{Ejercicios sobre continuidad y homeomorfismos}
\begin{ejercicio}[Lectura 5, ejercicio 1]\label{ej:lectura-5-1}
Demostrar que cada una de las siguientes condiciones equivale a continuidad de $f:X\to Y$:
\begin{enumerate}[label=(\roman*)]
\item Para todo $A\subset Y$, $f^{-1}(\operatorname{int}A)\subset\operatorname{int}(f^{-1}(A))$.
\item Para todo $A\subset Y$, $\cl{f^{-1}(A)}\subset f^{-1}(\cl A)$.
\end{enumerate}
En (ii), se corrige la letra $B$, no definida en el original, por $A$.
\end{ejercicio}
\begin{solucion}
\begin{enumerate}[label=(\roman*)]
\item Si $f$ es continua, $f^{-1}(\operatorname{int}A)$ es abierto y está contenido en $f^{-1}(A)$. El interior de un conjunto contiene a todos sus subconjuntos abiertos; de ahí la inclusión. Si la condición vale y $U$ es abierto, entonces $\operatorname{int}U=U$ y $f^{-1}(U)\subset\operatorname{int}(f^{-1}(U))$. La inclusión inversa es automática, de modo que $f^{-1}(U)$ es abierto y $f$ continua.
\item Si $f$ es continua, $f^{-1}(\cl A)$ es cerrado y contiene a $f^{-1}(A)$; por ello contiene a su clausura. Recíprocamente, para un cerrado $F$ tenemos $\cl F=F$, y la condición da $\cl{f^{-1}(F)}\subset f^{-1}(F)$. Por tanto todas las preimágenes de cerrados son cerradas y $f$ es continua.
\end{enumerate}
\end{solucion}

\begin{ejercicio}[Lectura 5, ejercicio 2]\label{ej:lectura-5-2}
Demostrar que una función es continua en el espacio si y solo si es continua en cada punto.
\end{ejercicio}
\begin{solucion}
Las dos direcciones están desarrolladas en el teorema~\ref{ch05:criterios}. El mecanismo es local: la preimagen de un abierto $U$ es abierta exactamente cuando contiene una bola alrededor de cada $x$ con $f(x)\in U$. La continuidad en ese $x$ produce esa bola, y la continuidad global garantiza que la preimagen de una bola alrededor de $f(x)$ es un abierto que contiene a $x$. Así se recupera en cada dirección la definición requerida.
\end{solucion}

\begin{ejercicio}[Lectura 5, ejercicio 3]\label{ej:lectura-5-3}
Probar que $f$ es continua si y solo si las preimágenes de todas las bolas abiertas son abiertas. Dar una función no continua cuyas preimágenes de todas las bolas cerradas sean cerradas.
\end{ejercicio}
\begin{solucion}
La necesidad sigue de que las bolas abiertas son abiertas. Para la suficiencia, todo abierto $U$ del codominio es una unión $\bigcup_iB_i$ de bolas abiertas. Como las preimágenes preservan uniones,
\[
f^{-1}(U)=\bigcup_i f^{-1}(B_i),
\]
que es abierto por hipótesis.

Para la segunda parte, tomemos $X=\R$ con la métrica usual, $Y=\R$ con la discreta y $f(t)=t$. Una bola cerrada de radio $r<1$ en $Y$ es un singleton, y una de radio $r\ge1$ es todo $Y$. Sus preimágenes son singletons o $\R$, todos cerrados en $X$. Pero $\{0\}$ es abierto en $Y$ y su preimagen no es abierta en $X$; así $f$ no es continua. El argumento con bolas abiertas no se transfiere a bolas cerradas, porque una unión arbitraria de cerrados no tiene por qué ser cerrada.
\end{solucion}

\begin{ejercicio}[Lectura 5, ejercicio 4]\label{ej:lectura-5-4}
Sean $f,g:X\to\R$ continuas. Probar que (i) $\{x:f(x)<g(x)\}$ es abierto y (ii) $\{x:f(x)\le g(x)\}$ es cerrado. Se corrige el dominio de $g$ a $X$, necesario para que ambas expresiones estén definidas.
\end{ejercicio}
\begin{solucion}
\begin{enumerate}[label=(\roman*)]
\item Si $f(x)<g(x)$, sea $\eta=(g(x)-f(x))/3>0$. La continuidad de ambas funciones permite elegir $\delta>0$ tal que $d(x,y)<\delta$ implique $|f(y)-f(x)|<\eta$ y $|g(y)-g(x)|<\eta$. Entonces
\[
f(y)<f(x)+\eta<g(x)-\eta<g(y).
\]
La bola está contenida en el conjunto, que es abierto.
\item El complemento es $\{x:g(x)<f(x)\}$, abierto por (i) aplicado con los papeles intercambiados. Por tanto el conjunto original es cerrado.
\end{enumerate}
\end{solucion}

\begin{ejercicio}[Lectura 5, ejercicio 5]\label{ej:lectura-5-5}\index{extension@extensión!por cero}\index{frontera}
Sea $A\subset X$ cerrado y $f:A\to\R$ continua, con $f=0$ en la frontera $\partial A=\cl A\cap\cl{X\setminus A}$. Demostrar que la extensión $g:X\to\R$, igual a $f$ en $A$ y a cero fuera de $A$, es continua.
\end{ejercicio}
\begin{solucion}
Distinguimos tres posiciones del punto $x$.

Si $x\notin A$, como $A$ es cerrado existe una bola alrededor de $x$ contenida en $X\setminus A$. Allí $g$ es constante y es continua en $x$.

Si $x\in\operatorname{int}A$, existe $r>0$ con $B_X(x,r)\subset A$. Dado $\eps>0$, la continuidad relativa de $f$ da $\delta_0>0$ tal que $y\in A$ y $d(x,y)<\delta_0$ implican $|f(y)-f(x)|<\eps$. Tomando $\delta=\min\{r,\delta_0\}$ se obtiene la misma desigualdad para $g$ y todo $y\in X$ con $d(x,y)<\delta$.

Finalmente, si $x\in A\setminus\operatorname{int}A$, todo entorno de $x$ intersecta al complemento de $A$, así que $x\in\partial A$ y $g(x)=f(x)=0$. Por continuidad relativa existe $\delta>0$ con $|f(y)|<\eps$ para $y\in A\cap B(x,\delta)$. Para $y$ de esa bola fuera de $A$, $g(y)=0$; para los puntos dentro de $A$, $|g(y)|=|f(y)|<\eps$. La misma tolerancia sirve para ambos lados de la frontera. Por tanto $g$ es continua en todos los puntos.
\end{solucion}

\begin{ejercicio}[Lectura 5, ejercicio 6: corrección del recíproco]\label{ej:lectura-5-6}\index{soporte}
Sea $S=\supp(f)=\cl{\{x:f(x)\ne0\}}$. Examinar la afirmación «$f$ es continua si y solo si $f=0$ en $\partial S$» y establecer una versión válida.
\end{ejercicio}
\begin{solucion}
La implicación directa es verdadera. Si $x\in\partial S$ y $f(x)\ne0$, por continuidad hay una bola en la que $|f(y)-f(x)|<|f(x)|/2$, así que $|f(y)|>|f(x)|/2>0$. Esa bola está contenida en $S$, lo que convierte a $x$ en punto interior de $S$ y contradice $x\in\partial S$.

El recíproco es falso. En $\R$, tomemos $f(x)=1$ para $x<0$ y $f(x)=2$ para $x\ge0$. La función no es continua en $0$, pero nunca se anula, de modo que $S=\R$ y $\partial S=\varnothing$. La condición en la frontera se satisface de forma vacía y no controla el salto interior.

Una versión exacta es: $f$ es continua en $X$ si y solo si $f|_S$ es continua en el subespacio $S$ y $f=0$ en $\partial S$. La necesidad se sigue de la continuidad de las restricciones y de la prueba anterior. Para la suficiencia, $S$ es cerrado por definición, $f$ es cero en $X\setminus S$ y el ejercicio~\ref{ej:lectura-5-5} aplica a la función continua $f|_S$. Su extensión por cero es precisamente $f$.
\end{solucion}

\begin{ejercicio}[Lectura 5, ejercicio 7]\label{ej:lectura-5-7}
Probar: (i) $f$ es abierta si y solo si las imágenes de las bolas abiertas son abiertas; (ii) si $f$ es biyectiva, es abierta si y solo si es cerrada; (iii) construir una función suprayectiva abierta que no sea cerrada y una cerrada que no sea abierta.
\end{ejercicio}
\begin{solucion}
\begin{enumerate}[label=(\roman*)]
\item La necesidad es inmediata porque una bola abierta es un abierto. Para la suficiencia, escribimos un abierto $U$ como unión de bolas $B_i$; entonces $f(U)=\bigcup_if(B_i)$ es una unión de abiertos.
\item Para una biyección se cumple $f(X\setminus A)=Y\setminus f(A)$. Si $f$ es abierta y $F$ cerrado, el conjunto $Y\setminus f(F)=f(X\setminus F)$ es abierto, así que $f(F)$ es cerrado. Si $f$ es cerrada y $U$ abierto, el conjunto $Y\setminus f(U)=f(X\setminus U)$ es cerrado, de modo que $f(U)$ es abierto. La identidad de complementos utiliza tanto inyectividad como suprayectividad.
\item La proyección $\pi:\R^2\to\R$, $\pi(x,y)=x$, es suprayectiva y abierta. Si $a\in\pi(U)$ con $U$ abierto, existe $b$ con $(a,b)\in U$ y un disco de radio $r>0$ centrado en $(a,b)$ contenido en $U$. Entonces $(a-r,a+r)\subset\pi(U)$, pues $(t,b)$ pertenece a ese disco cuando $|t-a|<r$. Sin embargo, el conjunto $F=\{(x,y):xy=1\}$ es cerrado, como preimagen de $\{1\}$ por la función continua producto, y $\pi(F)=\R\setminus\{0\}$ no es cerrado.

Para una función incluso suprayectiva y cerrada que no sea abierta, tomemos $h:\R\to[0,\infty)$, $h(t)=\max\{t,0\}$. Si $F$ es cerrado, su imagen es $F\cap[0,\infty)$, con el punto $0$ añadido cuando $F$ intersecta $(-\infty,0)$. En ambos casos es un cerrado de $[0,\infty)$. Pero $h((-2,-1))=\{0\}$ no es abierto en $[0,\infty)$: toda bola relativa alrededor de $0$ contiene puntos positivos.
\end{enumerate}
\end{solucion}

\begin{ejercicio}[Lectura 5, ejercicio 8]\label{ej:lectura-5-8}\index{pegado!sobre abiertos}
Sea $(A_i)_{i\in I}$ una cubierta abierta de $X$ y sean $f_i:A_i\to Y$ continuas que coinciden en todas las intersecciones. Probar que la función obtenida pegándolas es continua.
\end{ejercicio}
\begin{solucion}
Para $x\in X$, elijamos $i$ con $x\in A_i$ y definamos $g(x)=f_i(x)$. Si también $x\in A_j$, la hipótesis da $f_i(x)=f_j(x)$; por tanto la definición no depende de la elección. Cualquier función con las restricciones prescritas tiene necesariamente este valor en cada $x$, de modo que es única.

Sea $U$ abierto en $Y$. Entonces
\[
g^{-1}(U)=\bigcup_{i\in I}f_i^{-1}(U).
\]
Cada preimagen es abierta en $A_i$ por continuidad. Como $A_i$ es abierto en $X$, el ejercicio~\ref{ej:lectura-4-2} garantiza que también es abierta en $X$. Su unión es abierta y el criterio global prueba continuidad de $g$. No se requiere que la cubierta sea finita, porque las uniones arbitrarias de abiertos preservan apertura.
\end{solucion}

\begin{ejercicio}[Lectura 5, ejercicio 9]\label{ej:lectura-5-9}
Demostrar que ser homeomorfos es una relación de equivalencia entre espacios métricos.
\end{ejercicio}
\begin{solucion}
La identidad de $X$ es una biyección continua cuya inversa es ella misma; esto prueba reflexividad. Si $h:X\to Y$ es un homeomorfismo, su inversa $h^{-1}:Y\to X$ es biyectiva y continua y tiene inversa continua $h$; esto prueba simetría. Si $h:X\to Y$ y $k:Y\to Z$ son homeomorfismos, $k\circ h$ es una biyección continua y su inversa es $h^{-1}\circ k^{-1}$, continua por composición. Así se prueba transitividad. Cuando se habla de la clase de todos los espacios, «relación de equivalencia» se entiende en ese sentido de clases; sobre cualquier colección que sea un conjunto se obtiene la relación de equivalencia usual.
\end{solucion}

\section{Aproximar un espacio con pocos puntos}\label{lectura:06}
Los criterios de continuidad se han formulado mediante aproximaciones y entornos. Una sucesión solo utiliza un conjunto numerable de términos; cabe preguntar cuándo todo el espacio puede describirse, en términos de aproximación, a partir de un conjunto numerable. La densidad responderá a la primera parte de la pregunta, y la relación entre separabilidad y bases numerables permitirá controlar después familias de abiertos.

Un conjunto es \emph{a lo sumo numerable} si es finito o puede enumerarse mediante números naturales. Usaremos que una unión numerable de conjuntos numerables es numerable y que un producto finito de conjuntos numerables es numerable. Por ejemplo, los pares de naturales pueden enumerarse recorriendo las diagonales $m+n=2,3,\ldots$; una familia numerable de enumeraciones se reduce así a una lista de pares. Las elecciones necesarias se entienden en el marco usual de la teoría de conjuntos con el axioma de elección.

\begin{teorema}\label{ch06:denso}\index{conjunto!denso}
Para $A\subset X$, con $X$ métrico no vacío, son equivalentes:
\begin{enumerate}[label=(\roman*)]
\item $A\ne\varnothing$ y $d(x,A)=0$ para todo $x\in X$.
\item Todo abierto no vacío de $X$ intersecta a $A$.
\item $\cl A=X$.
\end{enumerate}
En este caso se dice que $A$ es denso en $X$.
\end{teorema}
\begin{proof}
Si (i) se cumple y $U$ es abierto no vacío, elijamos $x\in U$ y $r>0$ con $B(x,r)\subset U$. La igualdad $d(x,A)=0$ produce $a\in A$ con $d(x,a)<r$, así que $a\in U$. Si (ii) se cumple, cada bola $B(x,r)$ es un abierto no vacío e intersecta a $A$, lo que implica $x\in\cl A$ para cada $x$. Finalmente, si $\cl A=X$, entonces $A$ es no vacío porque $X$ lo es, y el ejercicio~\ref{ej:lectura-3-7} da $d(x,A)=0$.
\end{proof}

\begin{ejemplo}[Racionales y truncaciones]
En un espacio discreto, cada singleton es abierto. Por tanto un conjunto denso debe contener cada punto, y el único conjunto denso es el espacio entero.

En $\R$, $\Q$ es denso. Si $a<b$, elegimos $n$ con $1/n<b-a$. Existe un entero $m$ tal que $na<m\le na+1$; por ejemplo, $m=\lfloor na\rfloor+1$. Entonces $a<m/n\le a+1/n<b$. Así cada intervalo abierto no vacío contiene un racional.

En $\ell^1$, el subespacio $c_{00}=\operatorname{span}\{e_k:k\ge1\}$ de sucesiones con soporte finito es denso. Para $a\in\ell^1$, su truncación $a^{[N]}=(a_1,\ldots,a_N,0,\ldots)$ satisface
\[
\norm{a-a^{[N]}}_1=\sum_{k>N}|a_k|\longrightarrow0.
\]
La convergencia de la serie justifica que sus colas tiendan a cero. Sin embargo, $c_{00}\ne\ell^1$: la sucesión $(1/n^2)$ pertenece a $\ell^1$ y tiene infinitas coordenadas no nulas. Una familia puede generar un subespacio \emph{denso} sin generar algebraicamente todo el espacio.\index{c00@$c_{00}$}
\end{ejemplo}
Más adelante se demostrará que los polinomios son densos en $C([0,1])$ con la norma uniforme. Esa afirmación necesita un teorema de aproximación, y no se usará como premisa antes de demostrarlo.

\begin{definicion}\index{separabilidad}
Un espacio métrico es separable si posee un subconjunto denso a lo sumo numerable.
\end{definicion}
Así, $\R$ es separable aunque no sea numerable. Un espacio discreto es separable exactamente cuando es numerable, porque su único subconjunto denso es él mismo. No debe confundirse la cantidad de puntos del espacio con la cantidad de puntos necesaria para aproximarlos.

\section{Bases numerables}
\begin{definicion}\index{base!de una topologia@de una topología}\index{segunda numerabilidad}
Una base de la topología de $X$ es una familia $\mathcal B$ de abiertos tal que todo abierto de $X$ es unión de elementos de $\mathcal B$. El espacio es segundo numerable si posee una base a lo sumo numerable.
\end{definicion}
Una formulación equivalente, útil para las pruebas, es: para cada abierto $U$ y cada $x\in U$, existe $B\in\mathcal B$ con $x\in B\subset U$. Si todo abierto es unión de elementos de la base, uno de los que componen $U$ contiene a $x$. En sentido contrario, tomando esos elementos para todos los $x\in U$, su unión es exactamente $U$.

\begin{lema}\label{ch06:baseadenso}
De una base $\mathcal B$ puede obtenerse un conjunto denso cuya cardinalidad no exceda la de $\mathcal B$.
\end{lema}
\begin{proof}
Eliminemos el conjunto vacío de la base, si aparece: no aporta puntos a ninguna unión, por lo que esta sigue siendo base. Para cada $B\in\mathcal B$ elegimos $x_B\in B$, y ponemos $D=\{x_B:B\in\mathcal B\}$. Todo abierto no vacío contiene algún elemento no vacío de la base y, en consecuencia, contiene el punto elegido en él. Por el teorema~\ref{ch06:denso}, $D$ es denso. La aplicación $B\mapsto x_B$ es una sobreyección de la base sobre $D$, lo que da la cota cardinal.
\end{proof}

\begin{teorema}\label{ch06:separablebase}
Un espacio métrico es separable si y solo si es segundo numerable.
\end{teorema}
\begin{proof}
Una base numerable proporciona un conjunto denso numerable por el lema anterior. Para la otra dirección, sea $D=\{a_n:n\in\N\}$ denso; si $D$ es finito no vacío, se pueden repetir sus elementos en una enumeración. Definamos
\[
\mathcal B=\{B(a_n,2^{-k}):n,k\in\N\}.
\]
Esta familia es numerable porque está indexada por pares de naturales. Sea $U$ abierto y $x\in U$; elijamos $\eps>0$ con $B(x,\eps)\subset U$, y luego $k$ con $2^{1-k}<\eps$. La densidad permite elegir $a_n$ con $d(x,a_n)<2^{-k}$. Entonces $x\in B(a_n,2^{-k})$ y, para $y$ en esa bola,
\[
d(x,y)\le d(x,a_n)+d(a_n,y)<2\cdot2^{-k}<\eps.
\]
Por tanto $x\in B(a_n,2^{-k})\subset U$. Esta es la propiedad de base. En el espacio vacío, la familia vacía es una base y el conjunto vacío es denso en él, así que ambas propiedades también coinciden.
\end{proof}
La prueba revela por qué los centros densos son suficientes: se aproxima primero el centro deseado, reservando la mitad de la tolerancia para el radio de la bola. No basta elegir un centro cercano sin controlar también ese radio.

\begin{teorema}[Lindelöf]\label{ch06:lindelof}\index{Lindelof@Lindelöf, teorema de}
Sea $X$ métrico separable y $\mathcal U$ una familia de abiertos. Existe una subfamilia a lo sumo numerable con la misma unión que $\mathcal U$. En particular, toda cubierta abierta de $X$ admite una subcubierta numerable.
\end{teorema}
\begin{proof}
Sea $\mathcal B=\{B_n:n\in\N\}$ una base numerable. Consideremos los índices $n$ para los que $B_n$ está contenido en algún miembro de $\mathcal U$, y para cada uno elijamos un solo $U_n\in\mathcal U$ con $B_n\subset U_n$. La familia elegida es numerable.

Si $x\in\bigcup\mathcal U$, existe $U\in\mathcal U$ con $x\in U$. La propiedad de base proporciona un $B_n$ con $x\in B_n\subset U$. Este índice está entre los seleccionados; aunque el abierto elegido $U_n$ pueda ser distinto de $U$, contiene a $B_n$ y, por tanto, a $x$. Así $\bigcup\mathcal U\subset\bigcup_nU_n$. La inclusión contraria se sigue de que todos los $U_n$ pertenecen a $\mathcal U$. Si la unión es vacía, la subfamilia vacía sirve.
\end{proof}
El teorema no dice que la subcubierta pueda ser finita; esa exigencia define una propiedad más fuerte que estudiaremos bajo el nombre de compacidad.

\section{Ejercicios sobre densidad y numerabilidad}
\begin{ejercicio}[Lectura 6, ejercicio 1]\label{ej:lectura-6-1}
Demostrar que $\ell^\infty$ no es separable.
\end{ejercicio}
\begin{solucion}
Sea $E=\{0,1\}^{\N}\subset\ell^\infty$. Este conjunto no es numerable: si sus elementos se listaran como $a^{(1)},a^{(2)},\ldots$, la sucesión $b$ definida por $b_n=1-a^{(n)}_n$ diferiría de la sucesión número $n$ en la coordenada $n$ y no aparecería en la lista. Es el argumento diagonal de Cantor.

Para $a,b\in E$ distintos, existe una coordenada donde difieren en $1$, y ninguna diferencia tiene módulo mayor que $1$; por tanto $\norm{a-b}_\infty=1$. Las bolas $B(a,1/3)$, $a\in E$, son disjuntas dos a dos: un punto común daría por la triangular una distancia entre centros menor que $2/3$.

Si $D=\{d_n:n\in\N\}$ fuera denso, cada bola contendría algún $d_n$. Asignemos a cada $a\in E$ el primer índice de un punto de $D$ que pertenece a $B(a,1/3)$. Bolas disjuntas reciben índices distintos, lo que produciría una inyección del conjunto no numerable $E$ en $\N$. Es imposible; por tanto $\ell^\infty$ no es separable.
\end{solucion}

\begin{ejercicio}[Lectura 6, ejercicio 2]\label{ej:lectura-6-2}
Demostrar que un espacio normado $V$ es separable si y solo si existe un conjunto a lo sumo numerable $L\subset V$ cuyo subespacio generado sea denso.
\end{ejercicio}
\begin{solucion}
Si $D$ es denso y numerable, podemos tomar $L=D$, pues $D\subset\operatorname{span}D$ implica densidad del subespacio generado.

Para el recíproco, sea $\F_0=\Q$ en el caso real y $\F_0=\Q+i\Q$ en el complejo. Sea $D$ el conjunto de todas las combinaciones lineales finitas de elementos de $L$ con coeficientes en $\F_0$, incluyendo el vector cero. Para cada longitud de combinación hay solo un producto finito de conjuntos numerables; al tomar la unión sobre todas las longitudes, $D$ es numerable.

Probemos su densidad. Dado $v\in V$ y $\eps>0$, la densidad de $\operatorname{span}L$ permite escoger $w=\sum_{j=1}^m\alpha_jv_j$ con $\norm{v-w}<\eps/2$. Para cada $j$, elijamos $\beta_j\in\F_0$ tal que
\[
|\alpha_j-\beta_j|<\frac{\eps}{2m(1+\norm{v_j})}.
\]
Los racionales son densos en $\R$, y $\Q+i\Q$ es denso en $\C$ al aproximar por separado las partes real e imaginaria; por ello estas elecciones son posibles. Entonces $u=\sum_{j=1}^m\beta_jv_j\in D$ y
\[
\norm{w-u}\le\sum_{j=1}^m|\alpha_j-\beta_j|\norm{v_j}<\eps/2.
\]
La triangular da $\norm{v-u}<\eps$. Si la combinación es vacía, se toma $u=0$; en particular, $V=\{0\}$ queda cubierto. Así $D$ es denso.
\end{solucion}

\begin{ejercicio}[Lectura 6, ejercicio 3]\label{ej:lectura-6-3}
Si $(A_n)$ es una familia a lo sumo numerable de subespacios separables de un espacio métrico $X$, demostrar que $Y=\bigcup_nA_n$ es separable.
\end{ejercicio}
\begin{solucion}
Para cada $A_n$ no vacío, elijamos $D_n\subset A_n$ denso en su topología relativa y numerable; para el vacío tomemos $D_n=\varnothing$. El conjunto $D=\bigcup_nD_n$ es numerable. Si $y\in Y$ y $r>0$, existe $n$ con $y\in A_n$. Por densidad relativa, $B_X(y,r)\cap A_n$ contiene un punto de $D_n$, y por tanto $B_Y(y,r)$ contiene un punto de $D$. Esto demuestra $\cl D^{\,Y}=Y$. No se necesita que los $A_n$ sean abiertos ni cerrados.
\end{solucion}

\begin{ejercicio}[Lectura 6, ejercicio 4]\label{ej:lectura-6-4}
Demostrar que todo subespacio de un espacio métrico separable es separable.
\end{ejercicio}
\begin{solucion}
Sea $\{B_n\}$ una base numerable de $X$, cuya existencia asegura el teorema~\ref{ch06:separablebase}. Los conjuntos $Y\cap B_n$ forman una base de $Y$: si $U=Y\cap G$ es abierto relativo y $y\in U$, hay $B_n$ con $y\in B_n\subset G$, y entonces $y\in Y\cap B_n\subset U$. Esta base es numerable. Elegimos un punto de cada elemento no vacío de ella; por el lema~\ref{ch06:baseadenso}, esos puntos forman un denso numerable de $Y$.

No basta intersectar con $Y$ un denso numerable prefijado de $X$. Por ejemplo, $\Q$ es denso en $\R$, pero su intersección con $Y=\R\setminus\Q$ es vacía. La prueba con una base elige puntos dentro del subespacio y evita ese error.
\end{solucion}

\begin{ejercicio}[Lectura 6, ejercicio 5]\label{ej:lectura-6-5}
Sean $X_1,\ldots,X_m$ espacios métricos no vacíos. Demostrar que su producto finito, con una métrica producto, es separable si y solo si lo es cada factor.
\end{ejercicio}
\begin{solucion}
Podemos usar $d(x,y)=\max_{1\le j\le m}d_j(x_j,y_j)$. Su bola de radio $r$ es el producto de las bolas de radio $r$ de los factores. Si $D_j$ es denso numerable en $X_j$, el producto $D=\prod_{j=1}^mD_j$ es numerable. Dado $x$ y $r>0$, elegimos $a_j\in D_j$ con $d_j(x_j,a_j)<r$ para cada uno de los finitísimos índices. Entonces $a\in D\cap B(x,r)$ y $D$ es denso.

Recíprocamente, sea $D$ denso numerable en el producto. La proyección $\pi_j(D)$ es numerable. Dado $x_j\in X_j$ y $r>0$, elegimos puntos arbitrarios en los demás factores y formamos $x$ con coordenada $j$ igual a $x_j$. Una aproximación $a\in D$ con $d(a,x)<r$ proporciona $d_j(a_j,x_j)<r$, así que $\pi_j(D)$ es denso en $X_j$.

Las métricas producto finitas dadas por suma, máximo o norma $p$ inducen la misma topología: para números no negativos $t_j$, $\max_jt_j\le(\sum_jt_j^p)^{1/p}\le m^{1/p}\max_jt_j$ y $\max_jt_j\le\sum_jt_j\le m\max_jt_j$. Estas cotas comparan las bolas y trasladan el resultado a cualquiera de esas métricas. La hipótesis de factores no vacíos es necesaria: un factor vacío hace que el producto sea vacío, aunque otro factor sea no separable.
\end{solucion}

\begin{definicion}\index{denso en ninguna parte}
Un conjunto $A\subset X$ es denso en ninguna parte si $\operatorname{int}(\cl A)=\varnothing$. La clausura dentro de esta definición es esencial: tener interior vacío por sí solo no basta, como muestra $\Q\subset\R$.
\end{definicion}

\begin{ejercicio}[Lectura 6, ejercicio 6]\label{ej:lectura-6-6}
Demostrar que un conjunto denso en ninguna parte tiene complemento denso, y que un conjunto denso y cerrado tiene complemento denso en ninguna parte.
\end{ejercicio}
\begin{solucion}
Si $\operatorname{int}(\cl A)=\varnothing$, ningún abierto no vacío $U$ está contenido en $\cl A$. Por tanto $U\cap(X\setminus\cl A)\ne\varnothing$, y en particular $U\cap(X\setminus A)\ne\varnothing$. Así el complemento de $A$ es denso. Hemos probado algo más fuerte: el complemento de la clausura es abierto y denso.

Si $A$ es denso y cerrado, $A=\cl A=X$; por tanto su complemento es vacío, cuya clausura tiene interior vacío. La segunda afirmación de la lectura es correcta, aunque trivial con ambas hipótesis. Una formulación adicional útil es que el complemento de un abierto denso es cerrado y denso en ninguna parte: si $G$ es abierto y denso, $F=X\setminus G$ es cerrado y no contiene ningún abierto no vacío, porque todo abierto no vacío intersecta a $G$. Así $\operatorname{int}(\cl F)=\operatorname{int}F=\varnothing$.
\end{solucion}

\begin{ejercicio}[Lectura 6, ejercicio 7]\label{ej:lectura-6-7}
Demostrar que en un espacio métrico no vacío un conjunto no puede ser simultáneamente denso y denso en ninguna parte.
\end{ejercicio}
\begin{solucion}
Si $A$ es denso, $\cl A=X$. Entonces $\operatorname{int}(\cl A)=\operatorname{int}X=X\ne\varnothing$, lo que contradice ser denso en ninguna parte. La condición de espacio no vacío evita la excepción $X=A=\varnothing$.
\end{solucion}

\begin{ejercicio}[Lectura 6, ejercicio 8]\label{ej:lectura-6-8}\label{ch06:intervalos}
Demostrar que todo abierto de $\R$ es unión a lo sumo numerable de intervalos abiertos disjuntos dos a dos.
\end{ejercicio}
\begin{solucion}
Sea $U\subset\R$ abierto. Si es vacío, la unión vacía resuelve el problema. Para $x\in U$, definamos $I_x$ como la unión de todos los intervalos abiertos contenidos en $U$ que contienen a $x$. Hay al menos uno porque $U$ es abierto. La unión es abierta y está contenida en $U$.

Además, $I_x$ es un intervalo: si $a<b$ pertenecen a $I_x$, hay intervalos $J_a,J_b$ de la familia que contienen respectivamente a $a,b$ y ambos contienen a $x$. Como se intersectan, su unión es un intervalo. Para justificar este último hecho, un punto entre dos puntos de $J_a\cup J_b$ queda en uno de los intervalos, o entre ellos; en el segundo caso un punto común a ambos permite dividir el segmento en dos partes contenidas en los respectivos intervalos. Por ello todos los puntos entre $a,b$ pertenecen a $J_a\cup J_b\subset I_x$.

Así $I_x$ es un intervalo abierto, posiblemente una semirrecta o toda $\R$. Es maximal entre los intervalos de $U$ que contienen a $x$, porque su definición ya reúne a todos ellos. Si $I_x$ e $I_y$ se intersectan, su unión es un intervalo abierto contenido en $U$; contiene tanto a $x$ como a $y$. Por maximalidad está contenida tanto en $I_x$ como en $I_y$, y por tanto $I_x=I_y$. Los intervalos distintos de esta colección son disjuntos y cubren $U$ porque $x\in I_x$.

Cada intervalo no vacío contiene un racional. Fijemos una enumeración de $\Q$ y asignemos a cada intervalo el primer racional de la lista que contiene. Dos intervalos disjuntos reciben racionales distintos, de modo que la familia se inyecta en $\Q$ y es a lo sumo numerable. Esto concluye la descomposición, incluidos los intervalos no acotados.
\end{solucion}

\section{Una misma precisión para todos los puntos}\label{lectura:07}
La densidad explica qué puntos bastan para aproximar el espacio; ahora examinaremos qué control permite transportar esas aproximaciones a través de una función. La continuidad ordinaria permite elegir la precisión de entrada después de fijar un punto. La continuidad uniforme exige una elección que funcione simultáneamente en todo el dominio. Esta diferencia de cuantificadores explica por qué una función puede comportarse bien cerca de cada punto y, sin embargo, amplificar mucho las distancias en regiones cada vez más lejanas. Usaremos continuidad, sucesiones y la desigualdad triangular; no necesitaremos compacidad.

\begin{definicion}\label{ch07:def-uniforme}
Una aplicación $f:(X,d)\to(Y,\rho)$ es \emph{uniformemente continua} si
\[
 \forall\eps>0\ \exists\delta>0\ \forall x,y\in X,
 \qquad d(x,y)<\delta\Longrightarrow\rho(f(x),f(y))<\eps.
\]
Es \emph{Lipschitz} con constante $L\geq0$ si
\[
 \rho(f(x),f(y))\leq Ld(x,y)\qquad(x,y\in X).
\]
\end{definicion}
\index{continuidad!uniforme}\index{Lipschitz, funcion de@Lipschitz, función de}
La desigualdad de Lipschitz es no estricta: escribir $<$ para todo par de puntos produciría $0<0$ al tomar $x=y$. Una constante de Lipschitz no tiene por qué ser óptima; cualquier constante mayor también sirve.

\begin{proposicion}\label{ch07:lip-uniforme}
Toda aplicación Lipschitz es uniformemente continua, y toda aplicación uniformemente continua es continua.
\end{proposicion}
\begin{proof}
Si $L>0$, para una tolerancia $\eps>0$ tomamos $\delta=\eps/L$. Si $L=0$, todas las imágenes coinciden y sirve cualquier $\delta>0$. En ambos casos se cumple la definición. Para la segunda afirmación fijamos $x\in X$: el mismo $\delta$ que funciona para todos los pares funciona, en particular, para los pares $(x,y)$. Ésta es precisamente la continuidad en $x$.
\end{proof}

\begin{ejemplo}[Qué significa aquí discreto]
Si $d(x,y)=1$ para $x\ne y$, toda aplicación de $X$ a cualquier espacio métrico es uniformemente continua: con $\delta=1$ la condición $d(x,y)<\delta$ obliga a $x=y$. El argumento vale siempre que exista $c>0$ con $d(x,y)\geq c$ para puntos distintos; se dice entonces que $X$ es uniformemente discreto. Tener solamente topología discreta no basta. En $X=\{1/n:n\in\N\}\subset\R$, todos los puntos son aislados, pero $f(1/n)=n$ no es uniformemente continua: las distancias entre $1/n$ y $1/(n+1)$ tienden a cero, mientras que las imágenes distan $1$.
\end{ejemplo}

\begin{ejemplo}[Continuidad sin uniformidad]\label{ch07:cuadrado}
La función $f(x)=x^2$ es continua en $\R$, pero no uniformemente continua. Tomemos $x_n=n$ y $y_n=n+1/n$. Entonces
\[
 |x_n-y_n|=1/n\longrightarrow0,
 \qquad |f(x_n)-f(y_n)|=2+1/n^2\geq2.
\]
Si existiera un $\delta$ uniforme para $\eps=1$, un $n$ con $1/n<\delta$ contradiría la segunda desigualdad. Este razonamiento fija primero la tolerancia que fracasa y después produce pares arbitrariamente próximos. No es válido deducir una cota inferior proporcional a $\delta$ a partir de la sola desigualdad $|x-y|<\delta$.
\end{ejemplo}

\begin{proposicion}[Composición]\label{ch07:composicion}
Si $f:X\to Y$ y $g:Y\to Z$ son uniformemente continuas, también lo es $g\circ f$.
\end{proposicion}
\begin{proof}
Fijado $\eps>0$, la uniformidad de $g$ proporciona $\eta>0$ tal que $\rho(u,v)<\eta$ implica $\sigma(g(u),g(v))<\eps$. Aplicamos ahora la uniformidad de $f$ con tolerancia $\eta$: existe $\delta>0$ tal que $d(x,y)<\delta$ implica $\rho(f(x),f(y))<\eta$. Las dos implicaciones encadenadas dan el resultado. El orden de elección es $\eps$, después $\eta$ y finalmente $\delta$.
\end{proof}

\section{Sucesiones cuyos términos terminan estando próximos}
\begin{definicion}\label{ch07:cauchy}
Una sucesión $(x_n)$ de un espacio métrico es \emph{de Cauchy} si
\[
 \forall\eps>0\ \exists N\in\N\ \forall n,m\geq N,
 \qquad d(x_n,x_m)<\eps.
\]
\end{definicion}
\index{Cauchy, sucesion de@Cauchy, sucesión de}
La lectura original utiliza la condición $d(x_n,x_N)<\eps$ para $n\geq N$. Ambas formulaciones son equivalentes: la definición anterior implica esa condición al tomar $m=N$; a la inversa, aplicamos la condición con $\eps/2$ y usamos
\[
 d(x_n,x_m)\leq d(x_n,x_N)+d(x_N,x_m)<\eps.
\]
La condición de Cauchy compara términos de la sucesión entre sí; no presupone que su posible límite pertenezca al espacio.

\begin{teorema}\label{ch07:cauchy-imagen}
Las aplicaciones uniformemente continuas transforman sucesiones de Cauchy en sucesiones de Cauchy.
\end{teorema}
\begin{proof}
Sea $\eps>0$ y elijamos $\delta>0$ por la continuidad uniforme de $f$. Como $(x_n)$ es de Cauchy, existe $N$ tal que $n,m\geq N$ implica $d(x_n,x_m)<\delta$. Por tanto $\rho(f(x_n),f(x_m))<\eps$ para esos índices, que es la condición de Cauchy en $Y$.
\end{proof}
La continuidad ordinaria no basta: $x_n=1/n$ es de Cauchy en $(0,1]$, pero la función continua $x\mapsto1/x$ la transforma en $(n)$, que no es de Cauchy.

\section{Ejercicios sobre continuidad uniforme y extensiones}
\begin{ejercicio}[Lectura 7, ejercicio 1]\label{ej:lectura-7-1}
Demuestre que la composición de dos funciones Lipschitz es Lipschitz y determine una constante para la composición.
\end{ejercicio}
\begin{solucion}
Si $f:X\to Y$ tiene constante $L$ y $g:Y\to Z$ tiene constante $M$, entonces
\[
 \sigma(g(f(x)),g(f(y)))\leq M\rho(f(x),f(y))\leq MLd(x,y).
\]
Así, $ML$ es una constante válida, incluso cuando alguno de los factores es cero. No necesariamente es la mejor: si $g$ es constante en $f(X)$, la composición tiene constante óptima cero aunque $L,M$ sean positivos.
\end{solucion}

\begin{ejercicio}[Lectura 7, ejercicio 2]\label{ej:lectura-7-2}
Sean $f:X\to Y$ y $g:Y\to Z$ continuas. Determine si $g\circ f$ debe ser uniformemente continua cuando: \begin{enumerate}[label=(\roman*)]
\item $f$ es uniformemente continua;
\item $g$ es uniformemente continua.
\end{enumerate}
\end{ejercicio}
\begin{solucion}
En ambos casos la respuesta es negativa. En (i) tomamos $X=Y=Z=\R$, $f(x)=x$ y $g(x)=x^2$. La identidad es $1$-Lipschitz y la composición es el cuadrado, que falla por el ejemplo~\ref{ch07:cuadrado}. En (ii) intercambiamos los papeles: $f(x)=x^2$ y $g(x)=x$. Ahora $g$ es $1$-Lipschitz y la composición vuelve a ser $x^2$.
\end{solucion}

\begin{ejercicio}[Lectura 7, ejercicio 3]\label{ej:lectura-7-3}
Sea $A$ un cerrado no vacío de un espacio métrico $X$. Pruebe que $x\mapsto\dist(x,A)$ es uniformemente continua.
\end{ejercicio}
\begin{solucion}
Para $a\in A$, la desigualdad triangular da $d(x,a)\leq d(x,y)+d(y,a)$. Tomando ínfimos en $a$ obtenemos
\[
 \dist(x,A)\leq d(x,y)+\dist(y,A).
\]
Al intercambiar $x,y$ concluimos
\[
 |\dist(x,A)-\dist(y,A)|\leq d(x,y).
\]
La función es $1$-Lipschitz. La prueba no exige que $A$ sea cerrado; esa hipótesis solamente permite identificar $\dist(x,A)=0$ con $x\in A$. Se agrega la condición $A\ne\varnothing$ porque la distancia al vacío no es una función real finita.
\end{solucion}

\begin{ejercicio}[Lectura 7, ejercicio 4]\label{ej:lectura-7-4}
Dé una función uniformemente continua entre espacios métricos que no sea abierta.
\end{ejercicio}
\begin{solucion}
La función constante $f:\R\to\R$, $f(x)=0$, es $0$-Lipschitz. Sin embargo, la imagen del abierto $(-1,1)$ es $\{0\}$, que no es abierto en $\R$. La continuidad controla preimágenes de abiertos; ser abierta controla imágenes y es una condición distinta.
\end{solucion}

\begin{ejercicio}[Lectura 7, ejercicio 5]\label{ej:lectura-7-5}\label{ch07:ej-extension}
Sea $A$ abierto en un espacio métrico $X$ y sea $f:X\to\R$ uniformemente continua en $A$. Demuestre que $f|_A$ se extiende continuamente a $\cl A$. Dé un contraejemplo si se reemplaza $\R$ por un espacio métrico arbitrario.
\end{ejercicio}
\begin{solucion}
Probaremos la afirmación más fuerte en que $A$ es cualquier subconjunto y el espacio de llegada $Y$ es completo. Si $A=\varnothing$, la extensión tiene dominio vacío y no hay nada que definir. Para $x\in\cl A$ elegimos $a_n\in A$ con $d(a_n,x)<1/n$. Esta sucesión converge y es de Cauchy. Por el teorema~\ref{ch07:cauchy-imagen}, $(f(a_n))$ es de Cauchy; la completitud de $Y$ proporciona un límite. Definimos
\[
 g(x)=\lim_{n\to\infty}f(a_n).
\]
La elección no afecta al valor: si $b_n\to x$ en $A$, entonces $d(a_n,b_n)\leq d(a_n,x)+d(x,b_n)\to0$; la uniformidad implica $\rho(f(a_n),f(b_n))\to0$. Tomando límites, los dos valores propuestos coinciden. Si $x\in A$, podemos elegir la sucesión constante $a_n=x$, por lo que $g(x)=f(x)$.

Además, $g$ es uniformemente continua. Dado $\eps>0$, elijamos $\delta>0$ de modo que $d(a,b)<\delta$ implique $\rho(f(a),f(b))<\eps/2$ para $a,b\in A$. Si $x,y\in\cl A$ y $d(x,y)<\delta/3$, elegimos $a_n\to x$, $b_n\to y$. Para $n$ suficientemente grande, $d(a_n,b_n)<\delta$; al tomar límites queda $\rho(g(x),g(y))\leq\eps/2<\eps$. La extensión es única porque dos aplicaciones continuas que coinciden en un conjunto denso coinciden en su clausura: evaluamos ambas sobre una sucesión $a_n\to x$.

Para el contraejemplo tomamos $X=\R$, $A=(0,1)$ y $Y=(0,1)$ con la distancia usual. Definimos $f(x)=x$ en $A$ y $f(x)=1/2$ fuera de $A$. Su restricción a $A$ es $1$-Lipschitz. Una extensión continua $g:[0,1]\to(0,1)$ tendría $g(0)=\lim_{n\to\infty}g(1/(n+2))=0$, imposible en el espacio de llegada. Lo que falla es su completitud.
\end{solucion}

\begin{ejercicio}[Lectura 7, ejercicio 6]\label{ej:lectura-7-6}
Para subconjuntos no vacíos $A,B$ de $X$, defina $d(A,B)=\inf\{d(a,b):a\in A,b\in B\}$. Pruebe que $f:X\to Y$ es uniformemente continua si y sólo si
\[
 d(A,B)=0\Longrightarrow\rho(f(A),f(B))=0
\]
para todos esos conjuntos.
\end{ejercicio}
\begin{solucion}
Si $f$ es uniformemente continua, dado $\eps>0$ elegimos el correspondiente $\delta>0$. La igualdad $d(A,B)=0$ permite encontrar $a\in A,b\in B$ con $d(a,b)<\delta$. Entonces $\rho(f(A),f(B))\leq\rho(f(a),f(b))<\eps$. Un número no negativo menor que todo $\eps>0$ es cero.

Para el recíproco hay que cuidar las \emph{distancias cruzadas}: saber que $\rho(f(x_n),f(y_n))\geq\eps_0$ no basta por sí solo para separar los conjuntos de imágenes. Si $f$ no es uniformemente continua, existen $\eps_0>0$ y pares $x_n,y_n$ tales que
\[
 d(x_n,y_n)<1/n,\qquad \rho(p_n,q_n)\geq\eps_0,
 \quad p_n=f(x_n),\ q_n=f(y_n).
\]
Pongamos $r=\eps_0/8$. Si alguna bola $B(z,r)$ contiene infinitos $p_n$, retenemos esos índices $I$. Para $i,j\in I$, tenemos $\rho(q_j,z)\geq\eps_0-r$ y, por tanto, $\rho(p_i,q_j)>\eps_0-2r$. Si alguna bola contiene infinitos $q_n$, el argumento con los papeles intercambiados da la misma separación.

Queda el caso en que toda bola de radio $r$ contiene sólo un número finito de términos de cada sucesión, contando sus índices. Elegimos inductivamente índices crecientes $n_k$. Una vez elegidos los primeros $k$, descartamos todos los índices $n$ para los que $p_n\in B(q_{n_j},r)$ o $q_n\in B(p_{n_j},r)$ para algún $j\leq k$. Sólo descartamos una unión finita de conjuntos finitos; por ello queda un índice mayor que $n_k$. La sucesión retenida satisface $\rho(p_{n_i},q_{n_j})\geq r$ si $i\ne j$, y la misma desigualdad vale cuando $i=j$ porque $\eps_0\geq r$.

En cualquiera de los casos obtenemos un conjunto infinito de índices $I$ y una constante $c>0$ tales que $\rho(p_i,q_j)\geq c$ para todo $i,j\in I$. Para $A=\{x_i:i\in I\}$ y $B=\{y_i:i\in I\}$ se cumple $d(A,B)=0$, pues $I$ no está acotado y $d(x_i,y_i)<1/i$, pero $\rho(f(A),f(B))\geq c$. Esto contradice la propiedad supuesta.
\end{solucion}

\begin{ejercicio}[Lectura 7, ejercicio 7]\label{ej:lectura-7-7}
Suponga que $f:X\to Y$ satisface $\rho(f(x),f(y))\geq\alpha d(x,y)$ para alguna $\alpha>0$. Pruebe que es inyectiva, que su inversa sobre la imagen es Lipschitz y que $f$ es abierta si $f(X)$ es abierto en $Y$.
\end{ejercicio}
\begin{solucion}
Si $f(x)=f(y)$, la desigualdad da $0\geq\alpha d(x,y)$ y por tanto $x=y$. Para $u=f(x),v=f(y)$ obtenemos
\[
 d(f^{-1}(u),f^{-1}(v))\leq\alpha^{-1}\rho(u,v).
\]
Así, $f^{-1}:f(X)\to X$ es continua. Si $U\subset X$ es abierto, $f(U)=(f^{-1})^{-1}(U)$ es abierto en el subespacio $f(X)$. Cuando $f(X)$ es abierto en $Y$, todo abierto suyo es abierto en $Y$, y $f$ es una aplicación abierta. La hipótesis no afirma que $f$ sea continua; controla únicamente su inversa.
\end{solucion}

\section{Continuidad, aproximación y control uniforme}
El criterio de tolerancias permite obtener cotas concretas; el de abiertos simplifica composiciones y pegados; el secuencial permite detectar discontinuidades mediante aproximaciones explícitas. Un homeomorfismo conserva esta estructura de continuidad en ambas direcciones. No tiene por qué conservar distancias, acotación ni propiedades que dependan de la métrica concreta.

En un espacio métrico, una familia numerable de puntos de aproximación se convierte en una familia numerable de entornos básicos. Esa conversión permite reducir cubiertas abiertas a subfamilias numerables y heredar separabilidad a todos los subespacios. La condición «denso en ninguna parte» mide una limitación distinta: ni siquiera después de agregar puntos límite aparece una región abierta completa.

La continuidad uniforme controla simultáneamente todos los puntos, preserva la condición de Cauchy y permite extender funciones a puntos límite cuando el espacio de llegada contiene los límites necesarios. El capítulo~\ref{chap:08} estudia precisamente esa propiedad del espacio de llegada: la completitud.


\parte{Completitud y métodos de existencia}{%
Una sucesión de Cauchy describe una aproximación cuyos términos se acercan entre sí. La completitud garantiza que esa aproximación produce un punto del propio espacio. Estudiaremos cómo reconocer esta propiedad y cómo construir los puntos que faltan mediante la completación.

Dos consecuencias ofrecen métodos de existencia distintos. El teorema de Baire impide ciertas descomposiciones de un espacio completo; el principio de contracción construye una solución por iteración y controla el error. La continuidad uniforme y las sucesiones de Cauchy de la primera parte son las herramientas de entrada.}
\chapter{Espacios completos y completación}\label{cap:nuevo03}\label{chap:08}\label{chap:09}
Una sucesión de Cauchy contiene información interna: sus términos terminan siendo tan próximos como se quiera. La completitud asegura que esa información corresponde a un punto del propio espacio, y lo hace para todas las sucesiones de Cauchy. Comenzaremos por reconocer esta propiedad en espacios de sucesiones y en subespacios cerrados; después construiremos una completación para cualquier espacio métrico, conservando las distancias y haciendo que los puntos originales sean densos. La existencia de esa construcción y su unicidad relacionan los ejemplos concretos con una herramienta general. Usaremos la continuidad uniforme del capítulo~\ref{chap:07}, la propiedad arquimediana de $\R$ y las desigualdades elementales de las normas $\ell^p$.

\section{Límites que permanecen dentro del espacio}
\label{lectura:08}

\begin{proposicion}\label{ch08:cauchy-basico}
Toda sucesión convergente es de Cauchy. Toda sucesión de Cauchy es acotada. Si una sucesión de Cauchy tiene una subsucesión convergente a $x$, la sucesión completa converge a $x$.
\end{proposicion}
\begin{proof}
Si $x_n\to x$, dado $\eps>0$ elegimos $N$ tal que $d(x_n,x)<\eps/2$ para $n\geq N$. Entonces $d(x_n,x_m)<\eps$ para $n,m\geq N$. Si $(x_n)$ es de Cauchy, elegimos $N$ con $d(x_n,x_N)<1$ para $n\geq N$. El máximo de $1,d(x_1,x_N),\ldots,d(x_{N-1},x_N)$ es una cota de todas las distancias al centro $x_N$.

Por último, supongamos $x_{n_j}\to x$. Elegimos $N$ con $d(x_n,x_m)<\eps/2$ para $n,m\geq N$, y después $j$ con $n_j\geq N$ y $d(x_{n_j},x)<\eps/2$. Para cualquier $n\geq N$ se tiene $d(x_n,x)\leq d(x_n,x_{n_j})+d(x_{n_j},x)<\eps$.
\end{proof}

\begin{definicion}
Un espacio métrico es \emph{completo} si toda sucesión de Cauchy converge en él. Un espacio normado completo para la distancia $d(x,y)=\norm{x-y}$ se llama \emph{espacio de Banach}.
\end{definicion}
\index{completitud}\index{Banach, espacio de}

\begin{ejemplo}
En $\R\setminus\{0\}$, la sucesión $1/n$ es de Cauchy pero no converge en el espacio. De converger allí a $a$, también convergería a $a$ en $\R$, y la unicidad del límite obligaría a $a=0$. En cambio, todo espacio uniformemente discreto es completo: si los puntos distintos distan al menos $c>0$, la condición de Cauchy con tolerancia $c$ obliga a que la sucesión sea eventualmente constante. La mera topología discreta no garantiza completitud: $\{1/n:n\in\N\}$, con su distancia heredada, vuelve a ser un contraejemplo.
\end{ejemplo}

\section{Un ejemplo aritmético: la distancia p-ádica}
Fijemos un primo $p$. Todo racional no nulo se escribe como $r=p^k a/b$, donde $k\in\Z$, los enteros $a,b$ no son divisibles por $p$ y $b\ne0$. El exponente $k$ es único: es la potencia de $p$ del numerador menos la del denominador de cualquier fracción que represente a $r$. Los enteros $a,b$ no son únicos si no se exige además que la fracción esté reducida. Definimos $v_p(r)=k$, $|r|_p=p^{-k}$ y $|0|_p=0$. La distancia es
\[
 d_p(r,s)=|r-s|_p.
\]
Su propiedad decisiva, demostrada en el ejercicio~\ref{ej:lectura-8-1}, es
\[
 |r+s|_p\leq\max\{|r|_p,|s|_p\}.
\]
\index{metrica p-adica@métrica p-ádica}\index{ultrametrica@ultramétrica}
En esta geometría, ser divisible por una potencia grande de $p$ significa estar cerca de cero.

\begin{ejemplo}\label{ch08:padica-incompleta}
Para una sucesión de dígitos $a_j\in\{0,1,\ldots,p-1\}$ definamos $s_n=\sum_{j=0}^n a_jp^j$. Si $m>n$, la diferencia $s_m-s_n$ es divisible por $p^{n+1}$ y
\[
 |s_m-s_n|_p\leq p^{-(n+1)}\longrightarrow0.
\]
Se necesita una desigualdad, pues puede haber dígitos nulos adicionales o incluso diferencia cero. Por tanto $(s_n)$ es de Cauchy.

Hay incontables sucesiones de dígitos: si se enumeraran, podríamos cambiar el dígito de posición $n$ de la sucesión número $n$ y construir una que no aparece en la lista. Además, dos sucesiones distintas no pueden dar el mismo límite racional. Si su primer dígito distinto ocupa la posición $k$, para todo $n\geq k$ la diferencia entre sus sumas parciales tiene valoración exactamente $k$, pues el coeficiente de $p^k$ no es divisible por $p$. Las distancias entre esas sumas parciales son entonces constantemente $p^{-k}$; si ambas convergieran al mismo racional, esas distancias tenderían a cero. Como $\Q$ es numerable, no todas las sucesiones construidas pueden converger en $\Q$. Esto prueba que $(\Q,d_p)$ no es completo.

Algunas sí convergen: si $a_j=p-1$ para todo $j$, entonces $s_n=p^{n+1}-1$ y $|s_n-(-1)|_p=p^{-(n+1)}\to0$.
\end{ejemplo}

\section{Completitud de los espacios de sucesiones}
Recordemos que, para $1\leq p<\infty$,
\[
 \ell^p=\left\{x=(x_k):\sum_{k=1}^\infty|x_k|^p<\infty\right\},
 \qquad \norm{x}_p=\left(\sum_{k=1}^\infty|x_k|^p\right)^{1/p}.
\]
Las coordenadas pueden ser reales o complejas. En el segundo caso, la completitud de los escalares se obtiene aplicando la de $\R$ a partes reales e imaginarias.

\begin{teorema}\label{ch08:lp-completo}
Para cada $1\leq p<\infty$, el espacio $\ell^p$ es de Banach.
\end{teorema}
\begin{proof}
Sea $(x^{(n)})$ una sucesión de Cauchy en $\ell^p$, donde $x^{(n)}=(x^{(n)}_k)_{k\geq1}$. Distinguimos el índice $n$, que enumera vectores, del índice $k$, que enumera coordenadas. La estimación
\[
 |x^{(n)}_k-x^{(m)}_k|\leq\norm{x^{(n)}-x^{(m)}}_p
\]
muestra que cada sucesión escalar $(x^{(n)}_k)_n$ es de Cauchy. Definimos $x_k=\lim_nx^{(n)}_k$ y $x=(x_k)$.

Primero verificamos que $x$ pertenece a $\ell^p$. La sucesión de vectores es acotada, así que existe $M$ con $\norm{x^{(n)}}_p\leq M$ para todo $n$. Para cada entero $K$,
\[
 \sum_{k=1}^K|x_k|^p
 =\lim_{n\to\infty}\sum_{k=1}^K|x^{(n)}_k|^p\leq M^p.
\]
Aquí sólo se ha pasado el límite a través de una suma \emph{finita}. Las sumas parciales de la serie no negativa están acotadas y, por el axioma del supremo, convergen a un número no mayor que $M^p$.

Ahora probamos convergencia en norma. Dado $\eps>0$, existe $N$ tal que $\norm{x^{(n)}-x^{(m)}}_p<\eps/2$ siempre que $n,m\geq N$. Fijamos $n\geq N$ y $K\in\N$, y hacemos tender $m$ a infinito en
\[
 \sum_{k=1}^K|x^{(n)}_k-x^{(m)}_k|^p<(\eps/2)^p.
\]
Resulta $\sum_{k=1}^K|x^{(n)}_k-x_k|^p\leq(\eps/2)^p$. Tomando el supremo sobre $K$ obtenemos $\norm{x^{(n)}-x}_p\leq\eps/2<\eps$. El índice $N$ permanece fijo durante este argumento; no se hace tender a infinito un umbral que ya depende de $\eps$.
\end{proof}
La estrategia tiene dos partes que conviene distinguir: construir un candidato coordenada a coordenada y probar después una estimación uniforme en el número de coordenadas. La convergencia de cada coordenada, por sí sola, no implica convergencia en norma: los vectores canónicos $e_n$ convergen coordenada a coordenada a cero, pero $\norm{e_n}_p=1$.

\begin{teorema}[Subespacios completos]\label{ch08:cerrado-completo}
Sea $Y\subset X$ con la distancia heredada. Si $Y$ es completo, entonces es cerrado en $X$. Si $X$ es completo y $Y$ es cerrado en $X$, entonces $Y$ es completo.
\end{teorema}
\begin{proof}
Si $x\in\cl Y$, podemos elegir $y_n\in Y$ con $d(y_n,x)<1/n$. La sucesión es de Cauchy y, por completitud de $Y$, converge a algún $y\in Y$. También converge a $x$ en $X$, por lo que $x=y\in Y$. Esto prueba que $Y$ contiene su clausura. Si $Y$ es vacío, ya es cerrado.

Para la segunda afirmación, una sucesión de Cauchy en $Y$ lo es en $X$ porque las distancias son las mismas. Tiene, pues, un límite $x\in X$. Al ser $Y$ cerrado, todo límite de una sucesión suya pertenece a $Y$; así, $x\in Y$.
\end{proof}

\section{Ejercicios sobre completitud y espacios de sucesiones}
\begin{ejercicio}[Lectura 8, ejercicio 1]\label{ej:lectura-8-1}
Demuestre que $d_p$ es una métrica en $\Q$ y, más precisamente, una ultramétrica.
\end{ejercicio}
\begin{solucion}
La factorización prima garantiza la unicidad de $v_p(r)$ y $v_p(rs)=v_p(r)+v_p(s)$ para racionales no nulos. Sean $r,s\ne0$ y supongamos $v_p(r)=k\leq v_p(s)=l$. Escribiendo $r=p^ka/b$ y $s=p^lc/d$ con denominadores y numeradores no divisibles por $p$, obtenemos
\[
 r+s=p^k\frac{ad+p^{l-k}bc}{bd}.
\]
El denominador no es divisible por $p$, y el numerador es un entero. Si no es cero, su valoración es no negativa; de aquí $v_p(r+s)\geq k$. Si es cero, $|r+s|_p=0$. En ambos casos
\[
 |r+s|_p\leq p^{-k}=\max\{|r|_p,|s|_p\}.
\]
Si alguno de los sumandos es cero, la desigualdad se verifica directamente. También $|r|_p\geq0$, con igualdad exactamente para $r=0$, y $|-r|_p=|r|_p$. Por tanto $d_p$ es positiva, separa puntos y es simétrica. Aplicando la desigualdad anterior a $(r-t)+(t-s)$ obtenemos
\[
 d_p(r,s)\leq\max\{d_p(r,t),d_p(t,s)\}\leq d_p(r,t)+d_p(t,s).
\]
La primera desigualdad es la propiedad ultramétrica y la segunda da la desigualdad triangular usual.
\end{solucion}

\begin{ejercicio}[Lectura 8, ejercicio 2]\label{ej:lectura-8-2}
Demuestre la equivalencia entre el axioma del supremo y la completitud de $\R$ con su distancia usual.
\end{ejercicio}
\begin{solucion}
Trabajamos con las propiedades de cuerpo ordenado arquimediano de los reales; el contenido del ejercicio es comparar las dos formulaciones de su propiedad de completitud.

Supongamos el axioma del supremo y sea $(x_n)$ de Cauchy. Está acotada. Para cada $N$, definamos $a_N=\inf_{n\geq N}x_n$ y $b_N=\sup_{n\geq N}x_n$. El ínfimo existe aplicando el axioma del supremo al conjunto de los opuestos. La sucesión $(a_N)$ crece, $(b_N)$ decrece y $a_N\leq b_N$. Para cualquier $\eps>0$, la condición de Cauchy aplicada con $\eps/2$ implica que, a partir de cierto $N$, todos los términos de la cola están en un intervalo de longitud $\eps$; por ello $b_N-a_N\leq\eps$. Sea $x=\sup_Na_N$. Para todo $N$, se tiene $a_N\leq x\leq b_N$: la segunda desigualdad se debe a que $b_N$ es una cota superior de todos los $a_M$, también de los correspondientes a $M<N$. Si $n\geq N$, tanto $x_n$ como $x$ están en $[a_N,b_N]$, de modo que $|x_n-x|\leq b_N-a_N\to0$.

Recíprocamente, supongamos que toda sucesión real de Cauchy converge. Sea $A\ne\varnothing$ acotado superiormente. Si tiene máximo, ése es su supremo. En otro caso elegimos $a_0\in A$ y una cota superior $b_0>a_0$. La sucesión de intervalos se construye por bisección: si $m_n=(a_n+b_n)/2$ es cota superior de $A$, ponemos $(a_{n+1},b_{n+1})=(a_n,m_n)$; si no lo es, ponemos $(a_{n+1},b_{n+1})=(m_n,b_n)$. Aquí $a_n$ no tiene por qué pertenecer a $A$, pero nunca es cota superior, y $b_n$ siempre lo es. Los intervalos están anidados y su longitud es $(b_0-a_0)2^{-n}$. Esto hace que ambas sucesiones sean de Cauchy y que tengan el mismo límite $s$.

Todo $a\in A$ satisface $a\leq b_n$ para cada $n$, luego $a\leq s$. Si $t<s$, para algún $n$ se tiene $t<a_n$, y como $a_n$ no es cota superior existe $a\in A$ con $a>a_n>t$. Así, ningún número menor que $s$ es cota superior. Hemos probado $s=\sup A$. Si se desean puntos de $A$ como en la sugerencia original, elegimos $c_n\in A$ con $a_n<c_n\leq b_n$; entonces también $c_n\to s$.
\end{solucion}

\begin{ejercicio}[Lectura 8, ejercicio 3]\label{ej:lectura-8-3}
Demuestre que $\ell^\infty$ es completo.
\end{ejercicio}
\begin{solucion}
Sea $(x^{(n)})$ de Cauchy para $\norm{x}_\infty=\sup_k|x_k|$. Cada coordenada es de Cauchy, porque $|x^{(n)}_k-x^{(m)}_k|\leq\norm{x^{(n)}-x^{(m)}}_\infty$. Definamos $x_k=\lim_nx^{(n)}_k$. Existe $M$ con $\norm{x^{(n)}}_\infty\leq M$, y al pasar al límite para cada $k$ resulta $|x_k|\leq M$; así $x\in\ell^\infty$. Dado $\eps>0$, elegimos $N$ con $\norm{x^{(n)}-x^{(m)}}_\infty<\eps/2$ para $n,m\geq N$. Fijamos $n\geq N$ y dejamos $m\to\infty$: para cada $k$, $|x^{(n)}_k-x_k|\leq\eps/2$. Tomar el supremo en $k$ produce $\norm{x^{(n)}-x}_\infty\leq\eps/2<\eps$.
\end{solucion}

\begin{ejercicio}[Lectura 8, ejercicio 4]\label{ej:lectura-8-4}
Pruebe que los subespacios $c$ y $c_0$ de sucesiones convergentes y de sucesiones convergentes a cero son espacios de Banach con la norma uniforme.
\end{ejercicio}
\begin{solucion}
Ambos son subespacios lineales de $\ell^\infty$: las combinaciones lineales preservan la existencia del límite y, para $c_0$, preservan el límite cero. Por el teorema~\ref{ch08:cerrado-completo}, basta probar que son cerrados.

Sea $x^{(n)}\in c$ y $\norm{x^{(n)}-x}_\infty\to0$. Para $\eps>0$ elegimos $n$ con $\norm{x^{(n)}-x}_\infty<\eps/3$. Como la sucesión escalar $k\mapsto x^{(n)}_k$ converge, existe $K$ con $|x^{(n)}_k-x^{(n)}_j|<\eps/3$ para $k,j\geq K$. Por tanto
\[
 |x_k-x_j|\leq|x_k-x^{(n)}_k|+|x^{(n)}_k-x^{(n)}_j|+|x^{(n)}_j-x_j|<\eps.
\]
La sucesión de coordenadas de $x$ es de Cauchy en los escalares y converge; de ahí $x\in c$.

Si cada $x^{(n)}\in c_0$, elegimos $n$ con $\norm{x^{(n)}-x}_\infty<\eps/2$ y luego $K$ con $|x^{(n)}_k|<\eps/2$ para $k\geq K$. Entonces $|x_k|<\eps$ para $k\geq K$, luego $x\in c_0$.
\end{solucion}

\begin{ejercicio}[Lectura 8, ejercicio 5]\label{ej:lectura-8-5}
Para espacios métricos no vacíos, pruebe que el producto de dos espacios es completo si y sólo si cada factor lo es. Extienda el resultado a todo producto finito.
\end{ejercicio}
\begin{solucion}
Usamos la métrica máxima $D((x,y),(x',y'))=\max\{d(x,x'),\rho(y,y')\}$. Una sucesión es de Cauchy para $D$ exactamente cuando sus dos coordenadas lo son: una cota para $D$ acota ambas distancias, y dos umbrales para las coordenadas se reúnen tomando su máximo. Si ambos factores son completos, las coordenadas convergen a $x,y$, y el máximo de sus distancias a esos límites tiende a cero.

Si el producto es completo, fijamos $y_0\in Y$ y consideramos una sucesión de Cauchy $(x_n)$ de $X$. Los pares $(x_n,y_0)$ forman una sucesión de Cauchy que converge a algún $(x,y)$, y la primera coordenada da $x_n\to x$. Fijando $x_0$ se demuestra lo mismo para $Y$. Para $n$ factores, la métrica $\max_{1\leq j\leq n}d_j(x_j,y_j)$ permite repetir exactamente las dos implicaciones: tomamos el máximo de los $n$ umbrales y, para la necesidad, fijamos las otras $n-1$ coordenadas. La métrica suma y la euclídea de las distancias son equivalentes mediante cotas constantes a la métrica máxima; por ello tienen las mismas sucesiones de Cauchy. La no vacuidad evita que un producto vacío oculte un factor incompleto.
\end{solucion}

\begin{ejercicio}[Lectura 8, ejercicio 6]\label{ej:lectura-8-6}
Demuestre que el conjunto de Cantor, con la distancia heredada de $\R$, es completo y no tiene puntos aislados.
\end{ejercicio}
\begin{solucion}
Sea $C_0=[0,1]$ y sea $C_n$ el conjunto que queda después de eliminar los tercios medios abiertos durante las primeras $n$ etapas. Cada $C_n$ es una unión finita de intervalos cerrados de longitud $3^{-n}$, y $C=\bigcap_{n\geq0}C_n$ es cerrado en $\R$. Como $\R$ es completo, también lo es $C$.

Fijemos $x\in C$ y $\eps>0$. Elegimos $n$ con $3^{-n}<\eps$ y un intervalo componente $[a,b]$ de $C_n$ que contenga $x$. Ambos extremos pertenecen a todas las etapas posteriores: nunca se elimina un extremo al retirar un tercio medio abierto. Por tanto $a,b\in C$. Como $a\ne b$, al menos uno de ellos es distinto de $x$ y dista de $x$ a lo sumo $3^{-n}<\eps$. Ningún punto de $C$ es aislado.
\end{solucion}

\begin{ejercicio}[Lectura 8, ejercicio 7]\label{ej:lectura-8-7}
Demuestre que $C$ es incontable y que
\[
 C=\left\{\sum_{n=1}^\infty\frac{a_n}{3^n}:a_n\in\{0,2\}\right\}.
\]
\end{ejercicio}
\begin{solucion}
Para cada palabra $(a_1,\ldots,a_n)$ de ceros y doses, una componente de $C_n$ es
\[
 I(a_1,\ldots,a_n)=\left[\sum_{j=1}^n\frac{a_j}{3^j},\,
 \sum_{j=1}^n\frac{a_j}{3^j}+3^{-n}\right].
\]
Esto se verifica por inducción: escoger el tercio izquierdo agrega el dígito $0$ y escoger el derecho agrega el dígito $2$. Una serie con esos dígitos pertenece al intervalo indicado, pues su cola está entre cero y $\sum_{j>n}2/3^j=3^{-n}$. Pertenece entonces a todo $C_n$ y a $C$.

Recíprocamente, si $x\in C$, en cada etapa está en una componente y la componente de la etapa siguiente está contenida en la anterior. Sus elecciones izquierda/derecha proporcionan dígitos $a_n\in\{0,2\}$. La diferencia entre $x$ y la suma de sus primeros $n$ dígitos está entre cero y $3^{-n}$, por lo que $x$ es la serie correspondiente. No es necesario usar la expansión ternaria convencional en los extremos: por ejemplo, $1/3=0.02222\ldots$ usa sólo dígitos permitidos.

La correspondencia es inyectiva. Si dos sucesiones difieren por primera vez en $k$, la diferencia aportada allí tiene valor absoluto $2/3^k$, mientras que toda la cola posterior puede compensar a lo sumo $\sum_{j>k}2/3^j=1/3^k$. Las sumas difieren al menos en $1/3^k>0$. Finalmente, el conjunto de sucesiones de ceros y doses es incontable por el argumento diagonal: frente a una supuesta enumeración, cambiamos el dígito $n$ de la fila $n$. Su imagen inyectiva $C$ es, pues, incontable.
\end{solucion}

\section{Agregar límites sin alterar las distancias}
\label{lectura:09}
Los espacios $\ell^p$, $c$ y $c_0$ ofrecen modelos completos, mientras que los racionales p-ádicos muestran que una noción distinta de distancia puede exigir nuevos puntos. Buscamos ahora una construcción que agregue sistemáticamente los límites que faltan. Las distancias entre los puntos originales deben conservarse, y los puntos nuevos deben surgir como límites de puntos antiguos. La completación que construiremos satisface ambas exigencias y es única salvo una identificación que conserva distancias.

\begin{definicion}
Una aplicación $i:(X,d)\to(Y,\rho)$ es una \emph{isometría} si $\rho(i(x),i(y))=d(x,y)$ para todos $x,y\in X$. Una \emph{completación} de $X$ es un espacio completo $\widehat X$ junto con una isometría $i:X\to\widehat X$ cuya imagen es densa.
\end{definicion}
\index{isometria@isometría}\index{completacion@completación}
Una isometría es inyectiva: imágenes iguales dan distancia cero entre los puntos originales. Es $1$-Lipschitz y, por tanto, uniformemente continua; su inversa sobre la imagen también es una isometría. La inclusión de un subespacio con la distancia heredada es un ejemplo. Si $X$ ya es completo, $(X,\id_X)$ es una completación suya.

\section{Construcción mediante sucesiones de Cauchy}
\begin{teorema}\label{ch09:existencia}
Todo espacio métrico tiene una completación.
\end{teorema}
\begin{proof}
El espacio vacío es completo y se completa con él mismo. Supongamos $X\ne\varnothing$ y sea $\mathcal C(X)$ el conjunto de todas las sucesiones de Cauchy de $X$. Este conjunto no es vacío porque contiene todas las sucesiones constantes.

\textbf{Paso 1: una distancia provisional.} Para $u=(u_n)$ y $v=(v_n)$ ponemos
\[
 D(u,v)=\lim_{n\to\infty}d(u_n,v_n).
\]
Debemos demostrar que el límite existe. La desigualdad triangular, aplicada en los dos sentidos, produce
\[
 |d(u_n,v_n)-d(u_m,v_m)|\leq d(u_n,u_m)+d(v_n,v_m).
\]
Las dos sucesiones del lado derecho son de Cauchy; así, $(d(u_n,v_n))$ es una sucesión real de Cauchy. Su límite existe por la completitud de $\R$. Además, $D\geq0$, $D(u,u)=0$ y $D(u,v)=D(v,u)$. Si $w=(w_n)$, de $d(u_n,v_n)\leq d(u_n,w_n)+d(w_n,v_n)$ y la preservación de desigualdades al pasar al límite resulta $D(u,v)\leq D(u,w)+D(w,v)$. Por ello $D$ es una pseudométrica; puede anularse sobre sucesiones diferentes, por ejemplo sobre dos sucesiones que converjan al mismo punto.

\textbf{Paso 2: identificar las representaciones del mismo punto.} Definimos $u\sim v$ cuando $D(u,v)=0$. La reflexividad y la simetría provienen de las de $D$; la transitividad se sigue de $D(u,w)\leq D(u,v)+D(v,w)$. Sea
\[
 \widehat X=\mathcal C(X)/\!\sim,
 \qquad \widehat d([u],[v])=D(u,v).
\]
Esta fórmula no depende de representantes: si $u\sim u'$ y $v\sim v'$, entonces
\[
 |D(u,v)-D(u',v')|\leq D(u,u')+D(v,v')=0.
\]
Las propiedades de pseudométrica descienden al cociente y, además, $\widehat d([u],[v])=0$ implica $[u]=[v]$. Tenemos una métrica verdadera.

\textbf{Paso 3: insertar $X$ y probar densidad.} Definimos $i(x)$ como la clase de la sucesión constante $(x,x,\ldots)$. Entonces $\widehat d(i(x),i(y))=\lim_nd(x,y)=d(x,y)$, por lo que $i$ es una isometría. Si $u=[(u_n)]$, se cumple
\[
 \widehat d(i(u_n),u)=\lim_{k\to\infty}d(u_n,u_k)\longrightarrow0
 \quad(n\to\infty).
\]
En efecto, dado $\eps>0$, la condición de Cauchy permite que $d(u_n,u_k)<\eps/2$ para $n,k\geq N$; fijando $n\geq N$ y pasando $k$ al límite, queda una distancia no mayor que $\eps/2<\eps$. Esto prueba que $i(X)$ es denso.

\textbf{Paso 4: completar la construcción.} Sea $(U_n)$ una sucesión de Cauchy en $\widehat X$. Por la densidad ya probada, para cada $n$ elegimos $z_n\in X$ con
\[
 \widehat d(i(z_n),U_n)<2^{-n}.
\]
Estas elecciones se hacen una sola vez, con tolerancias predeterminadas. No dependen del $\eps$ que aparecerá después al comprobar la convergencia. Para $n,m\in\N$,
\[
 d(z_n,z_m)\leq 2^{-n}+\widehat d(U_n,U_m)+2^{-m}.
\]
Dado $\eps>0$, elegimos $N$ de modo que los dos términos geométricos sumen menos que $\eps/2$ para $n,m\geq N$ y que $\widehat d(U_n,U_m)<\eps/2$. Así $(z_n)$ es de Cauchy en $X$. Sea $Z=[(z_n)]\in\widehat X$. El paso 3 aplicado a esta sucesión da $i(z_n)\to Z$. Finalmente,
\[
 \widehat d(U_n,Z)\leq\widehat d(U_n,i(z_n))+\widehat d(i(z_n),Z)
 <2^{-n}+\widehat d(i(z_n),Z)\longrightarrow0.
\]
Toda sucesión de Cauchy de $\widehat X$ converge allí; por tanto es completo.
\end{proof}

\begin{observacion}[La elección diagonal y sus cuantificadores]
Si cada $U_n$ se representa por una sucesión $(u_{n,k})_k$, el punto $z_n$ del paso 4 puede tomarse como $u_{n,k_n}$ con $k_n$ suficientemente grande. Se puede imponer además $k_n>k_{n-1}$, porque la aproximación vale para todos los índices de una cola. Lo esencial es escoger $\widehat d(i(u_{n,k_n}),U_n)<2^{-n}$, y no un error fijo $\eps$ antes de definir la diagonal: una sucesión candidata al límite debe estar definida independientemente de la tolerancia con la que más tarde se compruebe que converge.
\end{observacion}

\section{Dos modelos y una precaución sobre las normas}
\begin{ejemplo}[Aproximación de funciones discontinuas]\leavevmode\par
En $C([0,1])$, la fórmula $\norm{f}_1=\int_0^1|f(t)|\,dt$ define una norma. Si una función continua no es cero, su valor absoluto es mayor que una constante positiva en algún intervalo de longitud positiva, y su integral es positiva. El valor absoluto es indispensable: $\int f$ por sí sola no es una norma.

Este espacio normado no es completo. Para $n\geq3$, definamos $f_n$ igual a cero en $[0,1/2]$, lineal de cero a uno en $[1/2,1/2+1/n]$ e igual a uno en el intervalo restante. Si $h$ es cero en $[0,1/2]$ y uno en $(1/2,1]$, entonces
\[
 \int_0^1|f_n-h|\leq1/n,
 \quad\text{y por tanto}\quad
 \norm{f_n-f_m}_1\leq1/n+1/m.
\]
Si existiera un límite continuo $f$ en esta norma, la misma desigualdad triangular daría $\int|f-h|=0$. En cualquier intervalo contenido en $(0,1/2)$, la continuidad obligaría a $f=0$: un valor no nulo produciría una integral positiva en una vecindad. De igual forma, $f=1$ en $(1/2,1)$. Esto contradice la continuidad en $1/2$.

El espacio de clases de funciones Lebesgue-integrables $L^1([0,1])$ proporcionará una realización de esta completación en el volumen de medida e integración. Aquí no se presupone esa teoría: la completación por sucesiones ya está construida. Para un compacto arbitrario $K\subset\R^n$, la integral puede ser sólo una seminorma sobre $C(K)$; por ejemplo, si $K=\{0\}$, todas las integrales respecto del volumen son cero. Por ello la afirmación sobre $C(K)$ requiere hipótesis sobre la medida y su soporte, o pasar primero al cociente por las funciones de seminorma cero.
\end{ejemplo}

\begin{ejemplo}[Los números p-ádicos]\label{ch09:qp}
La completación de $(\Q,d_p)$ se denota por $\Q_p$. Para una clase $r=[(r_n)]$, definimos $|r|_p=\widehat d(r,0)=\lim_n|r_n|_p$. La desigualdad ultramétrica pasa al límite, de modo que sigue siendo válida en la completación. Además de ser un espacio completo, conserva las operaciones de cuerpo. Veamos por qué. Si $(r_n)$ y $(s_n)$ son de Cauchy, sus valores absolutos p-ádicos están acotados, y
\[
 |r_ns_n-r_ms_m|_p\leq
 \max\{|r_n|_p|s_n-s_m|_p,\,|s_m|_p|r_n-r_m|_p\}.
\]
Esto prueba que los productos forman una sucesión de Cauchy; para la suma, $|(r_n+s_n)-(r_m+s_m)|_p$ está acotado por el máximo de las dos diferencias correspondientes. Las mismas estimaciones muestran que cambiar representantes por sucesiones equivalentes no cambia la clase de su suma o producto. Las identidades de cuerpo se conservan porque valen término a término. También $|rs|_p=|r|_p|s|_p$, pues la identidad racional $|r_ns_n|_p=|r_n|_p|s_n|_p$ pasa al límite de números reales.

Falta justificar la inversa. Para una clase $r\ne0$, su distancia a cero es un número $c>0$. Si $(r_n)$ la representa, entonces $|r_n|_p\to c$, de modo que $|r_n|_p\geq c/2$ en una cola. Allí
\[
 |r_n^{-1}-r_m^{-1}|_p
 =\frac{|r_n-r_m|_p}{|r_n|_p|r_m|_p}
 \leq \frac{4}{c^2}|r_n-r_m|_p.
\]
La sucesión de inversas es de Cauchy y su clase multiplica a $r$ para dar $1$. Los términos iniciales se pueden reemplazar sin alterar la clase. Así $\Q_p$ es un cuerpo.

Todo elemento tiene una expansión
\[
 r=\sum_{k\geq k_0}a_kp^k,
 \qquad k_0\in\Z,\quad a_k\in\{0,\ldots,p-1\},
\]
entendida como límite p-ádico de sus sumas parciales. Damos la razón de la existencia. Primero supongamos $|r|_p\leq1$ y tomemos racionales $r_n\to r$. Para cada $m\geq1$, elegimos uno con $|r-r_n|_p<p^{-m}$ y $|r_n|_p\leq1$. Podemos escribirlo $r_n=a/b$ con $p\nmid b$. Como $b$ tiene inversa módulo $p^m$, existe un único entero $c_m\in\{0,\ldots,p^m-1\}$ con $bc_m\equiv a\pmod{p^m}$. Por tanto $|r-c_m|_p\leq p^{-m}$. La unicidad de estos residuos, y la desigualdad ultramétrica, implican $c_{m+1}\equiv c_m\pmod{p^m}$; en consecuencia $c_{m+1}=c_m+a_mp^m$ para un dígito $a_m$. Partiendo de $c_1=a_0$, obtenemos la expansión y la cota $|r-c_m|_p\leq p^{-m}\to0$.

Para un $r$ arbitrario, su norma es finita y podemos multiplicarlo por una potencia positiva de $p$ hasta que su norma sea a lo sumo $1$; deshacer esa multiplicación produce el índice inicial $k_0$ posiblemente negativo. La unicidad de los dígitos se demuestra como en el ejemplo~\ref{ch08:padica-incompleta}: el primer dígito distinto determina la valoración de la diferencia y no puede ser cancelado por términos de mayor potencia. Este argumento muestra también cómo la completación agrega puntos que la aritmética racional no contenía.
\end{ejemplo}

\section{Unicidad y extensión: ejercicio resuelto}
\begin{ejercicio}[Lectura 9, ejercicio 1]\label{ej:lectura-9-1}\label{ch09:unicidad}
Sean $(Y,i)$ y $(Y',i')$ dos completaciones de $X$. Demuestre que existe una isometría sobreyectiva $T:Y\to Y'$ tal que $T\circ i=i'$.
\end{ejercicio}
\begin{solucion}
Si $X=\varnothing$, la densidad fuerza $Y=Y'=\varnothing$, y la afirmación es inmediata. En otro caso, para $y\in Y$ elegimos $x_n\in X$ con $i(x_n)\to y$. Como $i$ es isometría, $(x_n)$ es de Cauchy en $X$, y $(i'(x_n))$ lo es en el espacio completo $Y'$. Definimos
\[
 T(y)=\lim_{n\to\infty}i'(x_n).
\]
Si otra sucesión $z_n$ satisface $i(z_n)\to y$, entonces
\[
 d(x_n,z_n)=\rho(i(x_n),i(z_n))
 \leq\rho(i(x_n),y)+\rho(y,i(z_n))\longrightarrow0.
\]
La distancia entre $i'(x_n)$ e $i'(z_n)$ es la misma; sus límites coinciden y $T$ está bien definida. Para $y,z\in Y$, elegimos aproximaciones $i(x_n)\to y$, $i(z_n)\to z$ y obtenemos
\[
 \rho'(T(y),T(z))=\lim_nd(x_n,z_n)=\rho(y,z).
\]
Aquí la continuidad de la distancia se justifica por $|\rho(a,b)-\rho(c,d)|\leq\rho(a,c)+\rho(b,d)$. Por tanto $T$ es una isometría. La sucesión constante $x_n=x$ muestra $T(i(x))=i'(x)$.

Para probar sobreyectividad, sea $y'\in Y'$. Elegimos $x_n\in X$ con $i'(x_n)\to y'$. Entonces $(i(x_n))$ es de Cauchy y converge en $Y$ a cierto $y$. Por la definición anterior, $T(y)=y'$. Finalmente, la isometría que cumple la identidad requerida es única: es continua y queda determinada en la imagen densa $i(X)$; si $S$ es otra, para $i(x_n)\to y$ tenemos $S(y)=\lim_ni'(x_n)=T(y)$.
\end{solucion}

\begin{corolario}[Propiedad de extensión]\label{ch09:extension}
Toda aplicación uniformemente continua $f:X\to Z$, con $Z$ completo, se extiende de manera única a una aplicación uniformemente continua $\widehat f:\widehat X\to Z$ tal que $\widehat f\circ i=f$.
\end{corolario}
\begin{proof}
La aplicación $f\circ i^{-1}:i(X)\to Z$ es uniformemente continua porque $i^{-1}$ conserva distancias. Aplicamos la construcción y las estimaciones del ejercicio~\ref{ch07:ej-extension} al conjunto denso $i(X)\subset\widehat X$. Su clausura es todo $\widehat X$, por lo que obtenemos exactamente la extensión solicitada y su unicidad.
\end{proof}

\section{Completitud, construcción y unicidad}
Un punto nuevo es una clase de aproximaciones que terminan siendo indistinguibles en distancia. La densidad evita añadir puntos ajenos a esa idea, y la unicidad prueba que el resultado no depende de la realización concreta escogida. La completitud convierte después construcciones infinitas en puntos efectivos del espacio.

\chapter{Categoría de Baire y principio de contracción}\label{cap:nuevo04}\label{chap:10}\label{chap:11}
La completitud permite convertir procesos infinitos en resultados de existencia. Estudiaremos dos manifestaciones: el teorema de Baire, que hace compatibles familias numerables de condiciones abiertas y densas, y el principio de contracción, que construye una solución de $x=T(x)$ mediante iteraciones y controla su error. Comenzaremos con bolas cerradas anidadas cuyos radios tienden a cero; después desarrollaremos las consecuencias topológicas de Baire y las aplicaciones del punto fijo a sistemas, ecuaciones integrales y conjuntos autosimilares. En ambas partes será esencial justificar que los límites permanecen en los conjuntos donde se imponen las condiciones.

\section{De bolas cada vez menores a un punto}
\label{lectura:10}

\begin{teorema}[Bolas cerradas anidadas]\label{ch10:bolas}
Sea $X$ completo y sea $\overline B(x_n,r_n)$ una sucesión de bolas cerradas no vacías con
\[
 \overline B(x_{n+1},r_{n+1})\subset\overline B(x_n,r_n),
 \qquad r_n\longrightarrow0.
\]
Su intersección contiene exactamente un punto.
\end{teorema}
\begin{proof}
Para $m,k\geq n$, ambos centros pertenecen a $\overline B(x_n,r_n)$, y por tanto $d(x_m,x_k)\leq2r_n$. Esto prueba que $(x_n)$ es de Cauchy. Sea $x$ su límite. Fijado $n$, toda la cola $x_m$ con $m\geq n$ está en la bola cerrada número $n$; su límite también pertenece a ella. Como $n$ era arbitrario, $x$ pertenece a la intersección. Si $y$ perteneciera también, tendríamos $d(x,y)\leq2r_n$ para todo $n$, y dejar $n\to\infty$ da $x=y$.
\end{proof}
La bola cerrada $\overline B(x,r)=\{y:d(x,y)\leq r\}$ no debe confundirse automáticamente con la clausura de $B(x,r)$: en un espacio métrico general pueden ser diferentes. La construcción siguiente usará directamente bolas cerradas.

\section{Conjuntos pequeños en sentido topológico}
\begin{definicion}\index{F sigma@$F_\sigma$}\index{G delta@$G_\delta$}
Un conjunto es de tipo $F_\sigma$ si es unión numerable de cerrados, y de tipo $G_\delta$ si es intersección numerable de abiertos. Por las leyes de De Morgan, el complemento de un $F_\sigma$ es un $G_\delta$, y recíprocamente. Estas denominaciones se usarán para describir conjuntos de continuidad y de discontinuidad.
\end{definicion}


\begin{definicion}
Un subconjunto $A\subset X$ es \emph{denso en ninguna parte} si $\interior{\cl A}=\varnothing$. Es \emph{de primera categoría}, o \emph{magro}, si es unión numerable de conjuntos densos en ninguna parte. Es \emph{de segunda categoría} si no es de primera categoría. Un conjunto cuyo complemento es magro se llama \emph{residual}.
\end{definicion}
\index{Baire, teorema de}\index{categoria de Baire@categoría de Baire}\index{denso en ninguna parte}\index{conjunto!magro}\index{conjunto!residual}
El uso de la clausura es esencial. Los racionales tienen interior vacío en $\R$, pero su clausura es $\R$ y por ello no son densos en ninguna parte. Sí son magros: son unión numerable de conjuntos unipuntuales, cada uno cerrado y de interior vacío.

\begin{ejemplo}
En un espacio discreto, todo subconjunto es abierto y cerrado; por tanto, el único conjunto denso en ninguna parte es el vacío. Todo subconjunto no vacío es de segunda categoría. Así, $\Z$ es magro dentro de $\R$, porque es numerable y cada punto tiene interior vacío, pero es de segunda categoría dentro de sí mismo con la topología heredada, que es discreta. La categoría siempre se refiere a un espacio ambiente concreto.
\end{ejemplo}

\begin{teorema}[Baire, forma densa]\label{ch10:baire}
En un espacio métrico completo, la intersección de una familia numerable de abiertos densos es densa. Equivalentemente, el complemento de todo conjunto magro es denso.
\end{teorema}
\begin{proof}
Sean $G_1,G_2,\ldots$ abiertos densos, y sea $U$ un abierto no vacío. Vamos a construir un punto de $U\cap\bigcap_nG_n$. La densidad de $G_1$ hace no vacío al abierto $U\cap G_1$. Elegimos $x_1$ en él y $r_1>0$ suficientemente pequeño para que
\[
 \overline B(x_1,r_1)\subset U\cap G_1,\qquad r_1<1/2.
\]
La elección es posible porque, si $B(x_1,s)\subset U\cap G_1$, cualquier $0<r_1<s$ sirve. Suponiendo construida la bola número $n$, el conjunto $B(x_n,r_n)\cap G_{n+1}$ es abierto no vacío. Elegimos en él un punto $x_{n+1}$ y un radio $r_{n+1}>0$ con
\[
 \overline B(x_{n+1},r_{n+1})\subset B(x_n,r_n)\cap G_{n+1},
 \qquad r_{n+1}<2^{-(n+1)}.
\]
Las bolas cerradas son anidadas y sus radios tienden a cero. Por el teorema~\ref{ch10:bolas}, tienen un punto común $x$. La primera bola asegura $x\in U$, y la bola número $n$ asegura $x\in G_n$. Esto prueba densidad.

Si $A=\bigcup_nA_n$ con $\interior{\cl{A_n}}=\varnothing$, entonces $G_n=X\setminus\cl{A_n}$ es abierto y denso. La intersección de los $G_n$ está contenida en $X\setminus A$, que también será denso. A la inversa, para abiertos densos $G_n$, cada complemento $X\setminus G_n$ es cerrado de interior vacío, y el complemento de su unión es $\bigcap_nG_n$. Las dos formulaciones son, pues, equivalentes.
\end{proof}

\begin{corolario}\label{ch10:segunda}
Todo espacio métrico completo no vacío es de segunda categoría en sí mismo. Todo espacio métrico completo no vacío y sin puntos aislados es incontable.
\end{corolario}
\begin{proof}
Si $X$ fuera magro en sí mismo, el teorema daría que su complemento vacío es denso, imposible cuando $X\ne\varnothing$. Si además no tiene puntos aislados, cada singleton $\{x\}$ es cerrado de interior vacío. Si $X$ fuera numerable, sería unión numerable de esos conjuntos y resultaría magro, contradiciendo la primera afirmación.
\end{proof}

\begin{observacion}[Qué elecciones utiliza la demostración]
La construcción escoge una nueva bola a partir de la anterior. Trabajamos con el marco usual de teoría de conjuntos con elección, que permite estas elecciones sucesivas. En un espacio separable con una base numerable ya enumerada, se puede tomar en cada paso la primera bola básica que satisfaga las inclusiones y la cota de radio; así la regla queda determinada. La discusión de la equivalencia lógica entre versiones del teorema y principios de elección dependiente pertenece a los fundamentos de la teoría de conjuntos y no interviene en ninguna demostración del volumen. Lo que necesitaremos en análisis es el enunciado métrico que acabamos de demostrar.
\end{observacion}

\section{Ejercicios sobre categoría, cardinalidad y discontinuidades}
\begin{ejercicio}[Lectura 10, ejercicio 1]\label{ej:lectura-10-1}
Determine si es necesario que las bolas del teorema de bolas anidadas sean cerradas.
\end{ejercicio}


\begin{ejercicio}[Lectura 10, ejercicio 2]\label{ej:lectura-10-2}
Enuncie y pruebe el recíproco del teorema de bolas cerradas anidadas.
\end{ejercicio}


\begin{ejercicio}[Lectura 10, ejercicio 3]\label{ej:lectura-10-3}
Pruebe que $X$ es completo si y sólo si toda sucesión de cerrados \emph{no vacíos} $A_{n+1}\subset A_n$ con $\diam(A_n)\to0$ tiene intersección no vacía.
\end{ejercicio}


\begin{ejercicio}[Lectura 10, ejercicio 4]\label{ej:lectura-10-4}
Pruebe que $A$ es denso en ninguna parte si y sólo si $X\setminus\cl A$ es denso en $X$.
\end{ejercicio}


\begin{ejercicio}[Lectura 10, ejercicio 5]\label{ej:lectura-10-5}
Demuestre que todo subconjunto de un conjunto magro es magro.
\end{ejercicio}


\begin{ejercicio}[Lectura 10, ejercicio 6]\label{ej:lectura-10-6}
Sea $X$ completo no vacío y sea $X=\bigcup_{n\geq1}F_n$, donde los $F_n$ son cerrados. Pruebe que alguno tiene interior no vacío.
\end{ejercicio}


\begin{ejercicio}[Lectura 10, ejercicio 7]\label{ej:lectura-10-7}
Demuestre que un espacio de Banach de dimensión infinita posee una base algebraica y que toda base algebraica suya es incontable.
\end{ejercicio}


\begin{ejercicio}[Lectura 10, ejercicio 8]\label{ej:lectura-10-8}
Sea $(q_n)$ una enumeración de $\Q$, y para $\eps>0$ defina
\[
 A_\eps=\bigcup_{n\geq1}(q_n-\eps2^{-n-1},q_n+\eps2^{-n-1}).
\]
Determine si $E=\bigcap_{\eps>0}A_\eps$ contiene puntos irracionales además de los racionales.
\end{ejercicio}


\begin{ejercicio}[Lectura 10, ejercicio 9]\label{ej:lectura-10-9}
\begin{enumerate}[label=(\roman*)]
\item Examine la afirmación original: «No existe una función $f:\R\to\R$ discontinua en todo número irracional». Formule y demuestre la versión correcta referente al conjunto exacto de discontinuidades.
\item Para cada conjunto a lo sumo numerable $A\subset\R$, construya una función cuyas discontinuidades sean exactamente los puntos de $A$.
\end{enumerate}
\end{ejercicio}


\begin{ejercicio}[Lectura 10, ejercicio 10]\label{ej:lectura-10-10}
Construya funciones $\R\to\R$ discontinuas en todos los puntos de: \begin{enumerate}[label=(\roman*)]\item $\Q$;\item el conjunto de Cantor.\end{enumerate}
\end{ejercicio}


\section{Resolver una ecuación mediante iteraciones}
\label{lectura:11}
Baire permite satisfacer simultáneamente condiciones abiertas y densas. Esta herramienta explica la incontabilidad de ciertos conjuntos, impide bases algebraicas numerables en espacios de Banach infinitodimensionales y controla los conjuntos de discontinuidad; más adelante sostendrá los grandes principios sobre operadores lineales. Pasamos ahora a una segunda aplicación de la completitud, centrada en construir una solución a partir de una sucesión de aproximaciones.

Una ecuación puede escribirse en la forma $x=T(x)$. Si aplicar $T$ reduce las distancias por un factor fijo menor que uno, repetir la operación produce aproximaciones cada vez más estables. La completitud garantiza que esas aproximaciones desembocan en un punto del espacio. Además de existencia y unicidad, el argumento proporciona cotas explícitas del error.

\begin{definicion}
Una \emph{contracción} de un espacio métrico $X$ es una aplicación $T:X\to X$ para la que existe $0\leq q<1$ con
\[
 d(Tx,Ty)\leq qd(x,y)\qquad(x,y\in X).
\]
Un \emph{punto fijo} es un punto $x_*$ que satisface $T(x_*)=x_*$.
\end{definicion}
\index{contraccion@contracción}\index{punto fijo!teorema de Banach}
El caso $q=0$ incluye las aplicaciones constantes. La desigualdad estricta $d(Tx,Ty)<d(x,y)$ para $x\ne y$ no proporciona necesariamente una constante uniforme $q<1$; el teorema siguiente exige dicha constante.

\begin{teorema}[Punto fijo de Banach]\label{ch11:banach}
Si $X$ es completo y no vacío y $T:X\to X$ es una contracción con constante $q<1$, entonces existe un único punto fijo $x_*$. Para cualquier $x_0\in X$, las iteraciones $x_{n+1}=T(x_n)$ convergen a él. Si $0<q<1$, se tienen las cotas
\[
 d(x_n,x_*)\leq\frac{q^n}{1-q}d(x_1,x_0),
 \qquad d(x_n,x_*)\leq\frac{q}{1-q}d(x_n,x_{n-1})\quad(n\geq1).
\]
\end{teorema}
\begin{proof}
Si $q=0$, $T$ es constante, digamos $T(x)=c$, y $T(c)=c$; además ningún otro punto puede ser fijo. Supongamos $0<q<1$. Por inducción,
\[
 d(x_{n+1},x_n)\leq q^nd(x_1,x_0).
\]
Para $m>n$, sumamos distancias entre términos consecutivos:
\[
 d(x_m,x_n)\leq\sum_{k=n}^{m-1}d(x_{k+1},x_k)
 \leq d(x_1,x_0)\sum_{k=n}^{m-1}q^k
 \leq\frac{q^n}{1-q}d(x_1,x_0).
\]
El extremo derecho tiende a cero y no depende de $m$, así que $(x_n)$ es de Cauchy. Sea $x_*$ su límite en $X$. Una contracción es Lipschitz y continua; luego $T(x_n)\to T(x_*)$. Pero $T(x_n)=x_{n+1}\to x_*$, y la unicidad del límite da $T(x_*)=x_*$. Si $y_*$ fuese otro punto fijo,
\[
 d(x_*,y_*)=d(Tx_*,Ty_*)\leq qd(x_*,y_*),
\]
de donde $(1-q)d(x_*,y_*)\leq0$. Como $1-q>0$, los puntos coinciden.

Pasando $m\to\infty$ en la estimación de las colas obtenemos la primera cota. Para la segunda repetimos la suma desde $x_n$ y usamos $d(x_{n+j+1},x_{n+j})\leq q^{j+1}d(x_n,x_{n-1})$; la serie geométrica suma $q/(1-q)$. Esta última cota permite estimar el error a partir de dos iteraciones ya calculadas.
\end{proof}

\begin{ejemplo}[Existencia por continuidad y existencia por contracción]
Toda función continua $f:[0,1]\to[0,1]$ tiene al menos un punto fijo. En efecto, $h(x)=x-f(x)$ satisface $h(0)\leq0$ y $h(1)\geq0$; el teorema del valor intermedio proporciona un cero. La unicidad no se deduce: la identidad tiene todos los puntos fijos. Por otro lado, si $f:\R\to\R$ es derivable y $|f'(x)|\leq q<1$ para todo $x$, el valor medio da $|f(x)-f(y)|\leq q|x-y|$. Banach produce entonces un único punto fijo y convergencia desde cualquier valor inicial.
\end{ejemplo}

\section{Sistemas no lineales y elección de una norma}
\begin{proposicion}\label{ch11:jacobiano}
Sea $R\subset\R^n$ un rectángulo cerrado no vacío y $T:R\to R$ una aplicación de clase $C^1$ en una vecindad de $R$. Si
\[
 \sup_{x\in R}\max_i\sum_{j=1}^n
 \left|\frac{\partial T_i}{\partial x_j}(x)\right|\leq q<1,
\]
entonces $T$ es una contracción para la norma máxima y tiene un único punto fijo en $R$.
\end{proposicion}
\begin{proof}
El segmento entre $x,y\in R$ permanece en $R$ por convexidad. Aplicando el teorema fundamental del cálculo a $t\mapsto T_i(y+t(x-y))$,
\[
 T_i(x)-T_i(y)=\int_0^1\sum_j
 \frac{\partial T_i}{\partial x_j}(y+t(x-y))(x_j-y_j)\,dt.
\]
El valor absoluto está acotado por $q\norm{x-y}_\infty$. Tomar el máximo en $i$ demuestra la contracción. El rectángulo es cerrado en el espacio completo $\R^n$, por lo que es completo; ahora se aplica el teorema~\ref{ch11:banach}. Además de las derivadas, debe comprobarse que $T(R)\subset R$. Sin esta condición las iteraciones podrían salir del dominio.
\end{proof}
En particular, basta que cada derivada parcial tenga valor absoluto no mayor que $q/n$. La hipótesis de rectángulo cerrado corrige la posible pérdida de completitud en un rectángulo abierto.

\begin{ejemplo}[Un sistema de tres ecuaciones]\label{ch11:sistema3}
Consideremos
\begin{align*}
 3x-\cos(yz)-\tfrac12&=0,\\
 x^2-81(y+0.1)^2+\sin z+1.06&=0,\\
 e^{-xy}+20z+\tfrac{10\pi-3}{3}&=0.
\end{align*}
La rama $y+0.1\geq0$ conduce a
\[
 T(x,y,z)=\left(\frac{\cos(yz)}3+\frac16,
 \frac{\sqrt{x^2+\sin z+1.06}}9-0.1,
 -\frac{e^{-xy}}{20}-\frac{10\pi-3}{60}\right).
\]
Verificamos todas las condiciones en $R=[-1,1]^3$. El radicando $Q=x^2+\sin z+1.06$ satisface
\[
 Q\geq1.06-\sin1>0.21,\qquad Q<2.91<4.
\]
La raíz está bien definida y es diferenciable en una vecindad del cubo. La primera coordenada de $T$ está entre $-1/6$ y $1/2$; la segunda está entre $-0.1$ y $2/9-0.1$, y la tercera satisface
\[
 |T_3|\leq e/20+(10\pi-3)/60<1.
\]
Así $T(R)\subset R$.

Las derivadas no nulas son
\begin{align*}
 \partial_yT_1&=-z\sin(yz)/3,&\partial_zT_1&=-y\sin(yz)/3,\\
 \partial_xT_2&=x/(9\sqrt Q),&\partial_zT_2&=\cos z/(18\sqrt Q),\\
 \partial_xT_3&=ye^{-xy}/20,&\partial_yT_3&=xe^{-xy}/20.
\end{align*}
Las de $T_1$ están acotadas por $\sin1/3<0.3$; las de $T_2$ por $1/(9\sqrt{0.21})<0.3$ y $1/(18\sqrt{0.21})<0.3$; las de $T_3$ por $e/20<0.3$. Cada fila tiene suma menor que $0.9$, de modo que $q=0.9$ es una constante válida. Banach garantiza un único punto fijo de esta aplicación en el cubo.

En este caso el punto fijo puede darse exactamente:
\[
 (x_*,y_*,z_*)=(1/2,0,-\pi/6).
\]
La primera coordenada de $T$ es $1/3+1/6=1/2$; para la segunda el radicando vale $1/4-1/2+1.06=0.81$, cuya raíz es $0.9$; la tercera es $-1/20-(10\pi-3)/60=-\pi/6$. Se comprueban así las tres igualdades. La unicidad afirmada corresponde a la aplicación y a la rama elegida; despejar la raíz positiva no permite descartar sin examen las soluciones de la otra rama del sistema original.
\end{ejemplo}

\section{La distancia de Hausdorff y un fractal como punto fijo}
Sea $\mathcal F(X)$ la familia de subconjuntos \emph{no vacíos}, cerrados y acotados de $X$. Para $A,B\in\mathcal F(X)$ definimos
\[
 H(A,B)=\max\left\{\sup_{a\in A}\dist(a,B),
                         \sup_{b\in B}\dist(b,A)\right\}.
\]
\index{Hausdorff, distancia de}\index{fractal}\index{Cantor, conjunto de}
La no vacuidad evita distancias indefinidas y supremos sobre familias vacías. La acotación hace finitos los dos supremos: fijando $b_0\in B$, por ejemplo, $\dist(a,B)\leq d(a,b_0)$ queda acotada sobre $A$.

\begin{proposicion}\label{ch11:hausdorff-metrica}
La función $H$ es una métrica en $\mathcal F(X)$. Además,
\[
 |\dist(x,A)-\dist(x,B)|\leq H(A,B)\qquad(x\in X).
\]
\end{proposicion}
\begin{proof}
La simetría y la no negatividad están en la definición. Si $H(A,B)=0$, cada $a\in A$ tiene distancia cero a $B$ y, como $B$ es cerrado, pertenece a $B$. Al intercambiar los conjuntos obtenemos $A=B$. La vuelta es inmediata.

Para la desigualdad triangular observamos que, dados $a\in A$ y $\eta>0$, podemos elegir $b\in B$ con $d(a,b)<\dist(a,B)+\eta$ y después $c\in C$ con $d(b,c)<\dist(b,C)+\eta$. Entonces
\[
 \dist(a,C)\leq\dist(a,B)+H(B,C)+2\eta.
\]
Dejando $\eta\downarrow0$ y tomando el supremo en $a$ queda una cota por $H(A,B)+H(B,C)$. La otra distancia dirigida, de $C$ hacia $A$, se acota escogiendo primero un punto de $B$ próximo a un punto de $C$ y después uno de $A$ próximo a ese punto de $B$; da la misma suma. Tomar el máximo prueba $H(A,C)\leq H(A,B)+H(B,C)$.

Para la última afirmación, elegimos $a\in A$ con $d(x,a)<\dist(x,A)+\eta$ y un punto de $B$ a distancia menor que $H(A,B)+\eta$ de $a$. Resulta $\dist(x,B)\leq\dist(x,A)+H(A,B)+2\eta$. Al dejar $\eta\downarrow0$ e intercambiar $A,B$ obtenemos la estimación.
\end{proof}

El ejercicio~\ref{ej:lectura-11-5} demostrará que $\mathcal F(X)$ es completo exactamente cuando $X$ lo es. Para $X=\R$, definamos
\[
 \mathcal T(A)=\{a/3:a\in A\}\cup\{(a+2)/3:a\in A\}.
\]
Las transformaciones afines son homeomorfismos de $\R$, por lo que sus imágenes de cerrados son cerradas; también preservan no vacuidad y acotación. Así $\mathcal T$ actúa en $\mathcal F(\R)$. Para un punto $a/3$ de la primera pieza y un punto $b\in B$ próximo a $a$,
\[
 \dist(a/3,\mathcal T(B))\leq|a-b|/3.
\]
Tomando el ínfimo sobre $b$ se obtiene una cota $\dist(a,B)/3$. Para $(a+2)/3$ se usa la segunda pieza de $\mathcal T(B)$ y resulta la misma cota. Tomando supremos, y después intercambiando $A,B$, concluimos
\[
 H(\mathcal T(A),\mathcal T(B))\leq\tfrac13H(A,B).
\]
Por Banach existe un único conjunto fijo. Es el conjunto de Cantor: en su representación ternaria con dígitos $0$ y $2$, la primera pieza antepone el dígito $0$ y la segunda antepone el dígito $2$. Toda sucesión de esos dígitos comienza con uno de ellos, por lo que $C=C/3\cup(C+2)/3$. Las iteraciones desde cualquier cerrado acotado no vacío convergen a $C$ en distancia de Hausdorff.

\section{Ejercicios sobre contracciones, iteraciones y distancia de Hausdorff}
\begin{ejercicio}[Lectura 11, ejercicio 1]\label{ej:lectura-11-1}
Si $X$ es completo no vacío y $f:X\to X$ satisface que $f\circ f$ es una contracción, determine si $f$ tiene un punto fijo.
\end{ejercicio}


\begin{ejercicio}[Lectura 11, ejercicio 2]\label{ej:lectura-11-2}
Usando el ejercicio~\ref{ej:lectura-11-1}, pruebe que $x=\cos x$ tiene una única solución en $[0,1]$.
\end{ejercicio}


\begin{ejercicio}[Lectura 11, ejercicio 3]\label{ej:lectura-11-3}
Pruebe que $f(x)=x^2+x^{-1}$ no tiene puntos fijos en $[1,\infty)$.
\end{ejercicio}


\begin{ejercicio}[Lectura 11, ejercicio 4]\label{ej:lectura-11-4}
Sean $0<r<1$ y $\alpha\in\R$. En $C([0,r])$ defina
\[
 (\Phi f)(x)=\alpha+\int_0^x f(t)\,dt.
\]
Pruebe que $\Phi$ es contracción para la norma uniforme y encuentre su punto fijo.
\end{ejercicio}


\begin{ejercicio}[Lectura 11, ejercicio 5]\label{ej:lectura-11-5}\label{ch11:hiper-completo}
Demuestre que $(\mathcal F(X),H)$ es completo si y sólo si $X$ es completo.
\end{ejercicio}


\begin{ejercicio}[Lectura 11, ejercicio 6]\label{ej:lectura-11-6}
Calcule por iteración la solución del sistema
\[
 5x^2-y^2=0,\qquad y-\tfrac14(\sin x+\cos y)=0
\]
que se encuentra en la región positiva próxima a $(1/4,1/4)$.
\end{ejercicio}


\section{Existencia por densidad y por aproximaciones sucesivas}
Las bolas anidadas organizan dos formas de obtener existencia en espacios completos. En Baire, permiten conservar una cantidad numerable de condiciones densas; en el principio de contracción, las estimaciones geométricas conducen a un único punto fijo. Para iterar, la aplicación debe conservar el dominio, ese dominio debe ser completo y la constante de contracción debe ser estrictamente menor que uno. Verificados esos tres puntos, la misma demostración produce soluciones de ecuaciones, funciones definidas por ecuaciones integrales y conjuntos geométricos autosimilares, con cotas de error explícitas.


\parte{Compacidad y estructura topológica}{%
La compacidad permite obtener información global a partir de controles locales: subcubiertas finitas, subsucesiones convergentes y extremos de funciones continuas. En espacios métricos, su relación con la completitud y la acotación total permite reconocerla de varias maneras.

La conexidad estudia cuándo un espacio puede separarse en piezas abiertas y cómo intervienen los caminos. Con esas herramientas se abordan las compactificaciones y los modelos universales de espacios métricos. Se presupone el lenguaje topológico de la primera parte y la completitud desarrollada en la segunda.}
\chapter{Compacidad en espacios métricos}\label{cap:nuevo05}\label{chap:13}\label{chap:14}\label{chap:15}

\section{De los controles locales a la compacidad}\label{lectura:13}
Una función continua controla su variación cerca de cada punto, pero el tamaño de esa vecindad puede cambiar de un punto a otro. La compacidad permite reducir una familia de controles locales a una cantidad finita y, después, tomar el mínimo de los radios o el máximo de las cotas. Esta es la razón de su presencia en tantos argumentos de análisis. Usaremos abiertos relativos, continuidad, sucesiones de Cauchy y cerradura. Partiremos de las cubiertas abiertas, pasaremos a la extracción de subsucesiones y obtendremos después el criterio de completitud y acotación total. La compacidad relativa y la compacidad local mostrarán cómo usar estas ideas dentro de un espacio ambiente y alrededor de cada punto. Este capítulo integra las lecturas 13--15 de \cite{valencia}.

\section{Cubiertas y conjuntos compactos}
\begin{definicion}
Una \emph{cubierta abierta}\index{cubierta abierta} de $X$ es una familia $\mathcal U$ de abiertos de $X$ cuya unión es $X$. Una \emph{subcubierta} es una subfamilia que todavía cubre $X$. El espacio $X$ es \emph{compacto}\index{compacidad} si toda cubierta abierta tiene una subcubierta finita. Un subconjunto $K\subset X$ es compacto cuando lo es con su topología relativa. El conjunto vacío es compacto: la subfamilia vacía basta para cubrirlo.
\end{definicion}
Para un radio fijo $r>0$, las bolas $B(x,r)$, $x\in X$, cubren $X$. Si $D$ es denso, también las bolas con centros en $D$: dado $x\in X$, existe $a\in D\cap B(x,r)$ y la simetría de la distancia da $x\in B(a,r)$. En $\R$, los intervalos $(n,\infty)$, $n\in\Z$, cubren la recta pero ninguna cantidad finita la cubre. En un espacio discreto la familia pertinente es $\{\{x\}:x\in X\}$; sus elementos son conjuntos, mientras que los elementos de $X$ son puntos. Esta cubierta no admite quitar un solo miembro.

\begin{proposicion}\label{ch13:criterios}
Sean $\mathcal B$ una base de la topología de $X$ y $K\subset X$.
\begin{enumerate}[label=(\roman*)]
\item $X$ es compacto si y solo si toda cubierta de $X$ formada por miembros de $\mathcal B$ admite una subcubierta finita.
\item $K$ es compacto si y solo si toda familia de abiertos de $X$ que cubre $K$ admite una subfamilia finita que cubre $K$.
\end{enumerate}
\end{proposicion}
\begin{proof}
Para (i), una dirección es parte de la definición. Para la otra, dada una cubierta abierta $\mathcal U$, tomamos todos los $B\in\mathcal B$ que estén contenidos en algún $U\in\mathcal U$. Cubren $X$: si $x\in U$, la propiedad de base proporciona $x\in B\subset U$. Una subcubierta $B_1,\ldots,B_m$ permite elegir $U_i$ con $B_i\subset U_i$, y estos $U_i$ cubren $X$.

Para (ii), si $K$ es compacto, los abiertos relativos $K\cap U$ forman una cubierta de $K$ y por tanto basta una cantidad finita. Recíprocamente, toda cubierta relativa tiene la forma $\{K\cap U_\alpha\}$ con $U_\alpha$ abierto en $X$. Aplicar la hipótesis a $\{U_\alpha\}$ da la subcubierta relativa requerida.
\end{proof}

\begin{ejemplo}[Espacios finitos e intervalos cerrados]
Todo espacio finito es compacto: se elige un miembro de la cubierta para cada punto. Un espacio discreto infinito no lo es, por su cubierta de singletons.

El intervalo $[a,b]$ es compacto. Presentamos la prueba, en lugar de dar por conocida esta propiedad del cálculo. Sea $\mathcal U$ una cubierta y definamos
\[
E=\{t\in[a,b]:[a,t]\text{ se cubre con finitos miembros de }\mathcal U\}.
\]
El conjunto contiene $a$ y tiene supremo $c\in[a,b]$. Si $c<b$, un abierto $U\ni c$ contiene $(c-\delta,c+\delta)\cap[a,b]$ para algún $\delta>0$. Por la definición de supremo hay $t\in E$ con $t>c-\delta/2$; si $c=a$ podemos tomar $t=a$. Una cubierta finita de $[a,t]$, junto con $U$, cubre $[a,\min(b,c+\delta/2)]$, contradiciendo que $c$ sea cota superior. Así $c=b$. Tomando ahora un abierto que contenga $b$ y un $t\in E$ suficientemente próximo a $b$, obtenemos una cubierta finita de todo $[a,b]$.
\end{ejemplo}

\section{Imágenes, cerrados y continuidad uniforme}
\begin{teorema}\label{ch13:imagen}
La imagen continua de un espacio compacto es compacta.
\end{teorema}
\begin{proof}
Sea $f:X\to Y$ continua y $\mathcal U$ una familia de abiertos de $Y$ que cubre $f(X)$. Las preimágenes $f^{-1}(U)$ son abiertas y cubren $X$. Elegimos $U_1,\ldots,U_m$ cuyas preimágenes cubran $X$. Si $y=f(x)$, entonces $x\in f^{-1}(U_i)$ para algún $i$, de donde $y\in U_i$. La proposición \ref{ch13:criterios} termina la prueba.
\end{proof}

\begin{proposicion}\label{ch13:cerrados}
En un espacio métrico, todo compacto es cerrado y acotado. Todo subconjunto cerrado de un compacto es compacto.
\end{proposicion}
\begin{proof}
Sea $K$ compacto no vacío y $x\notin K$. Para $y\in K$ pongamos $r_y=d(x,y)/3>0$. Elegimos $y_1,\ldots,y_m$ tales que $K\subset\bigcup_iB(y_i,r_i)$ y ponemos $r=\min_i r_i>0$. Si $z\in B(x,r)$ y $y\in K$, elegimos $i$ con $d(y,y_i)<r_i$. La desigualdad triangular, escrita en sentido inverso, da
\[
d(z,y)\ge d(x,y_i)-d(x,z)-d(y,y_i)>3r_i-r-r_i\ge r.
\]
Así $z\notin K$ y $B(x,r)\subset X\setminus K$. Esto prueba que $K$ es cerrado. Para la acotación, fijemos $x_0\in X$; las bolas $B(x_0,n)$, $n\in\N$, cubren $K$. Una subcubierta finita queda contenida en la de mayor radio.

Si $F$ es cerrado en el compacto $K$, una cubierta relativa de $F$ por abiertos de $K$, junto con $K\setminus F$, cubre $K$. Extraemos una subcubierta finita y eliminamos $K\setminus F$ si aparece. Los miembros restantes cubren $F$.
\end{proof}

\begin{teorema}[Heine--Cantor]\label{ch13:heinecantor}
Toda función continua $f:X\to Y$ entre espacios métricos, con $X$ compacto, es uniformemente continua.
\end{teorema}
\begin{proof}
Fijemos $\eps>0$. Para cada $x\in X$ elegimos $r_x>0$ tal que $d(x,z)<r_x$ implique $d_Y(f(x),f(z))<\eps/2$. Extraemos una subcubierta finita $B(x_i,r_i/2)$ y ponemos $\delta=\frac12\min_i r_i$. Dados $u,v$ con $d(u,v)<\delta$, elegimos $i$ con $u\in B(x_i,r_i/2)$. Entonces
\[
d(v,x_i)\le d(v,u)+d(u,x_i)<\delta+r_i/2\le r_i.
\]
Las imágenes de $u$ y $v$ están a distancia menor que $\eps/2$ de $f(x_i)$, por lo que $d_Y(f(u),f(v))<\eps$. El radio $\delta$ se eligió antes que $u,v$: esta independencia es exactamente la continuidad uniforme.
\end{proof}

\begin{corolario}[Máximo y mínimo]\label{ch13:extremos}
Si $K$ es compacto no vacío y $f:K\to\R$ es continua, existen $x_-,x_+\in K$ tales que $f(x_-)\le f(x)\le f(x_+)$ para todo $x\in K$.
\end{corolario}
\begin{proof}
La imagen es compacta, luego cerrada y acotada. Si $M=\sup f(K)$, para cada $n$ existe $y_n\in f(K)$ con $M-1/n<y_n\le M$. Por cerradura, $M\in f(K)$. Aplicamos el mismo argumento a $-f$ para el mínimo.
\end{proof}

\section{Ejercicios propuestos: cubiertas e imágenes compactas}
\begin{ejercicio}[Lectura 13, ejercicio 1]\label{ej:lectura-13-1}
Una familia no vacía $\mathcal F$ tiene la propiedad de intersección finita cuando toda intersección de una cantidad finita de sus miembros es no vacía. Probar que $X$ es compacto si y solo si toda familia de cerrados con esta propiedad tiene intersección no vacía.
\end{ejercicio}


\begin{ejercicio}[Lectura 13, ejercicio 2]\label{ej:lectura-13-2}
Probar que una unión finita de subconjuntos compactos es compacta.
\end{ejercicio}


\begin{ejercicio}[Lectura 13, ejercicio 3]\label{ej:lectura-13-3}
En un compacto $X$, una sucesión decreciente de cerrados no vacíos $F_n$ tiene intersección no vacía. Deducir que $X$ es completo.
\end{ejercicio}


\begin{ejercicio}[Lectura 13, ejercicio 4]\label{ej:lectura-13-4}
Probar que la bola unitaria cerrada de $\ell^\infty$ es cerrada y acotada, pero no compacta.
\end{ejercicio}


\begin{ejercicio}[Lectura 13, ejercicio 5]\label{ej:lectura-13-5}
Si $X$ es compacto y $f:X\to Y$ es continua e inyectiva, probar que $f^{-1}:f(X)\to X$ es continua.
\end{ejercicio}


\begin{ejercicio}[Lectura 13, ejercicio 6]\label{ej:lectura-13-6}
Sean $X$ compacto, $f:X\to Y$ continua y $A\subset X$ tal que $f(X\setminus A)\subset Y\setminus f(A)$. Probar: (i) $f|_A:A\to f(A)$ es cerrada; (ii) si además es inyectiva, es un homeomorfismo.
\end{ejercicio}


\begin{ejercicio}[Lectura 13, ejercicio 7]\label{ej:lectura-13-7}
Si $X,Y$ son compactos, probar que $X\times Y$, con la métrica máximo, es compacto.
\end{ejercicio}


\begin{ejercicio}[Lectura 13, ejercicio 8]\label{ej:lectura-13-8}
Probar que toda cubierta abierta $\mathcal U$ de un compacto métrico admite un número de Lebesgue: existe $r>0$ tal que todo subconjunto de diámetro menor que $r$ está contenido en algún miembro de $\mathcal U$.
\end{ejercicio}


La compacidad no significa simplemente que un conjunto sea pequeño en diámetro. Significa que cualquier descripción local mediante abiertos puede reducirse a finitos datos. La extracción de esa subfamilia finita es el paso que produce controles uniformes, extremos y continuidad de inversas.

\section{Compacidad secuencial y puntos de acumulación}\label{lectura:14}
Las cubiertas abiertas expresan bien el significado de compacidad, pero en problemas concretos solemos disponer de sucesiones: aproximaciones, minimizantes o funciones que queremos hacer converger. En espacios métricos ambas perspectivas son equivalentes. La palabra \emph{métrico} es importante: fuera de esta categoría no se pueden intercambiar sin más compacidad y compacidad secuencial. Los resultados sobre cubiertas que acabamos de demostrar permiten ahora construir límites de subsucesiones. Usaremos también puntos de acumulación y bases numerables.

\begin{definicion}
Un espacio es \emph{secuencialmente compacto}\index{compacidad!secuencial} si toda sucesión tiene una subsucesión convergente a un punto del espacio. Es \emph{compacto por puntos de acumulación}\index{compacidad!por puntos de acumulacion@por puntos de acumulación} si todo subconjunto infinito tiene un punto de acumulación en el espacio.
\end{definicion}

\begin{lema}\label{ch14:acumulacion}
Si $a$ es punto de acumulación de $A$ en un espacio métrico, cada bola alrededor de $a$ contiene infinitos puntos de $A$ distintos de $a$.
\end{lema}
\begin{proof}
Si $B(a,r)\cap(A\setminus\{a\})=\{b_1,\ldots,b_m\}$ fuera finito, podríamos elegir $0<s<r$ menor que todas las distancias positivas $d(a,b_j)$. Entonces $B(a,s)$ no contendría ningún punto de $A$ distinto de $a$, contradiciendo la acumulación. Si la lista ya es vacía, la contradicción es directa.
\end{proof}

\begin{teorema}\label{ch14:equivalencias}
Para un espacio métrico $X$ son equivalentes:
\begin{enumerate}[label=(\roman*)]
\item $X$ es compacto;
\item $X$ es secuencialmente compacto;
\item todo subconjunto infinito de $X$ tiene un punto de acumulación en $X$.
\end{enumerate}
\end{teorema}
\begin{proof}
\emph{(i)$\Rightarrow$(ii).} Dada $(x_n)$, ponemos $F_n=\cl{\{x_k:k\ge n\}}$. Son cerrados no vacíos y decrecientes. El ejercicio \ref{ej:lectura-13-3} proporciona $x\in\bigcap_nF_n$. Elegimos $n_1$ con $d(x_{n_1},x)<1$. Una vez elegido $n_j$, la pertenencia $x\in F_{n_j+1}$ permite escoger $n_{j+1}>n_j$ con $d(x_{n_{j+1}},x)<1/(j+1)$. Esta es una subsucesión que converge a $x$. La barra de cerradura en $F_n$ es indispensable.

\emph{(ii)$\Rightarrow$(iii).} Si $A$ es infinito, elegimos puntos distintos $a_n\in A$. Una subsucesión converge a cierto $a\in X$. Toda bola alrededor de $a$ contiene términos de esa subsucesión distintos de $a$, porque los términos son distintos entre sí y a lo sumo uno puede ser igual a $a$. Así $a$ es punto de acumulación.

\emph{(iii)$\Rightarrow$(ii).} Dada $(x_n)$, si su recorrido es finito, algún valor aparece infinitas veces: una unión finita de conjuntos finitos de índices no podría ser $\N$. Se obtiene una subsucesión constante. Si el recorrido es infinito, tiene un punto de acumulación $x$. Elegido $n_j$, el lema \ref{ch14:acumulacion} permite hallar en $B(x,1/(j+1))$ un valor del recorrido distinto de $x_1,\ldots,x_{n_j}$. Algún índice que lo produzca debe ser mayor que $n_j$. Se construye así una subsucesión convergente.

\emph{(ii)$\Rightarrow$(i).} Primero demostraremos dos hechos. Para cada $r>0$, finitas bolas de radio $r$ cubren $X$. Si esto fallara, elegiríamos $x_1$ y, sucesivamente, $x_{n+1}\notin\bigcup_{j=1}^nB(x_j,r)$. Las distancias entre términos distintos serían al menos $r$, de modo que ninguna subsucesión sería de Cauchy y ninguna convergería.

Segundo, toda cubierta abierta $\mathcal U$ tiene un número de Lebesgue. En caso contrario, para cada $n$ habría un conjunto no vacío $A_n$ de diámetro menor que $1/n$ que no estuviera contenido en ningún miembro de $\mathcal U$. Elegimos $a_n\in A_n$ y una subsucesión $a_{n_j}\to a$. Tomamos $U\in\mathcal U$ y $r>0$ con $B(a,r)\subset U$. Para $j$ grande, $d(a_{n_j},a)<r/2$ y $1/n_j<r/2$. Todo $z\in A_{n_j}$ satisface $d(z,a)<r$, luego $A_{n_j}\subset U$, contradicción.

Sea $\lambda$ un número de Lebesgue. Cubrimos $X$ con finitas bolas de radio $\lambda/3$. Cada una tiene diámetro a lo sumo $2\lambda/3<\lambda$, por lo que está contenida en algún miembro de $\mathcal U$. Estos miembros forman una subcubierta finita. Esta prueba no usa el lema de Lebesgue para compactos para demostrar su propia hipótesis: acabamos de obtener el lema directamente de la compacidad secuencial.
\end{proof}

\begin{corolario}\label{ch14:separable}
Todo compacto métrico es separable y admite una base numerable.
\end{corolario}
\begin{proof}
Para cada $n$ elegimos un conjunto finito $D_n$ cuyos centros de bolas de radio $1/n$ cubran $X$. La unión $D=\bigcup_nD_n$ es a lo sumo numerable. Si $B(x,r)$ es una bola, elegimos $n$ con $1/n<r$ y un centro $a\in D_n$ con $d(x,a)<1/n$. Así $D$ es denso. Las bolas de centro en $D$ y radio racional positivo forman una base: si $B(x,r)$ está dado, elegimos $a\in D$ suficientemente próximo a $x$ y un racional $q$ con $d(a,x)<q<r-d(a,x)$. Entonces $x\in B(a,q)\subset B(x,r)$.
\end{proof}

\section{Ejercicios propuestos: sucesiones, productos y colas}
\begin{ejercicio}[Lectura 14, ejercicio 1]\label{ej:lectura-14-1}
Si cada $(X_n,d_n)$ es compacto, probar que $\prod_{n\ge1}X_n$ es compacto para la métrica
\[
D(x,y)=\sum_{n=1}^\infty 2^{-n}\frac{d_n(x_n,y_n)}{1+d_n(x_n,y_n)}.
\]
\end{ejercicio}


\begin{ejercicio}[Lectura 14, ejercicio 2]\label{ej:lectura-14-2}
Sea $X$ compacto y $U\subset X\times X$ abierto que contiene la diagonal $\Delta_X$. Probar que existe $r>0$ tal que $d(x,y)<r$ implica $(x,y)\in U$.
\end{ejercicio}


\begin{ejercicio}[Lectura 14, ejercicio 3]\label{ej:lectura-14-3}
Probar que todo espacio métrico no compacto contiene un subconjunto infinito numerable, cerrado y discreto.
\end{ejercicio}


\begin{ejercicio}[Lectura 14, ejercicio 4]\label{ej:lectura-14-4}
Probar que $X$ es compacto si y solo si toda función continua $X\to\R$ es acotada.
\end{ejercicio}


\begin{ejercicio}[Lectura 14, ejercicio 5]\label{ej:lectura-14-5}
Sea $1\le p<\infty$. Probar que un subconjunto cerrado y acotado $A\subset\ell^p$ es compacto si y solo si
\[
\forall\eps>0\ \exists N\ \forall x\in A:\quad
\sum_{k=N}^\infty |x_k|^p<\eps.
\]
Deducir que $K=\{x\in\ell^2:0\le x_n\le1/n\}$ es compacto.
\end{ejercicio}


\begin{ejercicio}[Lectura 14, ejercicio 6]\label{ej:lectura-14-6}
Sea $W$ una vecindad compacta de $x$ y $f:W\to[0,\infty)$ continua, con $f(x)=0$ y $f(y)>0$ para $y\in W\setminus\{x\}$. Probar que para toda vecindad $V$ de $x$ existe $r>0$ tal que $\{y\in W:f(y)<r\}\subset V$.
\end{ejercicio}


En un compacto métrico siempre podemos extraer una subsucesión convergente. Para reconocer compactos de sucesiones o funciones conviene buscar dos ingredientes: control de finitas coordenadas y una estimación uniforme de la cola. El primero por sí solo no garantiza convergencia en la norma del espacio.

\section{Acotación total y caracterización métrica}\label{lectura:15}
En $\R^n$, ser cerrado y acotado equivale a ser compacto. En $\ell^\infty$ ya encontramos una bola cerrada que contiene infinitos puntos separados por una distancia fija. La noción que excluye esta situación es la acotación total: a cualquier escala, el espacio debe poder describirse con finitos centros. La completitud garantiza, además, que los límites que esa descripción produce permanezcan dentro del espacio. La caracterización secuencial obtenida en la sección \ref{lectura:14} nos permite precisar ahora esta relación entre aproximación finita y existencia de límites.

\begin{definicion}
Un espacio métrico $X$ es \emph{totalmente acotado}\index{acotacion total@acotación total} si para cada $\eps>0$ existen $x_1,\ldots,x_m\in X$ tales que $X=\bigcup_{i=1}^mB(x_i,\eps)$. Los centros constituyen una \emph{red finita de radio $\eps$}. Aquí ``red'' significa conjunto de centros, no una familia indexada por un conjunto dirigido.
\end{definicion}
Un espacio totalmente acotado no vacío es acotado: basta una cubierta por finitas bolas de radio $1$ y estimar las distancias de los centros a uno de ellos. Todo subconjunto $A$ de un espacio totalmente acotado es totalmente acotado con su métrica relativa. Para probarlo, se cubre el espacio con bolas de radio $\eps/2$ y, de cada bola que toca $A$, se elige un centro nuevo en $A$; las bolas de radio $\eps$ con los nuevos centros cubren $A$. Este cambio de centros evita suponer que los centros originales pertenecen a $A$.

\begin{teorema}\label{ch15:criterio}
Un espacio métrico es compacto si y solo si es completo y totalmente acotado.
\end{teorema}
\begin{proof}
Si es compacto, es completo por el ejercicio \ref{ej:lectura-13-3}. La cubierta por todas las bolas de un radio fijo tiene una subcubierta finita, lo que prueba la acotación total.

Para el recíproco, sea $(x_n)$ una sucesión. Una cubierta finita por bolas de radio $1/2$ contiene una bola a la que pertenecen infinitos términos: de lo contrario, la unión de las cantidades finitas de índices sería finita. Sea $I_1$ un conjunto infinito de esos índices. Cubrimos ahora $X$ por finitas bolas de radio $1/4$ y elegimos un conjunto infinito $I_2\subset I_1$ cuyos términos estén en una sola bola. Continuamos con radios $2^{-j}$. Elegimos índices $n_j\in I_j$ con $n_{j+1}>n_j$; esto es posible porque cada $I_j$ es infinito. Si $k,l\ge j$, entonces $n_k,n_l\in I_j$ y
\[
d(x_{n_k},x_{n_l})<2\cdot2^{-j}.
\]
La subsucesión es de Cauchy. La completitud le proporciona un límite en $X$. Así toda sucesión admite una subsucesión convergente y el teorema \ref{ch14:equivalencias} implica compacidad.
\end{proof}

\begin{observacion}[Heine--Borel y su alcance]
En $\R^n$, un conjunto acotado está contenido en una caja $[-M,M]^n$. Para cada $\eps>0$, dividimos cada intervalo en finitos subintervalos de longitud menor que $\eps/\sqrt n$. Cada caja pequeña tiene diámetro menor que $\eps$; eligiendo un punto del conjunto en cada caja que lo toque obtenemos una cubierta finita por bolas de radio $\eps$. Por tanto el conjunto es totalmente acotado. Si además es cerrado, es completo como subespacio cerrado del espacio completo $\R^n$, y el teorema da compacidad. La necesidad de cerrado y acotado ya se probó. En dimensión infinita, acotado no implica totalmente acotado, como muestran los vectores $e_n$ de la bola de $\ell^\infty$.\index{teorema de Heine--Borel}
\end{observacion}

\begin{proposicion}\label{ch15:relativa}
Para $A\subset X$ son equivalentes:
\begin{enumerate}[label=(\roman*)]
\item la cerradura $\cl A^{\,X}$ es compacta;
\item toda sucesión en $A$ tiene una subsucesión convergente a algún punto de $X$;
\item $A$ es totalmente acotado y $\cl A^{\,X}$ es completo.
\end{enumerate}
Cuando se cumplen, $A$ es \emph{relativamente compacto en $X$}.\index{compacidad!relativa}
\end{proposicion}
\begin{proof}
De (i) se sigue (ii) aplicando compacidad secuencial a la cerradura. También se sigue (iii), porque la cerradura es completa y totalmente acotada, y sus subconjuntos son totalmente acotados.

Para (ii)$\Rightarrow$(i), sea $(z_n)$ una sucesión en $\cl A$. Si $A$ es vacío no hay nada que probar. En otro caso, la definición de cerradura permite elegir $a_n\in A$ con $d(z_n,a_n)<1/n$. Por (ii), una subsucesión $a_{n_j}$ converge a $a\in X$. Entonces
\[
d(z_{n_j},a)\le 1/n_j+d(a_{n_j},a)\longrightarrow0.
\]
Como $\cl A$ es cerrado, $a\in\cl A$. Así $\cl A$ es secuencialmente compacto. Para (iii)$\Rightarrow$(i), el ejercicio \ref{ej:lectura-15-1} demuestra que $\cl A$ es totalmente acotado; se aplica el teorema \ref{ch15:criterio}.
\end{proof}

\begin{ejemplo}[El espacio ambiente importa]
El intervalo $(0,1)$ es relativamente compacto en $\R$, pues su cerradura allí es $[0,1]$. No es relativamente compacto en sí mismo: su cerradura relativa es $(0,1)$ y la sucesión $1/n$ no tiene ninguna subsucesión que converja a un punto de ese espacio. En un ambiente completo, ser relativamente compacto equivale a ser totalmente acotado. En un ambiente incompleto hay que comprobar también la completitud de la cerradura.
\end{ejemplo}

\section{Compacidad alrededor de cada punto}
\begin{definicion}
Un espacio métrico $X$ es \emph{localmente compacto}\index{compacidad!local} si cada punto tiene una vecindad compacta: para cada $x$ existen un abierto $O$ y un compacto $K$ con $x\in O\subset K$. Equivalentemente, cada punto tiene una vecindad cuya cerradura es compacta. Una vecindad no tiene por qué ser abierta.
\end{definicion}
La equivalencia se verifica así: una vecindad compacta es cerrada, por lo que es su propia cerradura; si una vecindad tiene cerradura compacta, esa cerradura sigue siendo una vecindad del punto.

\begin{proposicion}\label{ch15:subespacios}
Todo subespacio abierto o cerrado de un espacio métrico localmente compacto es localmente compacto.
\end{proposicion}
\begin{proof}
Sea $x\in A$ y sea $K$ una vecindad compacta de $x$ en $X$. Si $A$ es cerrado, $K\cap A$ es compacto y contiene una vecindad relativa de $x$ en $A$.

Si $A$ es abierto, elegimos un abierto $O$ con $x\in O\subset K$. Existe $r>0$ con $\overline B(x,r)\subset A\cap O$: basta escoger primero una bola de radio $2r$ contenida en ese abierto. La bola cerrada es un cerrado de $K$, luego es compacta, y constituye una vecindad de $x$ tanto en $X$ como en $A$.
\end{proof}

\begin{proposicion}\label{ch15:engrosamiento}
Si $X$ es localmente compacto y $A\subset X$ es compacto, existe $r>0$ tal que $\cl{B(A,r)}$ es compacto, donde $B(A,r)=\{x:\dist(x,A)<r\}$.
\end{proposicion}
\begin{proof}
Para $A=\varnothing$, interpretamos $B(A,r)=\varnothing$. En otro caso, por el ejercicio \ref{ej:lectura-15-3} cada $a\in A$ admite $r_a>0$ con $\overline B(a,r_a)$ compacta. Elegimos una subcubierta finita $B(a_i,r_i/2)$ de $A$, ponemos $r=\frac14\min_i r_i$ y $K=\bigcup_i\overline B(a_i,r_i)$. El conjunto $K$ es compacto. Si $z\in B(A,r)$, hay $a\in A$ con $d(z,a)<r$, y después un $i$ con $d(a,a_i)<r_i/2$. Así $d(z,a_i)<r+r_i/2<r_i$ y $z\in K$. Como $K$ es cerrado, también $\cl{B(A,r)}\subset K$. Esta cerradura es un cerrado de un compacto, luego compacta. No se ha supuesto que todas las bolas cerradas de $X$ sean compactas: solo se emplean radios locales adecuados.
\end{proof}

\section{Ejercicios propuestos: compacidad local y funciones de corte}
\begin{ejercicio}[Lectura 15, ejercicio 1]\label{ej:lectura-15-1}
Si $A\subset X$ es totalmente acotado, probar que su cerradura también lo es.
\end{ejercicio}


\begin{ejercicio}[Lectura 15, ejercicio 2]\label{ej:lectura-15-2}
Probar que el recorrido de toda sucesión de Cauchy es totalmente acotado.
\end{ejercicio}


\begin{ejercicio}[Lectura 15, ejercicio 3]\label{ej:lectura-15-3}
Probar que $X$ es localmente compacto si y solo si cada $x$ admite $r_x>0$ con la bola cerrada $\overline B(x,r_x)$ compacta.
\end{ejercicio}


\begin{ejercicio}[Lectura 15, ejercicio 4]\label{ej:lectura-15-4}
Probar que $\Q$, con la métrica inducida de $\R$, no es localmente compacto.
\end{ejercicio}


\begin{ejercicio}[Lectura 15, ejercicio 5]\label{ej:lectura-15-5}
Sean $X$ localmente compacto, $A\subset X$ compacto y $U\subset X$ abierto con $A\subset U$. Probar:
\begin{enumerate}[label=(\roman*)]
\item existen $V$ abierto y $B$ compacto con $A\subset V\subset B\subset U$;
\item existe $f:X\to[0,1]$ continua, igual a $1$ en $A$, cuyo soporte es compacto y está contenido en $U$.
\end{enumerate}
\end{ejercicio}


\begin{ejercicio}[Lectura 15, ejercicio 6]\label{ej:lectura-15-6}
Si el compacto $A$ está cubierto por abiertos $U_1,\ldots,U_n$ de un espacio localmente compacto $X$, construir funciones continuas $f_k:X\to[0,1]$ de soporte compacto contenido en $U_k$, tales que $\sum_k f_k\le1$ en $X$ y $\sum_kf_k=1$ en $A$.
\end{ejercicio}


\section{De la extracción finita a la localización}
Las cubiertas, las subsucesiones y las redes finitas son tres formas complementarias de reconocer la compacidad métrica. Las primeras convierten datos locales en controles uniformes; las segundas permiten obtener límites de aproximaciones. La acotación total controla todas las escalas; la completitud retiene sus límites. La compacidad local permite construir funciones de corte y particiones finitas de la unidad, herramientas que más adelante harán posible localizar problemas de integración y aproximación.

\chapter{Conexidad, componentes y caminos}\label{cap:nuevo06}\label{chap:17}\label{chap:18}

\section{Separaciones, componentes y conexiones}\label{lectura:17}
La compacidad controla cubiertas; la conexidad controla separaciones. Un espacio conexo no puede dividirse en dos regiones abiertas disjuntas y no vacías. Esta propiedad explica el teorema del valor intermedio y se conserva bajo aplicaciones continuas, aunque no implique todavía que dos puntos puedan unirse mediante una curva. Usaremos topologías relativas, cerraduras y continuidad. Comenzaremos por las separaciones y sus formulaciones equivalentes, estudiaremos después las componentes maximales y distinguiremos finalmente los caminos de los arcos inyectivos. Esta distinción permitirá explicar tanto la curva del seno como la conexión mediante segmentos en abiertos de espacios normados. El capítulo integra las lecturas 17--18 de \cite{valencia}.

\section{Definiciones equivalentes y ejemplos}
\begin{proposicion}\label{ch17:def}
Para un espacio métrico $X$ son equivalentes:
\begin{enumerate}[label=(\roman*)]
\item los únicos subconjuntos simultáneamente abiertos y cerrados son $\varnothing$ y $X$;
\item no existen abiertos disjuntos no vacíos $U,V$ con $X=U\cup V$;
\item si $U,V$ son abiertos disjuntos y cubren $X$, uno de ellos es vacío.
\end{enumerate}
Si se cumplen, $X$ es \emph{conexo}; en caso contrario es \emph{disconexo}. Una pareja $U,V$ como en (ii) se llama una \emph{separación}. Un subconjunto es conexo si lo es con su topología relativa.\index{conexidad}\index{separacion@separación}
\end{proposicion}
\begin{proof}
Un subconjunto $A$ abierto y cerrado, no vacío y distinto de $X$, produce la separación $A,X\setminus A$: ambos son abiertos, disjuntos y no vacíos. En sentido inverso, en una separación $U,V$ tenemos $V=X\setminus U$, de modo que $U$ es abierto y cerrado. Las condiciones (ii) y (iii) son la misma afirmación expresada por negación y por implicación, respectivamente.
\end{proof}
El vacío y los singletons son conexos. Un espacio discreto no vacío es conexo si y solo si consta de un punto: si tiene al menos dos, cualquier singleton es abierto y cerrado propio.

\begin{lema}\label{ch17:intervalos}
Todo intervalo de $\R$ es conexo; se admiten intervalos abiertos, cerrados, semiabiertos, ilimitados, singletons y el vacío.
\end{lema}
\begin{proof}
Sea $I$ un intervalo y supongamos que tiene una separación $U,V$. Elegimos $a\in U$, $b\in V$; intercambiando los nombres si hace falta, podemos suponer $a<b$. El conjunto $E=U\cap[a,b]$ es no vacío y acotado; sea $c=\sup E$. Por ser $I$ intervalo, $[a,b]\subset I$ y $c\in I$. Existe una sucesión de puntos de $E$ que converge a $c$, por la propiedad del supremo; como $U$ es cerrado en $I$, tenemos $c\in U$. En particular $c\ne b$. Como $U$ es abierto relativo, contiene $I\cap(c-r,c+r)$ para algún $r>0$. Elegimos $t$ con $c<t<\min\{c+r,b\}$. Entonces $t\in E$, contradiciendo la definición de $c$. No existe tal separación.
\end{proof}

\begin{proposicion}\label{ch17:ambiental}
Un subconjunto $A\subset X$ es disconexo si y solo si existen abiertos disjuntos $U,V$ del ambiente $X$ tales que $A\subset U\cup V$ y ambos $A\cap U$, $A\cap V$ son no vacíos.
\end{proposicion}
\begin{proof}
La suficiencia se obtiene intersectando con $A$: aparecen dos abiertos relativos que lo separan. Para la necesidad, sea $A=A_1\cup A_2$ una separación relativa. Cada punto de $A_1$ tiene una bola que no toca $A_2$, porque $A_1$ es abierto en $A$; por ello $\dist(a,A_2)>0$ para $a\in A_1$. La afirmación correspondiente vale intercambiando los índices. Definimos
\[
U=\{x:\dist(x,A_1)<\dist(x,A_2)\},\qquad
V=\{x:\dist(x,A_2)<\dist(x,A_1)\}.
\]
Las distancias a conjuntos no vacíos son funciones continuas, así que $U,V$ son abiertos y disjuntos. En $A_1$ la primera distancia es cero y la segunda positiva; luego $A_1\subset U$. Del mismo modo $A_2\subset V$. Así los abiertos ambientales tienen todas las propiedades requeridas. Es útil construirlos explícitamente: los abiertos ambientales que representan inicialmente una separación relativa no tienen por qué ser disjuntos fuera de $A$.
\end{proof}

\section{Imágenes y familias de conjuntos conexos}
\begin{teorema}\label{ch17:imagen}
La imagen continua de un espacio conexo es conexa.
\end{teorema}
\begin{proof}
Si $f(X)$ tuviera una separación relativa $U,V$, las preimágenes $f^{-1}(U)$ y $f^{-1}(V)$ serían abiertas, disjuntas y cubrirían $X$. Serían no vacías porque ambos miembros de la separación contienen puntos de la imagen. Esto contradice la conexidad de $X$.
\end{proof}

\begin{corolario}\label{ch17:binario}
$X$ es conexo si y solo si toda aplicación continua $X\to\{0,1\}$ es constante. Equivalentemente, toda aplicación continua de $X$ a cualquier espacio discreto es constante.
\end{corolario}
\begin{proof}
Si $X$ es conexo, su imagen continua es conexa y una imagen discreta no puede tener dos puntos. Para la otra dirección, una separación $U,V$ define la función que vale $0$ en $U$ y $1$ en $V$. Es continua porque las preimágenes de los singletons son abiertas, y no es constante. El mismo argumento sirve con dos puntos distintos de cualquier espacio discreto de al menos dos elementos. Para el dominio vacío, usamos la convención habitual de que la función vacía es constante en el sentido de que no toma dos valores diferentes.
\end{proof}

\begin{definicion}
Una familia no vacía $\mathcal C$ de conjuntos no vacíos es \emph{concatenada}\index{familia concatenada} si para cualesquiera $A,B\in\mathcal C$ existe una cadena finita $A=C_1,\ldots,C_m=B$ de miembros de la familia con $C_i\cap C_{i+1}\ne\varnothing$.
\end{definicion}

\begin{proposicion}\label{ch17:union}
La unión de una familia concatenada de conjuntos conexos es conexa. En particular, una familia de conjuntos conexos con un punto común tiene unión conexa.
\end{proposicion}
\begin{proof}
Sea $f:\bigcup\mathcal C\to\{0,1\}$ continua. En cada miembro de $\mathcal C$ es constante, por el corolario \ref{ch17:binario}. Si dos miembros se intersectan, sus valores constantes coinciden al evaluar $f$ en un punto de intersección. Una cadena finita transmite la igualdad entre cualesquiera dos miembros. Así $f$ es constante en toda la unión, que es conexa por el mismo corolario. Si hay un punto común, cualquier pareja de miembros se intersecta y la familia es concatenada. Si una familia contiene conjuntos vacíos, se pueden eliminar sin alterar la unión.
\end{proof}

\begin{teorema}\label{ch17:producto}
El producto finito de espacios métricos no vacíos $X_1,\ldots,X_n$, con su topología producto, es conexo si y solo si cada factor es conexo.
\end{teorema}
\begin{proof}
Si el producto es conexo, la proyección sobre cada factor es continua y sobreyectiva; la imagen es conexa. La sobreyectividad requiere que los demás factores sean no vacíos.

Para la otra dirección, supongamos que todos los factores son conexos. Para $a=(a_1,\ldots,a_n)$ y $i$ definimos la fibra
\[
L_{a,i}=\{x\in\textstyle\prod_jX_j:x_j=a_j\text{ para todo }j\ne i\}.
\]
Aquí se exige igualdad en las coordenadas que no varían; escribir desigualdad cambiaría el conjunto y no permitiría identificarlo con $X_i$. La evaluación de la coordenada $i$ es un homeomorfismo de $L_{a,i}$ sobre $X_i$, con inversa que inserta esa coordenada y deja fijas las demás. Por tanto cada fibra es conexa.

Las fibras cubren el producto y forman una familia concatenada. En efecto, dados puntos $a,b$, las etapas
\[
(a_1,\ldots,a_n),\ (b_1,a_2,\ldots,a_n),\ldots,(b_1,\ldots,b_n)
\]
cambian una coordenada por vez. Cada etapa consecutiva está en una misma fibra, y las fibras sucesivas comparten el punto de transición. Para conectar dos fibras arbitrarias, tomamos un punto en cada una y añadimos al comienzo y al final esas dos fibras a la cadena construida. Aplicamos la proposición \ref{ch17:union}.
\end{proof}

\begin{ejemplo}[La intersección puede romper la conexidad]
Las semicircunferencias superior e inferior de la circunferencia unitaria son conexas: son imágenes continuas de $[0,\pi]$ y $[\pi,2\pi]$ por $t\mapsto(\cos t,\sin t)$. Su intersección es $\{(-1,0),(1,0)\}$, que es disconexa. Por tanto no hay una regla general de estabilidad de la conexidad bajo intersecciones.
\end{ejemplo}

\section{Ejercicios resueltos: separaciones y continuidad}
\begin{ejercicio}[Lectura 17, ejercicio 1: precisión del enunciado]\label{ej:lectura-17-1}
Probar que $X$ es conexo si y solo si todo subconjunto no vacío y \emph{propio} $A\subsetneq X$ tiene frontera no vacía.
\end{ejercicio}
\begin{solucion}
La palabra propio es necesaria, pues $\partial X=\varnothing$ en cualquier espacio. Recordemos $\partial A=\cl A\setminus\interior A$. Si esta frontera es vacía, de $\interior A\subset A\subset\cl A$ se deduce $\interior A=A=\cl A$; por tanto $A$ es abierto y cerrado. Recíprocamente, si $A$ es abierto y cerrado, esas tres expresiones coinciden y su frontera es vacía. Así la condición del ejercicio excluye exactamente los abiertos y cerrados no triviales, que caracterizan la desconexión.
\end{solucion}

\begin{ejercicio}[Lectura 17, ejercicio 2]\label{ej:lectura-17-2}
Probar: (i) los subconjuntos conexos de $\R$ son exactamente los intervalos; (ii) una función real continua sobre un espacio conexo toma todos los valores intermedios entre dos valores de su imagen; (iii) si toma un valor negativo y uno positivo, tiene un cero.
\end{ejercicio}
\begin{solucion}
Para (i), los intervalos son conexos por el lema \ref{ch17:intervalos}. Si $A\subset\R$ no es un intervalo, hay $a,b\in A$ y $c\notin A$ con $a<c<b$. Entonces $A\cap(-\infty,c)$ y $A\cap(c,\infty)$ son dos abiertos relativos disjuntos no vacíos que cubren $A$, de modo que $A$ no es conexo.

Para (ii), si $X$ es conexo y $f:X\to\R$ continua, el teorema \ref{ch17:imagen} muestra que $f(X)$ es conexo, y por (i) es un intervalo. Si $f(x)=a\le b=f(y)$, todo $t\in[a,b]$ pertenece a $f(X)$, lo que significa que existe $z\in X$ con $f(z)=t$. Para (iii), aplicamos (ii) a $t=0$, que está entre el valor negativo y el positivo. Los extremos de $X$ no necesitan existir: el argumento solo exige dos puntos donde se conozcan los signos.\index{teorema del valor intermedio}\index{teorema de Bolzano}
\end{solucion}

\begin{ejercicio}[Lectura 17, ejercicio 3]\label{ej:lectura-17-3}
Si $A$ es conexo y $A\subset B\subset\cl A$, probar que $B$ es conexo. Deducir que la cerradura de un conexo es conexa.
\end{ejercicio}
\begin{solucion}
Sea $f:B\to\{0,1\}$ continua. Su restricción a $A$ es constante, digamos de valor $0$. Si $b\in B$ tuviera $f(b)=1$, la preimagen $f^{-1}(\{1\})$ sería una vecindad abierta relativa de $b$ que no toca $A$. Pero $b\in\cl A$ y $A\subset B$, de modo que toda vecindad relativa de $b$ en $B$ toca $A$, contradicción. Así $f$ es constante y $B$ es conexo. Para $A=\varnothing$, también $B=\varnothing$ y el resultado es válido. Tomar $B=\cl A$ da la conclusión final.
\end{solucion}

\begin{ejercicio}[Lectura 17, ejercicio 4]\label{ej:lectura-17-4}
Probar que todo espacio normado es conexo.
\end{ejercicio}
\begin{solucion}
Para $x\in V$, el segmento $[0,x]=\{tx:0\le t\le1\}$ es imagen continua de $[0,1]$: $\norm{tx-sx}=|t-s|\norm{x}$. Es por tanto conexo. Todos estos segmentos contienen $0$, y su unión es $V$. La proposición \ref{ch17:union} prueba que $V$ es conexo. No se requiere completitud; el argumento vale también en un espacio normado complejo usando el parámetro real $t$.
\end{solucion}

\begin{ejercicio}[Lectura 17, ejercicio 5]\label{ej:lectura-17-5}
Sea $X$ conexo, $X=A\cup B$ y $A\cap B$ conexo. Probar que $A,B$ son conexos si (i) ambos son cerrados o (ii) ambos son abiertos. En (iii), decidir qué ocurre si $A$ es abierto y $B$ cerrado.
\end{ejercicio}
\begin{solucion}
Si $A\cap B=\varnothing$ y ambos son abiertos, forman una separación salvo que uno sea vacío. Si ambos son cerrados, cada uno es complemento del otro y por tanto también abierto; la misma conclusión se aplica. El miembro no vacío es entonces $X$ y ambos son conexos.

Supongamos ahora $A\cap B\ne\varnothing$. Si $A$ fuera disconexo, existiría una función continua no constante $g:A\to\{0,1\}$. Su restricción al conexo $A\cap B$ sería constante, de valor $c$. Definimos $f=g$ en $A$ y $f=c$ en $B$; las definiciones coinciden en la intersección. Si $A,B$ son abiertos, las preimágenes de abiertos por $f$ son uniones de conjuntos abiertos en $X$, así que $f$ es continua. Si $A,B$ son cerrados, para cada cerrado $F$ del codominio la preimagen es la unión de un cerrado relativo de $A$ y uno de $B$; ambos son cerrados en $X$, luego la unión lo es y $f$ también es continua. En cualquiera de los casos $f$ sería continua no constante sobre $X$, contradicción. Así $A$ es conexo. Intercambiando los papeles de $A$ y $B$, el mismo argumento de pegado prueba que $B$ es conexo.

Para (iii) no se garantiza la conexidad de ninguno de los dos por separado. En $X=\R$, tomemos $A=\R\setminus\{0\}$ y $B=[-1,0]$. El primero es abierto y disconexo, el segundo cerrado, su unión es $\R$ y su intersección $[-1,0)$ es conexa. Para que falle la conexidad de $B$, tomemos $X=[0,2]$, $A=[0,2)$, abierto relativo, y $B=\{1,2\}$, cerrado relativo. Cubren $X$ y su intersección es $\{1\}$, pero $B$ es disconexo. La condición uniforme abierto/abierto o cerrado/cerrado es la que permite el pegado anterior.
\end{solucion}

\begin{ejercicio}[Lectura 17, ejercicio 6]\label{ej:lectura-17-6}\label{ch17:ex6}
Sean $X$ conexo, $A\subset X$ conexo y $B\subset X\setminus A$ abierto y cerrado en $X\setminus A$. Probar que $A\cup B$ es conexo.
\end{ejercicio}
\begin{solucion}
Si $A$ es vacío, $B$ es abierto y cerrado en $X$ y por tanto es vacío o todo $X$, ambos conexos. Si $A\ne\varnothing$, sea $f:A\cup B\to\{0,1\}$ continua. En $A$ es constante; cambiando los nombres de $0$ y $1$ podemos suponer que vale $0$. Pongamos $D=f^{-1}(\{1\})\subset B$. Este conjunto es abierto y cerrado en $A\cup B$.

Demostraremos que $D$ es abierto y cerrado en $X$. Como $B$ es cerrado en $X\setminus A$, tenemos $\cl B\subset A\cup B$. Por ello $\cl D\subset A\cup B$ y, puesto que $D$ es cerrado en este subespacio, $\cl D=D$: es cerrado en $X$. Para la apertura, si $x\in D$, existe un abierto $O$ de $X$ con $x\in O$ y $O\cap(A\cup B)\subset D$. Como $B$ es abierto en $X\setminus A$, existe un abierto $P$ de $X$ con $x\in P$ y $P\cap(X\setminus A)\subset B$. Todo punto de $O\cap P$ que estuviera en $A$ o en $B$ estaría en $D$ por la primera condición; todo punto que no estuviera en $A$ estaría en $B$ por la segunda. Así $O\cap P\subset D$, y $D$ es abierto. La conexidad de $X$ fuerza $D=\varnothing$ o $D=X$; la segunda posibilidad es imposible porque $A\ne\varnothing$ y $D\cap A=\varnothing$. Entonces $f$ es constante y $A\cup B$ es conexo.

La sugerencia de la lectura es una versión de pegado que también podemos justificar. Si $f$ es continua por separado en $E$ y $F$, y $E\setminus F$ es abierto y cerrado en la diferencia simétrica $E\mathbin{\triangle}F$, entonces $f$ es continua en $E\cup F$, siempre que las definiciones coincidan en la intersección. Un punto de $E\setminus F$ tiene una vecindad relativa en la unión contenida en $E$, por la apertura en la diferencia simétrica; la cerradura da la afirmación correspondiente para $F\setminus E$. En un punto de $E\cap F$, intersectamos las dos vecindades proporcionadas por la continuidad de las restricciones. En los tres casos obtenemos la continuidad en la unión.
\end{solucion}

\begin{ejercicio}[Lectura 17, ejercicio 7]\label{ej:lectura-17-7}
Construir: (i) un espacio métrico conexo con una bola disconexa; (ii) un conjunto conexo cuyo interior sea disconexo.
\end{ejercicio}
\begin{solucion}
Para (i), en $\R^2$ con la distancia euclídea consideremos
\[
X=(\{-1,1\}\times[0,2])\cup([-1,1]\times\{2\}).
\]
Es la unión de tres segmentos conexos concatenados, luego es conexo. La bola relativa de centro $(-1,0)$ y radio $\sqrt5$ contiene el brazo izquierdo completo, el tramo superior con $-1\le x<0$, y el tramo derecho $\{1\}\times[0,1)$. No contiene ningún punto con primera coordenada entre $0$ y $1$. Sus intersecciones con los semiplanos $x<1/2$ y $x>1/2$ son abiertos relativos disjuntos no vacíos que la cubren. Por tanto la bola es disconexa.

Para (ii), tomemos la unión $K$ de los discos cerrados de radio $1$ centrados en $(-1,0)$ y $(1,0)$. Cada disco es conexo por convexidad, y ambos se tocan en $(0,0)$; su unión es conexa. Su interior es exactamente la unión de los dos discos abiertos, que son disjuntos. Los puntos de frontera distintos del punto de tangencia no son interiores porque tienen vecindades pequeñas que no tocan el otro disco y contienen puntos exteriores al primero. El punto de tangencia tampoco es interior, pues $(0,t)\notin K$ para todo $t\ne0$. Así el interior tiene dos partes abiertas no vacías separadas por el eje vertical y es disconexo.
\end{solucion}

La prueba de conexidad puede reducirse a excluir funciones continuas no constantes con valores $0$ y $1$. Esta perspectiva hace transparentes las imágenes continuas, las uniones concatenadas y los argumentos de pegado. La conexidad pertenece a la topología del espacio, mientras que la forma de sus bolas depende de la métrica elegida.

\section{Componentes y desconexión total}\label{lectura:18}
Un espacio disconexo no siempre se descompone en abiertos conexos. Las componentes proporcionan sus piezas conexas maximales, que son siempre cerradas pero pueden dejar de ser abiertas. Por otra parte, unir dos puntos mediante una curva es una exigencia más concreta que la conexidad. Conviene distinguir un camino continuo de un arco inyectivo, pues la concatenación conserva continuidad pero puede introducir repeticiones. Las propiedades de las imágenes, las uniones y las cerraduras conexas que acabamos de establecer serán las herramientas para describir esas piezas y compararlas con las conexiones mediante curvas.

\begin{proposicion}\label{ch18:componentes}
En un espacio métrico $X$, definamos $x\sim y$ si existe un subconjunto conexo que contiene a ambos. Esta es una relación de equivalencia. Sus clases, llamadas \emph{componentes conexas}, son precisamente los subconjuntos conexos no vacíos maximales por inclusión. Todas ellas son cerradas.\index{componente conexa}
\end{proposicion}
\begin{proof}
La reflexividad se obtiene del conexo $\{x\}$; la simetría es inmediata porque contener a $x,y$ no depende del orden. Si $x,y\in A$ y $y,z\in B$ con $A,B$ conexos, su intersección contiene $y$, así que $A\cup B$ es conexo y contiene $x,z$. Esto prueba transitividad.

La clase $C_x$ de $x$ es la unión de todos los conjuntos conexos que contienen a $x$. Todos comparten ese punto, luego la unión es conexa. Si un conexo $D$ intersecta $C_x$, la unión $D\cup C_x$ es conexa y contiene $x$; por la definición de la clase, toda la unión está en $C_x$, así que $D\subset C_x$. En particular, no puede ampliarse $C_x$ a un conexo mayor. Recíprocamente, todo conexo no vacío $D$ está contenido en la clase de cualquiera de sus puntos; si es maximal, coincide con esa clase.

Finalmente $\cl{C_x}$ es conexo por el ejercicio \ref{ej:lectura-17-3}. Como contiene a $C_x$, su maximalidad obliga a $\cl{C_x}=C_x$. Esto demuestra la cerradura.
\end{proof}

\begin{definicion}
Un espacio es \emph{totalmente disconexo}\index{desconexion total@desconexión total} si sus componentes son singletons; equivalentemente, todo subconjunto conexo no vacío consta de un punto.
\end{definicion}
Un espacio no vacío es conexo si y solo si tiene una componente. En un espacio discreto las componentes son puntos, pero el recíproco falla: $\Q$ es totalmente disconexo sin ser discreto. En efecto, si un subconjunto de $\Q$ contiene racionales $a<b$, elegimos un irracional $c\in(a,b)$; sus intersecciones con $(-\infty,c)$ y $(c,\infty)$ lo separan. Además ningún racional está aislado, porque cualquier intervalo a su alrededor contiene otros racionales.

\begin{proposicion}\label{ch18:abiertosEuclideos}
Todo abierto de $\R^n$ es unión de a lo sumo numerables abiertos conexos disjuntos, sus componentes.
\end{proposicion}
\begin{proof}
Sea $A$ abierto y $C$ una componente de $A$. Para $x\in C$, existe una bola $B(x,r)\subset A$. La bola es convexa, luego conexa: cada punto puede unirse con el centro por un segmento y la unión de esos segmentos es conexa. Como la bola intersecta $C$, su unión con $C$ es conexa y la maximalidad implica $B(x,r)\subset C$. Por tanto cada componente es abierta en $\R^n$. Cada una contiene algún punto de $\Q^n$, pues este conjunto es denso. Fijando una enumeración de $\Q^n$, asignamos a cada componente el primer racional de la enumeración que contenga; componentes distintas reciben racionales distintos, así que hay a lo sumo numerables. Son cerradas en $A$ por la proposición anterior, aunque no necesariamente cerradas en $\R^n$: el intervalo $(0,1)$ es una componente de sí mismo.
\end{proof}

\section{Caminos, arcos y el problema de las repeticiones}
\begin{definicion}
Un \emph{camino}\index{camino} en $X$ es una aplicación continua $\gamma:[0,1]\to X$. Un \emph{arco}\index{arco} es un camino inyectivo. Sus extremos son $\gamma(0),\gamma(1)$. Un espacio es \emph{conexo por caminos} si cada pareja de puntos es extremo de un camino. Es \emph{arcoconexo}\index{arcoconexidad} si cada pareja de puntos \emph{distintos} es extremo de un arco. Por convención, un punto está conectado consigo mismo; no se exige un mapa inyectivo con extremos iguales, pues eso sería imposible. El vacío y el singleton cumplen ambas condiciones con estas convenciones.
\end{definicion}
Un arco es un homeomorfismo de $[0,1]$ sobre su imagen, porque es una aplicación continua inyectiva desde un compacto a un espacio métrico. Si dos caminos tienen un extremo común, su concatenación, recorriendo cada uno en la mitad del tiempo, es continua por el lema de pegado sobre los intervalos cerrados $[0,1/2]$ y $[1/2,1]$. Pero incluso si ambos caminos son inyectivos, la concatenación puede visitar un punto dos veces. Para conservar un arco, necesitamos el siguiente recorte.

\begin{lema}[Recorte de dos arcos]\label{ch18:recorte}
Si hay un arco de $x$ a $y$ y otro de $y$ a $z$, entonces, cuando $x\ne z$, hay un arco de $x$ a $z$ contenido en la unión de sus imágenes.
\end{lema}
\begin{proof}
Sean $p,q:[0,1]\to X$ los arcos con $p(0)=x$, $p(1)=q(0)=y$ y $q(1)=z$. Si $z\in p([0,1])$, la inyectividad da un único $s\in(0,1]$ con $p(s)=z$, y $t\mapsto p(st)$ es el arco buscado.

Si $z\notin p([0,1])$, el conjunto
\[
E=\{t\in[0,1]:q(t)\in p([0,1])\}
\]
es cerrado, porque la imagen compacta de $p$ es cerrada. Es no vacío, pues contiene $0$, y no contiene $1$. Tiene máximo $t_0<1$. Existe un único $s_0\in[0,1]$ con $p(s_0)=q(t_0)$. Por maximalidad de $t_0$, la imagen de $q$ para $t>t_0$ no toca la imagen de $p$.

Si $s_0=0$, basta el arco $t\mapsto q(t_0+(1-t_0)t)$. Si $s_0>0$, usamos
\[
r(t)=\begin{cases}
p(2s_0t),&0\le t\le1/2,\\
q(t_0+(2t-1)(1-t_0)),&1/2\le t\le1.
\end{cases}
\]
Las dos expresiones coinciden en $t=1/2$, así que $r$ es continua. Cada tramo es inyectivo y sus imágenes solo se intersectan en el punto de unión, por la elección de $t_0$. Por ello $r$ es inyectiva y tiene extremos $x,z$. Este último control de la intersección es lo que una concatenación sin recorte no garantiza.
\end{proof}

\begin{proposicion}\label{ch18:arcoimplica}
Todo espacio conexo por caminos, y en particular todo espacio arcoconexo, es conexo.
\end{proposicion}
\begin{proof}
Si existiera una separación $U,V$, elegiríamos $x\in U$ y $y\in V$. Un camino $\gamma$ entre ambos tendría imagen conexa por ser imagen continua del intervalo $[0,1]$. Sin embargo, $\gamma([0,1])\cap U$ y $\gamma([0,1])\cap V$ serían una separación no trivial de esa imagen. La contradicción excluye la separación de $X$.
\end{proof}

\begin{ejemplo}[Una curva límite conexa sin conexión por caminos]\label{ch18:seno}
Consideremos
\[
S=\{(x,\sin(1/x)):x\ne0\}\cup(\{0\}\times[-1,1]).
\]
La gráfica $\Gamma_+$ para $x>0$ es conexa por ser imagen continua del intervalo $(0,\infty)$; la gráfica $\Gamma_-$ para $x<0$ también lo es. La cerradura de cada una incluye el segmento vertical $L=\{0\}\times[-1,1]$. Para verlo, dado $y\in[-1,1]$, elegimos $\alpha$ con $\sin\alpha=y$ y puntos $x_n=1/(2\pi n+\alpha)>0$ para $n$ grande. Entonces $(x_n,\sin(1/x_n))\to(0,y)$. Para la rama negativa usamos un ángulo $\beta$ con $-\sin\beta=y$ y $x_n=-1/(2\pi n+\beta)$. No hay otros puntos nuevos en las cerraduras: si la primera coordenada límite no es cero, la continuidad de $\sin(1/x)$ determina la segunda; si es cero, la segunda queda en $[-1,1]$. Por tanto $\cl{\Gamma_+}=\Gamma_+\cup L$ y $\cl{\Gamma_-}=\Gamma_-\cup L$. Son conexos y se intersectan en $L$, de modo que su unión $S$ es conexa.

No existe siquiera un camino en $S$ desde $(1,\sin1)$ hasta $(-1,-\sin1)$. Supongamos que $f(t)=(u(t),v(t))$ fuera uno. Por el valor intermedio, $u$ se anula; sea
\[
t_0=\min\{t\in[0,1]:u(t)=0\}.
\]
El mínimo existe porque el conjunto de ceros es cerrado no vacío en un compacto. Además $t_0>0$ y $u(t)>0$ para $0\le t<t_0$: no puede cambiar de signo sin anularse. Fijemos $0<\eta<t_0$. Tenemos $u(t_0-\eta)>0$. Existen números $a,b$ entre $0$ y $u(t_0-\eta)$ con $\sin(1/a)=1$ y $\sin(1/b)=-1$, tomando denominadores suficientemente grandes de la forma $\pi/2+2\pi n$ y $3\pi/2+2\pi n$. El valor intermedio proporciona $s,r\in(t_0-\eta,t_0)$ con $u(s)=a$, $u(r)=b$. Puesto que sus primeras coordenadas son positivas y $f$ toma valores en $S$, se sigue $v(s)=1$ y $v(r)=-1$.

Esto ocurre para todo $\eta>0$ suficientemente pequeño, así que $v$ no tiene límite por la izquierda en $t_0$, contradiciendo continuidad. El argumento no presupone que los tiempos de una rama y otra formen intervalos, ni identifica límites de composiciones sin justificar qué valores recorre la primera coordenada.
\end{ejemplo}

\section{Abiertos conexos de espacios normados}
\begin{definicion}
En un espacio vectorial real o complejo, el \emph{segmento}\index{segmento} entre $x,y$ es
\[
[x,y]=\{(1-t)x+ty:0\le t\le1\}.
\]
Un conjunto es convexo si contiene el segmento entre cualesquiera dos de sus puntos. Si $x\ne y$, la parametrización del segmento es inyectiva, porque $(t-s)(y-x)=0$ implica $t=s$.
\end{definicion}

\begin{teorema}\label{ch18:abiertoconexo}
Todo abierto conexo de un espacio normado es arcoconexo. En particular, la afirmación vale en un espacio de Banach, sin que se necesite su completitud.
\end{teorema}
\begin{proof}
El vacío y los singletons no presentan dificultad. Fijemos $x\in A$ y sea $P$ el conjunto de puntos que son iguales a $x$ o pueden unirse a $x$ mediante un arco contenido en $A$. Mostraremos que $P$ y $A\setminus P$ son abiertos relativos en $A$.

Si $y\in P$, elegimos $r>0$ con $B(y,r)\subset A$. Para $z\in B(y,r)$, el segmento $[y,z]$ está en esa bola, ya que
\[
\norm{(1-t)y+tz-y}=t\norm{z-y}<r.
\]
Si $y=x$, este segmento proporciona directamente la conexión a $z$; si $y\ne x$, combinamos el arco de $x$ a $y$ con el segmento de $y$ a $z$ y aplicamos el lema de recorte \ref{ch18:recorte}. Los casos de puntos iguales se resuelven por la convención de conexión consigo mismo. Así toda la bola está en $P$ y $P$ es abierto.

Si $y\in A\setminus P$, elegimos otra vez $r>0$ con $B(y,r)\subset A$. Si algún $z$ de esa bola perteneciera a $P$, podríamos unir $x$ a $z$ y después unir $z$ a $y$ por un segmento dentro de la bola; el lema de recorte produciría un arco de $x$ a $y$, contradicción. Por tanto la bola está en $A\setminus P$, y este conjunto también es abierto. Como $A$ es conexo y $x\in P$, debe ser $P=A$. Al haber elegido arbitrariamente el punto inicial $x$, cualesquiera dos puntos distintos de $A$ se unen por un arco.
\end{proof}

\section{Ejercicios resueltos: componentes, Cantor y arcos}
\begin{ejercicio}[Lectura 18, ejercicio 1]\label{ej:lectura-18-1}
Si un espacio métrico tiene una cantidad finita de componentes, probar que cada componente es abierta.
\end{ejercicio}
\begin{solucion}
Sean $C_1,\ldots,C_m$ las componentes. Cada una es cerrada, y para un índice $i$,
\[
X\setminus C_i=\bigcup_{j\ne i}C_j
\]
es una unión finita de cerrados, por tanto cerrada. Su complemento $C_i$ es abierto. La finitud es esencial para este argumento: en $\Q$ hay infinitas componentes unipuntuales y ninguna es abierta.
\end{solucion}

\begin{ejercicio}[Lectura 18, ejercicio 2]\label{ej:lectura-18-2}
Probar que todo espacio ultramétrico es totalmente disconexo.
\end{ejercicio}
\begin{solucion}
Recordemos que una ultramétrica satisface $d(x,z)\le\max\{d(x,y),d(y,z)\}$. Toda bola abierta $B(a,r)$ es también cerrada. En efecto, si $y\notin B(a,r)$, entonces $d(a,y)\ge r$. Para $z\in B(y,d(a,y))$, si $d(a,z)<d(a,y)$, la desigualdad ultramétrica daría
\[
d(a,y)\le\max\{d(a,z),d(z,y)\}<d(a,y),
\]
una contradicción. Así $d(a,z)\ge d(a,y)\ge r$, y la bola alrededor de $y$ queda en el complemento de $B(a,r)$. El complemento es abierto.

Dados $x\ne y$, la bola $B(x,d(x,y))$ contiene a $x$ y no a $y$, y es abierta y cerrada. Sus intersecciones con cualquier conjunto que contenga ambos puntos lo separan. Por tanto ningún conexo puede contener dos puntos distintos, que es la desconexión total.
\end{solucion}

\begin{ejercicio}[Lectura 18, ejercicio 3]\label{ej:lectura-18-3}
Un espacio es totalmente separado si para cualesquiera puntos \emph{distintos} $x,y$ hay una separación $U,V$ del espacio con $x\in U$, $y\in V$. Probar que Cantor es totalmente separado.
\end{ejercicio}
\begin{solucion}
Usamos la construcción por eliminación de tercios del ejercicio \ref{ej:lectura-8-6}. En la etapa $n$, el conjunto $C_n$ es unión de $2^n$ intervalos cerrados, disjuntos entre sí, de longitud $3^{-n}$; dos intervalos distintos están separados por un intervalo abierto de longitud positiva. Además $C\subset C_n$ y $C=\bigcap_n C_n$.

Dados $x,y\in C$ distintos, elegimos $n$ con $3^{-n}<|x-y|$. No pueden estar en el mismo intervalo de $C_n$, porque el diámetro de cada uno es $3^{-n}$. Sea $I$ el intervalo de esa etapa que contiene a $x$ y sea $J$ la unión de los demás. Definimos
\[
U=C\cap I,\qquad V=C\cap J.
\]
Ambos son cerrados relativos en $C$: $I$ es cerrado y $J$ es una unión finita de cerrados. Son disjuntos, su unión es $C$ y contienen respectivamente a $x$ y a $y$. Al ser complementarios, ambos son también abiertos relativos. Así constituyen la separación requerida. La precisión ``distintos'' es necesaria: un mismo punto no puede pertenecer a dos abiertos disjuntos.\index{separacion total@separación total}

Esta prueba también tendrá una interpretación por coordenadas cuando establezcamos el homeomorfismo $C\cong\{0,1\}^{\N}$ en la proposición \ref{ch16:cantor}. Dos puntos distintos corresponden entonces a sucesiones que difieren en alguna coordenada $k$. Los conjuntos $\{b:b_k=0\}$ y $\{b:b_k=1\}$ son abiertos en la topología producto, pues la proyección de coordenada es continua y los singletons de $\{0,1\}$ son abiertos. Son complementarios y no vacíos; transportados a $C$, separan los dos puntos. Si los dígitos aparecen en el orden contrario, se intercambian los conjuntos. Esta segunda descripción no se necesita para la prueba mediante intervalos.
\end{solucion}

\begin{ejercicio}[Lectura 18, ejercicio 4]\label{ej:lectura-18-4}
Si $A\subset X$ es abierto y cerrado, probar que es unión de componentes de $X$.
\end{ejercicio}
\begin{solucion}
Si una componente $C$ toca $A$, los conjuntos $C\cap A$ y $C\cap(X\setminus A)$ son abiertos y cerrados relativos, disjuntos, y cubren $C$. La conexidad de $C$ y el hecho de que $C\cap A$ sea no vacío obligan a que el otro sea vacío. Por tanto $C\subset A$. Cada punto de $A$ pertenece a una componente, y acabamos de probar que esa componente entera está en $A$. Así $A$ es la unión de las componentes que toca. No se afirma que toda unión de componentes sea abierta: en $\Q$ un singleton es una unión de componentes pero no es abierto.
\end{solucion}

\begin{ejercicio}[Lectura 18, ejercicio 5]\label{ej:lectura-18-5}
Probar que $A=\{(x,x\sin(1/x)):x\ne0\}\cup\{(0,0)\}$ es arcoconexo.
\end{ejercicio}
\begin{solucion}
Definimos $h:\R\to\R$ por $h(x)=x\sin(1/x)$ para $x\ne0$ y $h(0)=0$. Fuera del cero es continua, y en cero la estimación $|h(x)|\le|x|$ prueba continuidad. El mapa $F(x)=(x,h(x))$ es continuo e inyectivo, y su inversa sobre $A$ es la primera proyección, que es continua. Por tanto $A$ es homeomorfo a $\R$. Explícitamente, para dos puntos distintos $F(a),F(b)$, el mapa
\[
\gamma(t)=F((1-t)a+tb),\qquad0\le t\le1,
\]
es continuo e inyectivo porque $a\ne b$ y la primera coordenada es estrictamente monótona. Tiene los extremos requeridos y su imagen está en $A$. El factor $x$ amortigua las oscilaciones y permite atravesar el origen, a diferencia del ejemplo \ref{ch18:seno}.
\end{solucion}

\begin{ejercicio}[Lectura 18, ejercicio 6]\label{ej:lectura-18-6}
Sean $X$ conexo, $A\subset X$ conexo y $B$ una componente de $X\setminus A$. Probar que $X\setminus B$ es conexo.
\end{ejercicio}
\begin{solucion}
Si $A=\varnothing$, la única componente del espacio no vacío conexo $X$ es $X$, así que $X\setminus B=\varnothing$ es conexo. Supongamos $A\ne\varnothing$ y que $Y=X\setminus B$ admite una separación $U,V$. Puesto que $A\subset Y$ es conexo, debe estar enteramente en uno de sus miembros; renombramos para que $A\subset U$. El conjunto $V$ es no vacío, abierto y cerrado en $Y=X\setminus B$.

Aplicamos el ejercicio \ref{ej:lectura-17-6} al espacio conexo $X$, al conjunto conexo $B$ y al conjunto $V$, que es abierto y cerrado en su complemento. Obtenemos que $B\cup V$ es conexo. Como $V\cap A=\varnothing$ y $B\subset X\setminus A$, esta unión es un subconjunto conexo de $X\setminus A$. Contiene estrictamente a $B$, porque $V$ es no vacío y disjunto de $B$. Esto contradice la maximalidad de $B$ como componente de $X\setminus A$. Por tanto $Y$ no admite separación.
\end{solucion}

\begin{ejercicio}[Lectura 18, ejercicio 7: convención necesaria]\label{ej:lectura-18-7}
Definamos $x\approx y$ cuando $x=y$ o existe un arco de extremos $x,y$. Probar que es una relación de equivalencia y que sus clases, las componentes arcoconexas, están contenidas en componentes conexas.
\end{ejercicio}
\begin{solucion}
Sin la alternativa $x=y$, la relación escrita con arcos inyectivos no sería reflexiva: una aplicación inyectiva de $[0,1]$ no puede tener los mismos valores en $0$ y $1$. Con la convención corregida, la reflexividad forma parte de la definición. Para simetría, si $p$ es un arco de $x$ a $y$, $t\mapsto p(1-t)$ es un arco de $y$ a $x$. Para transitividad, los casos con puntos iguales son inmediatos por sustitución; si hay arcos de $x$ a $y$ y de $y$ a $z$, el lema \ref{ch18:recorte} produce un arco de $x$ a $z$ cuando son distintos, y si coinciden se aplica la primera alternativa. Así la relación es de equivalencia.

Además, si $p$ une $x$ y $y$, todo punto de su imagen está en la clase de $x$, porque la restricción de $p$ al intervalo que va de $0$ al parámetro de ese punto, reparametrizada linealmente, es un arco, salvo el punto inicial que ya está incluido por igualdad. En consecuencia, las clases son arcoconexas como subespacios, no solo pares de puntos conectables mediante arcos que podrían salir de la clase. Por la proposición \ref{ch18:arcoimplica}, cada clase es conexa; la proposición \ref{ch18:componentes} la sitúa dentro de una componente conexa. También se obtiene directamente de que la imagen de cualquier arco es un conjunto conexo que contiene sus dos extremos.\index{componente arcoconexa}
\end{solucion}

\section{Separar, descomponer y unir por curvas}
La conexidad impide las separaciones y se transmite por continuidad; las componentes son las piezas conexas maximales y siempre son cerradas. La conexión por curvas exige un argumento adicional; la curva del seno muestra por qué. En abiertos de espacios normados, la geometría local de las bolas permite construir conexiones y recortar sus repeticiones para obtener arcos. Conexidad, conexión por caminos e inyectividad no deben confundirse en las demostraciones.

\chapter{Compactificaciones y espacios universales}\label{cap:nuevo07}\label{chap:16}

\section{Añadir límites conservando la topología}\label{lectura:16}
Compactificar significa situar un espacio dentro de otro compacto, añadiendo puntos que pueden representar comportamientos límite. La inclusión debe conservar exactamente la topología original: una inyección continua por sí sola puede cambiarla. Desarrollamos la lectura 16 de \cite{valencia}, corrigiendo esa distinción y resolviendo sus ejercicios, incluso cuando el enunciado original requiere modificación. Necesitaremos productos numerables, compacidad y separabilidad, junto con la conexidad y los caminos estudiados en el capítulo \ref{cap:nuevo06}. El ejercicio \ref{ej:lectura-16-2} incluye una introducción autónoma al índice de un lazo para justificar las diferencias topológicas que solicita.

\section{La definición y dos modelos compactos}
\begin{definicion}
Una \emph{compactificación métrica}\index{compactificacion@compactificación} de $X$ es un par $(K,j)$, donde $K$ es compacto métrico y $j:X\to K$ es un \emph{encaje}: un homeomorfismo de $X$ sobre $j(X)$, con $j(X)$ denso en $K$.
\end{definicion}
\begin{observacion}[Por qué exigir un encaje]
La definición original solo exige que $j$ sea continua e inyectiva. Eso no permite deducir que $X$ sea separable. Por ejemplo, tomemos $\R$ con la métrica discreta y una biyección $j:\R\to(0,1)$. Toda función desde un espacio discreto es continua, de modo que $j$ es una inyección continua de imagen densa en $[0,1]$, aunque el dominio discreto no numerable no es separable. Su inversa no es continua. Con la definición de encaje, el argumento correcto sí funciona: un compacto métrico es segundo numerable, sus subespacios son segundo numerables y todo espacio segundo numerable es separable, eligiendo un punto en cada miembro no vacío de una base numerable. Por tanto todo espacio con compactificación métrica debe ser separable.
\end{observacion}

\begin{definicion}
El \emph{conjunto de Cantor}\index{conjunto de Cantor} es
\[
C=\left\{\sum_{n=1}^\infty\frac{2b_n}{3^n}:b_n\in\{0,1\}\right\}\subset[0,1].
\]
El \emph{cubo de Hilbert}\index{cubo de Hilbert} es $H=[0,1]^{\N}$ con la métrica
\[
D(x,y)=\sum_{n=1}^\infty2^{-n}\frac{|x_n-y_n|}{1+|x_n-y_n|}.
\]
\end{definicion}
El uso del nombre cubo de Hilbert para el compacto $K\subset\ell^2$ del ejercicio \ref{ej:lectura-14-5} es compatible: el ejercicio \ref{ej:lectura-16-4} construye un homeomorfismo entre ambos modelos.

\begin{proposicion}\label{ch16:cantor}
$C$ y $H$ son compactos. El mapa
\[
\theta:\{0,1\}^{\N}\longrightarrow C,\qquad
\theta(b)=\sum_{n=1}^\infty2b_n3^{-n},
\]
es un homeomorfismo. Por tanto el mapa $\theta(b)\mapsto b$ encaja $C$ en $H$ como el conjunto de sus vértices.
\end{proposicion}
\begin{proof}
Los espacios $\{0,1\}^{\N}$ y $H$ son compactos por el ejercicio \ref{ej:lectura-14-1}. Si dos sucesiones coinciden hasta la coordenada $N$, sus imágenes por $\theta$ difieren a lo sumo en $\sum_{n>N}2/3^n=3^{-N}$, lo que prueba continuidad en la topología producto. El mapa es sobreyectivo por la definición de $C$. Si el primer dígito distinto está en la posición $k$, la diferencia de los términos de esa posición es $2/3^k$, mientras que toda la cola puede compensar a lo sumo $1/3^k$; las imágenes difieren por lo menos en $1/3^k$. Así $\theta$ es inyectivo. Una biyección continua de un compacto a un espacio métrico es un homeomorfismo, por el ejercicio \ref{ej:lectura-13-5}. En particular, su imagen $C$ es compacta.

Para relacionar la fórmula con la construcción por tercios, un prefijo $b_1,\ldots,b_N$ determina el intervalo
\[
\left[\sum_{n=1}^N\frac{2b_n}{3^n},\ 
\sum_{n=1}^N\frac{2b_n}{3^n}+3^{-N}\right].
\]
La unión de estos $2^N$ intervalos se obtiene eliminando los tercios centrales durante $N$ etapas. Los intervalos de cada etapa son disjuntos, y cada punto de todas las etapas determina prefijos compatibles; su distancia a la suma correspondiente es a lo sumo $3^{-N}$ para todo $N$. Por tanto la intersección de las etapas coincide exactamente con $C$.
\end{proof}

\section{El encaje universal del espacio separable}
\begin{teorema}\label{ch16:encaje}
Todo espacio métrico separable se encaja en $H$. En consecuencia, un espacio métrico admite una compactificación métrica si y solo si es separable.
\end{teorema}
\begin{proof}
Si $X=\varnothing$, la aplicación vacía $X\to H$ es un encaje, y la clausura vacía de su imagen es una compactificación de $X$. Para $X\ne\varnothing$, sustituimos $d$ por $\rho=d/(1+d)$, que es una métrica equivalente, acotada por $1$. Su desigualdad triangular se demostró en el ejercicio \ref{ej:lectura-14-1}; la equivalencia se sigue de que $t/(1+t)$ es continua, estrictamente creciente y tiene inversa $s/(1-s)$ para $0\le s<1$.

Sea $(a_n)$ una sucesión cuyo recorrido es denso; si el conjunto denso es finito, repetimos sus puntos. Definimos
\[
J(x)=(\rho(x,a_n))_{n\ge1}\in H.
\]
Cada coordenada es continua y, de hecho,
\[
D(J(x),J(y))\le\sum_{n\ge1}2^{-n}|\rho(x,a_n)-\rho(y,a_n)|\le\rho(x,y).
\]
Por tanto $J$ es continua. Si $J(x)=J(y)$, dado $\eps>0$ elegimos $a_n$ con $\rho(x,a_n)<\eps/2$. La igualdad de coordenadas implica $\rho(y,a_n)=\rho(x,a_n)$, así que $\rho(x,y)<\eps$. Como $\eps$ es arbitrario, $x=y$.

Falta la continuidad de la inversa, que no se sigue solamente de la inyectividad. Fijemos $x\in X$ y $\eps>0$. Elegimos $a_n$ con $\rho(x,a_n)<\eps/4$. Si $J(y)$ está en una vecindad de $J(x)$ que haga variar la coordenada $n$ menos que $\eps/4$, entonces $\rho(y,a_n)<\eps/2$ y
\[
\rho(y,x)\le\rho(y,a_n)+\rho(a_n,x)<3\eps/4<\eps.
\]
Tal vecindad existe porque la evaluación de una coordenada es continua en la métrica producto. Esto prueba que $J^{-1}$ es continua sobre $J(X)$. Finalmente, $K=\cl{J(X)}^{\,H}$ es compacto como cerrado de $H$, y $J(X)$ es denso en $K$. La necesidad de separabilidad se probó después de la definición.
\end{proof}

\begin{ejemplo}[Dos extremos para la recta]
La función
\[
u(x)=\frac12\left(1+\frac{x}{1+|x|}\right)
\]
es un homeomorfismo de $\R$ sobre $(0,1)$. Es continua, estrictamente creciente y su inversa se obtiene poniendo $s=2t-1$ y $x=s/(1-|s|)$. Como $(0,1)$ es denso en $[0,1]$, este último compactifica la recta. Si añadimos los símbolos $-\infty,+\infty$ y ponemos $u(-\infty)=0$, $u(+\infty)=1$, la fórmula $\rho(x,y)=|u(x)-u(y)|$ metriza la recta extendida y la identifica con $[0,1]$. La métrica extendida induce en la recta original su topología usual, aunque no su distancia usual.

Se pueden adoptar convenciones de cálculo con infinitos, como $x+(+\infty)=+\infty$ para $x\in\R$ y $x(+\infty)=+\infty$ para $x>0$. Estas convenciones no convierten la recta extendida en un cuerpo; $+\infty-\infty$ queda sin definir. La convención $0\cdot\infty=0$, útil en ciertos contextos de medida, tampoco es una extensión continua de la multiplicación.
\end{ejemplo}

\begin{ejemplo}[El espacio de Baire dentro de Cantor]
Sea $\mathcal N=\N^{\N}$, con $\N=\{1,2,\ldots\}$, y distancia $0$ entre sucesiones iguales y $1/k$ si el primer desacuerdo aparece en la posición $k$. Codificamos $x=(x_n)$ por la sucesión binaria
\[
\underbrace{0\cdots0}_{x_1-1}1\,
\underbrace{0\cdots0}_{x_2-1}1\,
\underbrace{0\cdots0}_{x_3-1}1\cdots.
\]
Su imagen es el conjunto de sucesiones binarias con infinitos unos. La codificación es biyectiva sobre esa imagen porque las longitudes de los bloques recuperan $x_n-1$. Es continua: fijar los primeros $m$ números fija un prefijo binario de longitud al menos $m$. La inversa es continua: para recuperar los primeros $m$ números basta leer el prefijo hasta el $m$-ésimo uno. La imagen es densa en $\{0,1\}^{\N}$, ya que cualquier prefijo finito puede prolongarse con infinitos unos. Mediante $\theta$, obtenemos una compactificación de $\mathcal N$ homeomorfa a $C$. El ajuste $x_n-1$ evita que la convención $\N\ni1$ prohíba unos consecutivos; si se usan naturales desde cero, se pueden tomar directamente $x_n$ ceros.\index{espacio de Baire}
\end{ejemplo}

\section{Todo compacto métrico no vacío es imagen de Cantor}
\begin{teorema}\label{ch16:imagenCantor}
Un espacio métrico no vacío $X$ es compacto si y solo si existe una aplicación continua sobreyectiva $C\to X$.
\end{teorema}
\begin{proof}
La suficiencia es el teorema de imagen continua. Para la necesidad construiremos una dirección de aproximación para cada sucesión de dígitos, de modo que sus aproximaciones tengan un único límite.

Partimos de $K_{\varnothing}=X$. Cubrimos cada compacto no vacío $K_s$ ya construido por finitos conjuntos compactos no vacíos contenidos en él y de diámetro a lo sumo $2^{-n}$ en el nivel $n$. Esto se consigue intersectando $K_s$ con finitas bolas cerradas de radio $2^{-n-1}$ cuyos centros están en $K_s$ y cuyas bolas abiertas cubren $K_s$. En cada nivel hay solo finitos nodos: sea $m_n$ el mayor número de hijos de un nodo del nivel anterior. Si a un nodo le faltan hijos para llegar a $m_n$, repetimos uno de los existentes. Así los hijos se etiquetan uniformemente por $1,\ldots,m_n$.

Para $a=(a_n)\in\prod_n\{1,\ldots,m_n\}$, la cadena
\[
K_{a_1}\supset K_{a_1a_2}\supset\cdots
\]
está formada por compactos no vacíos. Tiene intersección no vacía por la propiedad de intersección finita dentro de $X$, y la intersección tiene a lo sumo un punto porque los diámetros tienden a cero. Llamemos $P(a)$ a ese único punto. Si dos sucesiones coinciden hasta el nivel $n$, sus imágenes están en un mismo conjunto de diámetro a lo sumo $2^{-n}$; esto prueba continuidad de $P$. Para cualquier $x\in X$, elegimos sucesivamente un hijo que contenga $x$, lo cual es posible porque los hijos cubren al padre. Se obtiene $a$ con $P(a)=x$, de modo que $P$ es sobreyectiva.

Falta obtener este producto finito a partir de dígitos binarios. Elegimos enteros $l_n\ge1$ con $2^{l_n}\ge m_n$ y una sobreyección $q_n:\{0,1\}^{l_n}\to\{1,\ldots,m_n\}$. Dividimos una sucesión binaria en bloques sucesivos de longitudes $l_1,l_2,\ldots$ y aplicamos $q_n$ al bloque $n$. El mapa resultante $Q$ es continuo porque cada coordenada depende de finitos dígitos, y es sobreyectivo porque podemos elegir por separado una preimagen para cada coordenada. La composición $P\circ Q\circ\theta^{-1}:C\to X$ es la aplicación buscada.
\end{proof}

\begin{observacion}[La aplicación binaria y un salto lógico que hay que evitar]
La fórmula correcta $g(b)=\sum_{n\ge1}b_n/2^n$ lleva continuamente $\{0,1\}^{\N}$ sobre $[0,1]$. La continuidad sigue del control de la cola; la sobreyectividad se obtiene de la expansión binaria: para $t<1$, el algoritmo $b_n=\lfloor2t_{n-1}\rfloor$, $t_n=2t_{n-1}-b_n$, comenzando con $t_0=t$, da $t=\sum_{n\le N}b_n2^{-n}+2^{-N}t_N$, y el resto tiende a cero. Para $t=1$ se toman todos los dígitos iguales a uno. Con denominador $2^{n+1}$ y el índice inicial $n=1$, la imagen solo alcanzaría $[0,1/2]$.

Una biyección entre $\N$ y $\N\times\N$ reordena las coordenadas y prueba que $(\{0,1\}^{\N})^{\N}$ es homeomorfo a $\{0,1\}^{\N}$. La continuidad en ambas direcciones sigue de que cada coordenada se convierte en otra coordenada. Aplicando $g$ fila por fila obtenemos una sobreyección continua $C\to H$. Sin embargo, de $X\subset H$ no se deduce automáticamente una sobreyección $C\to X$: un mapa sobre $H$ no tiene por qué tomar valores en $X$. La construcción del teorema \ref{ch16:imagenCantor} cubre exactamente ese paso.
\end{observacion}

\section{Ejercicios resueltos: compactificaciones y obstrucciones topológicas}
\begin{ejercicio}[Lectura 16, ejercicio 1]\label{ej:lectura-16-1}
Probar que $x\mapsto e^{2i\arctan x}$ encaja $\R$ densamente en $S^1$. Probar que $S^1$ y $[0,1]$ no son homeomorfos.
\end{ejercicio}
\begin{solucion}
La función $2\arctan$ es un homeomorfismo de $\R$ sobre $(-\pi,\pi)$, y $t\mapsto e^{it}$ es un homeomorfismo de ese intervalo sobre $S^1\setminus\{-1\}$. Su inversa es el argumento en $(-\pi,\pi)$; es continuo en la circunferencia cortada porque localmente puede elegirse un intervalo angular de longitud menor que $2\pi$. Así la composición es un encaje. Su imagen es densa: el único punto que falta, $-1$, es límite de $e^{2i\arctan n}$.

Al quitar cualquier punto de $S^1$ queda un espacio homeomorfo a $\R$, por una rotación seguida de la parametrización anterior; por tanto queda conexo. Al quitar $1/2$ de $[0,1]$, quedan los dos abiertos relativos no vacíos $[0,1/2)$ y $(1/2,1]$, que son disjuntos. Un homeomorfismo induciría otro entre los espacios con los puntos correspondientes eliminados, y preservaría conexidad. Esto es imposible. La conexidad de los intervalos quedó demostrada en el lema \ref{ch17:intervalos} a partir de la propiedad del supremo.
\end{solucion}

\begin{ejercicio}[Lectura 16, ejercicio 2]\label{ej:lectura-16-2}
Probar que $S^2$, el toro $S^1\times S^1$ y $[0,1]^2$ son compactificaciones de $\R^2$ y son dos a dos no homeomorfos.
\end{ejercicio}
\begin{solucion}
\emph{Construcción de las compactificaciones.} Los tres espacios son compactos: la esfera es cerrada y acotada en $\R^3$, y los otros dos son productos finitos de compactos. La proyección estereográfica inversa
\[
\sigma(u,v)=\left(\frac{2u}{1+u^2+v^2},
\frac{2v}{1+u^2+v^2},\frac{u^2+v^2-1}{1+u^2+v^2}\right)
\]
es un homeomorfismo de $\R^2$ sobre $S^2\setminus\{(0,0,1)\}$. Sustituir la fórmula verifica que la suma de cuadrados es uno, y la inversa es $(x,y,z)\mapsto(x/(1-z),y/(1-z))$. Ambos mapas son continuos en sus dominios. El punto norte es límite de $\sigma(n,0)$, de modo que la imagen es densa. El producto de dos copias del encaje del ejercicio \ref{ej:lectura-16-1} encaja $\R^2$ densamente en el toro. El producto del mapa $u:\R\to(0,1)$ del ejemplo sobre la recta extendida encaja $\R^2$ densamente en $[0,1]^2$. La densidad en los productos se verifica en rectángulos abiertos, que forman una base.

El producto $S^1\times S^1$ está naturalmente en $\R^4$, no literalmente en $\R^3$. Su modelo usual en $\R^3$ es la imagen de
\[
(e^{i\alpha},e^{i\beta})\longmapsto
((2+\cos\beta)\cos\alpha,(2+\cos\beta)\sin\alpha,\sin\beta).
\]
Es continuo e inyectivo: el radio cilíndrico $\rho=\sqrt{x^2+y^2}$ determina $\cos\beta=\rho-2$, la tercera coordenada determina $\sin\beta$, y $(x/\rho,y/\rho)$ determina $e^{i\alpha}$. Como el dominio es compacto, es un homeomorfismo sobre su imagen.

\emph{Herramienta elemental para distinguirlos.} Un \emph{lazo} en un espacio es una aplicación continua $\gamma:[0,1]\to X$ con $\gamma(0)=\gamma(1)$. Decimos que se contrae si existe un mapa continuo $F:[0,1]^2\to X$ con $F(t,0)=\gamma(t)$, $F(t,1)$ constante en $t$ y $F(0,s)=F(1,s)$ para todo $s$. Esta propiedad es preservada por homeomorfismos. Demostraremos los hechos que necesitamos, sin suponer teoría del grupo fundamental. Usaremos la conexidad del intervalo, demostrada mediante el supremo en el lema \ref{ch17:intervalos}, y las nociones de caminos del capítulo \ref{cap:nuevo06}.\index{lazo}\index{homotopia@homotopía}

Todo camino $\gamma:[0,1]\to S^1$ admite un argumento continuo $\vartheta:[0,1]\to\R$ con $\gamma(t)=e^{i\vartheta(t)}$. Para construirlo, por continuidad uniforme elegimos una partición tal que, en cada subintervalo $[t_j,t_{j+1}]$, $|\gamma(t)-\gamma(t_j)|<1$. El cociente $\gamma(t)/\gamma(t_j)$ evita $-1$ y tiene un argumento principal continuo en $(-\pi,\pi)$. Elegida $\vartheta(0)$, definimos en ese subintervalo
\[
\vartheta(t)=\vartheta(t_j)+\operatorname{Arg}(\gamma(t)/\gamma(t_j)).
\]
Las definiciones se pegan en los extremos. Dos argumentos continuos difieren en un múltiplo constante de $2\pi$, pues su diferencia es continua y toma valores en el conjunto discreto $2\pi\Z$.

Para un lazo, el entero
\[
\operatorname{ind}(\gamma)=\frac{\vartheta(1)-\vartheta(0)}{2\pi}
\]
no depende del argumento elegido. Si dos lazos $\alpha,\beta$ satisfacen $\sup_t|\alpha(t)-\beta(t)|<2$, entonces $q(t)=\operatorname{Arg}(\beta(t)/\alpha(t))$ es continuo y $q(0)=q(1)$. Un argumento de $\beta$ se obtiene sumando $q$ a uno de $\alpha$, así que tienen el mismo índice. En una deformación continua $F(t,s)$, la continuidad uniforme permite dividir el intervalo de $s$ en finitos pasos tales que lazos sucesivos estén uniformemente a distancia menor que $1$. Por lo anterior el índice es constante durante la deformación. El lazo $e^{2\pi it}$ tiene índice $1$; un lazo constante tiene índice $0$. En consecuencia, el primero no se contrae en $S^1$.\index{indice de un lazo@índice de un lazo}

\emph{El toro tiene un lazo que no se contrae.} Consideremos $\gamma(t)=(e^{2\pi it},1)$. Si se contrajera en el toro, al componer la contracción con la primera proyección obtendríamos una contracción de $e^{2\pi it}$ en $S^1$, que acabamos de excluir.

\emph{Todos los lazos de la esfera se contraen.} Este paso necesita justificación: una curva continua puede llenar una región, así que no debemos suponer sin prueba que omite algún punto de la esfera. Sea $\gamma$ un lazo en $S^2$. Elegimos una partición tan fina que en cada intervalo $[t_j,t_{j+1}]$ se cumpla $\norm{\gamma(t)-\gamma(t_j)}<1/4$. Unimos los valores consecutivos $a_j=\gamma(t_j)$ y $a_{j+1}$ mediante el arco
\[
\beta(t)=\frac{(1-r)a_j+ra_{j+1}}{\norm{(1-r)a_j+ra_{j+1}}},\qquad
r=\frac{t-t_j}{t_{j+1}-t_j}.
\]
El denominador es positivo porque los extremos están a distancia menor que $1/4$, de modo que no son antipodales. Si $v=(1-r)a_j+ra_{j+1}$, entonces $\norm{v-a_j}<1/4$ y
\[
\norm{v/\norm v-a_j}\le |1-\norm v|+\norm{v-a_j}<1/2.
\]
Por tanto $\norm{\beta(t)-\gamma(t)}<3/4$. Los dos valores nunca son antipodales, y
\[
F(t,s)=\frac{(1-s)\gamma(t)+s\beta(t)}{\norm{(1-s)\gamma(t)+s\beta(t)}}
\]
define una deformación continua de lazos entre $\gamma$ y $\beta$.

La imagen de $\beta$ está en una unión finita de planos propios que pasan por el origen, intersectados con $S^2$, porque cada arco está en el plano generado por sus extremos; si ambos son iguales elegimos cualquier plano que contenga ese punto. Hay un punto $q\in S^2$ fuera de todos esos planos. Para verlo sin un argumento de medida, escribamos sus ecuaciones como $A_jx+B_jy+C_jz=0$, con algún coeficiente no cero. El vector $(1,t,t^2)$ pertenece al plano $j$ solo si $A_j+B_jt+C_jt^2=0$, ecuación con finitas raíces. Elegimos $t$ fuera de la unión finita de raíces y normalizamos el vector. Entonces $q$ no pertenece a la imagen de $\beta$.

Una rotación seguida de la proyección estereográfica identifica $S^2\setminus\{q\}$ con $\R^2$. Allí todo lazo se contrae por interpolación lineal con un punto fijo. Transportamos esa contracción a la esfera y la concatenamos con $F$, usando la primera mitad del parámetro para $F$ y la segunda para la contracción. Así $\gamma$ se contrae. La esfera y el toro no son homeomorfos.

\emph{El cuadrado y el toro tampoco son homeomorfos.} Todo lazo del cuadrado se contrae mediante $F(t,s)=(1-s)\gamma(t)+s(1/2,1/2)$, pues el cuadrado es convexo; el toro tiene el lazo no contraíble ya construido.

\emph{La esfera y el cuadrado no son homeomorfos.} Al quitar cualquier punto de la esfera queda un espacio homeomorfo a $\R^2$, donde todos los lazos se contraen. En cambio, al quitar $a=(1/2,1/2)$ del cuadrado, el lazo
\[
\gamma(t)=a+\tfrac14(\cos2\pi t,\sin2\pi t)
\]
no se contrae: componer una supuesta contracción con $x\mapsto(x-a)/\norm{x-a}$ daría una contracción del lazo de índice uno en $S^1$. Si existiera un homeomorfismo de la esfera al cuadrado, al eliminar la preimagen de $a$ produciría un homeomorfismo entre estos dos espacios perforados. La contradicción termina las tres comparaciones.
\end{solucion}

\begin{ejercicio}[Lectura 16, ejercicio 3: examen y corrección del enunciado]\label{ej:lectura-16-3}
El enunciado propone que el conjunto de Cantor sea una compactificación de todo espacio discreto infinito separable. Decidir si es válido con la definición de encaje y construir una compactificación correcta.
\end{ejercicio}
\begin{solucion}
Un espacio discreto separable es numerable, porque todo subconjunto denso debe contener cada singleton abierto. Sin embargo, \emph{no puede encajarse densamente en Cantor}. En efecto, si $D$ es un subespacio denso y discreto de un espacio métrico $K$, para cada $d\in D$ existe un abierto $U\subset K$ con $U\cap D=\{d\}$. Si $U$ contuviera un punto $y\ne d$, el abierto no vacío $U\setminus\{d\}$ tendría que tocar $D$ por densidad, contradicción. Así $U=\{d\}$ y $d$ es un punto aislado de $K$. Pero Cantor no tiene puntos aislados: manteniendo un prefijo binario y cambiando algún dígito posterior obtenemos puntos distintos arbitrariamente próximos.

La sugerencia de enumerar un subconjunto denso de vértices proporciona una inyección continua desde el discreto numerable, pero su inversa sobre la imagen no es continua. Esta es precisamente la dificultad de la definición original. Una versión correcta es la siguiente: si $X=\{x_1,x_2,\ldots\}$ es discreto, el mapa $x_n\mapsto1/n$ lo encaja densamente en
\[
K=\{0\}\cup\{1/n:n\ge1\}.
\]
Cada $1/n$ es aislado en $K$, de modo que la imagen tiene la topología discreta; $0$ es su único punto límite añadido. El conjunto $K$ es cerrado y acotado en $\R$, luego compacto.

La afirmación adicional de que todo espacio ultramétrico tiene a Cantor como compactificación tampoco es válida en esa generalidad: los discretos anteriores son ultramétricos, y los no separables ni siquiera admiten compactificación métrica. Sí puede afirmarse, por el teorema \ref{ch16:encaje}, que todo espacio ultramétrico \emph{separable} admite alguna compactificación métrica.
\end{solucion}

\begin{ejercicio}[Lectura 16, ejercicio 4]\label{ej:lectura-16-4}
Probar que $T:K\to H$, $T(x)_n=nx_n$, es un homeomorfismo, donde $K=\{x\in\ell^2:0\le x_n\le1/n\}$.
\end{ejercicio}
\begin{solucion}
El mapa toma valores en $H$ y es inyectivo. Para $y\in H$, la sucesión $x_n=y_n/n$ satisface $\sum_n|x_n|^2\le\sum_n1/n^2<\infty$ y pertenece a $K$, de modo que es sobreyectivo. Si $x^{(j)}\to x$ en $\ell^2$, entonces $|x_n^{(j)}-x_n|\le\norm{x^{(j)}-x}_2\to0$ para cada $n$; por tanto $nx_n^{(j)}\to nx_n$ y la convergencia coordenada implica convergencia en $H$. Así $T$ es continua. Como $K$ es compacto y $H$ es métrico, la biyección continua $T$ es un homeomorfismo. Obsérvese que no estamos afirmando continuidad de este multiplicador en todo $\ell^2$ con valores en $\ell^2$: aquí el codominio tiene la topología producto.
\end{solucion}

\begin{ejercicio}[Lectura 16, ejercicio 5]\label{ej:lectura-16-5}
Probar que todo espacio métrico separable es homeomorfo a un subespacio del toro infinito $\mathbb T^\infty=(S^1)^{\N}$.
\end{ejercicio}
\begin{solucion}
Por el teorema \ref{ch16:encaje}, basta encajar $H$ en $\mathbb T^\infty$. El mapa $h:[0,1]\to S^1$, $h(t)=e^{i\pi t}$, es continuo e inyectivo; por compacidad es un homeomorfismo sobre el semicírculo superior. Aplicamos $h$ coordenada a coordenada:
\[
\Phi((x_n))=(e^{i\pi x_n})_{n\ge1}.
\]
Este mapa es continuo e inyectivo en las topologías producto. Como $H$ es compacto y el producto numerable de circunferencias es métrico, $\Phi$ es un encaje. La composición $\Phi\circ J$ da el encaje de $X$. Se utiliza un semicírculo, no $e^{2\pi it}$ en todo $[0,1]$, pues este último identificaría los dos extremos.
\end{solucion}

\section{Encajes e imágenes continuas de compactos}
Una compactificación conserva la topología y añade límites; una inyección continua puede hacer algo distinto. El cubo de Hilbert contiene modelos de todos los espacios métricos separables, mientras que Cantor se proyecta continuamente sobre todos los compactos métricos no vacíos. Encajar y proyectar son operaciones diferentes, y distinguirlas evita errores en las demostraciones y en los ejercicios.


\parte{Espacios normados y operadores}{%
La estructura lineal añade a la distancia operaciones, subespacios y transformaciones. Primero se comparan la dimensión finita y la infinita, se estudian las series y se precisa cuándo los subespacios son cerrados. Después, la norma de operador permite controlar simultáneamente todas las imágenes de la bola unitaria.

La acotación uniforme y los operadores compactos conectan esta teoría con Baire y la compacidad. Finalmente, Hahn--Banach proporciona funcionales que detectan vectores y separan conjuntos convexos. Se requieren álgebra lineal elemental y los resultados métricos de las partes anteriores; el dual y los operadores se construyen aquí.}
\chapter{Geometría normada y dimensión finita}\label{cap:nuevo08}\label{chap:19}\label{chap:22}

Un espacio normado reúne dos estructuras: podemos sumar y multiplicar vectores, y también medir distancias. Comenzaremos mostrando que ambas estructuras son compatibles; después estudiaremos series, subespacios cerrados y direcciones que permanecen alejadas de un subespacio propio. Estas herramientas permitirán reconocer qué fenómenos dependen realmente de la dimensión finita.

En $\R^n$ todas las normas describen la misma convergencia, toda transformación lineal es continua y los conjuntos cerrados acotados son compactos. No son propiedades de una norma aislada: el número finito de coordenadas es decisivo. Explicaremos tanto estas afirmaciones como las obstrucciones que aparecen en dimensión infinita.

Trabajaremos sobre $\F=\R$ o $\C$. Se requieren completitud, separabilidad, compacidad y el teorema de Baire. Una dimensión sin calificativo será algebraica: una base de Hamel permite únicamente combinaciones lineales finitas. No se presupone la teoría de operadores del capítulo siguiente; las afirmaciones sobre aplicaciones lineales en dimensión finita se demostrarán directamente mediante coordenadas.

\section{Operaciones, bolas y coordenadas}\label{lectura:19}

\begin{proposicion}\label{ch19:operaciones}
En un espacio normado $V$, la suma $V\times V\to V$ y el producto $\F\times V\to V$ son continuos. La norma es una función de Lipschitz de constante $1$.
\end{proposicion}
\begin{proof}
La desigualdad triangular, aplicada en ambos sentidos, da
\[
 \big|\norm{x}-\norm{y}\big|\leq\norm{x-y}.
\]
Para la suma, fijados $x_0,y_0$, tenemos
\[
 \norm{(x+y)-(x_0+y_0)}\leq\norm{x-x_0}+\norm{y-y_0}.
\]
Si ambos sumandos del lado derecho son menores que $\eps/2$, el resultado es menor que $\eps$. Para el producto, separamos los dos cambios mediante la identidad
\[
 \alpha x-\alpha_0x_0=\alpha(x-x_0)+(\alpha-\alpha_0)x_0.
\]
Si $|\alpha-\alpha_0|<\delta\leq1$, entonces $|\alpha|<|\alpha_0|+1$. Por tanto, si además $\norm{x-x_0}<\delta$, resulta
\[
 \norm{\alpha x-\alpha_0x_0}
 <\bigl(|\alpha_0|+1+\norm{x_0}\bigr)\delta.
\]
Basta elegir $\delta=\min\{1,\eps/(|\alpha_0|+1+\norm{x_0})\}$. La elección depende del punto de continuidad, como permite la definición.
\end{proof}

\begin{corolario}\label{ch19:limites}
Si $x_n\to x$, $y_n\to y$ y $\alpha_n\to\alpha$, entonces $x_n+y_n\to x+y$ y $\alpha_nx_n\to\alpha x$.
\end{corolario}
\begin{proof}
Aplicamos las dos estimaciones anteriores a los pares de términos de las sucesiones. En la segunda, $\alpha_n$ es acotada porque converge, de modo que $|\alpha_n|\norm{x_n-x}+|\alpha_n-\alpha|\norm{x}\to0$.
\end{proof}

\begin{observacion}
En el corolario de la lectura original aparece $\alpha_n+v_n$ donde corresponde $\alpha_nv_n$. En un espacio vectorial general ni siquiera está definida la suma de un escalar con un vector.
\end{observacion}

\begin{lema}[Coordenadas en dimensión finita]\label{ch19:coordenadas}
Si $v_1,\ldots,v_m$ son linealmente independientes en $V$, existen $c,C>0$ tales que
\[
 c\max_i|a_i|\leq\norm{\sum_{i=1}^m a_iv_i}\leq C\max_i|a_i|
 \qquad(a_1,\ldots,a_m\in\F).
\]
En particular, todo subespacio de dimensión finita es completo y cerrado en $V$.
\end{lema}
\begin{proof}
La cota superior vale con $C=\sum_i\norm{v_i}$. Definamos $L(a)=\sum_i a_iv_i$. Esa misma cota da $\norm{L(a)-L(b)}\leq C\norm{a-b}_\infty$, así que $a\mapsto\norm{L(a)}$ es continua. La esfera $S=\{a\in\F^m:\norm a_\infty=1\}$ es compacta por Heine--Borel en el espacio euclidiano $\R^m$ o $\R^{2m}$. La función alcanza allí un mínimo $c$. La independencia implica $L(a)\ne0$ en $S$, y por ello $c>0$. Para $a\ne0$, aplicamos esta cota a $a/\norm a_\infty$ y multiplicamos por $\norm a_\infty$; para $a=0$ la desigualdad es inmediata.

Sea ahora $x_n=L(a^{(n)})$ una sucesión de Cauchy en el subespacio generado. La cota inferior implica que cada sucesión de coordenadas es de Cauchy en $\F$. Si $a^{(n)}\to a$, la cota superior da $x_n\to L(a)$ en el subespacio. Este es completo. Si una sucesión del subespacio converge en $V$, es de Cauchy y el límite que acabamos de obtener pertenece al subespacio; la unicidad del límite demuestra que este es cerrado. El subespacio $\{0\}$ satisface ambas conclusiones directamente.
\end{proof}

Este lema se ha demostrado antes de aplicar Baire: sus consecuencias sobre cerradura y completitud se necesitan ya en las propiedades generales que siguen. La sección~\ref{sec:nuevo08:finita} desarrollará después el papel de la dimensión finita de manera sistemática.

\begin{proposicion}\label{ch19:geometria}
Sea $V$ un espacio normado y $r>0$.
\begin{enumerate}[label=(\roman*)]
\item $\interior{\overline B(x,r)}=B(x,r)$ y $\cl{B(x,r)}=\overline B(x,r)$.
\item $V$ es conexo por caminos; dos puntos distintos pueden unirse por un arco.
\item Si $V\ne\{0\}$, entonces $V$ no es compacto.
\item $V$ es separable si y solo si contiene un subespacio denso de dimensión a lo sumo numerable.
\item Si $V$ es de Banach, su dimensión es finita o infinita no numerable.
\end{enumerate}
\end{proposicion}
\begin{proof}
(i) La continuidad de la norma muestra que la bola abierta es abierta y la cerrada es cerrada. Por ello $B(x,r)\subseteq\interior{\overline B(x,r)}$ y $\cl{B(x,r)}\subseteq\overline B(x,r)$. Si $\norm{y-x}=r$, los puntos $x+(1+s)(y-x)$, con $s>0$ tendiendo a cero, están fuera de la bola cerrada y convergen a $y$. Así, $y$ no es interior. En cambio, $x+(1-s)(y-x)$ pertenece a la bola abierta para $0<s<1$ y converge a $y$, lo que prueba la otra igualdad. Si $V=\{0\}$ no hay puntos de frontera y las igualdades siguen siendo ciertas.

(ii) Dados $u,v$, la función $\gamma(t)=(1-t)u+tv$, $0\leq t\leq1$, es continua por la proposición~\ref{ch19:operaciones}. Une $u$ con $v$. Si $u\ne v$, es inyectiva, pues $\gamma(s)=\gamma(t)$ implica $(s-t)(v-u)=0$ y, por tanto, $s=t$. Como su dominio es compacto y el codominio es Hausdorff, es un homeomorfismo sobre su imagen. La conexión por caminos implica conexión: una separación del espacio induciría una separación del intervalo a través de un camino que una puntos de las dos partes.

(iii) Elija $v\ne0$. La sucesión $nv$ tiene normas $n\norm v$ no acotadas. Por tanto $V$ no es acotado, mientras que todo compacto métrico sí lo es. La excepción $V=\{0\}$ es necesaria: ese espacio es compacto.

(iv) Si $D$ es numerable y denso, $M=\spanop D$ contiene a $D$, es denso y tiene una base a lo sumo numerable. Recíprocamente, supongamos $M=\spanop\{v_1,v_2,\ldots\}$ y $\cl M=V$. Sea $\F_0=\Q$ en el caso real y $\F_0=\Q+i\Q$ en el complejo. El conjunto $D$ de todas las combinaciones finitas de los $v_j$ con coeficientes en $\F_0$ es numerable: es una unión numerable de conjuntos numerables. Si $x=\sum_{j=1}^m a_jv_j$, elegimos $q_j\in\F_0$ tan cercanos a $a_j$ que $\sum_{j=1}^m|q_j-a_j|\norm{v_j}<\eps$. Entonces $\norm{x-\sum q_jv_j}<\eps$. Así $M\subseteq\cl D$, y al tomar cierres obtenemos $V=\cl D$.

(v) Si existiera una base infinita numerable $v_1,v_2,\ldots$, tendríamos $V=\bigcup_{n\geq1}M_n$, con $M_n=\spanop\{v_1,\ldots,v_n\}$. Cada $M_n$ es cerrado por el lema~\ref{ch19:coordenadas}. Tiene interior vacío: si $B(x,r)\subseteq M_n$, entonces $x\in M_n$, y restar $x$ da $B(0,r)\subseteq M_n$; al multiplicar por escalares se obtiene $M_n=V$, contradicción. Baire prohíbe que el espacio completo $V$ sea esa unión de cerrados con interior vacío.
\end{proof}

\index{separabilidad!espacios normados}
\index{base de Hamel}
\begin{observacion}
Separabilidad y dimensión algebraica numerable no son equivalentes. Por ejemplo, un espacio de Banach separable e infinito dimensional puede tener un subespacio denso generado por una sucesión, pero esa sucesión no constituye una base de Hamel de todo el espacio. La diferencia consiste en permitir límites además de sumas finitas.
\end{observacion}

\section{Series y completitud}

\begin{definicion}
Dada una sucesión $(v_n)$ de $V$, la serie $\sum_{n=1}^{\infty}v_n$ converge a $s$ cuando sus sumas parciales $s_N=\sum_{n=1}^N v_n$ convergen a $s$. Se llama \emph{absolutamente convergente} cuando $\sum_{n=1}^{\infty}\norm{v_n}<\infty$.
\end{definicion}
\index{serie!absolutamente convergente}

\begin{teorema}\label{ch19:series}
Un espacio normado $V$ es completo si y solo si toda serie absolutamente convergente en $V$ converge en $V$.
\end{teorema}
\begin{proof}
Si $V$ es completo y $\sum\norm{v_n}<\infty$, para $m>n$ tenemos
\[
 \norm{s_m-s_n}\leq\sum_{j=n+1}^m\norm{v_j}.
\]
Las colas de una serie real convergente de términos no negativos tienden a cero. Por ello $(s_n)$ es de Cauchy, y la completitud proporciona su límite en $V$.

Para el recíproco, sea $(x_n)$ de Cauchy. Elegimos inductivamente enteros $n_1<n_2<\cdots$ tales que
\[
 p,q\geq n_j\quad\Longrightarrow\quad\norm{x_p-x_q}<2^{-j}.
\]
La condición de Cauchy da un umbral para cada $j$; aumentarlo permite exigir además $n_{j+1}>n_j$. Definimos $v_1=x_{n_1}$ y $v_{j+1}=x_{n_{j+1}}-x_{n_j}$. Entonces
\[
 \sum_{j=1}^{\infty}\norm{v_j}
 \leq\norm{x_{n_1}}+\sum_{j=1}^{\infty}2^{-j}<\infty.
\]
Por hipótesis la serie converge a algún $x\in V$. Su suma parcial número $j$ es $x_{n_j}$, porque todos los términos intermedios se cancelan. Tenemos, pues, una subsucesión convergente.

Falta justificar que converge toda la sucesión. Dado $\eps>0$, elija $N$ tal que $\norm{x_p-x_q}<\eps/2$ para $p,q\geq N$. Después elija $j$ con $n_j\geq N$ y $\norm{x_{n_j}-x}<\eps/2$. Para todo $n\geq N$,
\[
 \norm{x_n-x}\leq\norm{x_n-x_{n_j}}+\norm{x_{n_j}-x}<\eps.
\]
Así $V$ es completo.
\end{proof}

La subsucesión anterior no se elige para simplificar una notación: convierte diferencias pequeñas en una serie numérica sumable. Es una técnica que reaparecerá en construcciones por aproximaciones sucesivas.

\section{Subespacios cerrados y lema de Riesz}

\begin{proposicion}\label{ch19:cierre-subespacio}
El cierre de un subespacio vectorial $M\subseteq V$ es un subespacio vectorial.
\end{proposicion}
\begin{proof}
Contiene al cero. Si $x,y\in\cl M$, existen $x_n,y_n\in M$ con $x_n\to x$ e $y_n\to y$. Para $\alpha,\beta\in\F$, cada $\alpha x_n+\beta y_n$ pertenece a $M$ y converge a $\alpha x+\beta y$ por el corolario~\ref{ch19:limites}. Este límite pertenece a $\cl M$.
\end{proof}

\begin{lema}[Riesz]\label{ch19:riesz}
Si $M$ es un subespacio propio cerrado de $V$, para todo $0<\theta<1$ existe $v\in V$ con $\norm v=1$ y $\dist(v,M)>\theta$. En particular, para todo $\eps>0$ puede conseguirse $\dist(v,M)\geq1-\eps$.
\end{lema}
\index{Riesz, lema de}
\begin{proof}
Tome $u\notin M$ y $d=\dist(u,M)$. El número $d$ es positivo porque $M$ es cerrado: si fuera cero, habría puntos de $M$ que convergen a $u$. Como $d/\theta>d$, la definición de ínfimo permite elegir $w\in M$ tal que $\norm{u-w}<d/\theta$. Sea $v=(u-w)/\norm{u-w}$. Para cualquier $z\in M$, el vector $w+\norm{u-w}z$ también pertenece a $M$, y por ello
\[
 \norm{v-z}
 =\frac{\norm{u-(w+\norm{u-w}z)}}{\norm{u-w}}
 \geq\frac d{\norm{u-w}}>\theta.
\]
La cota inferior es independiente de $z$, de modo que también vale para el ínfimo. Si $0<\eps<1$, tome $\theta=1-\eps$; si $\eps\geq1$, cualquier vector unitario cumple la desigualdad solicitada. Tal vector existe porque $M$ es propio.
\end{proof}

Las operaciones continuas y la caracterización de completitud mediante series describen cualquier espacio normado. El lema de Riesz añade una restricción geométrica: un subespacio cerrado propio no puede aproximar uniformemente todas las direcciones unitarias. Combinaremos ahora esa restricción con las cotas de coordenadas para caracterizar la dimensión finita.

\section{Identificar vectores: espacios cocientes}\label{ampl:cocientes}
Antes de volver a la dimensión finita estudiaremos otra operación con subespacios: declarar equivalentes dos vectores cuya diferencia pertenece a un subespacio dado. Esta construcción permite ignorar direcciones que no interesa distinguir. También resuelve una dificultad ya encontrada: una integral puede medir cero para una función que no es idénticamente nula. En ese caso hay que identificar primero los objetos que la seminorma no distingue.

La organización de esta ampliación se apoya en los apuntes y la presentación de Maximenko \cite{max-cociente,max-seminormado}. Desarrollamos las verificaciones algebraicas y los pasos con ínfimos; los ejercicios indican cuáles adaptan propuestas de esa fuente. No se requiere todavía la teoría general de operadores.
\index{espacio!cociente}\index{norma!cociente}\index{seminorma!nucleo@núcleo}

\begin{definicion}
Si $M$ es un subespacio de un espacio vectorial $V$, escribimos $x\sim y$ cuando $x-y\in M$. La clase de $x$ es $[x]=x+M$, y el conjunto de clases se denota $V/M$. Sus operaciones son
\[
 [x]+[y]=[x+y],\qquad \lambda[x]=[\lambda x].
\]
\end{definicion}
La relación es reflexiva porque $0\in M$, simétrica porque $-M=M$ y transitiva porque $M+M=M$. Las operaciones no dependen de los representantes: si $x-x'$, $y-y'\in M$, entonces $(x+y)-(x'+y')\in M$ y $\lambda x-\lambda x'\in M$. Las identidades vectoriales se heredan de $V$. El vector cero del cociente es la clase $[0]=M$, que generalmente contiene muchos vectores de $V$. La aplicación $q:V\to V/M$, $q(x)=[x]$, es lineal, sobreyectiva y tiene núcleo $M$.

\begin{proposicion}[Seminorma del cociente]\label{ampl:cociente-seminorma}
Sea $p$ una seminorma en $V$ y $M\subset V$ un subespacio. La expresión
\[
 \bar p([x])=\inf_{m\in M}p(x-m)
\]
define una seminorma en $V/M$. Es una norma si y solo si $M$ es cerrado para la pseudométrica $d_p(x,y)=p(x-y)$. En particular, si $V$ es normado y $M$ es cerrado,
\[
 \norm{[x]}_{V/M}=\dist(x,M)
\]
es una norma. La proyección $q$ satisface $\bar p(qx-qy)\le p(x-y)$.
\end{proposicion}
\begin{proof}
Los representantes de $[x]$ son exactamente $x-m$, $m\in M$; cambiar $x$ por otro representante no cambia ese conjunto. El ínfimo es finito y no negativo porque $m=0$ es una elección permitida. Para $\eps>0$ elegimos $m,n\in M$ tales que
\[
 p(x-m)<\bar p([x])+\eps/2,\qquad
 p(y-n)<\bar p([y])+\eps/2.
\]
Entonces $m+n\in M$ y
\[
 \bar p([x+y])\le p(x+y-m-n)
 <\bar p([x])+\bar p([y])+\eps.
\]
Dejar $\eps\downarrow0$ demuestra subaditividad. Para $\lambda\ne0$, al escribir $m=\lambda n$ obtenemos
\[
 \bar p([\lambda x])=\inf_{n\in M}|\lambda|p(x-n)
 =|\lambda|\bar p([x]).
\]
Para $\lambda=0$ ambos miembros son cero. Además $\bar p([x])=0$ equivale a que para todo $\eps>0$ exista $m\in M$ con $p(x-m)<\eps$, esto es, a que $x\in\cl M$ en la topología de $d_p$. Por tanto la única clase de seminorma cero es $M$ exactamente cuando $\cl M=M$. Finalmente, elegir $m=0$ en el ínfimo para $[x-y]$ prueba la última desigualdad.
\end{proof}

\begin{observacion}[El ínfimo no exige un representante óptimo]
La definición permite encontrar representantes de norma menor que $\norm{[x]}_{V/M}+\eps$ para cada $\eps>0$. No afirma que alguno tenga norma exactamente igual al ínfimo. Esta distinción será esencial en la prueba de completitud y en el ejercicio~\ref{ej:ampl-cocientes-5}.
\end{observacion}

\begin{proposicion}[Eliminar el núcleo de una seminorma]\label{ampl:nucleo-seminorma}
Sea $N_p=\{x:p(x)=0\}$. Es un subespacio cerrado para $d_p$, y
\[
 \norm{[x]}_{V/N_p}=p(x)
\]
está bien definida y es una norma.
\end{proposicion}
\begin{proof}
Si $u,v\in N_p$, entonces $p(\alpha u+\beta v)\le|\alpha|p(u)+|\beta|p(v)=0$; así $N_p$ es subespacio. La desigualdad $|p(x)-p(y)|\le p(x-y)$ prueba continuidad de $p$ respecto de $d_p$, por lo que $N_p=p^{-1}(\{0\})$ es cerrado. Si $n\in N_p$, las desigualdades triangulares en ambos sentidos dan $p(x+n)=p(x)$. Todos los representantes de una clase tienen, por tanto, el mismo valor de $p$, y su ínfimo es ese valor. La proposición anterior demuestra que la seminorma resultante es una norma.
\end{proof}

\begin{ejemplo}[Una dirección invisible]
En $\R^2$, $p(s,t)=|s|$ no distingue puntos que difieren verticalmente. Su núcleo es $\{0\}\times\R$. La aplicación $[(s,t)]\mapsto s$ identifica isométricamente el cociente con $\R$. La coordenada vertical se pierde de manera deliberada; no se está afirmando que dos pares originalmente distintos sean iguales en $\R^2$.
\end{ejemplo}

\begin{teorema}[Completitud del cociente]\label{ampl:cociente-banach}
Si $V$ es Banach y $M$ es un subespacio cerrado, $V/M$ es Banach con la norma cociente.
\end{teorema}
\begin{estrategia}
Una sucesión de clases no viene con representantes de Cauchy elegidos de antemano. Es más cómodo usar el criterio de series del teorema~\ref{ch19:series}: levantaremos una serie absolutamente convergente a otra en $V$, permitiendo un error sumable al elegir cada representante.
\end{estrategia}
\begin{proof}
Sea $\sum_{n\ge1}A_n$ una serie en $V/M$ con $\sum_n\norm{A_n}_{V/M}<\infty$. Por la propiedad del ínfimo elegimos $x_n\in A_n$ tales que
\[
 \norm{x_n}<\norm{A_n}_{V/M}+2^{-n}.
\]
Así $\sum_n\norm{x_n}<\infty$, y la completitud de $V$ da $x=\sum_nx_n\in V$. Como $q$ es lineal y $1$-Lipschitz,
\[
 \left\|\sum_{n=1}^N A_n-q(x)\right\|_{V/M}
 \le\left\|\sum_{n=1}^N x_n-x\right\|\longrightarrow0.
\]
Toda serie absolutamente convergente del cociente converge. El mismo criterio de series prueba su completitud.
\end{proof}

\begin{proposicion}[Qué ocurre con las bolas]\label{ampl:cociente-bolas}
Para $V$ normado, $M$ cerrado y $r>0$,
\[
 q(B(x,r))=B(qx,r).
\]
En particular, $q$ envía abiertos en abiertos.
\end{proposicion}
\begin{proof}
Si $y\in B(x,r)$, entonces $\norm{qy-qx}\le\norm{y-x}<r$. Recíprocamente, si $\norm{[z]-[x]}<r$, existe $m\in M$ tal que $\norm{z-x-m}<r$. El vector $z-m$ pertenece a $B(x,r)$ y representa la clase $[z]$. Para un abierto $U$ y una clase $qx\in q(U)$, con $x\in U$, existe $r>0$ tal que $B(x,r)\subset U$. La igualdad de bolas da una bola alrededor de $qx$ contenida en $q(U)$; así $q(U)$ es abierto.
\end{proof}

\section{Ejercicios sobre cocientes y seminormas}

\begin{ejercicio}[Ampliación: una recta como cociente]\label{ej:ampl-cocientes-1}
En $\R^2$ con la norma euclidiana sea $M=\{(t,t):t\in\R\}$. Calcule la norma de $[(a,b)]$, encuentre el representante de norma mínima y construya una isometría lineal de $\R^2/M$ sobre $\R$. Este ejercicio y su desarrollo son propios de esta edición.
\end{ejercicio}


\begin{ejercicio}[Ampliación: el límite como coordenada del cociente]\label{ej:ampl-cocientes-2}
Sean $c$ y $c_0$ los espacios de sucesiones escalares convergentes y convergentes a cero, con la norma uniforme. Demuestre que $c_0$ es cerrado en $c$ y que $[x]\mapsto\lim_nx_n$ es un isomorfismo isométrico de $c/c_0$ sobre $\F$. Adaptación del ejercicio de la presentación \cite{max-seminormado}; la solución se desarrolla aquí.
\end{ejercicio}


\begin{ejercicio}[Ampliación: un cociente que pierde toda distancia]\label{ej:ampl-cocientes-3}
En $\ell^2$ sea $c_{00}$ el subespacio de sucesiones de soporte finito. Pruebe que $\ell^2/c_{00}$ no es el espacio vectorial cero, aunque su seminorma cociente vale cero en todas las clases. Explique por qué esto no contradice el teorema~\ref{ampl:cociente-banach}. Adaptación del contraejemplo propuesto en \cite{max-seminormado}.
\end{ejercicio}


\begin{ejercicio}[Ampliación: observar sólo los extremos]\label{ej:ampl-cocientes-4}
En $C([0,1],\F)$ defina $p(f)=|f(0)|+|f(1)|$. Describa su núcleo y pruebe que el cociente por ese núcleo es isométrico a $\F^2$ con la norma $\norm{(a,b)}_1=|a|+|b|$. Determine si el cociente es completo aunque $p$ no sea una norma. Ejercicio propio de esta edición.
\end{ejercicio}


\begin{ejercicio}[Ampliación: norma cociente sin mínimo]\label{ej:ampl-cocientes-5}
Sea $\varphi:V\to\F$ lineal y no nula, y suponga
$C=\sup_{\norm v\le1}|\varphi(v)|<\infty$. Pruebe que $M=\ker\varphi$ es cerrado y
\[
 \norm{[x]}_{V/M}=\frac{|\varphi(x)|}{C}.
\]
Después considere $V=\ell^1$ y $\varphi(x)=\sum_{n\ge1}\frac{n}{n+1}x_n$. Demuestre que $C=1$ y que la clase cuyos representantes satisfacen $\varphi(x)=1$ tiene norma $1$, pero no contiene ningún vector de norma $1$. Ejercicio propio de esta edición.
\end{ejercicio}


Los cocientes permiten medir distancia a un subespacio y las seminormas obligan a precisar qué vectores se identifican. Volvemos ahora a la dimensión finita: allí las coordenadas proporcionan controles que pueden perderse en los espacios de los últimos ejemplos.

\section{Subespacios de dimensión finita y continuidad}\label{lectura:22}\label{sec:nuevo08:finita}

\begin{proposicion}\label{ch22:suma-cerrada}
Si $M$ es un subespacio cerrado de un espacio normado $V$ y $F$ es un subespacio de dimensión finita, entonces $M+F$ es cerrado.
\end{proposicion}
\begin{proof}
Empecemos por $F=\F u$. Si $u\in M$, la suma es $M$ y no hay nada que probar. Si $u\notin M$, sea $d=\dist(u,M)>0$. Para cualquier escalar $t$, la invariancia de $M$ bajo multiplicación da
\[
 \dist(tu,M)=|t|d;
\]
para $t=0$ ambos lados son cero, y para $t\ne0$ se obtiene cambiando $m$ por $m/t$ en el ínfimo.

Sea $z_n=m_n+t_nu$ una sucesión de $M+\F u$ que converge a $z$ en $V$. Como $m_n-m_k\in M$,
\[
 |t_n-t_k|d\leq\norm{z_n-z_k}.
\]
La convergencia de $(z_n)$ implica que es de Cauchy, de modo que $(t_n)$ es de Cauchy en $\F$ y converge a algún $t$. Entonces $m_n=z_n-t_nu\to z-tu$. Por cerradura de $M$, $z-tu\in M$, y así $z\in M+\F u$. Esto prueba que la suma es cerrada.

Para $F=\spanop\{u_1,\ldots,u_k\}$, aplique el caso anterior sucesivamente a
\[
 M,\quad M+\F u_1,\quad M+\F u_1+\F u_2,\quad\ldots.
\]
Cada paso añade una dirección a un subespacio ya cerrado. Después de un número finito de pasos se obtiene $M+F$ cerrado.
\end{proof}

\begin{corolario}\label{ch22:finito-cerrado}
Todo subespacio de dimensión finita de un espacio normado es cerrado.
\end{corolario}
\begin{proof}
Tome $M=\{0\}$ en la proposición anterior. Esta conclusión coincide con la obtenida por coordenadas en el lema~\ref{ch19:coordenadas}; la nueva prueba resalta la estabilidad de la cerradura al añadir un número finito de direcciones.
\end{proof}



\begin{corolario}\label{ch22:lineal-finito}
Si $V$ tiene dimensión finita, toda transformación lineal $T:V\to W$ es continua. Todas las normas en $V$ son equivalentes y $V$ es de Banach.
\end{corolario}
\index{dimension finita@dimensión finita!equivalencia de normas}
\begin{proof}
Si $V=\{0\}$, todas las conclusiones son inmediatas. Supongamos $\dim V=k>0$ y elija una base $v_1,\ldots,v_k$. Por el lema~\ref{ch19:coordenadas} existe $c>0$ tal que $\norm{\sum_j a_jv_j}\geq c\max_j|a_j|$. Para una transformación lineal $T:V\to W$ y $v=\sum_j a_jv_j$, se obtiene la cota explícita
\[
 \norm{Tv}\leq\sum_{j=1}^k|a_j|\norm{Tv_j}
 \leq\frac{\sum_{j=1}^k\norm{Tv_j}}{c}\norm v=:C_T\norm v.
\]
Por linealidad, $\norm{Tx-Ty}\leq C_T\norm{x-y}$. Si $C_T>0$, una distancia de entrada menor que $\eps/C_T$ produce una distancia de salida menor que $\eps$; si $C_T=0$, $T$ es el operador cero. Esto demuestra continuidad sin utilizar todavía la teoría general de la norma de operador.

Si $\norm{\cdot}_1$ y $\norm{\cdot}_2$ son dos normas en $V$, la identidad de $(V,\norm{\cdot}_1)$ en $(V,\norm{\cdot}_2)$ y la identidad inversa son lineales. Por lo recién probado, ambas son continuas. Aplicando la cota explícita anterior a ambas identidades, existen $c,C>0$ tales que $c\norm v_1\leq\norm v_2\leq C\norm v_1$. Estas desigualdades son la equivalencia de normas.

Elija una base de $V$ y la norma de coordenadas $\norm{\sum a_jv_j}_0=\max_j|a_j|$. Una sucesión de Cauchy para esa norma es de Cauchy en cada coordenada; la completitud de $\F$ produce sus límites y, por tanto, un límite en $V$. Así $\norm{\cdot}_0$ es completa. Las desigualdades de equivalencia muestran que una sucesión de Cauchy en la norma original lo es en $\norm{\cdot}_0$, y que la convergencia en esta implica convergencia en la original. Por tanto $V$ es de Banach.
\end{proof}

\begin{proposicion}[Estabilidad de la independencia]\label{ch22:independencia}
Si $v_1,\ldots,v_k$ son linealmente independientes en $V$, existe $\delta>0$ tal que $\norm{w_j-v_j}<\delta$ para todo $j$ implica que $w_1,\ldots,w_k$ también son linealmente independientes.
\end{proposicion}
\begin{proof}
El lema~\ref{ch19:coordenadas} proporciona $c>0$ con
\[
 \norm{\sum_{j=1}^k a_jv_j}\geq c\max_j|a_j|
 \geq\frac ck\sum_{j=1}^k|a_j|.
\]
Tome $\delta=c/(2k)$. Si $\sum_j a_jw_j=0$, entonces
\begin{align*}
 \frac ck\sum_j|a_j|
 &\leq\norm{\sum_j a_jv_j}
 =\norm{\sum_j a_j(v_j-w_j)}\\
 &\leq\sum_j|a_j|\norm{v_j-w_j}
 \leq\frac c{2k}\sum_j|a_j|.
\end{align*}
Esto fuerza $\sum_j|a_j|=0$, y por ello todos los coeficientes son cero.
\end{proof}

\section{Compacidad y dimensión}

\begin{teorema}[Heine--Borel en dimensión finita]\label{ch22:heine-borel}
En un espacio normado de dimensión finita, un conjunto es compacto si y solo si es cerrado y acotado.
\end{teorema}
\begin{proof}
En todo espacio métrico, los compactos son cerrados y acotados. Para el recíproco, use una base para identificar $V$ linealmente con $\F^k$. El corolario~\ref{ch22:lineal-finito} muestra que la identificación y su inversa son continuas y que las normas se acotan mutuamente. Por tanto la imagen de un cerrado acotado de $V$ es cerrada y acotada en $\F^k$. Heine--Borel euclidiano la hace compacta, y su imagen por la identificación inversa es compacta en $V$.
\end{proof}

\begin{teorema}\label{ch22:caracterizacion}
Para un espacio normado $V$, las siguientes condiciones son equivalentes:
\begin{enumerate}[label=(\roman*)]
\item $V$ tiene dimensión finita;
\item la bola cerrada $\overline B(0,1)$ es compacta;
\item $V$ es localmente compacto.
\end{enumerate}
\end{teorema}
\index{compacidad!local en espacios normados}
\begin{proof}
Si la dimensión es finita, Heine--Borel da la compacidad de todas las bolas cerradas. Cada punto tiene una de esas bolas como vecindad compacta, por lo que (i) implica (ii) y (iii).

Si $V$ es localmente compacto, el origen tiene una vecindad compacta $K$. Existe $r>0$ con $B(0,r)\subseteq K$. La bola cerrada $\overline B(0,r/2)$ es un cerrado de $V$ contenido en $K$, luego es compacta. La homotecia $v\mapsto (2/r)v$ muestra que $\overline B(0,1)$ es compacta. Esto prueba (iii)$\Rightarrow$(ii).

Supongamos (ii). La compacidad implica acotación total, así que la bola unitaria está cubierta por un número finito de bolas abiertas de radio $1/2$, con centros $v_1,\ldots,v_k\in V$. Sea $M=\spanop\{v_1,\ldots,v_k\}$. Es cerrado por ser de dimensión finita. Si $M\ne V$, el lema de Riesz daría $v$ con $\norm v=1$ y $\dist(v,M)>1/2$. Sin embargo, algún centro satisface $\norm{v-v_j}<1/2$, contradicción. Por tanto $M=V$ y $V$ tiene dimensión finita.
\end{proof}

\begin{corolario}\label{ch22:no-sigma-compacto}
Un espacio de Banach infinito dimensional no es unión numerable de conjuntos compactos.
\end{corolario}
\begin{proof}
Si $V=\bigcup_nK_n$ con cada $K_n$ compacto, todos son cerrados. Baire da un $K_n$ con interior no vacío. Existe una bola $B(x,r)\subseteq K_n$. La bola cerrada $\overline B(x,r/2)$ es un cerrado contenido en el compacto $K_n$, y por ello es compacta. Trasladar y reescalar hace compacta la bola unitaria de $V$. El teorema anterior contradice la dimensión infinita.
\end{proof}

\begin{ejemplo}
En $\ell^2$, los vectores canónicos $e_n$ tienen norma $1$ y $\norm{e_n-e_m}=\sqrt2$ si $n\ne m$. La sucesión no tiene ninguna subsucesión de Cauchy y, por tanto, ninguna convergente. La bola unitaria es cerrada y acotada, pero no compacta. El teorema anterior demuestra que siempre existe una obstrucción de este tipo en dimensión infinita, aunque no haya una base canónica disponible.
\end{ejemplo}

\section{Ejercicios propuestos: operaciones, series y subespacios}

\begin{ejercicio}[Lectura 19, ejercicio 1]\label{ej:lectura-19-1}
Para $u,v\ne0$, ponga $u'=u/\norm u$ y $v'=v/\norm v$. Demuestre que
\[\norm{u'-v'}\leq\frac{2\norm{u-v}}{\norm u}.\]
\end{ejercicio}


\begin{ejercicio}[Lectura 19, ejercicio 2]\label{ej:lectura-19-2}
Para $A,B\subseteq V$, sea $A+B=\{a+b:a\in A,b\in B\}$. Pruebe: (i) si $A$ es abierto, $A+B$ es abierto; (ii) si $A,B$ son compactos, $A+B$ es compacto; (iii) si $A$ es cerrado y $B$ compacto, $A+B$ es cerrado. Dé además dos conjuntos cerrados cuya suma no sea cerrada.
\end{ejercicio}


\begin{ejercicio}[Lectura 19, ejercicio 3]\label{ej:lectura-19-3}
Demuestre que todo subespacio propio $M$ de $V$ tiene interior vacío.
\end{ejercicio}


\begin{ejercicio}[Lectura 19, ejercicio 4]\label{ej:lectura-19-4}
Pruebe que $V$ es separable si y solo si $\overline B(0,1)$ es separable.
\end{ejercicio}


\begin{ejercicio}[Lectura 19, ejercicio 5]\label{ej:lectura-19-5}
Si $(v_n)$ es de Cauchy, $v_n\ne0$, y no converge a $0$, pruebe que $(v_n/\norm{v_n})$ es de Cauchy. Deduzca que $V$ es de Banach si y solo si $\overline B(0,1)$ es completa.
\end{ejercicio}


\begin{ejercicio}[Lectura 19, ejercicio 6]\label{ej:lectura-19-6}
Sean $M_1,M_2\leq V$ tales que $\cl{M_1}\cap M_2=\{0\}$ y $M_1+M_2$ es cerrado. Demuestre que $M_1$ es cerrado.
\end{ejercicio}


\begin{ejercicio}[Lectura 19, ejercicio 7]\label{ej:lectura-19-7}
Si $v_n\to v$, demuestre que las medias de Cesàro $u_n=(v_1+\cdots+v_n)/n$ convergen a $v$.
\end{ejercicio}


\section{Ejercicios propuestos: dimensión y mejores aproximantes}

\begin{ejercicio}[Lectura 22, ejercicio 4]\label{ej:lectura-22-4}
Si $M$ tiene dimensión finita y $v\notin M$, demuestre que existe $m_0\in M$ con $\dist(v,M)=\norm{v-m_0}$. Deduzca que existe $u$ con $\norm u=\dist(u,M)=1$.
\end{ejercicio}


\begin{ejercicio}[Lectura 22, ejercicio 5]\label{ej:lectura-22-5}
Sea $c_{00}=\spanop\{e_n:n\in\N\}$, el espacio de sucesiones de soporte finito. Demuestre que no es completo con ninguna norma.
\end{ejercicio}


En dimensión finita, las coordenadas controlan continuidad, completitud y compacidad. El capítulo siguiente estudia cómo cambian estas propiedades al considerar operadores entre espacios normados, especialmente cuando la dimensión es infinita.

\chapter{Operadores acotados y compactos}\label{cap:nuevo09}\label{chap:20}\label{chap:23}

Una transformación lineal respeta sumas y escalares, pero el análisis exige además controlar cómo una perturbación de la entrada afecta a la salida. Para un operador lineal, la continuidad equivale a una desigualdad global. La norma de operador cuantifica ese control y permite estudiar sucesiones de transformaciones como puntos de un nuevo espacio normado.

La acotación por sí sola no obliga a que las imágenes de sucesiones acotadas tengan subsucesiones convergentes. Los operadores compactos satisfacen precisamente esa condición más fuerte y recuperan parte del comportamiento de la dimensión finita. Desarrollaremos primero las cotas y los límites de operadores, incluido el principio de acotación uniforme, y después la compacidad y su estabilidad.

Usaremos las propiedades de los espacios normados y de la dimensión finita del capítulo~\ref{cap:nuevo08}, el teorema de Baire y los criterios métricos de compacidad. La completitud del dominio será esencial para la acotación uniforme; la del codominio, para completar el espacio de operadores y conservar la compacidad bajo límites en norma. La topología del espacio de operadores será siempre la de su norma, salvo indicación expresa.

\section{Continuidad y norma de operador}\label{lectura:20}

\begin{teorema}\label{ch20:acotado}
Para una transformación lineal $T:V\to W$ entre espacios normados, son equivalentes:
\begin{enumerate}[label=(\roman*)]
\item $T$ es continua en todo $V$;
\item $T$ es continua en $0$;
\item existe $C\geq0$ tal que $\norm{Tv}\leq C\norm v$ para todo $v\in V$.
\end{enumerate}
Cuando se cumplen, $T$ es una función de Lipschitz y, por tanto, uniformemente continua.
\end{teorema}
\begin{proof}
La primera afirmación implica la segunda. Supongamos continuidad en cero. Como $T0=0$, existe $\delta>0$ tal que $\norm u<\delta$ implica $\norm{Tu}<1$. Si $v\ne0$, tome $u=\delta v/(2\norm v)$. Entonces $\norm u=\delta/2<\delta$ y, por linealidad,
\[
 \frac\delta{2\norm v}\norm{Tv}<1.
\]
De aquí $\norm{Tv}\leq (2/\delta)\norm v$. Para $v=0$ la misma desigualdad vale porque ambos lados son cero.

Finalmente, si se cumple (iii), para $x,y\in V$ obtenemos
\[
 \norm{Tx-Ty}=\norm{T(x-y)}\leq C\norm{x-y}.
\]
Si $C>0$, elegir $\norm{x-y}<\eps/C$ garantiza $\norm{Tx-Ty}<\eps$, con una elección independiente de $x,y$. Si $C=0$, $T$ es el operador cero y la conclusión es inmediata.
\end{proof}

\begin{definicion}
Denotamos por $\Bcal(V,W)$ el espacio de operadores lineales continuos de $V$ en $W$. También se llaman \emph{operadores acotados}. Sus operaciones son puntuales:
\[
 (S+T)v=Sv+Tv,\qquad (\alpha T)v=\alpha(Tv).
\]
Definimos
\[
 \norm T=\sup_{\norm v\leq1}\norm{Tv}.
\]
\end{definicion}
\index{operador!acotado}\index{norma!de operador}

El término ``acotado'' no significa que $T(V)$ sea un conjunto acotado. Un operador lineal no nulo tiene imagen no acotada, porque $T(nv)=nTv$. Significa que manda conjuntos acotados en conjuntos acotados.

\begin{proposicion}\label{ch20:norma}
Las operaciones anteriores hacen de $\Bcal(V,W)$ un espacio vectorial, y $\norm{\cdot}$ es una norma. Se tiene
\[
 \norm T=\inf\{C\geq0:\norm{Tv}\leq C\norm v\text{ para todo }v\in V\}.
\]
Si $V\ne\{0\}$, también
\[
 \norm T=\sup_{v\ne0}\frac{\norm{Tv}}{\norm v}
 =\sup_{\norm v=1}\norm{Tv}.
\]
En particular, $\norm{Tv}\leq\norm T\norm v$ para todo $v$.
\end{proposicion}
\begin{proof}
Las identidades vectoriales se verifican evaluando en cada $v$ y usando las correspondientes identidades en $W$. Además, si $S,T$ están acotados por $C_S,C_T$, entonces
\[
 \norm{(S+T)v}\leq(C_S+C_T)\norm v,
 \qquad\norm{(\alpha T)v}\leq|\alpha|C_T\norm v.
\]
Esto prueba que las operaciones permanecen dentro de $\Bcal(V,W)$.

El supremo de la definición es finito por el teorema~\ref{ch20:acotado}. Es no negativo. Si vale cero, para $v\ne0$ tenemos $\norm{T(v/\norm v)}=0$, y por linealidad $Tv=0$; también $T0=0$, así que $T=0$. El operador cero tiene norma cero. Para cualquier escalar,
\[
 \norm{\alpha T}=\sup_{\norm v\leq1}|\alpha|\norm{Tv}=|\alpha|\norm T.
\]
Y para $\norm v\leq1$,
\[
 \norm{(S+T)v}\leq\norm{Sv}+\norm{Tv}\leq\norm S+\norm T.
\]
Tomar el supremo da la desigualdad triangular.

Si $v\ne0$, la normalización $u=v/\norm v$ muestra
\[
 \norm{Tv}=\norm v\norm{Tu}\leq\norm T\norm v.
\]
Por ello $\norm T$ es una constante admisible en el ínfimo. Toda constante admisible $C$ satisface $\norm{Tv}\leq C$ en la bola unitaria, de modo que $\norm T\leq C$. Queda probada la igualdad con el ínfimo.

Finalmente, cada cociente $\norm{Tv}/\norm v$ es el valor de $\norm{Tu}$ en un vector unitario. Y si $0<\norm v\leq1$, entonces $\norm{Tv}\leq\norm{T(v/\norm v)}$. Estos dos hechos dan las igualdades entre los tres supremos. Cuando $V=\{0\}$ solo existe el operador cero; su norma se define mediante la bola, que contiene al cero, y vale cero. No es necesario asignar una norma mediante el supremo sobre una esfera vacía.
\end{proof}

\begin{ejemplo}
En $C([a,b])$ con norma uniforme, la evaluación $E_t(f)=f(t)$ tiene norma $1$: $|f(t)|\leq\norm f_\infty$ da la cota superior y la función constante $1$ da la inferior. El operador $Tf(x)=\int_a^x f(s)\,ds$ satisface $\norm T=b-a$ si $a<b$, porque $|Tf(x)|\leq(b-a)\norm f_\infty$ y la misma función constante alcanza la cota. Estos cálculos siguen un patrón: primero una estimación para todo vector, después un vector que la alcanza o una sucesión que se aproxima a ella.
\end{ejemplo}

\section{Rango finito y continuidad}\label{sec:nuevo09:rango-finito}

Una funcional lineal es una aplicación lineal $\varphi:V\to\F$. Si es continua, su norma es la norma de operador con codominio $\F$, es decir, $\sup_{\norm v\leq1}|\varphi(v)|$. Antes de estudiar sistemáticamente el dual en el capítulo~\ref{cap:nuevo10}, necesitamos un criterio elemental que convierte una hipótesis sobre el núcleo en una cota de continuidad.

\begin{proposicion}\label{ch21:kernel-cerrado}
Una funcional lineal $\varphi$ en un espacio normado es continua si y solo si $\ker\varphi$ es cerrado.
\end{proposicion}
\begin{proof}
Si es continua, su núcleo es la preimagen del cerrado $\{0\}$. Para la otra dirección, si $\varphi=0$ ya es continua. Suponga $\varphi\ne0$, $M=\ker\varphi$ cerrado, y elija $v_0$ con $\varphi(v_0)=1$. Entonces $d=\dist(v_0,M)>0$. Para $v$ con $\varphi(v)\ne0$,
\[
 \frac{v}{\varphi(v)}-v_0\in M,
\]
y, como $M$ es un subespacio,
\[
 d\leq\norm{\frac v{\varphi(v)}}=\frac{\norm v}{|\varphi(v)|}.
\]
Resulta $|\varphi(v)|\leq d^{-1}\norm v$. Si $\varphi(v)=0$, la misma desigualdad vale directamente. Así $\varphi$ es acotada y continua.
\end{proof}

\begin{teorema}\label{ch22:rango-finito-kernel}
Sea $T:V\to W$ lineal, con $T(V)$ de dimensión finita. Entonces $T$ es continuo si y solo si $\ker T$ es cerrado.
\end{teorema}
\begin{proof}
La necesidad se sigue de que $\ker T=T^{-1}(\{0\})$. Para la suficiencia, si $T(V)=\{0\}$, $T$ es el operador cero. En otro caso, elija una base $w_1,\ldots,w_k$ de $T(V)$ y vectores $v_j\in V$ con $Tv_j=w_j$. Para cada $v$ existen coordenadas únicas $\varphi_j(v)$ tales que
\[
 Tv=\sum_{j=1}^k\varphi_j(v)w_j.
\]
La unicidad de coordenadas y la linealidad de $T$ implican que cada $\varphi_j$ es lineal. Afirmamos
\[
 \ker\varphi_i=\ker T+\spanop\{v_j:j\ne i\}.
\]
Si $\varphi_i(v)=0$, el vector $v-\sum_{j\ne i}\varphi_j(v)v_j$ tiene imagen cero por $T$, de modo que pertenece a $\ker T$. Esto prueba una inclusión. En la otra dirección, la imagen de $m+\sum_{j\ne i}a_jv_j$, con $m\in\ker T$, es $\sum_{j\ne i}a_jw_j$, cuya coordenada $i$ es cero.

La proposición~\ref{ch22:suma-cerrada} muestra que $\ker\varphi_i$ es cerrado. Por la proposición~\ref{ch21:kernel-cerrado}, $\varphi_i$ es continua. Finalmente,
\[
 \norm{Tv}\leq\sum_{j=1}^k|\varphi_j(v)|\norm{w_j}
 \leq\left(\sum_{j=1}^k\norm{\varphi_j}\norm{w_j}\right)\norm v.
\]
Esta cota demuestra continuidad de $T$.
\end{proof}

El último teorema también proporciona otra demostración de la continuidad automática en dominio finito del corolario~\ref{ch22:lineal-finito}: el núcleo de una transformación lineal con dominio finito dimensional es un subespacio finito dimensional y, por ello, cerrado. La imagen también tiene dimensión finita, como se ve tomando las imágenes de una base del dominio. El criterio de rango finito concluye entonces que la transformación es continua. La prueba por coordenadas del capítulo anterior y esta prueba por el núcleo establecen el mismo resultado desde dos perspectivas complementarias.

\section{Límites de operadores y completitud}

\begin{lema}\label{ch20:limite-puntual}
Supongamos que $T_n\in\Bcal(V,W)$, que $\sup_n\norm{T_n}=M<\infty$ y que $T_nv\to Tv$ en $W$ para cada $v\in V$. Entonces $T$ es lineal y continuo, con $\norm T\leq M$.
\end{lema}
\begin{proof}
Para $x,y\in V$ y $\alpha,\beta\in\F$, la continuidad de las operaciones en $W$ da
\[
 T(\alpha x+\beta y)=\lim_nT_n(\alpha x+\beta y)
 =\alpha\lim_nT_nx+\beta\lim_nT_ny=\alpha Tx+\beta Ty.
\]
La continuidad de la norma permite además pasar al límite en $\norm{T_nv}\leq M\norm v$ y obtener $\norm{Tv}\leq M\norm v$. El teorema~\ref{ch20:acotado} concluye la prueba.
\end{proof}

\begin{teorema}\label{ch20:banach-operadores}
Si $W$ es de Banach, entonces $\Bcal(V,W)$ es de Banach, aunque $V$ no sea completo.
\end{teorema}
\begin{proof}
Sea $(T_n)$ de Cauchy en norma de operador. Para cada $v$,
\[
 \norm{T_nv-T_mv}\leq\norm{T_n-T_m}\norm v,
\]
por lo que $(T_nv)$ es de Cauchy en $W$. La completitud de $W$ permite definir $Tv=\lim_nT_nv$. Como toda sucesión de Cauchy es acotada, $\sup_n\norm{T_n}<\infty$, y el lema anterior muestra que $T\in\Bcal(V,W)$.

Aún debemos demostrar convergencia en norma, que es más fuerte que la convergencia puntual utilizada para construir $T$. Dado $\eps>0$, elija $N$ tal que $\norm{T_n-T_m}<\eps/2$ para $n,m\geq N$. Fije $n\geq N$ y $v$ con $\norm v\leq1$. Al hacer $m\to\infty$,
\[
 \norm{(T_n-T)v}=\lim_m\norm{(T_n-T_m)v}\leq\eps/2.
\]
La cota no depende de $v$, así que el supremo sobre la bola da $\norm{T_n-T}\leq\eps/2<\eps$ para todo $n\geq N$.
\end{proof}

La convergencia en norma equivale, por definición, a la convergencia uniforme sobre la bola unitaria. La convergencia puntual solo controla cada vector por separado; la diferencia entre ambas reaparecerá en el ejercicio~\ref{ej:lectura-20-7}.

\section{El principio de acotación uniforme}

\begin{teorema}[Banach--Steinhaus]\label{ch20:uniforme}
Sea $V$ de Banach, sea $W$ normado y sea $\mathcal T\subseteq\Bcal(V,W)$. Si para cada $v\in V$ el conjunto $\{Tv:T\in\mathcal T\}$ es acotado, entonces
\[\sup_{T\in\mathcal T}\norm T<\infty.\]
\end{teorema}
\index{acotacion uniforme@acotación uniforme, principio de}
\index{Banach--Steinhaus, teorema de}
\begin{proof}
Si $\mathcal T$ es vacía no hay nada que demostrar. Para cada entero $n\geq1$, defina
\[
 E_n=\{v\in V:\norm{Tv}\leq n\text{ para todo }T\in\mathcal T\}.
\]
Cada $E_n$ es cerrado: es la intersección, sobre $T\in\mathcal T$, de las preimágenes de la bola cerrada de radio $n$. La hipótesis puntual implica $V=\bigcup_{n\geq1}E_n$, porque para cada $v$ existe un entero que supera todas las normas $\norm{Tv}$.

Por Baire, algún $E_{n_0}$ tiene interior no vacío. Existen $v_0\in V$ y $r>0$ con $B(v_0,r)\subseteq E_{n_0}$. Si $\norm h<r$, los vectores $v_0+h$ y $v_0-h$ pertenecen a $E_{n_0}$. Para cada $T\in\mathcal T$,
\begin{align*}
 2\norm{Th}
 &=\norm{T(v_0+h)-T(v_0-h)}\\
 &\leq\norm{T(v_0+h)}+\norm{T(v_0-h)}\leq2n_0.
\end{align*}
Si $\norm u=1$ y $0<s<r$, aplicamos esta estimación a $h=su$ y obtenemos $\norm{Tu}\leq n_0/s$. Al dejar $s\uparrow r$, resulta $\norm{Tu}\leq n_0/r$. Tomando el supremo en $u$, $\norm T\leq n_0/r$ para todo $T\in\mathcal T$. Si $V=\{0\}$ la conclusión vale directamente.
\end{proof}

\begin{observacion}
La hipótesis intercambia los cuantificadores de una manera muy precisa:
\[
 \forall v\ \exists C_v\ \forall T:\norm{Tv}\leq C_v
 \quad\Longrightarrow\quad
 \exists C\ \forall v\ \forall T:\norm{Tv}\leq C\norm v.
\]
No se hace ese intercambio por álgebra; Baire es la herramienta que lo justifica. Sin completitud puede fallar. En $c_{00}$ con norma uniforme, las funcionales $T_n(x)=nx_n$ son continuas y tienen norma $n$. Para cada sucesión de soporte finito $x$, sin embargo, $T_n(x)=0$ salvo para un número finito de índices. La familia está puntualmente acotada y sus normas no lo están.
\end{observacion}

Hasta aquí, las estimaciones controlan el tamaño de las imágenes y la convergencia de las transformaciones. La geometría de la bola en dimensión infinita muestra que ese control no basta para extraer límites de una sucesión acotada de imágenes. La siguiente noción incorpora esa posibilidad de extracción sin exigir que el rango sea finito dimensional.

\section{Definición y criterios de compacidad}\label{lectura:23}

\begin{definicion}
Un operador lineal $T:V\to W$ es \emph{compacto} si para toda sucesión acotada $(v_n)$ de $V$, la sucesión $(Tv_n)$ tiene una subsucesión convergente en $W$. Denotamos por $\Kcal(V,W)$ el conjunto de estos operadores.
\end{definicion}
\index{operador!compacto}

La frase ``en $W$'' forma parte de la definición. Si $W$ no es completo, una sucesión de imágenes podría ser de Cauchy sin converger en el codominio; eso no basta para la compacidad.

\begin{proposicion}\label{ch23:compacto-acotado}
Todo operador compacto es continuo.
\end{proposicion}
\begin{proof}
Si $T$ no fuera continuo, el criterio del teorema~\ref{ch20:acotado} permitiría elegir $v_n$ con $\norm{v_n}=1$ y $\norm{Tv_n}>n$. La sucesión de entradas es acotada, pero ninguna subsucesión de las imágenes puede converger: a lo largo de cualquier subsucesión sus normas siguen tendiendo a infinito, mientras que toda sucesión convergente es acotada. Esto contradice la definición.
\end{proof}

\begin{proposicion}\label{ch23:bola-compacta}
Un operador lineal $T:V\to W$ es compacto si y solo si $\cl{T(\overline B(0,1))}$ es compacto en $W$.
\end{proposicion}
\begin{proof}
Suponga primero que $T$ es compacto. Sea $(y_n)$ una sucesión en el cierre de la imagen de la bola. Para cada $n$ elija $v_n$ con $\norm{v_n}\leq1$ y $\norm{Tv_n-y_n}<1/n$. Existe una subsucesión $Tv_{n_j}\to y$ en $W$. Entonces
\[
 \norm{y_{n_j}-y}\leq1/n_j+\norm{Tv_{n_j}-y}\to0.
\]
El límite pertenece al cierre de la imagen porque este es cerrado. Hemos probado que ese cierre es secuencialmente compacto y, por ser métrico, compacto.

Recíprocamente, si el cierre es compacto y $(v_n)$ es acotada, elija $R>0$ con $\norm{v_n}\leq R$ para todo $n$. La sucesión $T(v_n/R)$ pertenece a un conjunto de cierre compacto y tiene una subsucesión convergente. Multiplicar por $R$ da una subsucesión convergente de $(Tv_n)$.
\end{proof}

\begin{proposicion}\label{ch23:totalmente-acotado}
Si $W$ es de Banach, un operador lineal $T:V\to W$ es compacto si y solo si manda conjuntos acotados en conjuntos totalmente acotados.
\end{proposicion}
\begin{proof}
Si $T$ es compacto y $A$ es acotado, existe $R>0$ con $A\subseteq R\overline B(0,1)$. La proposición anterior muestra que $\cl{T(A)}$ está contenido en un compacto, por lo que $T(A)$ es totalmente acotado.

Si $T$ manda conjuntos acotados en totalmente acotados, la imagen de la bola unitaria es totalmente acotada. Su cierre también lo es. Por ser un cerrado del espacio completo $W$, ese cierre es completo. Un espacio métrico completo y totalmente acotado es compacto, por lo que la proposición anterior da la compacidad de $T$. La completitud solo se utilizó en esta segunda dirección.
\end{proof}

\section{Ejemplos y dimensión finita}

\begin{proposicion}\label{ch23:rango-finito}
Todo operador lineal \emph{acotado} de rango finito es compacto. En particular, un operador acotado es compacto cuando su dominio o su codominio tiene dimensión finita.
\end{proposicion}
\begin{proof}
Si $F=T(V)$ es finito dimensional, la imagen de una sucesión acotada es acotada en $F$ porque $T$ es acotado. Una bola cerrada suficientemente grande de $F$ es compacta por Heine--Borel. Existe por ello una subsucesión de imágenes que converge en $F$, y por tanto en $W$.

Si el dominio es finito dimensional, todo operador lineal es continuo y su rango es finito dimensional. Si el codominio es finito dimensional, el rango también lo es, pero la continuidad todavía debe suponerse.
\end{proof}

\begin{observacion}[Corrección de una hipótesis]
La lectura original afirma que el rango finito basta por sí solo. Es falso. En $c_{00}$ con norma uniforme, la funcional
\[
 \varphi(x)=\sum_{n=1}^{\infty}n x_n
\]
está bien definida porque cada $x$ tiene soporte finito, es lineal y tiene rango de dimensión $1$. Sin embargo, $\norm{e_n}_\infty=1$ y $\varphi(e_n)=n$. No es continua ni compacta. La hipótesis ``acotado'' en la proposición anterior corrige exactamente esta dificultad.
\end{observacion}

\begin{ejemplo}[Un operador compacto de rango infinito]\label{ch23:diagonal}
En $\ell^2$, defina $T(x)_k=x_k/k$. La serie que define su norma converge porque
\[
 \sum_{k=1}^{\infty}\left|\frac{x_k}{k}\right|^2\leq\sum_{k=1}^{\infty}|x_k|^2.
\]
Así $T$ lleva $\ell^2$ en sí mismo. Es lineal coordenada a coordenada, $\norm T\leq1$ y, al evaluar en $e_1$, $\norm T=1$.

Para probar su compacidad directamente, sea $(x^{(n)})$ acotada, con $\norm{x^{(n)}}_2\leq R$. Si $R=0$ todas las entradas son cero. Supongamos $R>0$. Para cada coordenada $k$, la sucesión de escalares $x^{(n)}_k/k$ es acotada. Mediante extracciones sucesivas y una diagonal obtenemos una subsucesión $(Tx^{(n_j)})$ cuyas coordenadas convergen todas. La extracción se hace de modo que la subsucesión de la etapa $k+1$ esté contenida en la de la etapa $k$; tomar el término $j$ de la etapa $j$ conserva la convergencia de cada coordenada desde un índice en adelante.

La convergencia de coordenadas todavía no implica convergencia en $\ell^2$. Controlamos las colas uniformemente:
\[
 \sum_{k>N}|(Tx^{(n)})_k|^2
 \leq\frac1{(N+1)^2}\sum_{k>N}|x^{(n)}_k|^2
 \leq\frac{R^2}{(N+1)^2}.
\]
Dado $\eps>0$, elija $N$ con $2R/(N+1)<\eps/2$. Para dos términos de la subsucesión, la norma de la diferencia de sus colas es a lo sumo $2R/(N+1)$. Como las primeras $N$ coordenadas convergen, la norma euclidiana de la diferencia de esas coordenadas es menor que $\eps/2$ cuando ambos índices son grandes. Sumando ambas estimaciones, la subsucesión es de Cauchy en $\ell^2$. La completitud de $\ell^2$ da su convergencia. Por tanto $T$ es compacto. Su rango es infinito dimensional porque contiene $T(e_k)=e_k/k$, una familia linealmente independiente.
\end{ejemplo}

\begin{teorema}\label{ch23:identidad}
La identidad de un espacio normado $V$ es compacta si y solo si $V$ tiene dimensión finita.
\end{teorema}
\begin{proof}
Si la dimensión es finita, la identidad es acotada y de rango finito, luego compacta. Recíprocamente, si es compacta, la proposición~\ref{ch23:bola-compacta} dice que el cierre de $I(\overline B(0,1))=\overline B(0,1)$ es compacto. La bola ya es cerrada, así que es compacta. El teorema~\ref{ch22:caracterizacion} implica dimensión finita.
\end{proof}

\section{Separabilidad y límites en norma}

\begin{teorema}\label{ch23:separable}
Si $T:V\to W$ es compacto, tanto $T(V)$ como $\cl{T(V)}$ son separables. No hace falta que $W$ sea completo.
\end{teorema}
\begin{proof}
Por la proposición~\ref{ch23:bola-compacta}, los conjuntos
\[
 K_m=\cl{T(\overline B(0,m))}\qquad(m\in\N)
\]
son compactos. Todo compacto métrico es separable: para cada $n$ se puede elegir una red finita de radio $1/n$, y la unión de todas esas redes es numerable y densa. Por tanto cada subespacio métrico $T(\overline B(0,m))$ es separable; el argumento de base numerable del ejercicio~\ref{ej:lectura-19-4} justifica esta herencia a subconjuntos.

Elija un conjunto numerable $D_m$ denso en $T(\overline B(0,m))$. Como
\[
 T(V)=\bigcup_{m=1}^{\infty}T(\overline B(0,m)),
\]
la unión $D=\bigcup_mD_m$ es numerable y densa en $T(V)$. En efecto, cada punto de $T(V)$ pertenece a alguna de esas imágenes y puede aproximarse con puntos de su $D_m$. Tomando cierres, $\cl D=\cl{T(V)}$, de modo que el mismo conjunto es denso en el cierre.
\end{proof}

\begin{proposicion}
El conjunto $\Kcal(V,W)$ es un subespacio de $\Bcal(V,W)$.
\end{proposicion}
\begin{proof}
El operador cero es compacto. Si $S,T$ son compactos y $(v_n)$ es acotada, extraiga primero una subsucesión en la que $Sv_n$ converja. De esa subsucesión extraiga otra en la que $Tv_n$ converja. En la segunda, ambas imágenes convergen, y por la continuidad de las operaciones también converge $\alpha Sv_n+\beta Tv_n$. Esto vale para cualesquiera escalares $\alpha,\beta$, por lo que $\alpha S+\beta T$ es compacto.
\end{proof}

\begin{teorema}\label{ch23:cierre}
Si $W$ es de Banach y $T_n\in\Kcal(V,W)$ converge en norma de operador a $T\in\Bcal(V,W)$, entonces $T$ es compacto. En consecuencia, $\Kcal(V,W)$ es un subespacio cerrado y completo de $\Bcal(V,W)$.
\end{teorema}
\begin{proof}
Demostraremos que $T(\overline B(0,1))$ es totalmente acotado. Dado $\eps>0$, elija $n$ con $\norm{T-T_n}<\eps/3$. Como $T_n$ es compacto, su imagen de la bola unitaria es totalmente acotada. Existen $y_1,\ldots,y_k\in W$ tales que cada $T_nv$, con $\norm v\leq1$, está a distancia menor que $\eps/3$ de alguno de ellos. Para ese centro,
\[
 \norm{Tv-y_j}\leq\norm{(T-T_n)v}+\norm{T_nv-y_j}<2\eps/3<\eps.
\]
La familia finita de bolas $B(y_j,\eps)$ cubre la imagen por $T$. Reescalar muestra lo mismo para la imagen de cualquier bola, y todo conjunto acotado está contenido en una bola. Por la proposición~\ref{ch23:totalmente-acotado}, $T$ es compacto.

La cerradura se sigue del criterio de sucesiones en el espacio normado de operadores. Como $\Bcal(V,W)$ es completo por el teorema~\ref{ch20:banach-operadores}, su subespacio cerrado $\Kcal(V,W)$ también lo es.
\end{proof}

\begin{observacion}
El operador del ejemplo~\ref{ch23:diagonal} también se aproxima en norma por sus truncaciones $T_N(x)=(x_1,x_2/2,\ldots,x_N/N,0,\ldots)$. Estas son de rango finito y
\[
 \norm{T-T_N}=\frac1{N+1}.
\]
La cota superior se obtiene de la estimación de colas, y el vector $e_{N+1}$ da la cota inferior. El teorema de cerradura proporciona así una segunda prueba de compacidad. No se está afirmando que todo operador compacto entre espacios de Banach arbitrarios admita aproximación por operadores de rango finito; el resultado probado va en la dirección de la aproximación hacia la compacidad.
\end{observacion}

\begin{ejemplo}[Por qué importa el codominio completo]
Considere $V=W=c_{00}$ con norma uniforme y $Tx=(x_k/k)_k$. Sus truncaciones $T_N$ tienen rango finito, son acotadas y compactas. Además $\norm{T-T_N}=1/(N+1)\to0$. Sin embargo, $T$ no es compacto con valores en $c_{00}$. Las entradas $x^{(m)}=(1,\ldots,1,0,\ldots)$, con $m$ unos, tienen norma $1$. Sus imágenes convergen uniformemente en $c_0$ a $(1,1/2,1/3,\ldots)$, que no tiene soporte finito. Ninguna subsucesión puede converger en $c_{00}$, porque cualquier límite tendría necesariamente esas mismas coordenadas. Así la hipótesis de completitud en el teorema~\ref{ch23:cierre} no es decorativa.
\end{ejemplo}

\section{Ejercicios propuestos: acotación y convergencia de operadores}

\begin{ejercicio}[Lectura 20, ejercicio 1]\label{ej:lectura-20-1}
Si $T\in\Bcal(V_1,V_2)$ y $S\in\Bcal(V_2,V_3)$, demuestre que $ST\in\Bcal(V_1,V_3)$ y $\norm{ST}\leq\norm S\norm T$.
\end{ejercicio}


\begin{ejercicio}[Lectura 20, ejercicio 2]\label{ej:lectura-20-2}
Pruebe que un operador lineal es continuo si y solo si manda sucesiones acotadas en sucesiones acotadas.
\end{ejercicio}


\begin{ejercicio}[Lectura 20, ejercicio 3]\label{ej:lectura-20-3}
Si $T:V\to W$ es lineal continuo y $F\leq W$ tiene dimensión finita, demuestre que $T^{-1}(F)$ es un subespacio cerrado. Deduzca que $\ker T$ es cerrado.
\end{ejercicio}


\begin{ejercicio}[Lectura 20, ejercicio 4]\label{ej:lectura-20-4}
Demuestre que un operador lineal es continuo si y solo si es continuo en algún punto de su dominio.
\end{ejercicio}


\begin{ejercicio}[Lectura 20, ejercicio 5]\label{ej:lectura-20-5}
Si $T_n\to T$ en $\Bcal(V_1,V_2)$ y $S_n\to S$ en $\Bcal(V_2,V_3)$, demuestre que $S_nT_n\to ST$ en norma.
\end{ejercicio}


\begin{ejercicio}[Lectura 20, ejercicio 6]\label{ej:lectura-20-6}
En $\R^2$ con norma uniforme, sea $T$ el operador de matriz $\begin{pmatrix}a&b\\c&d\end{pmatrix}$. Demuestre que
\[\norm T=\max\{|a|+|b|,|c|+|d|\}.\]
\end{ejercicio}


\begin{ejercicio}[Lectura 20, ejercicio 7]\label{ej:lectura-20-7}
Sea $(T_n)$ una sucesión acotada en $\Bcal(V,W)$ que converge puntualmente a $T$. Si $A\subseteq V$ es totalmente acotado, pruebe que la convergencia es uniforme en $A$.
\end{ejercicio}


\section{Ejercicios propuestos: operadores y dimensión finita}\label{sec:nuevo09:ejercicios-finito}

\begin{ejercicio}[Lectura 22, ejercicio 1]\label{ej:lectura-22-1}
Demuestre que toda transformación lineal sobreyectiva de un espacio normado sobre un espacio normado de dimensión finita es abierta.
\end{ejercicio}


\begin{ejercicio}[Lectura 22, ejercicio 2]\label{ej:lectura-22-2}
Sea $V$ finito dimensional con base $v_1,\ldots,v_k$, donde $k\geq1$. Pruebe que existe $C>0$ tal que
\[
 \norm T\leq C\max_j\norm{Tv_j}
\]
para todo operador lineal $T$ de $V$. Deduzca que si $T_n$ converge puntualmente a cero, entonces $\norm{T_n}\to0$. Compruebe por separado esta última conclusión cuando $V=\{0\}$.
\end{ejercicio}


\begin{ejercicio}[Lectura 22, ejercicio 3]\label{ej:lectura-22-3}
Sea $\mathcal A_n=\{T\in\Bcal(V,W):\dim T(V)\leq n\}$. Demuestre que $\mathcal A_n$ es cerrado en norma de operador.
\end{ejercicio}


\begin{ejercicio}[Lectura 22, ejercicio 6]\label{ej:lectura-22-6}
Sea $V\subseteq\ell^1$ de dimensión finita. Demuestre que para todo $\eps>0$ existe $N\in\N$ tal que, para todo $x\in V$ y todo $\lambda\in\F$,
\[
 \norm{x+\lambda e_N}_1\geq|\lambda|+(1-\eps)\norm x_1.
\]
\end{ejercicio}


\section{Ejercicios propuestos: compacidad de operadores}

\begin{ejercicio}[Lectura 23, ejercicio 1]\label{ej:lectura-23-1}
Si $f:V\to W$ es uniformemente continua en todo el espacio normado $V$ y $A\subseteq V$ es acotado, demuestre que $f(A)$ es acotado.
\end{ejercicio}


\begin{ejercicio}[Lectura 23, ejercicio 2]\label{ej:lectura-23-2}
Si $W$ es de Banach, pruebe que un operador lineal $T:V\to W$ es compacto si y solo si manda conjuntos acotados en conjuntos totalmente acotados.
\end{ejercicio}


\begin{ejercicio}[Lectura 23, ejercicio 3]\label{ej:lectura-23-3}
Sean $S:V\to W$ y $T:W\to X$ operadores continuos. Si uno de ellos es compacto, demuestre que $TS$ es compacto.
\end{ejercicio}


\begin{ejercicio}[Lectura 23, ejercicio 4]\label{ej:lectura-23-4}
Sea $T:V\to V$ compacto y $V$ infinito dimensional. Pruebe que no existe un operador continuo $S$ con $ST=TS=I$.
\end{ejercicio}

\chapter{Dualidad, Hahn--Banach y geometría convexa}\label{cap:nuevo10}\label{chap:21}\label{chap:24}

Una funcional lineal asigna un escalar a cada vector y puede interpretarse como una medición: su núcleo registra los vectores que esa medición no detecta. El dual reúne todas las mediciones continuas. Estudiaremos cómo describen hiperplanos y recuperan la norma de un vector, y cómo las cotas inferiores permiten reconocer isomorfismos sobre la imagen.

El puente hacia la geometría será el teorema de Hahn--Banach: extender una funcional conservando una cota permite separar conjuntos convexos. Sus puntos de máximo identificarán después caras extremas, y el teorema de Krein--Milman reconstruirá un compacto convexo a partir de sus puntos extremos mediante combinaciones convexas y límites. Las demostraciones se presentan aquí en el orden de sus dependencias.

Se requieren la norma de operador y su completitud cuando el codominio es completo, las propiedades de los compactos métricos y el lema de Zorn. Este último afirma que un conjunto parcialmente ordenado en el que toda cadena tiene una cota superior posee un elemento maximal. Lo aplicaremos primero a extensiones compatibles de funcionales y después a subconjuntos extremos compactos. En cada aplicación comprobaremos expresamente la condición sobre las cadenas.

\section{Funcionales e hiperplanos}\label{lectura:21}

\begin{definicion}
El \emph{dual continuo} de $V$ es $V^*=\Bcal(V,\F)$. Sus elementos son funcionales lineales continuas y llevan la norma de operador
\[\norm\varphi=\sup_{\norm v\leq1}|\varphi(v)|.\]
El dual algebraico contiene todas las funcionales lineales, sin exigir continuidad, y puede ser mayor que $V^*$.
\end{definicion}
\index{dual!continuo}\index{funcional lineal}

\begin{proposicion}
El espacio $V^*$ es de Banach, aunque $V$ no sea completo.
\end{proposicion}
\begin{proof}
Los cuerpos $\R$ y $\C$ son completos. Aplicamos el teorema~\ref{ch20:banach-operadores} con $W=\F$.
\end{proof}

\begin{proposicion}\label{ch21:kernel-proporcional}
Dos funcionales lineales no nulas tienen el mismo núcleo si y solo si son múltiplos escalares no nulos una de la otra.
\end{proposicion}
\begin{proof}
Si $\psi=c\varphi$ y $c\ne0$, sus conjuntos de ceros coinciden. Recíprocamente, si sus núcleos coinciden, elija $v_0$ con $\varphi(v_0)=1$. Esto es posible: primero elija un vector donde $\varphi$ no se anule y divídalo entre su valor. Para cualquier $v$, el vector $v-\varphi(v)v_0$ pertenece a $\ker\varphi$, pues
\[
 \varphi(v-\varphi(v)v_0)=\varphi(v)-\varphi(v)\cdot1=0.
\]
También pertenece a $\ker\psi$, de modo que $\psi(v)=\psi(v_0)\varphi(v)$. El escalar $\psi(v_0)$ no es cero, porque $v_0$ no está en el núcleo común. Si se admiten funcionales nulas, la igualdad de núcleos igual a $V$ significa que ambas son cero.
\end{proof}

\begin{teorema}\label{ch21:hiperplanos}
Los subespacios propios maximales de un espacio vectorial $V$ son exactamente los núcleos de funcionales lineales no nulas. Cada uno tiene codimensión $1$. La correspondencia es biyectiva si se identifican funcionales que difieren por un factor escalar no nulo.
\end{teorema}
\index{hiperplano}
\begin{proof}
Si $M=\ker\varphi$ y $\varphi(v_0)=1$, todo $v$ admite la descomposición
\[
 v=\underbrace{v-\varphi(v)v_0}_{\in M}+\varphi(v)v_0.
\]
Es única, porque $M\cap\F v_0=\{0\}$. Así $V=M\oplus\F v_0$. Todo subespacio que contiene estrictamente a $M$ contiene algún $w$ con $\varphi(w)\ne0$; la misma descomposición usando $w/\varphi(w)$ muestra que ese subespacio es $V$. Por ello $M$ es maximal.

Recíprocamente, si $M$ es propio maximal, elija $v_0\notin M$. El subespacio $M+\F v_0$ contiene estrictamente a $M$, por lo que es $V$. Como la suma es directa, la regla $\varphi(m+\alpha v_0)=\alpha$ está bien definida, es lineal, no nula y tiene núcleo $M$. La proposición~\ref{ch21:kernel-proporcional} prueba que la única indeterminación es un factor no nulo.
\end{proof}

El criterio de continuidad por cerradura del núcleo ya se demostró en la proposición~\ref{ch21:kernel-cerrado}, al estudiar operadores de rango finito. Aplicado a las funcionales de esta sección, enlaza su descripción algebraica con la topología.

En consecuencia, los hiperplanos cerrados corresponden a las clases de funcionales continuas no nulas. Además, el núcleo de una funcional discontinua es denso: su cierre es un subespacio que contiene al núcleo; por maximalidad, ese cierre es el propio núcleo o todo $V$. La primera posibilidad lo haría cerrado y forzaría continuidad.

\section{Isomorfismos y cotas inferiores}

\begin{definicion}
Un \emph{isomorfismo de espacios normados} $T:V\to W$ es una transformación lineal biyectiva para la cual $T$ y $T^{-1}$ son continuas. Una aplicación lineal es una \emph{isometría lineal} si $\norm{Tv}=\norm v$ para todo $v$; no se exige sobreyectividad. Un isomorfismo isométrico es, además, biyectivo. Una aplicación puede ser un isomorfismo de $V$ sobre su imagen sin ser sobreyectiva sobre $W$.
\end{definicion}
\index{isomorfismo!de espacios normados}

\begin{teorema}\label{ch21:cota-inferior}
Un operador lineal $T:V\to W$ es un isomorfismo de $V$ sobre $T(V)$ si y solo si existen $a,b>0$ tales que
\[
 a\norm v\leq\norm{Tv}\leq b\norm v\qquad(v\in V).
\]
Para concluir que es un isomorfismo de $V$ sobre $W$ hay que exigir además $T(V)=W$.
\end{teorema}
\begin{proof}
Si $T$ es un isomorfismo sobre su imagen y $V\ne\{0\}$, la continuidad da
\[
 \norm{Tv}\leq\norm T\norm v,
 \qquad\norm v=\norm{T^{-1}Tv}\leq\norm{T^{-1}}\norm{Tv}.
\]
Por tanto sirven $a=1/\norm{T^{-1}}$ y $b=\norm T$. Para el espacio nulo pueden elegirse cualesquiera constantes positivas.

Recíprocamente, la cota superior hace continuo a $T$. Si $Tv=0$, la inferior obliga a que $v=0$, por lo que $T$ es inyectivo. Tiene una inversa definida en $T(V)$, y para $w=Tv$,
\[
 \norm{T^{-1}w}=\norm v\leq a^{-1}\norm w.
\]
Esta inversa es continua. No hemos probado que cada $w\in W$ sea de la forma $Tv$, porque tal afirmación no se deduce de las desigualdades.
\end{proof}

\begin{observacion}[Corrección de la lectura original]
La inclusión $T:\R\to\R^2$, $T(t)=(t,0)$, es una isometría y satisface ambas cotas con $a=b=1$, pero no es sobreyectiva. Es un contraejemplo al enunciado original si ``isomorfismo'' se entiende sobre todo el codominio. La formulación anterior distingue la imagen del codominio.
\end{observacion}

La inversa de un isomorfismo es un isomorfismo, y la composición de dos isomorfismos también: la inversa de $ST$ es $T^{-1}S^{-1}$, y la composición de operadores continuos es continua. Cuando $V\ne\{0\}$,
\[
 1=\norm{I_V}=\norm{T^{-1}T}\leq\norm{T^{-1}}\norm T.
\]
El factor $2I$ es un isomorfismo pero no una isometría. Por otra parte, si $V$ es completo y $T$ tiene cota inferior positiva, $T(V)$ es cerrado en $W$: de $Tv_n\to w$ se deduce que $(v_n)$ es de Cauchy; si $v_n\to v$, la continuidad da $w=Tv$.

\section{La extensión de Hahn--Banach}

La construcción de funcionales a partir de una base de Hamel produce funcionales algebraicas, pero no controla su norma. El siguiente argumento conserva una cota durante la extensión. Se presenta con detalle para justificar tanto la identificación isométrica con una parte del bidual como los teoremas de separación de conjuntos convexos que siguen.

\begin{definicion}
En un espacio vectorial real, $p:V\to\R$ es \emph{sublineal} si
\[
 p(x+y)\leq p(x)+p(y),\qquad p(tx)=tp(x)\quad(t\geq0).
\]
No se exige $p(-x)=p(x)$. Una norma es un ejemplo, pero los funcionales asociados a conjuntos convexos pueden no ser simétricos.
\end{definicion}

\begin{teorema}[Hahn--Banach real, forma dominada]\label{ch21:hb-real}
Sea $p$ sublineal en un espacio vectorial real $V$, sea $M\leq V$ y sea $f:M\to\R$ lineal con $f(m)\leq p(m)$ para todo $m\in M$. Existe una funcional lineal $F:V\to\R$ que extiende a $f$ y satisface $F(v)\leq p(v)$ para todo $v\in V$.
\end{teorema}
\index{Hahn--Banach, teorema de}
\begin{proof}
Primero extendemos una sola dimensión. Si $x_0\notin M$, cualquier extensión a $M+\R x_0$ tiene la forma
\[
 f_1(m+tx_0)=f(m)+tc,
\]
donde $c$ es el valor que asignaremos a $x_0$. La representación es única porque $x_0\notin M$. Para que $f_1\leq p$ necesitamos
\[
 \sup_{y\in M}\{f(y)-p(y-x_0)\}\leq c
 \leq\inf_{z\in M}\{p(z+x_0)-f(z)\}.
\]
El intervalo no es vacío. En efecto, para $y,z\in M$,
\[
 f(y)+f(z)=f(y+z)\leq p(y+z)
 \leq p(y-x_0)+p(z+x_0).
\]
Reordenando obtenemos que cada término del extremo izquierdo es menor o igual que cada término del derecho. Además ambos extremos son finitos: tomar $y=0$ proporciona una cota inferior finita para el supremo, y tomar $z=0$ proporciona una cota superior finita; la desigualdad cruzada controla las otras dos cotas. Podemos elegir un número real $c$ en ese intervalo.

Si $t>0$, la cota superior aplicada a $z=m/t$ da
\[
 f(m)+tc\leq t\,p(m/t+x_0)=p(m+tx_0).
\]
Si $t=-s<0$, la cota inferior aplicada a $y=m/s$ da
\[
 f(m)-sc\leq s\,p(m/s-x_0)=p(m-sx_0).
\]
Para $t=0$ se trata de la hipótesis inicial. Hemos construido la extensión de una dimensión.

Considere ahora todos los pares $(N,g)$ donde $M\subseteq N\subseteq V$ es un subespacio, $g$ extiende a $f$ y $g\leq p$ en $N$. Ordénelos por extensión. La familia no es vacía porque contiene $(M,f)$. En una cadena, las funciones coinciden en los dominios comunes. Por ello la unión de sus dominios es un subespacio y la unión de sus funciones es una funcional lineal bien definida dominada por $p$. Es una cota superior de la cadena. Zorn proporciona un par maximal $(N,g)$. Si $N\ne V$, el paso de una dimensión lo extendería a un subespacio mayor, contradiciendo maximalidad. Así $N=V$.
\end{proof}

\begin{corolario}[Extensión con conservación de la norma]\label{ch21:hb-norma}
Si $M$ es un subespacio de un espacio normado real o complejo $V$ y $f\in M^*$, existe $F\in V^*$ con $F|_M=f$ y $\norm F=\norm f$.
\end{corolario}
\begin{proof}
En el caso real aplique el teorema a $p(v)=\norm f\norm v$. Obtenemos $F(v)\leq\norm f\norm v$. Aplicar la misma desigualdad a $-v$ da $-F(v)\leq\norm f\norm v$, luego $|F(v)|\leq\norm f\norm v$ y $\norm F\leq\norm f$. La restricción da la desigualdad contraria.

En el caso complejo considere $V$ y $M$ como espacios reales. La función $g=\operatorname{Re}f$ es real lineal y está dominada por $p(v)=\norm f\norm v$. Extiéndala a una funcional real lineal $G$ dominada por $p$, y defina
\[
 F(v)=G(v)-iG(iv).
\]
Es aditiva y real homogénea. Además,
\[
 F(iv)=G(iv)-iG(-v)=G(iv)+iG(v)=iF(v),
\]
lo que implica homogeneidad compleja. Para $m\in M$, $G(m)=\operatorname{Re}f(m)$ y $G(im)=\operatorname{Re}(if(m))=-\operatorname{Im}f(m)$, así que $F(m)=f(m)$.

Si $F(v)\ne0$, elija un escalar $\lambda$ de módulo $1$ con $\lambda F(v)=|F(v)|$. Entonces
\[
 |F(v)|=\operatorname{Re}F(\lambda v)=G(\lambda v)
 \leq\norm f\norm{\lambda v}=\norm f\norm v.
\]
Si $F(v)=0$, la desigualdad es inmediata. Así $\norm F\leq\norm f$; la restricción vuelve a dar igualdad.
\end{proof}

\begin{corolario}\label{ch21:funcional-normante}
Para cada $v\ne0$ existe $\varphi\in V^*$ con $\norm\varphi=1$ y $\varphi(v)=\norm v$. En consecuencia,
\[\norm v=\sup_{\norm\varphi\leq1}|\varphi(v)|.\]
\end{corolario}
\begin{proof}
En $M=\F v$ defina $f(\alpha v)=\alpha\norm v$. Esta funcional tiene norma $1$, porque $|f(\alpha v)|=\norm{\alpha v}$. Extiéndala por el corolario anterior. La desigualdad $|\varphi(v)|\leq\norm\varphi\norm v$ da una cota superior para el supremo, y la extensión construida alcanza esa cota. Para $v=0$, todos los valores son cero.
\end{proof}

La extensión con conservación de la norma asegura que las funcionales continuas distinguen vectores; la forma dominada de Hahn--Banach permite algo más: adaptar una funcional a la geometría de un conjunto que no tiene por qué ser simétrico. Pasamos ahora de hiperplanos y mediciones a conjuntos convexos, antes de construir sus funcionales de separación.

\section{Conjuntos y combinaciones convexas}\label{lectura:24}

Toda convexidad de estas secciones es \emph{real}. Si el espacio es complejo, trabajamos con su espacio vectorial real subyacente. Los coeficientes de las combinaciones convexas son reales no negativos; las desigualdades de separación se escriben para funcionales reales, o para la parte real de una funcional compleja. No existe un orden de los números complejos que permita escribir sin más $\varphi(a)<\varphi(b)$. La convexidad permite interpolar entre puntos sin salir de un conjunto; una funcional que alcanza allí un máximo permitirá identificar sus partes extremas.



\begin{definicion}
El segmento entre $x$ e $y$ es
\[
 [x,y]=\{(1-t)x+ty:0\leq t\leq1\}.
\]
Un conjunto $A$ es \emph{convexo} si $[x,y]\subseteq A$ para todos $x,y\in A$. Un conjunto es un \emph{cono convexo} si $\alpha x+\beta y$ pertenece a él siempre que $x,y$ pertenezcan a él y $\alpha,\beta\geq0$.
\end{definicion}
\index{convexidad}\index{cono convexo}

\begin{ejemplo}
El vacío y todo subespacio vectorial son convexos. Toda bola abierta de un espacio normado es convexa: si $\norm{y-x}<r$, $\norm{z-x}<r$ y $0\leq t\leq1$, entonces
\[
 \norm{(1-t)y+tz-x}\leq(1-t)\norm{y-x}+t\norm{z-x}<r.
\]
La última desigualdad vale también en $t=0,1$, donde queda una sola de las dos hipótesis. Para la bola cerrada, las mismas estimaciones con $\leq r$ dan $\norm{(1-t)y+tz-x}\leq r$. Un cono convexo es convexo tomando $\alpha=1-t$ y $\beta=t$. Si un cono no es vacío, contiene al cero, porque se permiten ambos coeficientes iguales a cero.
\end{ejemplo}

\begin{proposicion}\label{ch24:interior-cierre}
El interior y el cierre de un conjunto convexo de un espacio normado son convexos.
\end{proposicion}
\begin{proof}
Si $x,y\in\interior A$, elija $r>0$ tal que $B(x,r)$ y $B(y,r)$ estén contenidos en $A$. Para $0\leq t\leq1$ y $\norm h<r$, los puntos $x+h,y+h$ pertenecen a $A$. Por convexidad,
\[
 (1-t)(x+h)+t(y+h)=(1-t)x+ty+h\in A.
\]
Así la bola de radio $r$ alrededor de $(1-t)x+ty$ está en $A$, y ese punto es interior. Si el interior es vacío, su convexidad se cumple por definición.

Para el cierre, si $x,y\in\cl A$, elija sucesiones $x_n,y_n\in A$ convergentes a ellos. Para cada $t\in[0,1]$, los puntos $(1-t)x_n+ty_n$ están en $A$ y convergen a $(1-t)x+ty$ por continuidad de las operaciones. Por ello el límite pertenece a $\cl A$.
\end{proof}

\begin{definicion}
La \emph{envolvente convexa}, o cápsula convexa, de $A$, denotada por $\operatorname{co}(A)$, es la intersección de todos los conjuntos convexos que contienen a $A$. La \emph{envolvente convexa cerrada} es la intersección de todos los conjuntos cerrados convexos que contienen a $A$.
\end{definicion}
\index{envolvente convexa}\index{envolvente convexa!cerrada}

Estas familias no son vacías, porque contienen a todo el espacio. La intersección de conjuntos convexos es convexa: si dos puntos están en todos los conjuntos de la familia, todo su segmento está en cada uno y por tanto en la intersección. La intersección también es cerrada cuando todos los conjuntos lo son. Así las definiciones producen los menores conjuntos con las propiedades indicadas.

\begin{proposicion}\label{ch24:combinaciones}
La envolvente convexa de $A$ es el conjunto de todas sus combinaciones convexas finitas:
\[
 \operatorname{co}(A)=
 \left\{\sum_{j=1}^n\alpha_ja_j:
 n\geq1,\ a_j\in A,\ \alpha_j\geq0,\ \sum_{j=1}^n\alpha_j=1\right\}.
\]
Además, su envolvente convexa cerrada es $\cl{\operatorname{co}(A)}$.
\end{proposicion}
\begin{proof}
Denote por $D$ el conjunto del lado derecho. Contiene a $A$ tomando una sola variable con coeficiente $1$. Si $x=\sum_{j=1}^n\alpha_ja_j$ e $y=\sum_{k=1}^m\beta_kb_k$, entonces
\[
 (1-t)x+ty=\sum_{j=1}^n(1-t)\alpha_ja_j+\sum_{k=1}^m t\beta_kb_k.
\]
Los coeficientes son no negativos y su suma es $(1-t)+t=1$. Por tanto $D$ es convexo, y la minimalidad de $\operatorname{co}(A)$ da $\operatorname{co}(A)\subseteq D$.

Para la otra inclusión, pruebe por inducción que todo conjunto convexo que contiene a $A$ contiene esas combinaciones. Para un término es inmediato. En una combinación de $n+1$ términos, si $\alpha_{n+1}=1$, el resultado es $a_{n+1}$. Si $\alpha_{n+1}<1$, ponga $s=1-\alpha_{n+1}>0$ y
\[
 y=\sum_{j=1}^n\frac{\alpha_j}{s}a_j.
\]
Por la hipótesis inductiva, $y$ pertenece al conjunto convexo, y la combinación original es $sy+\alpha_{n+1}a_{n+1}$, que también pertenece a él. Esto prueba $D\subseteq\operatorname{co}(A)$.

Por la proposición~\ref{ch24:interior-cierre}, $\cl{\operatorname{co}(A)}$ es cerrado, convexo y contiene a $A$. Cualquier cerrado convexo que contenga a $A$ contiene primero a $\operatorname{co}(A)$ y después a su cierre. Así ese cierre es exactamente la envolvente convexa cerrada. Si $A=\varnothing$, todas las envolventes son vacías y las afirmaciones siguen siendo válidas.
\end{proof}

\section{Separación: construcción de la funcional}

Para separar conjuntos convexos, primero convertimos la geometría de un abierto convexo en una función sublineal. La simetría no será necesaria; por eso utilizamos la versión real dominada de Hahn--Banach, no solamente su versión para normas.

\begin{lema}[Funcional de Minkowski]\label{ch24:minkowski}
Sea $U$ abierto convexo en un espacio normado real y suponga $0\in U$. Defina
\[
 p_U(x)=\inf\{t>0:x\in tU\}.
\]
Entonces $p_U$ es finito, no negativo y sublineal. Si $B(0,r)\subseteq U$, entonces $p_U(x)\leq\norm x/r$. Además,
\[U=\{x:p_U(x)<1\}.\]
\end{lema}
\index{Minkowski, funcional de}
\begin{proof}
Como $0$ es interior, existe $r>0$ con $B(0,r)\subseteq U$. Para todo $t>\norm x/r$, el vector $x/t$ está en esa bola y, por tanto, $x\in tU$. Esto prueba finitud y $p_U(x)\leq\norm x/r$. La no negatividad procede de que el ínfimo se toma entre números positivos, y $p_U(0)=0$ porque $0\in tU$ para todo $t>0$.

Para $s>0$, cambiar $t$ por $su$ en la definición da $p_U(sx)=sp_U(x)$; para $s=0$ la igualdad ya está probada. Observe que $0\in U$ y convexidad implican $aU\subseteq bU$ cuando $0<a<b$: si $u\in U$, entonces $(a/b)u+(1-a/b)0\in U$.

Si $a>p_U(x)$, la definición de ínfimo permite elegir $t<a$ con $x\in tU$, y la inclusión anterior da $x\in aU$. Si también $b>p_U(y)$, los puntos $x/a$ e $y/b$ están en $U$. Su combinación convexa
\[
 \frac{x+y}{a+b}=\frac a{a+b}\frac xa+\frac b{a+b}\frac yb
\]
pertenece a $U$. Así $p_U(x+y)\leq a+b$. Haciendo $a\downarrow p_U(x)$ y $b\downarrow p_U(y)$ obtenemos subaditividad.

Si $p_U(x)<1$, existe $0<t<1$ con $x\in tU$, y por convexidad $tU\subseteq U$. Por tanto $x\in U$. Recíprocamente, si $x\in U$ y $x\ne0$, la apertura de $U$ y la continuidad de $s\mapsto sx$ permiten elegir $s>1$ suficientemente próximo a $1$ con $sx\in U$. Entonces $x\in s^{-1}U$, de modo que $p_U(x)\leq s^{-1}<1$. Para $x=0$ ya se sabe que $p_U(0)=0$.
\end{proof}

\begin{observacion}
No hemos afirmado que $p_U(-x)=p_U(x)$. Por ejemplo, para $U=(-1,2)\subseteq\R$, $p_U(1)=1/2$ y $p_U(-1)=1$. La sublinealidad exige homogeneidad solo para escalares no negativos. Esta precisión es esencial para utilizar correctamente la convexidad sin añadir una simetría que el conjunto no tiene.
\end{observacion}

\begin{lema}\label{ch24:separacion-abierto}
Si $D$ es abierto, convexo, no vacío y $0\notin D$, existe una funcional real lineal continua $\ell\ne0$ tal que $\ell(d)>0$ para todo $d\in D$.
\end{lema}
\begin{proof}
Elija $a\in D$ y ponga $U=D-a$. Es abierto convexo y contiene al origen. El vector $x_0=-a$ no pertenece a $U$, porque de lo contrario $0\in D$. Por el lema anterior, $p_U(x_0)\geq1$.

En la recta $\R x_0$ defina $f(tx_0)=t p_U(x_0)$. Para $t\geq0$, la homogeneidad de $p_U$ da $f(tx_0)=p_U(tx_0)$. Para $t<0$, $f(tx_0)<0\leq p_U(tx_0)$. Así $f\leq p_U$ en su dominio. Hahn--Banach real, teorema~\ref{ch21:hb-real}, proporciona una extensión $F$ a todo el espacio con $F\leq p_U$.

Si $B(0,r)\subseteq U$, aplicamos la cota del lema tanto a $x$ como a $-x$ y obtenemos
\[
 F(x)\leq\norm x/r,\qquad -F(x)=F(-x)\leq\norm x/r.
\]
Por ello $F$ es continua. Además $F(x_0)=p_U(x_0)\geq1$, de modo que no es cero. Si $d\in D$, entonces $d-a\in U$ y
\[
 F(d)=F(d-a)+F(a)<1-p_U(x_0)\leq0.
\]
La funcional $\ell=-F$ satisface la conclusión.
\end{proof}

\begin{teorema}[Separación estricta de un compacto y un cerrado]\label{ch24:separacion}
Sean $A,B$ no vacíos, convexos y disjuntos en un espacio normado real, con $A$ compacto y $B$ cerrado. Existe $\ell\in V^*$ tal que
\[
 \sup_{a\in A}\ell(a)<\inf_{b\in B}\ell(b).
\]
En un espacio complejo, existe una funcional compleja continua $\varphi$ para la cual la misma desigualdad vale con $\operatorname{Re}\varphi$ en lugar de $\ell$.
\end{teorema}
\index{separacion@separación!de conjuntos convexos}
\begin{proof}
El conjunto $C=B-A$ es convexo. Es cerrado porque es la suma del cerrado $B$ y del compacto $-A$, como se demostró en el ejercicio~\ref{ej:lectura-19-2}. La disjunción implica $0\notin C$. Por cerradura existe $d=\dist(0,C)>0$. Sea $r=d/2$ y
\[
 D=C+B(0,r).
\]
Es convexo y abierto: la convexidad se comprueba sumando las combinaciones correspondientes, y la apertura procede de que es una unión de traslaciones de la bola abierta. No contiene a cero, porque $c+h=0$ con $\norm h<r$ implicaría $\norm c<d$, imposible.

El lema anterior da $\ell\ne0$ con $\ell(c+h)>0$ para $c\in C$ y $\norm h<r$. La simetría de la bola y la definición de norma de una funcional real implican
\[
 \inf_{\norm h<r}\ell(h)=-r\norm\ell.
\]
Para justificarlo, la cota $\ell(h)\geq-\norm\ell\norm h$ da un sentido; en el otro, elija vectores unitarios donde $|\ell|$ se aproxime a $\norm\ell$, cambie su signo para obtener valores negativos y multiplíquelos por escalares menores que $r$ que tiendan a $r$. De $\ell(c)+\ell(h)>0$ para todos esos $h$ se deduce
\[
 \ell(c)\geq r\norm\ell>0.
\]
En particular, para todo $a\in A$ y $b\in B$,
\[
 \ell(b)-\ell(a)\geq r\norm\ell.
\]
Como $A$ es compacto, $\ell$ alcanza en él un máximo finito. Tomándolo en la desigualdad y después tomando el ínfimo sobre $B$, obtenemos
\[
 \inf_B\ell\geq\max_A\ell+r\norm\ell>\max_A\ell.
\]

En el caso complejo aplique la prueba al espacio real subyacente. A partir de la funcional real $\ell$, defina $\varphi(v)=\ell(v)-i\ell(iv)$. El cálculo del corolario~\ref{ch21:hb-norma} muestra que es compleja lineal y continua, y su parte real es precisamente $\ell$. Esto da la versión compleja sin ordenar números complejos.
\end{proof}

\section{Puntos extremos y el teorema de Krein--Milman}

\begin{definicion}
Un punto $x\in A$ es \emph{extremo} si, cuando $x=(1-t)y+tz$ con $y,z\in A$ y $0<t<1$, necesariamente $y=z=x$. El conjunto de puntos extremos se denota $\operatorname{ext}(A)$.

Un subconjunto $F\subseteq A$ es \emph{extremo en $A$} si toda combinación $(1-t)y+tz\in F$, con $y,z\in A$ y $0<t<1$, obliga a que $y,z\in F$. Si además $F$ es convexo, se denomina una \emph{cara} de $A$.
\end{definicion}
\index{punto extremo}\index{cara!de un convexo}

La definición equivalente de punto extremo dice que no puede escribirse como combinación convexa estricta de dos puntos \emph{distintos} de $A$. Esa palabra falta en la formulación original: sin ella, $x=(1-t)x+tx$ excluiría a todos los puntos. Una singleton $\{x\}$ es extrema en $A$ exactamente cuando $x$ es un punto extremo.

\begin{lema}\label{ch24:caras}
Sea $K$ compacto no vacío y sea $\ell$ una funcional real continua. El conjunto
\[
 F=\{x\in K:\ell(x)=\max_K\ell\}
\]
es no vacío, compacto y extremo en $K$. Si $K$ es convexo, $F$ también lo es. Además, si $G$ es extremo en $F$ y $F$ es extremo en $A$, entonces $G$ es extremo en $A$.
\end{lema}
\begin{proof}
El máximo $m$ existe por compacidad y continuidad. El conjunto $F=K\cap\ell^{-1}(\{m\})$ es un cerrado no vacío de $K$, luego compacto. Si $(1-t)x+ty\in F$ con $x,y\in K$ y $0<t<1$, entonces
\[
 0=m-\ell((1-t)x+ty)
 =(1-t)(m-\ell(x))+t(m-\ell(y)).
\]
Los dos sumandos son no negativos y sus coeficientes son estrictamente positivos. Por ello ambos déficits son cero, de modo que $x,y\in F$. Si $K$ es convexo y $x,y\in F$, el segmento pertenece a $K$ y la linealidad da $\ell((1-t)x+ty)=m$, por lo que está en $F$.

Para la transitividad, si $(1-t)x+ty\in G\subseteq F$ y $x,y\in A$, la extremalidad de $F$ fuerza $x,y\in F$; después la extremalidad de $G$ dentro de $F$ fuerza $x,y\in G$.
\end{proof}

\begin{teorema}[Krein--Milman para compactos en norma]\label{ch24:krein-milman}
Si $A$ es no vacío, compacto y convexo en un espacio normado, entonces $\operatorname{ext}(A)\ne\varnothing$ y
\[
 A=\cl{\operatorname{co}(\operatorname{ext}(A))}.
\]
\end{teorema}
\index{Krein--Milman, teorema de}
\begin{proof}
Primero encontraremos un punto extremo. Sea $\mathcal F$ la familia de los subconjuntos \emph{no vacíos}, compactos y extremos de $A$. No es vacía porque $A\in\mathcal F$. Ordénela por inclusión inversa: $F\preceq G$ significa $G\subseteq F$. Buscamos así un elemento maximal para ese orden, que será un conjunto mínimo por inclusión.

Sea $\mathcal C$ una cadena no vacía de $\mathcal F$. Sus miembros son compactos y, por tanto, cerrados en el compacto $A$. Cualquier intersección finita de miembros de una cadena coincide con el más pequeño de esos miembros y no es vacía. La propiedad de intersección finita de los compactos implica
\[
 K=\bigcap_{F\in\mathcal C}F\ne\varnothing.
\]
Es compacto por ser cerrado en $A$. Es extremo porque, si una combinación estricta de puntos de $A$ pertenece a $K$, pertenece a cada $F\in\mathcal C$; la extremalidad de cada $F$ obliga a que ambos extremos pertenezcan a todos ellos y, por tanto, a $K$. Así $K\in\mathcal F$ y es una cota superior de la cadena en el orden inverso. Para la cadena vacía, cualquier miembro de $\mathcal F$ es una cota superior. Zorn proporciona un miembro minimal por inclusión $F_0$.

Si $F_0$ contuviera dos puntos distintos $u,v$, Hahn--Banach real aplicado al vector $v-u$ proporcionaría una funcional real continua $\ell$ con $\ell(v-u)=\norm{v-u}>0$. Por tanto $\ell(u)<\ell(v)$. El subconjunto de máximos de $\ell$ en $F_0$ es compacto, no vacío y extremo en $F_0$ por el lema anterior. Por transitividad, es extremo en $A$. Es un subconjunto propio de $F_0$ porque no contiene a $u$, contradiciendo la minimalidad. Así $F_0=\{e\}$ y $e\in\operatorname{ext}(A)$.

Ponga $E=\operatorname{ext}(A)$ y $C=\cl{\operatorname{co}(E)}$. Como $A$ es cerrado convexo y contiene a $E$, se tiene $C\subseteq A$. En particular, $C$ es compacto por ser cerrado en $A$; no estamos usando una afirmación general de compacidad de envolventes convexas.

Si existiera $x\in A\setminus C$, el teorema de separación aplicado a $C$ y $\{x\}$ daría una funcional real continua $\ell$ con
\[
 \sup_C\ell<\ell(x)\leq\max_A\ell=:m.
\]
La cara $F=\{a\in A:\ell(a)=m\}$ es no vacía, compacta y convexa, y es extrema en $A$. La primera parte de la demostración, aplicada a $F$, proporciona un punto extremo $e$ de $F$. La transitividad del lema hace que $e$ sea extremo en $A$, así que $e\in E\subseteq C$. Pero $\ell(e)=m>\sup_C\ell$, contradicción. En consecuencia $A\subseteq C$, y queda probada la igualdad.
\end{proof}

El cierre en el enunciado permite límites de combinaciones convexas finitas; no equivale a permitir una suma infinita arbitraria sin comprobar su convergencia. La versión aquí demostrada usa compacidad para la norma. Otras topologías conducen a versiones más generales, pero requieren hipótesis adicionales que no se están utilizando en este capítulo.

\section{Ejercicios propuestos: dualidad, distancia e invertibilidad}

\begin{ejercicio}[Lectura 21, ejercicio 1]\label{ej:lectura-21-1}
Demuestre que el espacio de todas las funcionales lineales de $\F^n$ es isomorfo a $\F^n$.
\end{ejercicio}


\begin{ejercicio}[Lectura 21, ejercicio 2]\label{ej:lectura-21-2}
Sea $\varphi\in V^*$ y $M=\ker\varphi$. Pruebe que
\[|\varphi(v)|=\norm\varphi\,\dist(v,M)\qquad(v\in V).\]
\end{ejercicio}


\begin{ejercicio}[Lectura 21, ejercicio 3]\label{ej:lectura-21-3}
Sea $A\subseteq V$ no vacío y $U=\overline B(0,1)$. Pruebe
\[
 \cl{A+U}=\{v\in V:\dist(v,A)\leq1\}.
\]
Deduzca que $A+U$ es cerrado si y solo si todo $v$ con $\dist(v,A)=1$ admite $a\in A$ con $\norm{v-a}=1$.
\end{ejercicio}


\begin{ejercicio}[Lectura 21, ejercicio 4]\label{ej:lectura-21-4}
Sea $\norm\varphi=1$, $M=\ker\varphi$ y $U=\overline B(0,1)$. Demuestre
\[
 \cl{M+U}=\{v:|\varphi(v)|\leq1\}.
\]
Pruebe que $M+U$ es cerrado si y solo si existe $u\in U$ con $|\varphi(u)|=1$.
\end{ejercicio}


\begin{ejercicio}[Lectura 21, ejercicio 5]\label{ej:lectura-21-5}
Sea $V$ de Banach y $T\in\Bcal(V,V)$ con $\sum_{n=0}^{\infty}\norm{T^n}<\infty$. Demuestre que $I-T$ es un isomorfismo y calcule su inversa.
\end{ejercicio}

\index{Neumann, serie de}

\begin{ejercicio}[Lectura 21, ejercicio 6]\label{ej:lectura-21-6}
Sean $V,W$ de Banach, con $V\ne\{0\}$, y $T:V\to W$ un isomorfismo. Si $S\in\Bcal(V,W)$ satisface $\norm{S-T}<1/\norm{T^{-1}}$, demuestre que $S$ es un isomorfismo. Deduzca que los isomorfismos forman un abierto de $\Bcal(V,W)$, incluyendo por separado el caso $V=W=\{0\}$.
\end{ejercicio}


\begin{ejercicio}[Lectura 21, ejercicio 7]\label{ej:lectura-21-7}
Sea $T$ un isomorfismo de $V$ sobre un subespacio de $W$. Pruebe que, si $\norm{T'-T}$ es suficientemente pequeña, $T'$ es un isomorfismo sobre su imagen.
\end{ejercicio}


\begin{ejercicio}[Lectura 21, ejercicio 8]\label{ej:lectura-21-8}
Pruebe que la evaluación canónica
\[
 J:V\longrightarrow V^{**},\qquad J(v)(\varphi)=\varphi(v),
\]
es una isometría lineal.
\end{ejercicio}


\section{Ejercicios propuestos: envolventes y separación}

\begin{ejercicio}[Lectura 24, ejercicio 1]\label{ej:lectura-24-1}
Demuestre que la intersección de una familia no vacía de conjuntos convexos es convexa.
\end{ejercicio}


\begin{ejercicio}[Lectura 24, ejercicio 2]\label{ej:lectura-24-2}
Si $A,B\subseteq V$ y $C\subseteq W$ son convexos, $\alpha\in\R$ y $T:V\to W$ es lineal, pruebe que $\alpha A$, $A+B$, $T(A)$ y $T^{-1}(C)$ son convexos.
\end{ejercicio}


\begin{ejercicio}[Lectura 24, ejercicio 3]\label{ej:lectura-24-3}
Sean $A_1,\ldots,A_n$ conjuntos convexos no vacíos. Demuestre que $\operatorname{co}(A_1\cup\cdots\cup A_n)$ está formado por las combinaciones $\sum_{j=1}^n\alpha_jx_j$, con $x_j\in A_j$, $\alpha_j\geq0$ y $\sum_j\alpha_j=1$.
\end{ejercicio}


\begin{ejercicio}[Lectura 24, ejercicio 4: revisión y solución]\label{ej:lectura-24-4}
La lectura propone, para conjuntos cerrados $A,B$, las afirmaciones: (i) si ambos son compactos, $\operatorname{co}(A\cup B)$ es compacto; (ii) si $A$ es compacto y $B$ acotado, $\operatorname{co}(A\cup B)$ es cerrado. Examine ambas afirmaciones y pruebe sus versiones correctas. Demuestre además que la envolvente convexa de $\{(0,0)\}\cup\{(1,y):y\in\R\}$ no es cerrada.
\end{ejercicio}


\begin{ejercicio}[Lectura 24, ejercicio 5]\label{ej:lectura-24-5}
Si $A$ es convexo, demuestre que $C=\bigcup_{s\geq0}sA$ es un cono convexo.
\end{ejercicio}


\begin{ejercicio}[Lectura 24, ejercicio 6]\label{ej:lectura-24-6}
Demuestre que la envolvente convexa cerrada de un compacto no vacío $A$ es el cierre de su envolvente convexa. Determine si la compacidad es necesaria.
\end{ejercicio}


\begin{ejercicio}[Lectura 24, ejercicio 7]\label{ej:lectura-24-7}
Sean $A,B$ convexos y disjuntos, y sea $x\in V$. Pruebe que
\[
 \operatorname{co}(A\cup\{x\})\cap B=\varnothing
 \quad\text{o}\quad
 \operatorname{co}(B\cup\{x\})\cap A=\varnothing.
\]
\end{ejercicio}


\parte{Espacios de funciones y aproximación}{%
Considerar funciones como puntos de un espacio permite estudiar límites y aproximaciones con las herramientas anteriores. La métrica uniforme proporciona un control común sobre todo el dominio; analizaremos cuándo ese control conserva continuidad, integrales y derivadas.

Arzelà--Ascoli caracteriza la compacidad de familias de funciones mediante equicontinuidad y acotación. Tietze permite extender funciones continuas, y Stone--Weierstrass explica cuándo un álgebra de funciones puede aproximar a todas las funciones continuas. Se usan compacidad, completitud y normas, además del cálculo elemental de una variable.}
\chapter{Espacios de funciones y convergencia uniforme}\label{cap:nuevo11}\label{chap:25}\label{chap:26}
\index{metrica uniforme@métrica uniforme}\index{convergencia!puntual}\index{convergencia!uniforme}

Una sucesión de funciones puede aproximar bien cada valor y, sin embargo, hacerlo mal en una región que cambia con el índice. La métrica uniforme mide el mayor error en todo el dominio y convierte las funciones en puntos de un nuevo espacio métrico. A partir de esa construcción estudiaremos la conservación de la continuidad, las series de funciones y el paso al límite en integrales y derivadas.

Usaremos completitud, continuidad, compacidad, propiedades del supremo y cálculo diferencial e integral de una variable. Cuando construyamos límites de funciones con valores en un espacio normado $V$, exigiremos que $V$ sea Banach. Para comparar funciones, su dominio no necesita inicialmente ninguna estructura: basta una distancia en el espacio de llegada. Salvo indicación contraria, los conjuntos considerados son no vacíos; si el dominio es vacío, existe una sola función con ese dominio y su espacio de funciones es el espacio métrico de un punto.

\section{Funciones acotadas y distancia del supremo}
\label{lectura:25}
\begin{definicion}
Sean $S$ un conjunto y $(Y,d)$ un espacio métrico. Una función $f:S\to Y$ es \emph{acotada} si $f(S)$ está contenido en alguna bola de radio finito. Denotamos por $B(S,Y)$ el conjunto de esas funciones. Para $f,g\in B(S,Y)$ definimos
\[
 d_\infty(f,g)=\sup_{s\in S}d(f(s),g(s)).
\]
\end{definicion}

Antes de verificar los axiomas de una métrica debemos comprobar que el supremo es finito. Si $f(S)\subset B(y_0,R)$ y $g(S)\subset B(y_1,T)$, la desigualdad triangular da
\[
 d(f(s),g(s))\le d(f(s),y_0)+d(y_0,y_1)+d(y_1,g(s))
 <R+d(y_0,y_1)+T.
\]
Por tanto, el conjunto cuyo supremo tomamos es no vacío, está acotado superiormente y está formado por números no negativos.

\begin{proposicion}\label{ch25:metrica}
La función $d_\infty$ es una métrica en $B(S,Y)$.
\end{proposicion}
\begin{proof}
La no negatividad se hereda de $d$. Si $d_\infty(f,g)=0$, entonces
$0\le d(f(s),g(s))\le 0$ para cada $s$, de modo que $f(s)=g(s)$ en todo $S$. Recíprocamente, si $f=g$, todas las distancias son cero y el supremo es cero. La simetría se obtiene de $d(f(s),g(s))=d(g(s),f(s))$.

Para la desigualdad triangular, fijamos $s\in S$ y escribimos
\[
 d(f(s),h(s))\le d(f(s),g(s))+d(g(s),h(s))
 \le d_\infty(f,g)+d_\infty(g,h).
\]
La expresión de la derecha no depende de $s$ y es una cota superior para todos los términos de la izquierda. La propiedad que caracteriza al supremo implica
$d_\infty(f,h)\le d_\infty(f,g)+d_\infty(g,h)$.
\end{proof}

\begin{observacion}[Un supremo no tiene que alcanzarse]
Si $f(s)=s$ y $g(s)=0$ en $S=(0,1)$, entonces $d_\infty(f,g)=1$, aunque $|f(s)-g(s)|<1$ para cada punto. Así, demostrar una desigualdad estricta punto por punto suele producir solamente una desigualdad no estricta después de tomar supremos. Para obtener finalmente un error menor que $\eps$ trabajaremos, por ejemplo, con $\eps/2$.
\end{observacion}

\begin{lema}\label{ch25:unionacotada}
Si una familia no vacía $\mathcal A\subset B(S,Y)$ está acotada para $d_\infty$, entonces $\bigcup_{f\in\mathcal A}f(S)$ está acotado en $Y$.
\end{lema}
\begin{proof}
Elegimos $f_0\in\mathcal A$. Como $\mathcal A$ está acotada, existe $M<\infty$ tal que $d_\infty(f,f_0)\le M$ para toda $f\in\mathcal A$. Como $f_0$ es una función acotada, existen $y_0\in Y$ y $R<\infty$ con $d(f_0(s),y_0)\le R$ para todo $s\in S$. En consecuencia,
\[
 d(f(s),y_0)\le d(f(s),f_0(s))+d(f_0(s),y_0)\le M+R
\]
para toda $f\in\mathcal A$ y todo $s\in S$. La misma bola, de centro $y_0$ y radio $M+R+1$, contiene todas las imágenes.
\end{proof}

Esta prueba proporciona directamente el control necesario. No hace falta construir una subsucesión con diámetros crecientes. Hay que distinguir dos afirmaciones: cada función de una familia puede ser acotada sin que la familia sea acotada en $d_\infty$. Las funciones constantes $f_n(s)=n$ ilustran esa diferencia.

\section{Por qué la completitud se transmite a las funciones}
\index{completitud!espacios de funciones}
\begin{teorema}\label{ch25:completitud}
Si $Y$ es completo, entonces $(B(S,Y),d_\infty)$ es completo.
\end{teorema}
\begin{estrategia}
La condición de Cauchy uniforme permite construir primero un límite en cada punto. Después debemos resolver dos problemas diferentes: demostrar que esa función límite es acotada y demostrar que la convergencia hacia ella es uniforme. Construir únicamente los límites puntuales no concluye la prueba.
\end{estrategia}
\begin{proof}
Sea $(f_n)$ una sucesión de Cauchy en $B(S,Y)$. Para $s\in S$ fijo tenemos
\[
 d(f_n(s),f_m(s))\le d_\infty(f_n,f_m).
\]
Por ello $(f_n(s))$ es de Cauchy en $Y$. La completitud de $Y$ proporciona un único punto $f(s)$ al cual converge esa sucesión. Al realizar esta construcción para cada $s$ obtenemos una función $f:S\to Y$.

Tomemos $N$ tal que $d_\infty(f_n,f_m)<1$ para $n,m\ge N$. Fijando $n=N$ y pasando al límite $m\to\infty$, la continuidad de la distancia da
\[
 d(f_N(s),f(s))\le 1 \qquad(s\in S).
\]
Si $f_N(S)\subset B(y_0,R)$, se sigue que $d(f(s),y_0)\le 1+R$ para todo $s$. Por tanto, $f\in B(S,Y)$.

Ahora sea $\eps>0$. Existe $N_\eps$ tal que $d_\infty(f_n,f_m)<\eps/2$ para $n,m\ge N_\eps$. Para $n\ge N_\eps$ y $s\in S$ fijos, al hacer $m\to\infty$ obtenemos
\[
 d(f_n(s),f(s))\le\eps/2.
\]
El índice $N_\eps$ fue elegido antes que $s$ y sirve para todos los puntos. Tomando supremos, $d_\infty(f_n,f)\le\eps/2<\eps$. Esto demuestra la convergencia en la métrica uniforme.
\end{proof}

\begin{ejemplo}[Por qué la hipótesis de completitud es necesaria]
Si $S=\{s_0\}$, la aplicación $f\mapsto f(s_0)$ es una isometría de $B(S,Y)$ sobre $Y$. Así, un espacio de llegada incompleto ya produce un espacio de funciones incompleto con un dominio de un solo punto. Más generalmente, para cualquier $S$ no vacío, las funciones constantes asociadas a una sucesión de Cauchy sin límite en $Y$ no pueden converger en $B(S,Y)$.
\end{ejemplo}

\section{Continuidad y límite uniforme}
\index{convergencia uniforme!continuidad del limite@continuidad del límite}
Si $X$ es un espacio métrico, escribimos $C_b(X,Y)$ para el conjunto de funciones continuas y acotadas de $X$ a $Y$. El subíndice $b$ recuerda la condición de acotación. Cuando $Y=\F$, usamos $C_b(X)$.

\begin{teorema}[Continuidad del límite uniforme]\label{ch25:cerradocb}\label{ch12:limite-continuo}
Sean $X,Y$ espacios métricos y $f_n:X\to Y$ funciones continuas. Si $f_n$ converge uniformemente a $f:X\to Y$, entonces $f$ es continua, sin necesidad de que las funciones sean acotadas. En particular, el conjunto $C_b(X,Y)$ es cerrado en $B(X,Y)$.
\end{teorema}
\begin{proof}
La convergencia uniforme significa que, dada una tolerancia positiva, todos los valores de $f_n$ están a esa distancia de los de $f$ cuando $n$ es suficientemente grande; el índice se elige antes que el punto. Para funciones acotadas equivale a convergencia en $d_\infty$. Fijemos $x\in X$ y $\eps>0$. Elegimos un índice $n$ tal que $d(f_n(z),f(z))<\eps/3$ para todo $z\in X$. Mantenemos fijo ese índice. La continuidad de $f_n$ en $x$ proporciona $\delta>0$ tal que
\[
 d_X(x,y)<\delta\quad\Longrightarrow\quad d(f_n(x),f_n(y))<\eps/3.
\]
Para esos puntos $y$, la desigualdad triangular da
\[
 \begin{split}
 d(f(x),f(y))&\le d(f(x),f_n(x))+d(f_n(x),f_n(y))+d(f_n(y),f(y))\\
 &<\eps/3+\eps/3+\eps/3=\eps.
 \end{split}
\]
Por tanto, $f$ es continua en $x$. Como $x$ era arbitrario, es continua en $X$. La elección uniforme del índice permite controlar a la vez el error en $x$ y en los puntos vecinos, sin exigir acotación global. Para la última afirmación, una sucesión de $C_b(X,Y)$ que converge en $B(X,Y)$ tiene un límite ya acotado por pertenecer a ese espacio ambiente y continuo por lo que acabamos de probar. La caracterización secuencial de cerrado permite concluir que $C_b(X,Y)$ es cerrado.
\end{proof}

El orden de las elecciones es decisivo: primero se fija una función $f_n$ que aproxima a $f$ en \emph{todos} los puntos y luego se usa la continuidad de esa función. No se necesita que el mismo $\delta$ funcione para todas las funciones de la sucesión.

\begin{corolario}
Si $Y$ es completo, entonces $C_b(X,Y)$ es completo con la métrica uniforme.
\end{corolario}
\begin{proof}
Una sucesión de Cauchy en $C_b(X,Y)$ converge en $B(X,Y)$ por el teorema~\ref{ch25:completitud}. Su límite pertenece a $C_b(X,Y)$ por el teorema~\ref{ch25:cerradocb}. Este es precisamente el significado de completitud del subespacio.
\end{proof}

\section{La norma uniforme y tres espacios que no deben confundirse}
\index{norma!uniforme}\index{funciones!de soporte compacto}\index{funciones!que se anulan en el infinito}
Sea $V$ un espacio normado sobre $\F$. Las operaciones puntuales
\[
 (f+g)(s)=f(s)+g(s),\qquad (\lambda f)(s)=\lambda f(s)
\]
convierten $B(S,V)$ en un espacio vectorial: las identidades vectoriales se verifican en cada punto; además, la suma y el producto por escalar conservan la acotación porque
$\norm{f(s)+g(s)}\le\norm{f(s)}+\norm{g(s)}$ y $\norm{\lambda f(s)}=|\lambda|\norm{f(s)}$.

\begin{proposicion}
La expresión $\norm f_\infty=\sup_{s\in S}\norm{f(s)}$ es una norma en $B(S,V)$ y su métrica asociada es $d_\infty$. Si $V$ es Banach, también lo son $B(S,V)$ y $C_b(X,V)$.
\end{proposicion}
\begin{proof}
La norma es no negativa y se anula exactamente cuando todos los valores $f(s)$ son cero. Para $\lambda\ne0$, multiplicar todas las cantidades de un conjunto por $|\lambda|$ multiplica su supremo por ese factor; así $\norm{\lambda f}_\infty=|\lambda|\norm f_\infty$. Para $\lambda=0$ la igualdad también se cumple. Finalmente,
\[
 \norm{f(s)+g(s)}\le\norm{f(s)}+\norm{g(s)}\le\norm f_\infty+\norm g_\infty;
\]
tomando supremos se obtiene la desigualdad triangular. La identidad $d_\infty(f,g)=\norm{f-g}_\infty$ se sigue de las definiciones. La completitud ya quedó demostrada. El conjunto $C_b(X,V)$ es un subespacio vectorial porque las operaciones de $V$ son continuas, y es cerrado por el teorema~\ref{ch25:cerradocb}.
\end{proof}

Para funciones escalares en un espacio métrico $X$ usaremos las notaciones
\[
 \begin{split}
 C_b(X)&=\{f:X\to\F:f\text{ es continua y acotada}\},\\
 C_c(X)&=\{f:X\to\F:f\text{ es continua y }\supp f\text{ es compacto}\},\\
 C_0(X)&=\{f:X\to\F:f\text{ es continua y }
 \{x:|f(x)|\ge\eps\}\text{ es compacto para cada }\eps>0\},
 \end{split}
\]
donde $\supp f=\cl{\{x:f(x)\ne0\}}$. Estos espacios satisfacen $C_c(X)\subset C_0(X)\subset C_b(X)$. Para justificar la segunda inclusión, el conjunto $K=\{|f|\ge1\}$ es compacto y $f$ es acotada en $K$ por continuidad; fuera de $K$, $|f|<1$. La primera inclusión se obtiene porque cada conjunto $\{|f|\ge\eps\}$ es cerrado y está contenido en el soporte compacto.

\begin{observacion}[Corrección de notación de la lectura original]
El espacio de funciones continuas de soporte compacto no es, en general, $C_b(X)$. En $X=\R$, la función constante $1$ pertenece a $C_b(\R)$ y no pertenece a $C_0(\R)$; la función $e^{-x^2}$ pertenece a $C_0(\R)$ y no pertenece a $C_c(\R)$. Si $X$ es compacto, en cambio, los tres espacios coinciden con $C(X)$.
\end{observacion}

\begin{proposicion}
Para un espacio métrico $X$, el espacio $C_0(X)$ es cerrado en $C_b(X)$. Por tanto, es Banach. Si, además, $X$ es localmente compacto, $C_c(X)$ es denso en $C_0(X)$ para la norma uniforme.
\end{proposicion}
\begin{proof}
Supongamos que $f_n\in C_0(X)$ y $\norm{f_n-f}_\infty\to0$. Dado $\eps>0$, elegimos $n$ con $\norm{f_n-f}_\infty<\eps/2$. Entonces
\[
 \{|f|\ge\eps\}\subset\{|f_n|\ge\eps/2\}.
\]
El conjunto de la izquierda es cerrado en $X$ y el de la derecha es compacto; por ello el primero es compacto. Esto demuestra $f\in C_0(X)$.

Para la densidad, fijemos $f\in C_0(X)$ y $\eps>0$. El conjunto $K=\{|f|\ge\eps/2\}$ es compacto. Si es vacío, la función cero aproxima a $f$ con error a lo sumo $\eps/2$. Si no es vacío, la compacidad local y una subcubierta finita proporcionan un abierto $U$ tal que $K\subset U$ y $\cl U$ es compacto. En efecto, cubrimos $K$ por finitos abiertos cuyas clausuras son compactas y tomamos su unión. Si $U=X$, podemos tomar $\chi=1$. En otro caso, los conjuntos cerrados disjuntos $K$ y $X\setminus U$ permiten definir
\[
 \chi(x)=\frac{\dist(x,X\setminus U)}{\dist(x,X\setminus U)+\dist(x,K)}.
\]
Las funciones distancia son continuas; el denominador no se anula, pues ello colocaría a $x$ en ambos conjuntos cerrados. Además, $0\le\chi\le1$, $\chi=1$ en $K$ y $\chi=0$ fuera de $U$. Por ello $h=\chi f$ es continua y su soporte está contenido en $\cl U$, luego es compacto. En $K$ tenemos $h=f$; fuera de $K$,
$|f-h|=|1-\chi|\,|f|<\eps/2$. Así $\norm{f-h}_\infty\le\eps/2<\eps$.
\end{proof}

\begin{observacion}[Verificaciones de la lectura 25]
La lectura 25 de la colección de partida no contiene una lista de ejercicios. Sus verificaciones están desarrolladas en las secciones anteriores: los axiomas de la métrica y de la norma, el control común de las imágenes, la completitud y la conservación de la continuidad. En particular, la prueba de completitud incluye explícitamente la convergencia uniforme hacia el límite construido.
\end{observacion}

\section{Convergencia puntual y uniforme: el orden de los cuantificadores}
\label{lectura:26}
La convergencia puntual responde una pregunta cada vez: después de fijar $x$, ¿se aproxima $f_n(x)$ a $f(x)$? La convergencia uniforme exige controlar simultáneamente todos los puntos con un único índice. Esa diferencia permite conservar continuidad, intercambiar integrales y límites y, con hipótesis adicionales, derivar límites.

La métrica anterior describe este control cuando las funciones son acotadas. Conviene formular ahora los cuantificadores de manera que también cubran funciones no acotadas y permitan precisar qué operaciones conservan los límites.

\begin{definicion}
Sean $f_n,f:S\to V$. Decimos que $f_n$ converge \emph{puntualmente} a $f$ si
\[
 \forall x\in S\ \forall\eps>0\ \exists N=N(x,\eps)\ \forall n\ge N:
 \norm{f_n(x)-f(x)}<\eps.
\]
Decimos que converge \emph{uniformemente}, y escribimos $f_n\rightrightarrows f$, si
\[
 \forall\eps>0\ \exists N=N(\eps)\ \forall n\ge N\ \forall x\in S:
 \norm{f_n(x)-f(x)}<\eps.
\]
\end{definicion}

Esta definición admite funciones no acotadas. Lo que acaba siendo acotado es el error $f_n-f$. Para funciones de $B(S,V)$ equivale a $\norm{f_n-f}_\infty\to0$. Exigir que \emph{todos} los errores, incluso los de los primeros términos, sean acotados añade una restricción innecesaria: modificar un número finito de términos no debe cambiar la convergencia. Adoptaremos la formulación mediante cuantificadores.

\begin{proposicion}
La convergencia uniforme implica convergencia puntual. El recíproco es falso.
\end{proposicion}
\begin{proof}
El índice que funciona para todos los puntos funciona en particular para cualquier punto fijado. Para refutar el recíproco tomamos $f_n(x)=x^n$ en $[0,1]$. Si $0\le x<1$, la sucesión geométrica $x^n$ tiende a cero; si $x=1$, permanece igual a uno. Por tanto, el límite puntual es
\[
 f(x)=\begin{cases}0,&0\le x<1,\\1,&x=1.\end{cases}
\]
Para cada $n$, $\sup_{0\le x<1}x^n=1$, de modo que $\norm{f_n-f}_\infty=1$. El error no tiende a cero. También podemos detectar el problema porque el límite es discontinuo, mientras que un límite uniforme de funciones continuas es continuo.
\end{proof}

\begin{observacion}[Corrección de la lectura original]
La observación 2 de la lectura 26 invierte esta implicación. La dirección válida es uniforme $\Rightarrow$ puntual; el ejemplo que la acompaña confirma precisamente que la otra dirección no vale.
\end{observacion}

\section{El teorema de Dini}
\index{Dini, teorema de}
\begin{teorema}[Dini]\label{ch26:dini}
Sean $X$ compacto y $(f_n)$ una sucesión de funciones continuas $X\to\R$ que converge puntualmente a una función continua $f$. Si la sucesión es creciente en todos los puntos, o decreciente en todos los puntos, entonces $f_n\rightrightarrows f$.
\end{teorema}
\begin{estrategia}
La convergencia puntual produce una cubierta por regiones donde el error es pequeño. La compacidad reduce esa cubierta a un número finito de regiones. La monotonía permite reemplazar sus distintos índices por uno solo.
\end{estrategia}
\begin{proof}
Supongamos primero que $f_n(x)\le f_{n+1}(x)$ para todo $x$. Al pasar al límite se obtiene $f_n(x)\le f(x)$. Así, las funciones continuas $g_n=f-f_n$ son no negativas, decrecientes y convergen puntualmente a cero.

Fijemos $\eps>0$ y definamos $U_n=\{x:g_n(x)<\eps/2\}$. Cada $U_n$ es abierto por continuidad y $U_n\subset U_{n+1}$ por monotonía. Para cada $x$, la convergencia $g_n(x)\to0$ asegura que $x\in U_n$ para algún $n$; por tanto, $(U_n)$ cubre $X$. Una subcubierta finita $U_{n_1},\ldots,U_{n_r}$ cubre $X$. Si $N=\max_j n_j$, el carácter creciente de la cubierta implica $X=U_N$, y también $X=U_n$ para $n\ge N$. Por ello
\[
 0\le f(x)-f_n(x)<\eps/2\quad(n\ge N,\ x\in X),
\]
lo que implica $\norm{f-f_n}_\infty\le\eps/2<\eps$. Si $(f_n)$ es decreciente, aplicamos la prueba anterior a $(-f_n)$, que es creciente y converge a $-f$.
\end{proof}

Las hipótesis tienen funciones diferentes. La sucesión $x^n$ en $[0,1]$ muestra que no basta con un límite puntual cualquiera: se necesita continuidad del límite. En $(0,1)$, la misma sucesión tiene límite continuo cero y es decreciente, pero no converge uniformemente; aquí falta compacidad. Sin monotonía podemos construir picos continuos de altura uno con soportes disjuntos que se acercan a cero; convergen puntualmente a cero, pero su norma uniforme permanece igual a uno. Por ejemplo, sirven los triángulos de altura uno centrados en $2^{-n}$ y de semiancho $2^{-n-2}$, extendidos por cero en $[0,1]$.

\section{Series de funciones y el criterio M}
\index{series de funciones}\index{Weierstrass!criterio M}
\begin{definicion}
La serie $\sum_{n=1}^\infty u_n$ converge puntual o uniformemente cuando lo hace su sucesión de sumas parciales $s_N=\sum_{n=1}^N u_n$. Converge \emph{absolutamente en cada punto} si $\sum_n\norm{u_n(x)}<\infty$ para cada $x$. Se dice que converge \emph{normalmente} si $\sum_n\norm{u_n}_\infty<\infty$.
\end{definicion}

La convergencia normal es una condición de control uniforme de las normas. No debe confundirse con la convergencia absoluta punto por punto, que permite que los controles dependan de $x$.

\begin{teorema}[Criterio M de Weierstrass]\label{ch26:criteriom}\label{ch12:weierstrass}
Sean $V$ Banach y $u_n:S\to V$. Supongamos que existen $M_n\ge0$ tales que $\norm{u_n(x)}\le M_n$ para todo $x\in S$ y que $\sum_n M_n<\infty$. Entonces $\sum_nu_n$ converge uniformemente y absolutamente en cada punto. Si $s$ es su suma,
\[
 \norm{s-s_N}_\infty\le\sum_{n>N}M_n.
\]
\end{teorema}
\begin{proof}
Para $q>p$ y cualquier $x$,
\[
 \norm{s_q(x)-s_p(x)}
 \le\sum_{n=p+1}^q\norm{u_n(x)}\le\sum_{n=p+1}^qM_n.
\]
La última expresión tiende a cero uniformemente respecto de $q>p$ cuando $p\to\infty$, pues la serie numérica converge. Para cada $x$, $(s_N(x))$ es de Cauchy y, por completitud, tiene límite $s(x)\in V$. Fijando $p=N$ y haciendo $q\to\infty$ obtenemos
$\norm{s(x)-s_N(x)}\le\sum_{n>N}M_n$. Tomar el supremo prueba simultáneamente la convergencia uniforme y la estimación del error. Por otra parte, $\sum_n\norm{u_n(x)}\le\sum_nM_n<\infty$, lo que da la convergencia absoluta.
\end{proof}

\begin{observacion}[La completitud no es decorativa]
La lectura original formula este criterio para un espacio normado sin explicitar completitud. En el espacio $c_{00}$ de sucesiones de soporte finito, con norma $\ell^1$, la serie de vectores $\sum_n 2^{-n}e_n$ satisface $\sum_n\norm{2^{-n}e_n}=1$, pero su supuesto límite tendría infinitas coordenadas no nulas y no pertenecería a $c_{00}$. Un dominio $S$ de un solo punto convierte este ejemplo en una serie de funciones. Por eso formulamos el criterio en un espacio de Banach.
\end{observacion}

\section{Integración de límites uniformes}
\index{integracion@integración!limite uniforme@límite uniforme}
\begin{teorema}\label{ch26:integracion}
Si $f_n:[a,b]\to\F$ son Riemann-integrables y $f_n\rightrightarrows f$, entonces $f$ es Riemann-integrable y
\[
 \int_a^bf(x)\,dx=\lim_{n\to\infty}\int_a^bf_n(x)\,dx.
\]
\end{teorema}
\begin{proof}
El caso $a=b$ es inmediato, de modo que supongamos $a<b$. Comenzamos con funciones reales. Elegimos un índice $n_0$ tal que $\norm{f-f_{n_0}}_\infty<1$. Como una función Riemann-integrable es acotada, $f$ también lo es.

Para una partición $P:a=x_0<\cdots<x_r=b$, denotamos por $U(h,P)$ y $L(h,P)$ las sumas superior e inferior de Darboux. Si $\norm{f-f_n}_\infty\le\delta$, en cada subintervalo $I$ se cumple
\[
 \sup_I f\le\sup_I f_n+\delta,
 \qquad \inf_I f\ge\inf_I f_n-\delta.
\]
Al multiplicar las diferencias por las longitudes y sumar, resulta
\[
 U(f,P)-L(f,P)\le U(f_n,P)-L(f_n,P)+2\delta(b-a).
\]
Dado $\eta>0$, elegimos primero $n$ con $\norm{f-f_n}_\infty<\eta/[4(b-a)]$. Después, por integrabilidad de $f_n$, elegimos una partición $P$ con $U(f_n,P)-L(f_n,P)<\eta/2$. La desigualdad anterior proporciona $U(f,P)-L(f,P)<\eta$. El criterio de Darboux demuestra que $f$ es integrable.

Para funciones complejas, las partes real e imaginaria de $f_n$ son integrables y convergen uniformemente a las partes correspondientes de $f$, pues $|\operatorname{Re}z|,|\operatorname{Im}z|\le|z|$. El argumento real demuestra la integrabilidad de ambas componentes y, por definición, la de $f$.

Finalmente, para una función integrable $h$ real o compleja, cada suma de Riemann satisface
\[
 \left|\sum_j h(\xi_j)(x_j-x_{j-1})\right|
 \le\sum_j|h(\xi_j)|(x_j-x_{j-1})
 \le(b-a)\norm h_\infty.
\]
El paso al límite da $|\int_a^b h|\le(b-a)\norm h_\infty$. Aplicado a $h=f_n-f$, obtenemos
\[
 \left|\int_a^bf_n-\int_a^bf\right|
 \le(b-a)\norm{f_n-f}_\infty\longrightarrow0.
\]
\end{proof}

\begin{corolario}[Integración término a término]
Si cada $u_n$ es Riemann-integrable en $[a,b]$ y $\sum_nu_n$ converge uniformemente a $s$, entonces
$\int_a^bs=\sum_n\int_a^bu_n$.
\end{corolario}
\begin{proof}
Las sumas parciales son integrables por linealidad. Aplicamos el teorema~\ref{ch26:integracion} a esas sumas y usamos $\int_a^b\sum_{n=1}^Nu_n=\sum_{n=1}^N\int_a^bu_n$.
\end{proof}

\section{Derivación de límites: la hipótesis correcta}
\index{derivacion@derivación!limite uniforme@límite uniforme}
La convergencia uniforme de funciones no controla sus derivadas. Por ejemplo, $f_n(x)=\sin(nx)/n$ converge uniformemente a cero en $\R$, pero $f_n'(0)=1$ para todo $n$, mientras que la derivada del límite es cero. El resultado válido comienza con convergencia uniforme de las \emph{derivadas} y fija además una constante de integración mediante el valor en un punto.

\begin{lema}[Desigualdad del valor medio para funciones complejas]\label{ch26:tvmdesigualdad}
Si $h:[u,v]\to\C$ es continua y es derivable, respecto de la variable real, en $(u,v)$, y $|h'(t)|\le M$, entonces $|h(v)-h(u)|\le M(v-u)$.
\end{lema}
\begin{proof}
Si $h(v)=h(u)$, la desigualdad es inmediata. En otro caso elegimos $\lambda\in\C$ con $|\lambda|=1$ y $\lambda(h(v)-h(u))=|h(v)-h(u)|$. La función real $\phi(t)=\operatorname{Re}(\lambda h(t))$ satisface el teorema del valor medio real. Existe $\xi\in(u,v)$ tal que
\[
 |h(v)-h(u)|=\phi(v)-\phi(u)=(v-u)\phi'(\xi)
 \le(v-u)|h'(\xi)|\le(v-u)M.
\]
\end{proof}

En general no existe un punto $\xi$ que satisfaga la igualdad compleja $h(v)-h(u)=h'(\xi)(v-u)$. Por ejemplo, para $h(t)=e^{it}$ en $[0,2\pi]$ el lado izquierdo es cero y la derivada nunca se anula. La desigualdad anterior es el sustituto adecuado para la demostración.

\begin{teorema}\label{ch26:derivacion}
Sean $a<b$ finitos y $f_n:(a,b)\to\C$ derivables. Supongamos que para algún $x_0\in(a,b)$ converge $(f_n(x_0))$ y que $f_n'\rightrightarrows g$ en $(a,b)$. Entonces existe $f$ con $f_n\rightrightarrows f$ en $(a,b)$, la función $f$ es derivable y $f'=g$.
\end{teorema}
\begin{proof}
La convergencia uniforme de $f_n'$ implica que, para $n,m$ suficientemente grandes, $\sup_{(a,b)}|f_n'-f_m'|$ es finito y pequeño. Aplicando el lema anterior a $h=f_n-f_m$ entre $x_0$ y $x$, obtenemos
\[
 |f_n(x)-f_m(x)|
 \le |f_n(x_0)-f_m(x_0)|+|x-x_0|\sup_{(a,b)}|f_n'-f_m'|
 \le |f_n(x_0)-f_m(x_0)|+(b-a)\sup_{(a,b)}|f_n'-f_m'|.
\]
La última cota no depende de $x$ y tiende a cero cuando $n,m\to\infty$. El criterio de Cauchy uniforme, demostrado en el ejercicio~\ref{ej:lectura-26-1}, produce $f_n\rightrightarrows f$.

Para demostrar la derivabilidad fijamos $c\in(a,b)$ y definimos los cocientes prolongados
\[
 q_n(x)=\begin{cases}
 \dfrac{f_n(x)-f_n(c)}{x-c},&x\ne c,\\[4pt]
 f_n'(c),&x=c.
 \end{cases}
\]
La derivabilidad de $f_n$ implica exactamente que $q_n$ es continua en $c$. Para $x\ne c$, el lema aplicado a $f_n-f_m$ proporciona
\[
 |q_n(x)-q_m(x)|\le\sup_{(a,b)}|f_n'-f_m'|.
\]
La misma desigualdad vale para $x=c$ por la definición del cociente prolongado. Por tanto, $(q_n)$ es de Cauchy uniformemente y converge uniformemente a alguna función $q$.

El límite $q$ es continuo en $c$. Para verificarlo sin exigir continuidad de $q_n$ en otros puntos, dado $\eps>0$ elegimos $n$ con $\sup|q-q_n|<\eps/3$ y luego $\delta>0$ con $|q_n(x)-q_n(c)|<\eps/3$ si $|x-c|<\delta$. La desigualdad triangular da $|q(x)-q(c)|<\eps$.

Si $x\ne c$, la convergencia puntual de $f_n$ implica
$q(x)=(f(x)-f(c))/(x-c)$. En $x=c$ tenemos
$q(c)=\lim_n f_n'(c)=g(c)$. Por continuidad de $q$ en $c$,
\[
 \lim_{x\to c}\frac{f(x)-f(c)}{x-c}=q(c)=g(c).
\]
Así $f'(c)=g(c)$, y $c$ era arbitrario.
\end{proof}

No hemos supuesto que las derivadas sean continuas. Tampoco se obtiene una afirmación sobre valores arbitrariamente asignados en los extremos $a,b$: el teorema se refiere al intervalo abierto. La finitud de su longitud interviene en el factor $b-a$.

\section{Ejercicios sobre límites, productos y series uniformes}
\begin{ejercicio}[Lectura 26, ejercicio 1]\label{ej:lectura-26-1}\label{ch26:ej1}
Sea $V$ Banach y $f_n:S\to V$. Demuestre que $(f_n)$ converge uniformemente si y solo si para todo $\eps>0$ existe $N$ tal que $\norm{f_n(x)-f_m(x)}<\eps$ para todos $n,m\ge N$ y $x\in S$.
\end{ejercicio}


\begin{ejercicio}[Lectura 26, ejercicio 2]\label{ej:lectura-26-2}
Sean $(f_n)$ y $(g_n)$ sucesiones de funciones escalares sobre $S$, acotadas en $B(S,\F)$, que convergen uniformemente. Demuestre que $(f_ng_n)$ converge uniformemente. Dé un contraejemplo al suprimir el control de acotación.
\end{ejercicio}


\begin{ejercicio}[Lectura 26, ejercicio 3: enunciado corregido]\label{ej:lectura-26-3}
Sean $V,W$ Banach, $f_n\rightrightarrows f$ y $\norm{f_n}_\infty\le\alpha$. La lectura propone que la continuidad de $g$ en la bola de radio $\alpha$ basta para que $g\circ f_n\rightrightarrows g\circ f$. Examine la afirmación y pruebe una versión válida.
\end{ejercicio}


\begin{ejercicio}[Lectura 26, ejercicio 4]\label{ej:lectura-26-4}
Sean $V$ Banach y $f_n:S\to V$. Suponga que $\norm{f_{n+1}(x)-f_n(x)}\le M_n$ para todo $x$, donde $\sum_nM_n<\infty$. Demuestre que $(f_n)$ converge uniformemente.
\end{ejercicio}


\begin{ejercicio}[Lectura 26, ejercicio 5]\label{ej:lectura-26-5}
Estudie la relación entre convergencia absoluta y uniforme:
\begin{enumerate}[label=(\roman*)]
\item En el dominio $[0,1]$, para $x>0$ sea $u_n(x)=\sin(\pi/x)\mathbf1_{[1/(n+1),1/n]}(x)$ y defina $u_n(0)=0$. Analice la convergencia absoluta puntual de $\sum_nu_n$ y la convergencia uniforme.
\item Demuestre que $\sum_{n=1}^\infty(-1)^n(x^2+n)/n^2$ converge uniformemente en $[0,1]$, pero no absolutamente en ningún punto.
\end{enumerate}
\end{ejercicio}


\begin{ejercicio}[Lectura 26, ejercicio 6]\label{ej:lectura-26-6}
Si $f_n:X\to V$ son continuas y convergen uniformemente a $f$, demuestre que $f_n(x_n)\to f(x)$ siempre que $x_n\to x$. ¿Vale el recíproco?
\end{ejercicio}


\begin{ejercicio}[Lectura 26, ejercicio 7: Dirichlet uniforme]\label{ej:lectura-26-7}
Sean $V$ Banach, $u_n:S\to V$ y $v_n:S\to\R$. Suponga que, para algún $M<\infty$,
\[
 \sup_{x\in S}\left\|\sum_{k=1}^nu_k(x)\right\|\le M\quad(n\ge1),
\]
que $v_n\rightrightarrows0$ y que $v_{n+1}(x)\le v_n(x)$ para todo $x$. Demuestre que $\sum_nv_nu_n$ converge uniformemente.
\end{ejercicio}


\begin{ejercicio}[Lectura 26, ejercicio 8: condensación uniforme de Cauchy]\label{ej:lectura-26-8}
Sea $(u_n)$ una sucesión de funciones reales sobre $S$, decreciente respecto de $n$ en cada punto. Demuestre que $\sum_{n=1}^\infty u_n$ y $\sum_{k=1}^\infty 2^ku_{2^k}$ convergen uniformemente o dejan de hacerlo simultáneamente.
\end{ejercicio}


\section{El problema de Cauchy como ecuación entre funciones}\label{ampl:picard}
Hasta ahora hemos construido límites de funciones y comprobado qué propiedades conservan. Podemos emplear esos resultados para construir una función desconocida a partir de una ecuación. En un problema de Cauchy se prescribe una velocidad que depende del tiempo y del estado, junto con el estado inicial:
\[
 u'(t)=f(t,u(t)),\qquad u(t_0)=\xi.
\]
La idea es integrar antes de buscar la solución. Así la incógnita aparece como punto fijo de una transformación sobre un espacio completo de funciones, y el teorema~\ref{ch11:banach} del capítulo~\ref{cap:nuevo04} se vuelve aplicable.

Las dos demostraciones siguientes retoman las estrategias de Maximenko \cite{max-picard}. Aquí trabajamos en un intervalo compacto y permitimos valores en $\R^d$. No exigimos que $f$ sea acotada en todo el cilindro $I\times\R^d$: la continuidad basta para que su composición con cada curva continua sea acotada. Esta precisión permite incluir ecuaciones lineales como $u'=u$.
\index{Picard, teorema de}\index{problema de Cauchy}\index{norma!ponderada}\index{aproximaciones sucesivas}

\begin{teorema}[Picard en un intervalo compacto]\label{ampl:picard-teorema}
Sea $I=[a,b]$, $a<b$, y sea $f:I\times\R^d\to\R^d$ continua. Suponga que existe $L\ge0$ tal que, para todo $t\in I$ y todo $x,y\in\R^d$,
\[
 \norm{f(t,x)-f(t,y)}_\infty\le L\norm{x-y}_\infty.
\]
Dados $t_0\in I$ y $\xi\in\R^d$, existe una única solución $u\in C^1(I,\R^d)$ del problema de Cauchy. En los extremos del intervalo se entienden derivadas laterales. Desde cualquier $u_0\in C(I,\R^d)$, las aproximaciones
\[
 u_{n+1}(t)=\xi+\int_{t_0}^t f(s,u_n(s))\,ds
\]
convergen uniformemente a esa solución.
\end{teorema}

\begin{proof}[Preparación común a las dos demostraciones]
El espacio $X=C(I,\R^d)$ es Banach para
\[
 \norm u_\infty=\sup_{t\in I}\norm{u(t)}_\infty
\]
por los teoremas~\ref{ch25:completitud} y~\ref{ch25:cerradocb}; $\R^d$ es completo y toda curva continua sobre $I$ es acotada. La integral vectorial se define integrando cada coordenada. La desigualdad escalar para integrales y el máximo de las coordenadas dan
\[
 \left\|\int_c^d v(s)\,ds\right\|_\infty
 \le\int_{\min(c,d)}^{\max(c,d)}\norm{v(s)}_\infty\,ds.
\]
Para $u\in X$, la función $s\mapsto f(s,u(s))$ es continua. Por el teorema fundamental del cálculo,
\[
 (Su)(t)=\xi+\int_{t_0}^t f(s,u(s))\,ds
\]
pertenece a $C^1(I,\R^d)$ y, en particular, a $X$. Un punto fijo de $S$ es una solución del problema de Cauchy; recíprocamente, integrar la ecuación de cualquier solución muestra que es punto fijo. Es suficiente probar existencia, unicidad y convergencia de las iteraciones de $S$. Si $T=\max\{t_0-a,b-t_0\}$, la estimación directa sólo da
\[
 \norm{Su-Sv}_\infty\le LT\norm{u-v}_\infty.
\]
Cuando $LT\ge1$ esta cota no permite aplicar Banach directamente. No significa que la solución deje de existir; significa que necesitamos una estimación mejor.
\end{proof}

\begin{proof}[Primera demostración: una norma ponderada]
Elijamos $\beta>0$ con $\beta\ge L$ y definamos
\[
 \norm u_\beta=\sup_{t\in I}e^{-\beta|t-t_0|}\norm{u(t)}_\infty.
\]
Es una norma: la positividad del peso asegura que sólo se anula en la función cero, y homogeneidad y desigualdad triangular se verifican punto por punto antes de tomar el supremo. Además,
\begin{equation}\label{ampl:picard-equivalencia}
 e^{-\beta T}\norm u_\infty\le\norm u_\beta\le\norm u_\infty.
\end{equation}
Así una sucesión de Cauchy para la norma ponderada lo es para la uniforme, y su límite uniforme también es límite en la ponderada. Por tanto $X$ sigue siendo completo.

Para $r=|t-t_0|$, parametrizamos el segmento entre $t_0$ y $t$ por su distancia $s$ al punto inicial; esto evita cambiar signos cuando $t<t_0$. Obtenemos
\[
 \begin{split}
 e^{-\beta r}\norm{Su(t)-Sv(t)}_\infty
 &\le L e^{-\beta r}\norm{u-v}_\beta\int_0^r e^{\beta s}\,ds\\
 &=\frac L\beta(1-e^{-\beta r})\norm{u-v}_\beta.
 \end{split}
\]
Al tomar supremos, $S$ es una contracción con constante
\begin{equation}\label{ampl:picard-constante}
 q=\frac L\beta(1-e^{-\beta T})<1.
\end{equation}
Si $L=0$, $q=0$; si $L>0$, usamos $L/\beta\le1$ y $1-e^{-\beta T}<1$. Banach da un único punto fijo y convergencia de todas las iteraciones en $\norm{\cdot}_\beta$. La equivalencia~\eqref{ampl:picard-equivalencia} convierte esa convergencia en uniforme. La preparación anterior identifica el punto fijo con la solución requerida.
\end{proof}

\begin{proof}[Segunda demostración: una potencia contractiva]
Demostraremos por inducción que, para todo $n\ge0$,
\begin{equation}\label{ampl:picard-factorial}
 \norm{S^nu(t)-S^nv(t)}_\infty
 \le\frac{L^n|t-t_0|^n}{n!}\norm{u-v}_\infty.
\end{equation}
Para $n=0$, la afirmación es la definición de norma uniforme; interpretamos $L^0=1$. Supongamos cierta la cota para $n$. Al aplicar $S$ una vez más y parametrizar de nuevo el segmento desde $t_0$,
\[
 \norm{S^{n+1}u(t)-S^{n+1}v(t)}_\infty
 \le L\int_0^{|t-t_0|}\frac{L^ns^n}{n!}\norm{u-v}_\infty\,ds,
\]
que es la cota deseada con $n+1$. En consecuencia,
\[
 \norm{S^nu-S^nv}_\infty\le\frac{(LT)^n}{n!}\norm{u-v}_\infty.
\]
Si $L=0$, $S$ ya es constante. Si $L>0$, los números $(LT)^n/n!$ tienden a cero: a partir de un índice, el cociente entre dos consecutivos es $LT/(n+1)\le1/2$. Existe, pues, $N\ge1$ para el cual $S^N$ es contracción.

Sea $u$ su único punto fijo. Como $S$ conmuta con $S^N$,
$S^N(Su)=S(S^Nu)=Su$; así $Su$ es otro punto fijo de $S^N$ y, por unicidad, $Su=u$. Cualquier punto fijo de $S$ lo es de $S^N$, por lo que también es único. Esto desarrolla el mecanismo del ejercicio~\ref{ej:lectura-11-1}.

Finalmente, para cualquier $u_0$, Banach da $S^{kN}u_0\to u$. Para cada residuo $0\le j<N$, la aplicación $S^j$ es continua por~\eqref{ampl:picard-factorial}; luego
$S^{kN+j}u_0\to S^ju=u$. Hay sólo un número finito de residuos, de modo que las $N$ subsucesiones juntas forman una sucesión que converge a $u$. La convergencia es uniforme y concluye la segunda prueba.
\end{proof}

\begin{proposicion}[Un error que puede calcularse]\label{ampl:picard-error}
En el teorema anterior, inicie con $u_0(t)=\xi$ y ponga
$M=\max_{t\in I}\norm{f(t,\xi)}_\infty$. Para $n\ge0$,
\[
 \norm{u-u_n}_\infty
 \le MT\sum_{k=n}^{\infty}\frac{(LT)^k}{(k+1)!}
 \le MT e^{LT}\frac{(LT)^n}{(n+1)!}.
\]
Cuando $L=0$, las iteraciones son exactas desde $u_1$; para $n=0$ se interpreta $(LT)^0=1$.
\end{proposicion}
\begin{proof}
La primera diferencia satisface $\norm{u_1(t)-u_0(t)}_\infty\le M|t-t_0|$. El mismo argumento inductivo por integración da
\[
 \norm{u_{k+1}(t)-u_k(t)}_\infty
 \le\frac{ML^k|t-t_0|^{k+1}}{(k+1)!}.
\]
Sumar desde $k=n$ hasta $m-1$ y dejar $m\to\infty$ produce la primera cota, pues $u_m\to u$ uniformemente. Para la segunda, escribimos $k=n+j$ y usamos
$(n+j+1)!\ge(n+1)!j!$. La suma restante queda acotada por $\sum_{j\ge0}(LT)^j/j!=e^{LT}$. La cota muestra por qué las integraciones repetidas pueden mejorar mucho el control: aparece un factorial en el denominador.
\end{proof}

\section{Ejercicios sobre existencia, iteración y estabilidad}
Los cinco ejercicios de esta sección son desarrollos propios de esta edición a partir del método explicado en \cite{max-picard}; no son una lista reproducida de ese apunte.

\begin{ejercicio}[Ampliación: iteraciones que revelan la solución]\label{ej:ampl-picard-1}
Resuelva $u'=u+t$, $u(0)=1$, en $[0,2]$ mediante Picard. Partiendo de $u_0=1$, encuentre una fórmula para $u_n$ y una cota uniforme explícita del error. Explique por qué la estimación directa del operador no basta para aplicar Banach en la norma uniforme.
\end{ejercicio}


\begin{ejercicio}[Ampliación: una ecuación no lineal en un intervalo largo]\label{ej:ampl-picard-2}
Para $u'=\sin u+t$, $u(0)=0$, en $[0,2]$, verifique las hipótesis de Picard, calcule $u_1,u_2$ desde $u_0=0$ y obtenga cotas del error mediante la norma $\norm v_1=\sup_{0\le t\le2}e^{-t}|v(t)|$ y mediante el factorial.
\end{ejercicio}


\begin{ejercicio}[Ampliación: estabilidad frente a datos perturbados]\label{ej:ampl-picard-3}
Sean $u,v$ las soluciones en $I$ de $u'=f(t,u)$, $v'=g(t,v)$ con $u(t_0)=\xi$, $v(t_0)=\eta$. Suponga las hipótesis de Picard para ambas ecuaciones, que $f$ tiene constante Lipschitz $L$, y que $\norm{f(t,x)-g(t,x)}_\infty\le\delta$ para todo $t,x$. Demuestre, con $r=|t-t_0|$, que para $L>0$,
\[
 \norm{u(t)-v(t)}_\infty
 \le e^{Lr}\norm{\xi-\eta}_\infty
       +\delta\frac{e^{Lr}-1}{L}.
\]
Obtenga la versión para $L=0$ y deduzca dependencia continua respecto del dato inicial.
\end{ejercicio}


\begin{ejercicio}[Ampliación: versión local y explosión de una solución]\label{ej:ampl-picard-4}
Sea $f$ continua en $[t_0-a,t_0+a]\times\{x:\norm{x-\xi}_\infty\le b\}$, con $a,b>0$, acotada allí por $M$ y Lipschitz en $x$ con constante $L$. Si $0<h\le a$, $Mh\le b$ y $Lh<1$, pruebe existencia y unicidad en $[t_0-h,t_0+h]$ entre las soluciones que permanecen en esa bola. Aplique el resultado a $u'=u^2$, $u(0)=1$, y explique por qué no puede concluirse existencia en todo $[0,2]$.
\end{ejercicio}


\begin{ejercicio}[Ampliación: continuidad sin unicidad]\label{ej:ampl-picard-5}
Para el problema $u'=2\sqrt{|u|}$, $u(0)=0$, en $[0,1]$, pruebe que para cada $c\in[0,1]$,
\[
 u_c(t)=\begin{cases}0,&0\le t\le c,\\(t-c)^2,&c\le t\le1\end{cases}
\]
es una solución. Identifique exactamente qué hipótesis de unicidad falta, incluso si se restringe el estado a un intervalo compacto alrededor de cero.
\end{ejercicio}


\section{Cómo elegir una herramienta de convergencia}
La convergencia uniforme conserva la continuidad y permite pasar integrales al límite en intervalos compactos. Para derivar, se necesita controlar las derivadas y fijar un valor. El criterio de Cauchy construye el límite sin conocer su fórmula; Dini añade orden y compacidad a la convergencia puntual. El criterio M usa una cota sumable, mientras Dirichlet aprovecha cancelaciones. Estas diferencias determinan qué hipótesis deben comprobarse. Los capítulos~\ref{cap:nuevo12} y~\ref{cap:nuevo13} aplicarán este control a la compacidad, la extensión y la aproximación de funciones.

\chapter{Equicontinuidad y teorema de Arzelà--Ascoli}\label{cap:nuevo12}\label{chap:27}
\index{equicontinuidad}\index{Arzela--Ascoli, teorema de@Arzelà--Ascoli, teorema de}

\section{Compacidad en un espacio de dimensión infinita}
\label{lectura:27}
En $\R^n$, una sucesión acotada tiene una subsucesión convergente. En el espacio $C([0,1])$ con la norma uniforme, la acotación por sí sola no basta: las funciones $x^n$ tienen norma uno y su único posible límite puntual es discontinuo. Ninguna subsucesión puede converger uniformemente. Necesitamos impedir que las funciones desarrollen cambios bruscos en regiones cada vez más pequeñas.

La condición que controla esos cambios es la \emph{equicontinuidad}. El teorema de Arzelà--Ascoli combina ese control con acotación puntual y proporciona una caracterización de compacidad relativa. Usaremos compacidad métrica, conjuntos densos numerables, Bolzano--Weierstrass y la completitud de $C(X)$ demostrada en el capítulo~\ref{chap:25}.

\section{Un control común de continuidad}
\begin{definicion}
Una familia $\mathcal A$ de funciones $X\to\F$ es \emph{equicontinua en $x$} si para cada $\eps>0$ existe $\delta>0$ tal que
\[
 d(x,y)<\delta\quad\Longrightarrow\quad
 |f(x)-f(y)|<\eps\quad\text{para toda }f\in\mathcal A.
\]
Es \emph{equicontinua} si lo es en cada punto. Es \emph{uniformemente equicontinua} si el mismo $\delta$ puede elegirse para todos los puntos $x$.
\end{definicion}

En esta definición $\delta$ puede depender de $\eps$ y, para la equicontinuidad puntual, de $x$, pero no de la función $f$. Cada miembro de una familia equicontinua es continuo: basta usar el $\delta$ común para ese miembro particular. La definición equivalente con $\sup_{f\in\mathcal A}|f(x)-f(y)|<\eps$ se obtiene ajustando la tolerancia, por ejemplo utilizando primero $\eps/2$.

\begin{ejemplo}[Potencias en el intervalo unidad]
La familia $f_n(t)=t^n$ es equicontinua en todo $x\in[0,1)$ y no lo es en $1$. Para el primer punto, elegimos $r$ con $x<r<1$ y una vecindad relativa de $x$ contenida en $[0,r]$. Dado $\eps>0$, elegimos $N$ con $r^N<\eps/2$. Si $n\ge N$ y $y$ está en esa vecindad, tanto $x^n$ como $y^n$ pertenecen a $[0,\eps/2)$, de modo que $|x^n-y^n|<\eps/2$. Para cada uno de los índices $1\le n<N$, la continuidad de $t^n$ proporciona una vecindad donde $|x^n-y^n|<\eps$. La intersección de esas finitísimas vecindades y de la primera sigue siendo una vecindad de $x$ y funciona para toda la familia. El argumento incluye $x=0$ al trabajar con vecindades relativas.

En cambio, $y_n=2^{-1/n}$ converge a $1$ y
$|f_n(1)-f_n(y_n)|=1/2$. Si hubiera equicontinuidad en $1$ con tolerancia $1/3$, un único radio $\delta$ controlaría a todas las funciones. Para $n$ suficientemente grande, $|1-y_n|<\delta$, lo que produciría la contradicción $1/2<1/3$.
\end{ejemplo}

\begin{proposicion}\label{ch27:equicontuniforme}
Si $X$ es compacto y $\mathcal A\subset C(X,\F)$ es equicontinua, entonces es uniformemente equicontinua.
\end{proposicion}
\begin{proof}
Dado $\eps>0$, para cada $z\in X$ elegimos $r_z>0$ tal que
$d(z,w)<r_z$ implica $|f(z)-f(w)|<\eps/2$ para toda $f\in\mathcal A$. Las bolas $B(z,r_z/2)$ cubren $X$, y la compacidad proporciona una subcubierta finita de centros $z_1,\ldots,z_m$. Sea $\delta=\frac12\min_jr_{z_j}>0$.

Si $d(x,y)<\delta$, elegimos $j$ con $d(x,z_j)<r_{z_j}/2$. Entonces
\[
 d(y,z_j)\le d(y,x)+d(x,z_j)<\delta+r_{z_j}/2\le r_{z_j}.
\]
Para toda $f\in\mathcal A$ tenemos
$|f(x)-f(y)|\le|f(x)-f(z_j)|+|f(z_j)-f(y)|<\eps$. El radio $\delta$ es independiente tanto de $x,y$ como de $f$.
\end{proof}

\begin{proposicion}\label{ch27:acotacioncomun}
Si $X$ es compacto, $\mathcal A$ es equicontinua y cada conjunto de valores $\mathcal A(x)=\{f(x):f\in\mathcal A\}$ es acotado, entonces $\mathcal A$ está acotada en la norma uniforme.
\end{proposicion}
\begin{proof}
Por equicontinuidad, para cada $x$ existe una vecindad $U_x$ donde $|f(y)-f(x)|<1$ para toda $f\in\mathcal A$. Elegimos finitos centros $x_1,\ldots,x_m$ cuyas vecindades cubran $X$. Por acotación puntual existen $M_j<\infty$ tales que $|f(x_j)|\le M_j$ para toda $f$. Para cualquier $y\in X$ elegimos $j$ con $y\in U_{x_j}$ y obtenemos
$|f(y)|\le|f(y)-f(x_j)|+|f(x_j)|<1+\max_jM_j$. Esta cota no depende de $y$ ni de $f$.
\end{proof}

La compacidad del dominio, por sí sola, no convierte la acotación puntual de una familia en acotación uniforme. La prueba utiliza expresamente la equicontinuidad. Por ejemplo, picos continuos con soportes disjuntos alrededor de $2^{-n}$ y alturas $n$ dan una familia puntualmente acotada en $[0,1]$ cuyas normas no están acotadas.

\section{De un conjunto denso a todo el espacio}
\begin{proposicion}\label{ch27:denso}
Sea $(f_n)$ una sucesión equicontinua de funciones $X\to\F$. Si $f_n(y)$ converge para cada punto de un conjunto denso $D\subset X$, entonces $f_n(x)$ converge para todo $x\in X$.
\end{proposicion}
\begin{proof}
Fijemos $x\in X$ y $\eps>0$. Por equicontinuidad elegimos $\delta>0$ con $|f_n(x)-f_n(y)|<\eps/3$ para todo $n$ cuando $d(x,y)<\delta$. Como $D$ es denso, existe $y\in D\cap B(x,\delta)$. La sucesión $(f_n(y))$ converge y, por tanto, es de Cauchy. Existe $N$ con $|f_n(y)-f_m(y)|<\eps/3$ para $n,m\ge N$. Así,
\[
 |f_n(x)-f_m(x)|\le|f_n(x)-f_n(y)|+|f_n(y)-f_m(y)|+|f_m(y)-f_m(x)|<\eps.
\]
La completitud de $\F$ da un límite en $x$. Obsérvese que la diferencia central debe evaluarse en $y$, donde conocemos convergencia. Este resultado no necesita compacidad del dominio.
\end{proof}

\begin{proposicion}\label{ch27:puntualuniforme}
Si $X$ es compacto y $(f_n)$ es equicontinua y converge puntualmente a $f$, entonces $f_n\rightrightarrows f$ y $f$ es continua.
\end{proposicion}
\begin{proof}
Dado $\eps>0$, la proposición~\ref{ch27:equicontuniforme} proporciona $\delta>0$ con $|f_n(x)-f_n(y)|<\eps/4$ para todo $n$ si $d(x,y)<\delta$. El límite puntual en esa desigualdad da $|f(x)-f(y)|\le\eps/4$ para esos puntos.

Por compacidad, existen $x_1,\ldots,x_m$ con $X=\bigcup_jB(x_j,\delta)$. La convergencia en esos finitos puntos permite elegir un único $N$ tal que $|f_n(x_j)-f(x_j)|<\eps/4$ para todo $j$ y $n\ge N$. Para cualquier $x$ elegimos un centro $x_j$ a distancia menor que $\delta$ y obtenemos
\[
 |f_n(x)-f(x)|
 \le |f_n(x)-f_n(x_j)|+|f_n(x_j)-f(x_j)|+|f(x_j)-f(x)|
 \le3\eps/4<\eps.
\]
Esto prueba la convergencia uniforme. La continuidad de $f$ se sigue del teorema~\ref{ch12:limite-continuo}; también se deduce directamente del control obtenido al pasar al límite.
\end{proof}

\section{La extracción diagonal y Arzelà--Ascoli}
\begin{teorema}[Arzelà--Ascoli real]\label{ch27:ascoli}
Sea $X$ un espacio métrico compacto. Una familia $\mathcal A\subset C(X,\R)$ tiene clausura compacta en la norma uniforme si y solo si es equicontinua y $\mathcal A(x)$ es acotado en $\R$ para todo $x\in X$.
\end{teorema}
\begin{estrategia}
La acotación permite extraer una subsucesión convergente en un punto. Un procedimiento diagonal hace que una misma subsucesión converja en un conjunto denso numerable. La equicontinuidad extiende ese control a todos los puntos, y la compacidad del dominio lo convierte en control uniforme. Finalmente hay que comprobar la compacidad de la clausura, no solamente la existencia de subsucesiones para sucesiones de la familia original.
\end{estrategia}
\begin{proof}
Supongamos primero equicontinuidad y acotación puntual. Si $\mathcal A$ es vacía no hay nada que probar. Como $X$ es compacto métrico, para cada $r\ge1$ tiene una red finita $D_r$ de radio $1/r$. El conjunto $D=\bigcup_rD_r$ es a lo sumo numerable y es denso: toda bola de radio positivo contiene un punto de $D_r$ cuando $1/r$ es suficientemente pequeño. Podemos enumerarlo como $d_1,d_2,\ldots$, repitiendo puntos si es finito.

Sea $(f_n)$ una sucesión de $\mathcal A$. La acotación de $\mathcal A(d_1)$ y Bolzano--Weierstrass producen una subsucesión $(f_{n^{(1)}_j})_j$ cuyos valores en $d_1$ convergen. De esa subsucesión extraemos otra $(f_{n^{(2)}_j})_j$ cuyos valores en $d_2$ convergen. La nueva sucesión conserva la convergencia en $d_1$. Por inducción construimos sucesiones de índices crecientes anidadas, de modo que la sucesión de nivel $r$ converge en $d_1,\ldots,d_r$.

Elegimos ahora la diagonal $h_j=f_{n^{(j)}_j}$. Los índices diagonales son estrictamente crecientes: al ser la sucesión del nivel $j+1$ una subsucesión de la del nivel $j$, su término $j+1$ tiene índice al menos $n^{(j)}_{j+1}>n^{(j)}_j$. Para cada $r$ fijo, la cola $(h_j)_{j\ge r}$ está tomada de la sucesión del nivel $r$, en orden creciente. Por ello $(h_j(d_r))$ converge. La subsucesión diagonal converge puntualmente en todo $D$.

La proposición~\ref{ch27:denso} proporciona convergencia puntual en $X$ y la proposición~\ref{ch27:puntualuniforme} proporciona convergencia uniforme a una función continua $h$. La completitud de los escalares interviene al construir los límites de las sucesiones de Cauchy; equivalentemente, podemos ver el último paso como una construcción de una sucesión de Cauchy en el Banach $C(X,\R)$.

Ahora sea $(g_n)$ una sucesión en $\cl{\mathcal A}$. Para cada $n$ elegimos $f_n\in\mathcal A$ con $\norm{g_n-f_n}_\infty<1/n$. Por lo demostrado, una subsucesión $f_{n_j}$ converge uniformemente a $h$. Entonces
\[
 \norm{g_{n_j}-h}_\infty
 \le1/n_j+\norm{f_{n_j}-h}_\infty\longrightarrow0.
\]
Como $\cl{\mathcal A}$ es cerrado, $h\in\cl{\mathcal A}$. Su clausura es secuencialmente compacta y, por ser un espacio métrico, es compacta.

Para la necesidad, supongamos que $K=\cl{\mathcal A}$ es compacto. La evaluación en un punto $x$ es continua, pues $|f(x)-g(x)|\le\norm{f-g}_\infty$. Su imagen $\{f(x):f\in K\}$ es compacta en $\R$, luego acotada, y contiene $\mathcal A(x)$.

Fijemos $x\in X$ y $\eps>0$. Por compacidad de $K$ existen $h_1,\ldots,h_r\in K$ tales que las bolas de radio $\eps/3$ y esos centros cubren $K$. Cada $h_j$ es continua en $x$, así que existe $\delta_j>0$ con $|h_j(x)-h_j(y)|<\eps/3$ si $d(x,y)<\delta_j$. Sea $\delta=\min_j\delta_j$. Para $f\in\mathcal A$ escogemos $h_j$ con $\norm{f-h_j}_\infty<\eps/3$. Si $d(x,y)<\delta$, obtenemos
\[
 |f(x)-f(y)|
 \le |f(x)-h_j(x)|+|h_j(x)-h_j(y)|+|h_j(y)-f(y)|<\eps.
\]
El radio $\delta$ es independiente de $f$, y esto demuestra la equicontinuidad.
\end{proof}

\begin{observacion}[Dónde vive el conjunto compacto]
La compacidad relativa del teorema corresponde al espacio de funciones $C(X,\R)$ con su norma uniforme. La familia $\mathcal A$ no es un subconjunto del dominio $X$. Esta distinción corrige una imprecisión de la formulación original.
\end{observacion}

\section{Ejercicios sobre familias equicontinuas y aplicaciones compactas}
En estos ejercicios, $X$ es un espacio métrico compacto, salvo indicación contraria.
\begin{ejercicio}[Lectura 27, ejercicio 1]\label{ej:lectura-27-1}
Demuestre que $\mathcal A\subset C(X,\R)$ es compacta si y solo si es cerrada, acotada en norma y equicontinua.
\end{ejercicio}


\begin{ejercicio}[Lectura 27, ejercicio 2]\label{ej:lectura-27-2}
Sea $v:C(X,\R)\times X\to\R$ la evaluación $v(f,x)=f(x)$.
\begin{enumerate}[label=(\roman*)]
\item Demuestre que $v$ es continua.
\item Para $\mathcal A\subset C(X,\R)$ y $\eps>0$, demuestre que
\[
 U_\eps=\{((f,x),(g,y)):|f(x)-g(y)|<\eps\}
\]
es abierto en $(\mathcal A\times X)^2$ y contiene la diagonal. Si $\mathcal A$ es compacta, deduzca la necesidad de Arzelà--Ascoli.
\end{enumerate}
\end{ejercicio}


\begin{ejercicio}[Lectura 27, ejercicio 3]\label{ej:lectura-27-3}
En $C^1([a,b],\R)$ consideramos $\norm f_{C^1}=\norm f_\infty+\norm{f'}_\infty$. Demuestre que la inclusión $J:C^1([a,b])\to C([a,b])$, $Jf=f$, es lineal y compacta.
\end{ejercicio}


\begin{ejercicio}[Lectura 27, ejercicio 4]\label{ej:lectura-27-4}
Sea $\mathcal A\subset C([a,b],\R)$ una familia cuyas funciones son derivables en $(a,b)$. Si $|f'(x)|\le M$ para todas las funciones y todos los puntos interiores, demuestre que $\mathcal A$ es equicontinua.
\end{ejercicio}


\begin{ejercicio}[Lectura 27, ejercicio 5: versión compleja]\label{ej:lectura-27-5}
Sea $\mathcal A\subset C(X,\C)$ y ponga $\mathcal A_1=\{\operatorname{Re}f:f\in\mathcal A\}$, $\mathcal A_2=\{\operatorname{Im}f:f\in\mathcal A\}$.
\begin{enumerate}[label=(\roman*)]
\item Pruebe que $\mathcal A$ es equicontinua si y solo si lo son $\mathcal A_1$ y $\mathcal A_2$.
\item Demuestre que $\Phi:C(X,\R)\times C(X,\R)\to C(X,\C)$, $\Phi(u,v)=u+iv$, es continua.
\item Deduzca Arzelà--Ascoli complejo: $\mathcal A$ es relativamente compacta si y solo si es equicontinua y cada $\mathcal A(x)$ es acotado.
\end{enumerate}
\end{ejercicio}


\section{La información que se vuelve finita}
Arzelà--Ascoli convierte dos controles sobre una familia infinita en una propiedad de extracción: de cualquier sucesión puede elegirse una subsucesión convergente uniformemente. El mecanismo combina un conjunto numerable de puntos para la extracción diagonal y una cubierta finita para controlar el error en todo el dominio. Las aplicaciones a operadores compactos, como la inclusión de $C^1$ en $C$, descansan en esa reducción y no en una supuesta extensión de Heine--Borel a todos los espacios normados.

\chapter{Extensión y aproximación de funciones continuas}\label{cap:nuevo13}\label{chap:12}\label{chap:28}

Una función definida en una parte de un espacio puede contener información que queremos conservar al ampliar el dominio. La continuidad uniforme permite extenderla a una clausura cuando el espacio de llegada es completo. Tietze resuelve otro problema: extender una función real continua desde un cerrado a todo un espacio métrico, sin exigir continuidad uniforme. La construcción combina separación de cerrados con una serie de correcciones de tamaño decreciente.

Una vez estudiada la extensión, buscaremos aproximaciones uniformes mediante funciones de una forma prescrita. Stone--Weierstrass conecta esa pregunta con las operaciones de un álgebra y la capacidad de distinguir puntos. Los resultados de convergencia uniforme del capítulo~\ref{cap:nuevo11} permiten desarrollar ambas construcciones con control explícito del error.

\section{Separar cerrados mediante funciones}
\label{lectura:12}
\begin{teorema}[Lema de Urysohn para espacios métricos]\label{ch12:urysohn}
Si $A,B$ son cerrados disjuntos de un espacio métrico $X$, existe una función continua $u:X\to[0,1]$ tal que $u|_A=0$ y $u|_B=1$.
\end{teorema}
\begin{proof}
Si $A=\varnothing$, sirve $u=1$; si $B=\varnothing$ y $A\ne\varnothing$, sirve $u=0$. Cuando ambos son no vacíos, definimos
\[
 u(x)=\frac{\dist(x,A)}{\dist(x,A)+\dist(x,B)}.
\]
Las funciones de distancia son continuas, incluso $1$-Lipschitz. El denominador no se anula: si ambas distancias fueran cero, el carácter cerrado de los conjuntos obligaría a $x\in A\cap B$. Por tanto el cociente es continuo y está entre cero y uno. En $A$, el numerador vale cero; en $B$, la distancia a $B$ es cero y la de $A$ es positiva, por lo que el cociente vale uno. Esto también permite separar los cerrados por los abiertos $u^{-1}([0,1/3))$ y $u^{-1}((2/3,1])$, entendiendo los intervalos como abiertos relativos de $[0,1]$.
\end{proof}
\index{Urysohn, lema de}\index{separacion de cerrados@separación de cerrados}
Los cerrados pueden tener distancia mutua cero; no se exige una separación uniforme. Lo que se necesita es que, para cada punto individual, no se anulen simultáneamente las dos distancias.

\begin{lema}[Una corrección que reduce el error]\label{ch12:correccion}
Sea $V$ cerrado en $X$, y sea $g:V\to\R$ continua con $|g|\leq\alpha$, donde $\alpha>0$. Existe una función continua $h:X\to\R$ tal que
\[
 |h(x)|\leq\alpha/3\quad(x\in X),\qquad
 |g(x)-h(x)|\leq2\alpha/3\quad(x\in V).
\]
\end{lema}
\begin{proof}
Los conjuntos
\[
 A=\{x\in V:g(x)\leq-\alpha/3\},\qquad
 B=\{x\in V:g(x)\geq\alpha/3\}
\]
son cerrados en $V$ por continuidad, y por tanto cerrados en $X$ porque $V$ lo es. Son disjuntos. Sea $u$ la función del lema de Urysohn con $u=0$ en $A$ y $u=1$ en $B$, y pongamos $h=(\alpha/3)(2u-1)$. Entonces $|h|\leq\alpha/3$, $h=-\alpha/3$ en $A$ y $h=\alpha/3$ en $B$.

Si $x\in A$, los valores de $g(x)$ están entre $-\alpha$ y $-\alpha/3$, de modo que $|g(x)+\alpha/3|\leq2\alpha/3$. En $B$, $g(x)$ está entre $\alpha/3$ y $\alpha$, y $|g(x)-\alpha/3|\leq2\alpha/3$. En $V\setminus(A\cup B)$ se tiene $|g(x)|<\alpha/3$ y $|h(x)|\leq\alpha/3$, así que la desigualdad triangular da el mismo resultado. Estos tres casos cubren todo $V$.
\end{proof}

Para sumar las correcciones usaremos los resultados del capítulo~\ref{cap:nuevo11}. El criterio M de Weierstrass, teorema~\ref{ch12:weierstrass}, garantiza convergencia uniforme si las cotas de los términos forman una serie numérica convergente, y controla el error por la cola de esa serie. El teorema~\ref{ch12:limite-continuo} asegura después que el límite de las sumas parciales continuas es continuo. En ambos pasos es esencial que el índice que controla el error no dependa del punto. Estos resultados ya están demostrados, también sin exigir acotación global en el teorema de continuidad.

\section{Demostración completa del teorema de extensión}
\begin{teorema}[Tietze]\label{ch12:tietze}
Sea $V$ un cerrado de un espacio métrico $X$. Toda función continua $g:V\to\R$ tiene una extensión continua $f:X\to\R$. Si $g(V)\subset[a,b]$, puede elegirse una extensión con valores en el mismo intervalo.
\end{teorema}
\index{Tietze, teorema de}\index{extension de funciones@extensión de funciones}
\begin{proof}
Si $V=\varnothing$, una función constante resuelve el problema; para la afirmación del intervalo elegimos su valor en $[a,b]$. Si $X=V$, sirve $g$. Tratamos el caso restante.

\textbf{Caso acotado y simétrico.} Supongamos $|g|\leq\alpha$ para algún $\alpha>0$. Aplicamos el lema~\ref{ch12:correccion} a $g$ y obtenemos $h_1$ con $|h_1|\leq\alpha/3$ en $X$ y error a lo sumo $2\alpha/3$ en $V$. El residuo $r_1=g-h_1|_V$ sigue siendo continuo. Aplicamos el mismo lema al residuo con cota $\alpha_1=2\alpha/3$ y obtenemos $h_2$ con
\[
 |h_2|\leq\frac\alpha3\frac23,
 \qquad |g-h_1-h_2|\leq\alpha(2/3)^2\text{ en }V.
\]
Inductivamente, si $h_1,\ldots,h_{n-1}$ ya están construidas, aplicamos el lema al residuo continuo $g-\sum_{k=1}^{n-1}h_k|_V$, cuya cota es $\alpha(2/3)^{n-1}$. Esto produce $h_n$ continua en $X$ con
\[
 |h_n(x)|\leq\frac\alpha3(2/3)^{n-1}\quad(x\in X),
 \qquad
 \left|g(x)-\sum_{k=1}^nh_k(x)\right|\leq\alpha(2/3)^n\quad(x\in V).
\]
La suma del error es \emph{finita} en cada etapa; ése es el residuo al que se aplica el siguiente paso. Como
\[
 \sum_{n=1}^\infty\frac\alpha3(2/3)^{n-1}=\alpha,
\]
el criterio M, teorema~\ref{ch12:weierstrass}, asegura que $f=\sum_nh_n$ converge uniformemente en $X$. Sus sumas parciales son continuas y el teorema~\ref{ch12:limite-continuo} muestra que $f$ es continua. La segunda estimación, al dejar $n\to\infty$, da $f|_V=g$. La primera, sumada, da $|f|\leq\alpha$ en todo $X$.

\textbf{Valores en un intervalo.} Si $a=b$, $g$ es constante y la misma constante la extiende. Si $a<b$, aplicamos el caso anterior a
\[
 g_0=g-\frac{a+b}{2},\qquad |g_0|\leq\frac{b-a}{2}.
\]
Su extensión $f_0$ con la misma cota produce $f=f_0+(a+b)/2$, cuyos valores están en $[a,b]$ y que coincide con $g$ en $V$.

\textbf{Caso no acotado.} La función $h=\arctan g$ es continua en $V$ y toma valores estrictamente entre $-\pi/2$ y $\pi/2$. Por el caso acotado tiene una extensión continua $H:X\to[-\pi/2,\pi/2]$. No podemos aplicar directamente $\tan H$: fuera de $V$, la extensión podría alcanzar los extremos, donde la tangente no está definida. Para evitarlo, definimos
\[
 w(x)=\frac1{1+\dist(x,V)},\qquad \widetilde H(x)=w(x)H(x).
\]
Como $V$ es cerrado y no vacío, $w$ es continua, vale $1$ en $V$ y satisface $0<w(x)<1$ fuera de $V$. En $V$, $|H|<\pi/2$ porque $H=h$; fuera de $V$, el factor $w<1$ garantiza $|\widetilde H|<\pi/2$. Por ello
\[
 f(x)=\tan(\widetilde H(x))
\]
es una función real continua en todo $X$, y en $V$ cumple $f=\tan(\arctan g)=g$. Esto completa el caso que requiere especial cuidado con los extremos del intervalo.
\end{proof}

\begin{ejemplo}[Por qué se necesita un cerrado]
La función $g:(0,\infty)\to\R$, $g(x)=\cos(1/x)$, es continua y está acotada por uno. Sin embargo, para
\[
 x_n=1/(2\pi n),\qquad y_n=1/((2n+1)\pi)
\]
ambas sucesiones tienden a cero y $g(x_n)=1$, $g(y_n)=-1$. Una extensión continua a $\R$ tendría que tener en cero dos límites distintos. La condición de que el dominio original sea cerrado no puede simplemente omitirse.
\end{ejemplo}

\section{Ejercicios sobre separación, soporte y particiones de la unidad}
\begin{ejercicio}[Lectura 12, ejercicio 1]\label{ej:lectura-12-1}
Sin usar el lema de Urysohn, separe dos cerrados disjuntos $A,B\subset X$ por abiertos disjuntos que los contengan.
\end{ejercicio}


\begin{ejercicio}[Lectura 12, ejercicio 2]\label{ej:lectura-12-2}\label{ch12:soporte}
Defina el soporte de una función por
\[
 \supp(f)=\cl{\{x\in X:f(x)\ne0\}}.
\]
Sean $V$ cerrado y $U$ abierto con $V\subset U$.
\begin{enumerate}[label=(\roman*)]
\item Construya $\varphi:X\to[0,1]$ continua, igual a $1$ en $V$, y con $\supp(\varphi)\subset U$.
\item Examine la afirmación de que toda $g:V\to[a,b]$ continua tiene una extensión continua $h:X\to[a,b]$ con soporte en $U$, y resuelva la versión válida.
\end{enumerate}
\end{ejercicio}


\begin{ejercicio}[Lectura 12, ejercicio 3]\label{ej:lectura-12-3}
Sean $U_1,\ldots,U_n$ abiertos que cubren $X$. Construya funciones continuas $\varphi_i:X\to[0,1]$ con $\supp(\varphi_i)\subset U_i$ y $\sum_{i=1}^n\varphi_i=1$.
\end{ejercicio}

\index{particion de la unidad@partición de la unidad}\index{soporte de una funcion@soporte de una función}

Las funciones de separación permiten fijar valores sobre cerrados; las series uniformemente convergentes convierten aproximaciones en una extensión exacta; y las funciones de soporte controlado permiten localizar esa extensión. Una partición de la unidad combina estas ideas para repartir una construcción global entre distintos abiertos sin perder continuidad. Pasaremos ahora de conservar valores en una parte del dominio a aproximar una función en todo un compacto mediante una familia algebraicamente sencilla.

\index{Stone--Weierstrass, teorema de}\index{algebra@álgebra!de funciones}

\section{Cuándo bastan las funciones sencillas}
\label{lectura:28}
El teorema de aproximación de Weierstrass afirma que una función continua en un intervalo compacto puede aproximarse uniformemente por polinomios. Stone--Weierstrass identifica qué propiedades de los polinomios hacen posible esa aproximación: pueden sumarse, multiplicarse y escalarse; distinguen puntos distintos; y no se anulan todos en un mismo punto.

Estudiaremos primero funciones reales. La prueba transformará esas propiedades algebraicas en la posibilidad de construir máximos y mínimos aproximados. Dos aplicaciones de compacidad permitirán ensamblar aproximaciones locales hasta obtener una aproximación en todo el dominio. Para el caso complejo aparecerá una condición adicional: estabilidad bajo conjugación.

Usaremos la norma uniforme, la continuidad de sumas y productos, el teorema de Dini del capítulo~\ref{chap:26} y la compacidad. Construiremos explícitamente los polinomios que aproximan al valor absoluto para evitar utilizar el propio teorema que queremos demostrar.

\section{Álgebras y ejemplos fundamentales}
\begin{definicion}
Un \emph{álgebra asociativa} sobre $\F$ es un espacio vectorial $A$ con una aplicación bilineal $(a,b)\mapsto ab$ de $A\times A$ en $A$ que satisface $(ab)c=a(bc)$. La bilinealidad significa
\[
 (a+b)c=ac+bc,\quad a(b+c)=ab+ac,\quad
 (\lambda a)b=a(\lambda b)=\lambda(ab).
\]
El álgebra es \emph{conmutativa} si $ab=ba$. Es \emph{unitaria} si tiene un elemento $1$ tal que $1a=a1=a$ para todo $a$. Una norma es \emph{submultiplicativa} si $\norm{ab}\le\norm a\norm b$; un álgebra con esa norma es un \emph{álgebra normada}. Si además es completa, es un \emph{álgebra de Banach}.
\end{definicion}
\index{algebra@álgebra!normada}\index{algebra@álgebra!de Banach}

Pedimos que la unidad actúe por ambos lados. Una identidad solamente a la derecha no es, en general, suficiente en un álgebra no conmutativa. Si existen dos unidades bilaterales $e,e'$, entonces $e=ee'=e'$, de modo que la unidad es única.

\begin{ejemplo}[Escalares y matrices]
$\R$ y $\C$ son álgebras conmutativas unitarias sobre sí mismas; $\C$ es también un álgebra sobre $\R$. Las matrices $M_n(\F)$ forman un álgebra con la suma, el producto por escalar y el producto matricial. La asociatividad se puede verificar calculando cada entrada:
\[
 ((AB)C)_{ij}=\sum_k\sum_\ell a_{i\ell}b_{\ell k}c_{kj}
 =\sum_\ell a_{i\ell}\sum_kb_{\ell k}c_{kj}=(A(BC))_{ij}.
\]
La identidad es la matriz $I$. Si $n\ge2$, las matrices elementales satisfacen $E_{12}E_{21}=E_{11}$ y $E_{21}E_{12}=E_{22}$, así que el producto no es conmutativo. Para $n=1$ sí lo es. La norma de operador es submultiplicativa porque $\norm{ABx}\le\norm A\norm B\norm x$.
\end{ejemplo}

\begin{ejemplo}[Cuaternios: un ejemplo asociativo no conmutativo]
\index{cuaternios}
Los cuaternios son las expresiones $q=a+bi+cj+dk$ con coeficientes reales y reglas
\[
 i^2=j^2=k^2=-1,\qquad ij=k,\ jk=i,\ ki=j,
 \qquad ji=-k,\ kj=-i,\ ik=-j.
\]
Estas reglas determinan el producto por bilinealidad. Para comprobar su consistencia y asociatividad sin multiplicar todas las combinaciones de tres elementos, escribimos $z=a+bi$, $w=c+di$ y representamos
\[
 \rho(q)=\begin{pmatrix}z&w\\-\overline w&\overline z\end{pmatrix}.
\]
El conjunto de estas matrices es un subespacio vectorial real y está cerrado bajo producto, pues
\[
 \begin{pmatrix}z&w\\-\overline w&\overline z\end{pmatrix}
 \begin{pmatrix}z'&w'\\-\overline {w'}&\overline {z'}\end{pmatrix}
 =\begin{pmatrix}
 zz'-w\overline {w'}&zw'+w\overline {z'}\\
 -\overline{zw'+w\overline {z'}}&\overline{zz'-w\overline {w'}}
 \end{pmatrix}.
\]
Las matrices de $i,j,k$ satisfacen las reglas indicadas y $\rho$ es inyectiva. La asociatividad del producto de matrices se transmite al producto de cuaternios. El cuaternio conjugado es $\overline q=a-bi-cj-dk$ y la multiplicación da
$q\overline q=\overline q q=(a^2+b^2+c^2+d^2)1$. Así, si $q\ne0$,
\[
 q^{-1}=\frac{\overline q}{a^2+b^2+c^2+d^2}.
\]
Todo elemento no nulo es invertible, pero $ij=-ji$, por lo que esta álgebra de división no es un campo conmutativo. La norma euclidiana es multiplicativa: como $\det\rho(q)=a^2+b^2+c^2+d^2$, la igualdad de determinantes de un producto implica $\norm{qr}_2^2=\norm q_2^2\norm r_2^2$.
\end{ejemplo}

\begin{ejemplo}[Operadores y grupo de invertibles]
Para un espacio normado $V$, los operadores lineales acotados $\mathcal B(V,V)$ forman un álgebra con el producto dado por composición. La asociatividad es la de la composición; las distributividades provienen de la linealidad, y
\[
 \norm{STv}\le\norm S\norm{Tv}\le\norm S\norm T\norm v
 \quad\Longrightarrow\quad\norm{ST}\le\norm S\norm T.
\]
La unidad es $I$. Si $V$ es Banach, el espacio de operadores es completo y constituye un álgebra de Banach.

En cambio,
\[
 \operatorname{GL}(V)=\{T\in\mathcal B(V,V):T\text{ tiene inverso acotado}\}
\]
es un \emph{grupo} bajo composición, no un álgebra cuando $V\ne\{0\}$. En efecto, $I$ y $-I$ son invertibles y su suma es el operador cero, que no es invertible. Ni siquiera es un subespacio vectorial. La cerradura bajo composición se verifica mediante $(ST)^{-1}=T^{-1}S^{-1}$, y ese inverso es acotado por submultiplicatividad. Esto corrige el ejemplo 2(iv) de la lectura original.
\end{ejemplo}

\begin{ejemplo}[Funciones continuas]
Para un espacio métrico $X$, las funciones continuas $C(X,\F)$ forman un álgebra con operaciones puntuales, pues todas las identidades se verifican en cada $x$. Sin embargo, la norma uniforme puede ser infinita si $X$ no es compacto; por ejemplo, $f(x)=x$ en $\R$. El álgebra normada adecuada es $C_b(X,\F)$, donde
\[
 |f(x)g(x)|\le\norm f_\infty\norm g_\infty,
 \qquad\norm{fg}_\infty\le\norm f_\infty\norm g_\infty.
\]
Es un álgebra de Banach por el capítulo~\ref{chap:25}; es conmutativa y tiene por unidad la función constante $1$. Si $X$ es compacto, toda función continua es acotada, y entonces $C(X,\F)=C_b(X,\F)$.
\end{ejemplo}

\begin{definicion}
Una \emph{subálgebra} $A$ de un álgebra $B$ es un subespacio vectorial cerrado bajo multiplicación. No exigimos que contenga la unidad de $B$. Una familia de funciones $A$ \emph{separa puntos} de $X$ si, para $x\ne y$, existe $h\in A$ con $h(x)\ne h(y)$. Decimos que $A$ \emph{no se anula en ningún punto} si para cada $x$ existe $h\in A$ con $h(x)\ne0$.
\end{definicion}
\index{separacion de puntos@separación de puntos}\index{subalgebra@subálgebra}

Los polinomios restringidos a un compacto de $\R$ forman una subálgebra, porque la suma y el producto de polinomios son polinomios. Incluyendo el término de grado cero tenemos la función constante $1$, que garantiza que no hay ceros comunes, y la función $x\mapsto x$ separa puntos. Si solamente se toman potencias de grado positivo y el compacto contiene a cero, todas las funciones de la subálgebra se anulan en ese punto: esa subálgebra no puede ser densa en todo $C(X)$.

\section{Aproximación del valor absoluto sin circularidad}
\begin{lema}\label{ch28:abspolinomios}
Para cada $\alpha\ge0$ existen polinomios reales $p_n$ con $p_n(0)=0$ que convergen uniformemente a $|t|$ en $[-\alpha,\alpha]$.
\end{lema}
\begin{proof}
Si $\alpha=0$, basta $p_n=0$. Construiremos primero polinomios en $[-1,1]$. Definimos
\[
 q_0(t)=0,\qquad q_{n+1}(t)=q_n(t)+\frac{t^2-q_n(t)^2}{2}.
\]
Por inducción, cada $q_n$ es un polinomio y $q_n(0)=0$. Fijemos $t\in[-1,1]$ y escribamos $r=|t|$. Probamos inductivamente que $0\le q_n(t)\le r$ y $q_{n+1}(t)\ge q_n(t)$. Si $0\le q_n(t)\le r$, entonces
\[
 q_{n+1}(t)-q_n(t)=\frac{r^2-q_n(t)^2}{2}\ge0,
\]
y
\[
 r-q_{n+1}(t)
 =(r-q_n(t))\left(1-\frac{r+q_n(t)}2\right)\ge0,
\]
porque $0\le r+q_n(t)\le2$. El caso inicial $q_0=0$ cumple las desigualdades.

Para cada $t$, la sucesión numérica $q_n(t)$ es creciente y está acotada por $r$, por lo que converge a algún $L(t)\in[0,r]$. Al pasar al límite en la recurrencia obtenemos
$L(t)=L(t)+(t^2-L(t)^2)/2$, de modo que $L(t)^2=t^2$. Como $L(t)\ge0$, se sigue que $L(t)=|t|$.

Hemos construido una sucesión creciente de funciones continuas en el compacto $[-1,1]$ cuyo límite $|t|$ es continuo. El teorema de Dini~\ref{ch26:dini} demuestra que la convergencia es uniforme. Para $\alpha>0$, ponemos $p_n(t)=\alpha q_n(t/\alpha)$. Entonces $p_n(0)=0$ y
\[
 \sup_{|t|\le\alpha}|p_n(t)-|t||
 =\alpha\sup_{|u|\le1}|q_n(u)-|u||\longrightarrow0.
\]
\end{proof}

La condición $p_n(0)=0$ tiene una finalidad algebraica: permite escribir $p_n(f)$ usando solo potencias positivas de $f$. Así permaneceremos dentro de una subálgebra aunque no contenga la función constante $1$.

\section{Interpolación en dos puntos}
\begin{lema}\label{ch28:interpolacion}
Si una subálgebra $A$ de funciones $S\to\F$ separa puntos y no se anula en ningún punto, entonces para $x\ne y$ y cualesquiera $\alpha,\beta\in\F$ existe $h\in A$ con $h(x)=\alpha$ y $h(y)=\beta$.
\end{lema}
\begin{proof}
Elegimos $g\in A$ con $g(x)\ne g(y)$, $a\in A$ con $a(x)\ne0$ y $b\in A$ con $b(y)\ne0$. Definimos
\[
 u=gb-g(x)b,\qquad v=ga-g(y)a.
\]
Ambas funciones pertenecen a $A$ porque se obtienen por productos, sumas y múltiplos escalares. No estamos suponiendo que las funciones constantes pertenezcan a $A$: $g(x)b$ es un múltiplo escalar de $b$. Calculando sus valores,
\[
 u(x)=0,\quad u(y)=(g(y)-g(x))b(y)\ne0,
\]
\[
 v(y)=0,\quad v(x)=(g(x)-g(y))a(x)\ne0.
\]
Por tanto,
$h=\alpha v/v(x)+\beta u/u(y)$ pertenece a $A$ y tiene los valores prescritos. No se requiere $\alpha\ne\beta$. Para un solo punto $x$ y un valor $\alpha$, basta elegir $a(x)\ne0$ y usar $\alpha a/a(x)$.
\end{proof}

\section{Demostración completa de Stone--Weierstrass real}
\begin{teorema}\label{ch28:stone}
Sea $X$ un espacio métrico compacto. Si $A$ es una subálgebra de $C(X,\R)$ que separa puntos y no se anula en ningún punto, entonces $\cl A=C(X,\R)$ en la norma uniforme.
\end{teorema}
\begin{proof}
Escribamos $B=\cl A$. Dividimos la prueba en pasos que tienen funciones distintas.

\emph{Paso 1: $B$ es una subálgebra cerrada.}
Para $f,g\in B$ existen sucesiones $f_n,g_n\in A$ con $f_n\to f$ y $g_n\to g$ uniformemente. Las estimaciones
\[
 \norm{(f_n+g_n)-(f+g)}_\infty
 \le\norm{f_n-f}_\infty+\norm{g_n-g}_\infty,
 \qquad\norm{\lambda f_n-\lambda f}_\infty=|\lambda|\norm{f_n-f}_\infty
\]
demuestran que $B$ es un subespacio vectorial. Como $(f_n)$ converge en norma, sus normas están acotadas por algún $M$. Así,
\[
 \norm{f_ng_n-fg}_\infty
 \le M\norm{g_n-g}_\infty+\norm g_\infty\norm{f_n-f}_\infty\longrightarrow0.
\]
Cada $f_ng_n$ pertenece a $A$, de modo que $fg\in B$. La cerradura de $B$ es parte de su definición.

\emph{Paso 2: $B$ contiene valores absolutos y máximos y mínimos finitos.}
Sea $f\in B$ y $\alpha=\norm f_\infty$. El lema~\ref{ch28:abspolinomios} proporciona polinomios $p_n$ sin término constante que aproximan uniformemente al valor absoluto en $[-\alpha,\alpha]$. Cada función $p_n(f)$ pertenece a $B$ porque es combinación lineal de potencias positivas de $f$. Además,
\[
 \norm{p_n(f)-|f|}_\infty
 \le\sup_{|t|\le\alpha}|p_n(t)-|t||\longrightarrow0.
\]
La cerradura de $B$ implica $|f|\in B$. Para $f,g\in B$ utilizamos las identidades puntuales
\[
 \max(f,g)=\frac{f+g+|f-g|}{2},\qquad
 \min(f,g)=\frac{f+g-|f-g|}{2}.
\]
Ambas pertenecen a $B$. Aplicando reiteradamente las identidades se obtienen máximos y mínimos de cualquier familia finita.

\emph{Paso 3: una aproximación global por abajo que coincide en un punto.}
Fijemos $f\in C(X,\R)$, una tolerancia $\eta>0$ y un punto $x\in X$. Para cada $y\in X$, el lema~\ref{ch28:interpolacion} proporciona $h_{x,y}\in A$ con
\[
 h_{x,y}(x)=f(x),\qquad h_{x,y}(y)=f(y).
\]
Si $y=x$ basta la versión de interpolación en un solo punto. El conjunto
\[
 U_y=\{z\in X:h_{x,y}(z)>f(z)-\eta\}
\]
es abierto por continuidad y contiene a $y$. Por compacidad, finitos de ellos, $U_{y_1},\ldots,U_{y_r}$, cubren $X$. Definimos
$g_x=\max\{h_{x,y_1},\ldots,h_{x,y_r}\}\in B$. Todos los términos valen $f(x)$ en $x$, por lo que $g_x(x)=f(x)$. Para cualquier $z\in X$, alguno de los abiertos elegidos contiene a $z$ y entonces
\[
 g_x(z)\ge h_{x,y_j}(z)>f(z)-\eta.
\]
La comparación es con $f(z)$, el valor de la función en el punto variable, y no con la constante $f(x)$.

\emph{Paso 4: controlar también el error por arriba.}
Como $g_x(x)=f(x)$, el conjunto
$V_x=\{z:g_x(z)<f(z)+\eta\}$ es abierto y contiene a $x$. Otra aplicación de compacidad produce una subcubierta $V_{x_1},\ldots,V_{x_m}$ de $X$. Ponemos
\[
 h=\min\{g_{x_1},\ldots,g_{x_m}\}\in B.
\]
Cada $g_{x_j}(z)$ es mayor que $f(z)-\eta$, luego también lo es su mínimo. Por otro lado, para cada $z$ existe $j$ con $z\in V_{x_j}$, y por ello
$h(z)\le g_{x_j}(z)<f(z)+\eta$. Hemos probado
\[
 f(z)-\eta<h(z)<f(z)+\eta\quad(z\in X),
 \qquad\norm{h-f}_\infty\le\eta.
\]

Para concluir, dado $\eps>0$ elegimos $\eta=\eps/3$ en la construcción. Como $h\in B=\cl A$, existe $a\in A$ con $\norm{a-h}_\infty<\eps/3$. Entonces
\[
 \norm{a-f}_\infty\le\norm{a-h}_\infty+\norm{h-f}_\infty
 <2\eps/3<\eps.
\]
Así cada función continua puede aproximarse arbitrariamente por elementos de $A$, como se quería.
\end{proof}

La prueba utiliza máximos y mínimos en la \emph{clausura} del álgebra. No afirma que el máximo de dos polinomios sea un polinomio. La última aproximación de $h\in\cl A$ por un elemento de $A$ devuelve el resultado a la familia original.

\begin{ejemplo}[Qué impide la densidad]
Si todas las funciones de $A$ satisfacen $h(x_0)=0$, todo límite uniforme tiene esa propiedad, pues la evaluación es continua. La función constante $1$ no puede estar en la clausura. Si todas satisfacen $h(x_0)=h(y_0)$ para dos puntos distintos, sus límites uniformes también satisfacen esa igualdad. En un espacio métrico, la función continua $z\mapsto d(z,x_0)$ toma valores distintos en esos dos puntos y, si el espacio es compacto, pertenece a $C(X)$. Por ello la falta de separación impide densidad. Estas observaciones explican las dos hipótesis algebraicas del teorema.
\end{ejemplo}

\section{Ejercicios sobre aproximación polinómica y álgebras complejas}
\begin{ejercicio}[Lectura 28, ejercicio 1: aproximación de Weierstrass]\label{ej:lectura-28-1}
Sea $f:[a,b]\to\F$ continua. Demuestre que existe una sucesión de polinomios con coeficientes en $\F$ que converge uniformemente a $f$.
\end{ejercicio}


\begin{ejercicio}[Lectura 28, ejercicio 2: determinación por momentos]\label{ej:lectura-28-2}
Sea $f:[0,1]\to\R$ continua y suponga que $\int_0^1 f(x)x^n\,dx=0$ para todo entero $n\ge0$. Demuestre que $f=0$.
\end{ejercicio}


\begin{ejercicio}[Lectura 28, ejercicio 3: Stone--Weierstrass complejo]\label{ej:lectura-28-3}
Sea $X$ un espacio métrico compacto. Una subálgebra compleja $A\subset C(X,\C)$ es \emph{autoadjunta} si $\overline f\in A$ para toda $f\in A$. Demuestre que, si $A$ es autoadjunta, separa puntos y no se anula en ningún punto, entonces es densa en $C(X,\C)$.
\end{ejercicio}


\begin{ejercicio}[Lectura 28, ejercicio 4: por qué hace falta conjugación]\label{ej:lectura-28-4}
Sea $\mathbb T=\{z\in\C:|z|=1\}$ y sea $A$ el álgebra de restricciones a $\mathbb T$ de los polinomios complejos en $z$:
\[
 p(e^{i\theta})=\sum_{k=0}^N a_ke^{ik\theta}.
\]
Demuestre que $A$ separa puntos y no se anula en ningún punto, pero no es densa en $C(\mathbb T,\C)$.
\end{ejercicio}


\section{Extender, localizar y aproximar}
Tietze construye una extensión exacta mediante correcciones cuyo error decrece geométricamente; las funciones de soporte controlado permiten localizarla. Stone--Weierstrass resuelve un problema complementario: aproximar una función continua por miembros de una subálgebra. Su argumento crea primero interpolaciones exactas en dos puntos y después las ensambla mediante máximos y mínimos. La aproximación polinómica del valor absoluto permite realizar ese ensamblaje dentro de la clausura de una subálgebra. En el caso complejo, la conjugación recupera partes reales e imaginarias y permite usar el teorema real. El ejercicio~\ref{ej:lectura-28-4} muestra que distinguir puntos no basta si se pierde esa operación. Esta relación entre álgebra y norma uniforme será una herramienta central para las álgebras de operadores y para el análisis armónico de los volúmenes siguientes.


\parte{Geometría de Hilbert}{%
El producto interno permite medir ángulos y expresar la ortogonalidad. En un espacio de Hilbert, la completitud convierte esa geometría en un método de mejor aproximación: la proyección sobre un subespacio cerrado se obtiene como un límite y se caracteriza mediante perpendicularidad.

Las bases ortonormales describen vectores mediante coeficientes de Fourier. La representación de Riesz relaciona esos vectores con los funcionales continuos y conduce al operador adjunto. Esta parte utiliza espacios normados, series y operadores de las partes anteriores, manteniendo la convención de linealidad del producto interno en el primer argumento.}
\chapter{Producto interno y proyección ortogonal}\label{cap:nuevo14}\label{chap:29}\label{chap:30}
\index{producto interno}\index{Hilbert, espacio de}
\index{proyeccion ortogonal@proyección ortogonal}\index{mejor aproximacion@mejor aproximación}

Una norma mide longitudes; un producto interno permite además medir ángulos y separar componentes perpendiculares. Esta información adicional convierte ciertos problemas de aproximación en problemas geométricos. En geometría euclidiana, el punto de una recta más cercano a un punto exterior se obtiene trazando una perpendicular. El mismo principio funciona en dimensión infinita, pero la existencia del punto más cercano exige atención.

Construiremos ese puente desde los axiomas del producto interno y la desigualdad de Cauchy--Schwarz hasta la proyección ortogonal. Incluiremos los espacios de sucesiones con un conjunto de índices no necesariamente numerable y distinguiremos la \emph{caracterización} del mejor aproximante de su \emph{existencia}, para identificar exactamente dónde interviene la completitud. Se presuponen las propiedades de las normas, la convergencia de series numéricas y la completitud. El capítulo reúne los temas de las lecturas 29 y 30 de \cite{valencia}, con las correcciones de hipótesis desarrolladas en sus ejercicios.

\section{Producto interno, norma inducida y completitud}
\label{lectura:29}
\begin{definicion}
Sea $V$ un espacio vectorial sobre $\F=\R$ o $\C$. Un \emph{producto interno} es una función $\ip{\cdot}{\cdot}:V\times V\to\F$ que cumple
\[
 \ip{x}{x}\geq0,\qquad \ip{x}{x}=0\Longleftrightarrow x=0,
\]
\[
 \ip{\alpha x+\beta y}{z}=\alpha\ip{x}{z}+\beta\ip{y}{z},
 \qquad \ip{x}{y}=\overline{\ip{y}{x}}.
\]
La última propiedad se llama \emph{simetría conjugada}. En todo este volumen el producto interno es lineal en la primera entrada.
\end{definicion}
La segunda entrada es antilineal, pues
\[
 \ip{x}{\alpha y+\beta z}
 =\overline{\ip{\alpha y+\beta z}{x}}
 =\bar\alpha\ip{x}{y}+\bar\beta\ip{x}{z}.
\]
Sobre $\R$ no aparecen conjugados. Es útil escribirlos siempre al principio: muchos errores sobre adjuntos proceden de cambiar inadvertidamente la convención.

\begin{teorema}[Cauchy--Schwarz]\label{ch29:cs}
\index{Cauchy--Schwarz, desigualdad de}
Si $\norm{x}=\sqrt{\ip{x}{x}}$, entonces
\[
 |\ip{x}{y}|\leq\norm{x}\norm{y}.
\]
\end{teorema}
\begin{proof}
Si $y=0$, ambos miembros son cero. Si $y\ne0$, para cualquier $\alpha\in\F$,
\[
0\leq\norm{x-\alpha y}^{2}
=\norm{x}^{2}-\bar\alpha\ip{x}{y}-\alpha\ip{y}{x}
 +|\alpha|^{2}\norm{y}^{2}.
\]
La elección $\alpha=\ip{x}{y}/\norm{y}^{2}$ cancela la componente de $x$ paralela a $y$ y da
\[
0\leq\norm{x}^{2}-\frac{|\ip{x}{y}|^{2}}{\norm{y}^{2}}.
\]
Multiplicar por $\norm{y}^{2}$ y tomar raíces produce la desigualdad.
\end{proof}

\begin{proposicion}\label{ch29:norma}
\index{paralelogramo, identidad del}
La función $x\mapsto\sqrt{\ip{x}{x}}$ es una norma. Además,
\begin{align}
 \norm{x+y}^{2}&=\norm{x}^{2}+2\operatorname{Re}\ip{x}{y}+\norm{y}^{2},\label{ch29:expansion}\\
 \norm{x+y}^{2}+\norm{x-y}^{2}&=2\norm{x}^{2}+2\norm{y}^{2}.\label{ch29:paralelogramo}
\end{align}
\end{proposicion}
\begin{proof}
La positividad y la separación del cero son axiomas. La sesquilinealidad da $\ip{\alpha x}{\alpha x}=|\alpha|^{2}\ip{x}{x}$ y, por tanto, homogeneidad absoluta de la norma. Desarrollar $\ip{x+y}{x+y}$ demuestra \eqref{ch29:expansion}; Cauchy--Schwarz implica
\[
 \norm{x+y}^{2}\leq \norm{x}^{2}+2\norm{x}\norm{y}+\norm{y}^{2}
 =(\norm{x}+\norm{y})^{2}.
\]
La raíz cuadrada preserva el orden de números no negativos, de modo que se obtiene la desigualdad triangular. Sumar las expansiones correspondientes a $x+y$ y $x-y$ cancela los términos cruzados y prueba \eqref{ch29:paralelogramo}.
\end{proof}

\begin{definicion}
Un espacio con producto interno es un \emph{espacio de Hilbert} cuando es completo para la norma inducida. Sin completitud se habla también de espacio \emph{prehilbertiano}.
\end{definicion}

\begin{ejemplo}[Tres modelos]
En $\F^{n}$ se utiliza $\ip{x}{y}=\sum_{j=1}^{n}x_j\bar y_j$. Es un espacio de Hilbert porque una sucesión de Cauchy converge coordenada a coordenada y, al haber solamente $n$ coordenadas, converge en norma.

En $\ell^{2}$, el producto $\ip{x}{y}=\sum_{j\geq1}x_j\bar y_j$ está bien definido: Cauchy--Schwarz para sumas finitas acota $\sum_{j=1}^{N}|x_jy_j|$ por $\norm{x}_2\norm{y}_2$. La completitud se demostrará, en mayor generalidad, en el ejercicio~\ref{ej:lectura-29-8}.

Si $a<b$, en $C([a,b])$ se tiene
\[
 \ip{f}{g}=\frac1{b-a}\int_a^b f(t)\overline{g(t)}\,dt.
\]
Si $f$ no es cero, existe un intervalo de longitud positiva donde $|f|$ está acotado inferiormente por una constante positiva. Por ello $\int|f|^{2}>0$. Sin embargo, este espacio no es completo. En $[0,1]$ considérese $f_n$ igual a cero hasta $1/2$, lineal de cero a uno entre $1/2$ y $1/2+1/n$, e igual a uno después. Si $s=\mathbf1_{(1/2,1]}$, entonces $\int|f_n-s|^{2}\leq1/n$, por lo que $(f_n)$ es de Cauchy. Si convergiera en la norma integral a una función continua $f$, la desigualdad triangular daría $\int|f-s|^{2}=0$. La continuidad obligaría a $f=0$ en $(0,1/2)$ y a $f=1$ en $(1/2,1)$, incompatible con la continuidad en $1/2$. Un cambio afín de variable traslada el ejemplo a cualquier $[a,b]$.
\end{ejemplo}

\section{Ortogonalidad, sumas y desigualdad de Bessel}
\index{ortogonalidad}\index{Bessel, desigualdad de}
Se escribe $x\perp y$ si $\ip{x}{y}=0$. Para $S\subset V$ se define
\[
 S^{\perp}=\{x\in V:\ip{x}{s}=0\text{ para todo }s\in S\}.
\]
La linealidad demuestra que $S^{\perp}$ es un subespacio. Es cerrado: si $x_n\in S^{\perp}$ y $x_n\to x$, entonces, para cada $s\in S$,
$|\ip{x}{s}|=|\ip{x-x_n}{s}|\leq\norm{x-x_n}\norm{s}\to0$.
El único vector perpendicular a todos los vectores de $V$ es cero, porque es entonces perpendicular a sí mismo. Dos vectores ortogonales \emph{no nulos} son linealmente independientes: de $\alpha x+\beta y=0$ se obtiene $\alpha\norm{x}^{2}=0$ y $\beta\norm{y}^{2}=0$. La condición de no nulidad es indispensable.

\begin{teorema}[Pitágoras]\label{ch29:pitagoras}
\index{Pitagoras@Pitágoras, teorema de}
Si $x_1,\dots,x_n$ son ortogonales a pares, entonces
\[
 \Big\|\sum_{j=1}^{n}x_j\Big\|^{2}=\sum_{j=1}^{n}\norm{x_j}^{2}.
\]
\end{teorema}
\begin{proof}
Por sesquilinealidad, el miembro izquierdo es $\sum_{j,k}\ip{x_j}{x_k}$. Los términos con $j\ne k$ son cero; los restantes son $\ip{x_j}{x_j}=\norm{x_j}^{2}$. Esta expansión también explica por qué la suma de varios vectores perpendiculares a $y$ sigue siendo perpendicular a $y$.
\end{proof}

\begin{definicion}\label{ch29:suma}
Para una familia $(a_s)_{s\in S}$ de números no negativos se define
\[
 \sum_{s\in S}a_s=\sup\Big\{\sum_{s\in F}a_s:F\subset S\text{ finito}\Big\}.
\]
Cada índice se cuenta una sola vez. Un conjunto $E\subset V$ es \emph{ortonormal} si sus vectores tienen norma uno y son ortogonales dos a dos.
\end{definicion}
\index{suma!familia no numerable}\index{conjunto ortonormal}

\begin{lema}\label{ch29:soporte}
Si $A=\sum_{s\in S}a_s<\infty$, entonces $\{s:a_s>0\}$ es a lo sumo numerable. Cualquier enumeración sin repeticiones de ese conjunto tiene suma $A$.
\end{lema}
\begin{proof}
Para cada $m\geq1$, el conjunto $S_m=\{s:a_s>1/m\}$ es finito: si contuviera $r$ elementos distintos, $r/m<\sum_{s\in F}a_s\leq A$; no pueden existir subconjuntos finitos de cardinal arbitrariamente grande. Como todo número positivo supera $1/m$ para algún $m$, el soporte es $\bigcup_m S_m$. Si lo enumeramos como $s_1,s_2,\dots$, las sumas parciales crecen y están acotadas por $A$, de modo que tienen límite $L\leq A$. Para cualquier conjunto finito $F$, sus índices con coeficiente no nulo aparecen entre los primeros $N$ términos para algún $N$, y entonces $\sum_{s\in F}a_s\leq L$. Tomar el supremo da $A\leq L$. El soporte finito o vacío se trata mediante una suma finita o la suma cero.
\end{proof}

\begin{teorema}[Bessel]\label{ch29:bessel}
Si $E$ es ortonormal, entonces para todo $x\in V$,
\[
 \sum_{e\in E}|\ip{x}{e}|^{2}\leq\norm{x}^{2}.
\]
\end{teorema}
\begin{proof}
Para $F\subset E$ finito, sea $p_F=\sum_{e\in F}\ip{x}{e}e$. Para cada $u\in F$,
\[
 \ip{x-p_F}{u}=\ip{x}{u}-\sum_{e\in F}\ip{x}{e}\ip{e}{u}=0.
\]
Por tanto, $x-p_F\perp p_F$ y Pitágoras da
\[
 \norm{x}^{2}=\norm{x-p_F}^{2}+\norm{p_F}^{2}
 =\norm{x-p_F}^{2}+\sum_{e\in F}|\ip{x}{e}|^{2}.
\]
Tomamos el supremo sobre $F$. El lema anterior demuestra que, para este $x$ fijo, solamente una cantidad a lo sumo numerable de coeficientes es distinta de cero. No demuestra que $E$ mismo sea numerable; en $\ell^{2}(S)$ puede haber un vector coordenado por cada elemento de un conjunto no numerable $S$.
\end{proof}

\section{Ejercicios de producto interno, polarización y completitud}

\begin{ejercicio}[Lectura 29, ejercicio 1]\label{ej:lectura-29-1}
Caracterizar la igualdad en Cauchy--Schwarz.
\end{ejercicio}


\begin{ejercicio}[Lectura 29, ejercicio 2: corrección del recíproco]\label{ej:lectura-29-2}
Analizar si $\norm{v_1+\cdots+v_n}^{2}=\sum_j\norm{v_j}^{2}$ implica ortogonalidad por pares en un espacio real.
\end{ejercicio}


\begin{ejercicio}[Lectura 29, ejercicio 3: identidad del paralelogramo]\label{ej:lectura-29-3}
Demostrar que una norma procede de un producto interno si y solo si
\[
 \norm{x+y}^{2}+\norm{x-y}^{2}=2\norm{x}^{2}+2\norm{y}^{2}.
\]
Aplicarlo a (i) $\ell^{1}$ y (ii) la norma uniforme en $B(S,W)$.
\end{ejercicio}


\begin{ejercicio}[Lectura 29, ejercicio 4]\label{ej:lectura-29-4}
Demostrar la continuidad conjunta del producto interno.
\end{ejercicio}


\begin{ejercicio}[Lectura 29, ejercicio 5: identidad polar]\label{ej:lectura-29-5}
\index{polarizacion@polarización, identidad de}
Si $\zeta$ es una raíz primitiva de la unidad de orden $n\geq3$, probar
\[
 \ip{x}{y}=\frac1n\sum_{k=0}^{n-1}\zeta^k\norm{x+\zeta^ky}^{2}.
\]
\end{ejercicio}


\begin{ejercicio}[Lectura 29, ejercicio 6]\label{ej:lectura-29-6}
Sea $h:[a,b]\to\F$ continua y sin ceros, con $a<b$. Determinar cuándo $\ip{f}{g}_h=\int_a^b f\bar g h$ es producto interno, y discutir el caso continuo a trozos.
\end{ejercicio}


\begin{ejercicio}[Lectura 29, ejercicio 7]\label{ej:lectura-29-7}
\index{ell2@\ensuremath{\ell^2}!ponderado}
Para $w_n>0$, sea $\ell^{2}_w=\{x:\sum_n w_n|x_n|^{2}<\infty\}$. Probar (i) que es espacio vectorial, (ii) que $\ip{x}{y}_w=\sum_nw_nx_n\bar y_n$ es producto interno y (iii) que es completo.
\end{ejercicio}


\begin{ejercicio}[Lectura 29, ejercicio 8]\label{ej:lectura-29-8}
\index{ell2@\ensuremath{\ell^2}!indices arbitrarios@índices arbitrarios}
Para cualquier conjunto $S$, definir $\ell^{2}(S)$ por $\sum_{s\in S}|f(s)|^{2}<\infty$, con la suma de la definición \ref{ch29:suma}. Probar (i) la estructura vectorial, (ii) la existencia del producto $\ip{f}{g}=\sum_{s\in S}f(s)\overline{g(s)}$ y (iii) la completitud.
\end{ejercicio}


\section{Caracterización geométrica de la mejor aproximación}
\label{lectura:30}
La identidad del paralelogramo caracteriza las normas que admiten ángulos, y Bessel cuantifica cuánta longitud se captura al conservar componentes ortogonales. Convertiremos ahora esta idea en un procedimiento de mejor aproximación: la proyección ortogonal. Primero probaremos qué propiedad debe satisfacer un mejor aproximante, cuando existe; después mostraremos que la completitud permite construirlo.
\begin{teorema}\label{ch30:caracterizacion}
Sea $V$ un espacio con producto interno y $M\subset V$ un subespacio. Para $x\in V$ y $m\in M$ son equivalentes:
\begin{enumerate}[label=(\roman*)]
 \item $x-m\in M^{\perp}$;
 \item $\norm{x-m}=\dist(x,M)$.
\end{enumerate}
Si se cumplen, $m$ es único y para todo $u\in M$,
\[
 \norm{x-u}^{2}=\norm{x-m}^{2}+\norm{m-u}^{2}.
\]
En particular, $\dist(x,M)^{2}=\norm{x}^{2}-\norm{m}^{2}$.
\end{teorema}
\begin{proof}
Si $x-m\perp M$, entonces $x-m\perp m-u$, de modo que Pitágoras da la igualdad anunciada. Para $u\ne m$, el segundo sumando es positivo; por ello $m$ es el único minimizador. Con $u=0$ obtenemos la fórmula de la distancia.

Recíprocamente, supongamos que $m$ minimiza. Fijemos $u\in M\setminus\{0\}$ y sea $r=x-m$. Como $m+\alpha u\in M$ para cada $\alpha\in\F$,
\[
0\leq\norm{r-\alpha u}^{2}-\norm{r}^{2}
=-2\operatorname{Re}(\bar\alpha\ip{r}{u})+|\alpha|^{2}\norm{u}^{2}.
\]
Elegimos $\alpha=\ip{r}{u}/\norm{u}^{2}$. El lado derecho se vuelve $-|\ip{r}{u}|^{2}/\norm{u}^{2}$. Debe ser no negativo, por lo que $\ip{r}{u}=0$. Para $u=0$ la igualdad ya se cumple. Así $r\in M^{\perp}$.
\end{proof}

\begin{definicion}
El vector del teorema se llama \emph{proyección ortogonal} de $x$ sobre $M$ y se escribe $P_Mx$. Un subespacio es \emph{ortogonalmente complementado}, o \emph{distinguido} en las lecturas originales, cuando dicha proyección existe para todo $x\in V$.
\end{definicion}
\index{subespacio!distinguido}\index{subespacio!ortogonalmente complementado}

\begin{ejemplo}[La distancia puede no alcanzarse]
En $\ell^{2}$, sea $M=c_{00}=\spanop\{e_n:n\geq1\}$. Los truncamientos de cualquier sucesión convergen en norma, de modo que $M$ es denso. Para $x=(1/n)_{n\geq1}$ tenemos $\dist(x,M)=0$, pero $x\notin M$. Si $u\in M$ y $u_n=0$ para $n\geq N$, entonces
\[
 \norm{x-u}^{2}\geq\sum_{n=N}^{\infty}\frac1{n^{2}}>0.
\]
Por ello ningún $u$ alcanza la distancia. El obstáculo es que el límite de una sucesión de aproximantes puede salir del subespacio.
\end{ejemplo}

\section{Existencia y descomposición ortogonal}
\begin{teorema}[Teorema de proyección]\label{ch30:proyeccion}
En un espacio con producto interno:
\begin{enumerate}[label=(\roman*)]
 \item todo subespacio ortogonalmente complementado es cerrado;
 \item todo subespacio completo es ortogonalmente complementado.
\end{enumerate}
Por tanto, en un espacio de Hilbert, un subespacio admite proyección ortogonal si y solo si es cerrado.
\end{teorema}
\begin{proof}
\emph{(i)} Si $x\in\cl M$, entonces $\dist(x,M)=0$. Si $m=P_Mx$ existe, $\norm{x-m}=0$, luego $x=m\in M$.

\emph{(ii)} Fijemos $x$ y escribamos $d=\dist(x,M)$. Por definición de ínfimo, para cada $n$ se puede elegir $m_n\in M$ con $d\leq\norm{x-m_n}<d+1/n$. El paralelogramo aplicado a $x-m_n$ y $x-m_k$ da
\[
 \norm{m_n-m_k}^{2}
 =2\norm{x-m_n}^{2}+2\norm{x-m_k}^{2}
 -4\left\|x-\frac{m_n+m_k}{2}\right\|^{2}.
\]
El punto medio pertenece a $M$, de modo que su distancia a $x$ es al menos $d$. Por ello
\[
0\leq\norm{m_n-m_k}^{2}
\leq2\norm{x-m_n}^{2}+2\norm{x-m_k}^{2}-4d^{2}\longrightarrow0.
\]
Esto prueba que $(m_n)$ es de Cauchy. Por completitud de $M$, converge a $m\in M$, y la continuidad de la norma implica $\norm{x-m}=d$. El teorema \ref{ch30:caracterizacion} termina la prueba.

Si el espacio ambiente es Hilbert, sus subespacios cerrados son completos: el límite ambiente de una sucesión de Cauchy permanece en el subespacio por cerradura. Esto explica la última equivalencia.
\end{proof}

\begin{teorema}\label{ch30:descomposicion}
\index{descomposicion ortogonal@descomposición ortogonal}
Un subespacio $M$ es ortogonalmente complementado si y solo si
\[
 V=M\oplus M^{\perp}.
\]
En ese caso $(M^{\perp})^{\perp}=M$, y $M^{\perp}$ también es ortogonalmente complementado.
\end{teorema}
\begin{proof}
Si existe $P_Mx$, escribimos $x=P_Mx+(x-P_Mx)$, con el primer sumando en $M$ y el segundo en $M^{\perp}$. Además $M\cap M^{\perp}=\{0\}$, pues un vector de la intersección es perpendicular a sí mismo. La descomposición es, por tanto, única.

Recíprocamente, si $x=m+n$ con $m\in M$ y $n\in M^{\perp}$, el teorema \ref{ch30:caracterizacion} demuestra que $m=P_Mx$. Es el componente en $M$, no el componente perpendicular.

Siempre $M\subset(M^{\perp})^{\perp}$. Si $x\in(M^{\perp})^{\perp}$, descompongamos $x=m+n$. Entonces $0=\ip{x}{n}=\ip{m}{n}+\ip{n}{n}=\norm{n}^{2}$, y $x=m\in M$. Así $V=M^{\perp}\oplus(M^{\perp})^{\perp}$; la equivalencia ya demostrada prueba que $M^{\perp}$ es complementado.
\end{proof}

\begin{corolario}\label{ch30:doble}
\index{complemento ortogonal!doble}
Si $H$ es Hilbert y $S\subset H$, entonces
\[
 (S^{\perp})^{\perp}=\cl{\spanop S}.
\]
\end{corolario}
\begin{proof}
Un vector perpendicular a $S$ es perpendicular a todas sus combinaciones lineales y a sus límites, por continuidad del producto interno. Por tanto, con $M=\cl{\spanop S}$ se tiene $S^{\perp}=M^{\perp}$. Como $M$ es cerrado en Hilbert, es complementado y el teorema anterior da $(M^{\perp})^{\perp}=M$.
\end{proof}

\section{Calcular una proyección sin ortonormalizar}\label{ampl:gram}
El teorema de proyección garantiza un mejor aproximante, pero no siempre lo expresa en las coordenadas que tenemos disponibles. Si un subespacio está generado por monomios, datos experimentales u otras funciones sencillas, conviene poder trabajar directamente con esa lista. La matriz de Gram convierte la condición de ortogonalidad del error en un sistema de ecuaciones lineales.

Esta ampliación utiliza \cite{max-gram,max-proyeccion} como orientación. Las demostraciones se escriben con la convención del volumen: el producto interno es lineal en el primer argumento. Las fórmulas complejas dependen de esa convención; no basta con copiar una matriz sin comprobar dónde aparecen los conjugados.
\index{Gram, matriz de}\index{ecuaciones normales}\index{minimos cuadrados@mínimos cuadrados}

\begin{definicion}
Para una lista $a_1,\ldots,a_m$ de vectores de un espacio con producto interno, su \emph{matriz de Gram} será
\[
 G=(G_{ij})_{i,j=1}^m,\qquad G_{ij}=\ip{a_j}{a_i}.
\]
Para una matriz compleja, $A^*$ denota su transpuesta conjugada. Los vectores de coeficientes se escriben como columnas.
\end{definicion}

\begin{proposicion}\label{ampl:gram-positiva}
La matriz $G$ es hermitiana y semidefinida positiva. Es definida positiva, y por tanto invertible, si y solo si $a_1,\ldots,a_m$ son linealmente independientes. Para cada columna $c=(c_j)$,
\[
 c^*Gc=\left\|\sum_{j=1}^mc_ja_j\right\|^2.
\]
\end{proposicion}
\begin{proof}
La simetría conjugada da $\overline{G_{ji}}=\overline{\ip{a_i}{a_j}}=\ip{a_j}{a_i}=G_{ij}$, luego $G^*=G$. Además,
\[
 c^*Gc=\sum_{i,j}\overline{c_i}c_j\ip{a_j}{a_i}
 =\ip{\sum_jc_ja_j}{\sum_ic_ia_i}\ge0.
\]
Si la lista es independiente y $c\ne0$, el vector dentro de la norma no es cero, por lo que la expresión es positiva. Si la lista es dependiente, existe $c\ne0$ que produce el vector cero, y la expresión no es estrictamente positiva. Finalmente, para una lista independiente, $Gc=0$ implicaría $c^*Gc=0$ y por tanto $c=0$; una matriz cuadrada con núcleo cero es invertible.
\end{proof}

\begin{teorema}[Ecuaciones normales de la proyección]\label{ampl:gram-proyeccion}
Sean $a_1,\ldots,a_m$ linealmente independientes y $M=\spanop\{a_1,\ldots,a_m\}$. Para $u$ en el espacio ambiente defina la columna $\alpha$ por $\alpha_i=\ip{u}{a_i}$. El sistema
\[
 Gc=\alpha
\]
tiene solución única, y
\[
 P_Mu=\sum_{j=1}^mc_ja_j,\qquad
 \dist(u,M)^2=\norm u^2-\alpha^*G^{-1}\alpha.
\]
No se requiere que el espacio ambiente sea completo.
\end{teorema}
\begin{proof}
Por la proposición anterior existe $c=G^{-1}\alpha$. Sea $p=\sum_jc_ja_j$. Para cada $i$,
\[
 \ip{u-p}{a_i}=\alpha_i-\sum_jc_j\ip{a_j}{a_i}
 =\alpha_i-(Gc)_i=0.
\]
Si $v=\sum_id_ia_i\in M$, entonces
$\ip{u-p}{v}=\sum_i\overline{d_i}\ip{u-p}{a_i}=0$. El residuo es ortogonal a todo $M$, de modo que el teorema~\ref{ch30:caracterizacion} prueba que $p$ es el único mejor aproximante. El argumento sólo usó sumas finitas y un sistema invertible, por lo que no necesita completitud ambiente.

Como $G$ es hermitiana, también lo es $G^{-1}$: al tomar adjuntos en $GG^{-1}=I$ se obtiene $(G^{-1})^*G=I$. Por ello,
\[
 \norm p^2=c^*Gc=(G^{-1}\alpha)^*G(G^{-1}\alpha)
 =\alpha^*G^{-1}\alpha.
\]
La descomposición $u=p+(u-p)$ es ortogonal. Pitágoras da la fórmula del error.
\end{proof}

En la práctica se resuelve $Gc=\alpha$; escribir $G^{-1}$ en una fórmula no obliga a calcular todas las entradas de la inversa. Cuando la lista ya es ortonormal, $G=I$ y recuperamos $c_i=\ip{u}{a_i}$. La matriz de Gram registra precisamente lo que cambia cuando las direcciones no son ortogonales.

\begin{ejemplo}[Matrices y ajuste por mínimos cuadrados]
En $\F^n$ con el producto interno usual, reúna los vectores $a_j$ como columnas de una matriz $A$ de tamaño $n\times m$. Entonces
\[
 G=A^*A,\qquad \alpha=A^*u.
\]
En efecto, la entrada $(i,j)$ de $A^*A$ es $\sum_k\overline{(a_i)_k}(a_j)_k=\ip{a_j}{a_i}$. Las ecuaciones normales son $A^*Ac=A^*u$ y minimizan $\norm{u-Ac}_2^2$. Si las columnas son independientes, la proyección queda expresada por
\[
 P_Mu=A(A^*A)^{-1}A^*u.
\]
La fórmula es una aplicación finita del teorema anterior; no presupone todavía la teoría del adjunto de un operador entre espacios de Hilbert.
\end{ejemplo}

\section{Ejercicios sobre Gram y mejores aproximaciones}
Los ejercicios siguientes son elaboraciones propias de esta edición. El último desarrolla también el problema sobre rango de Gram propuesto en \cite{max-gram}; las soluciones y las aplicaciones concretas se presentan aquí completas.

\begin{ejercicio}[Ampliación: aproximar una parábola por una recta]\label{ej:ampl-gram-1}
En $C([0,1],\R)$ con $\ip fg=\int_0^1f(t)g(t)\,dt$, encuentre el polinomio $a+bt$ que minimiza $\int_0^1(t^2-a-bt)^2\,dt$. Calcule el error mínimo y verifique directamente la ortogonalidad del residuo.
\end{ejercicio}


\begin{ejercicio}[Ampliación: dónde va el conjugado]\label{ej:ampl-gram-2}
En $\C^2$, con $\ip xy=x_1\overline{y_1}+x_2\overline{y_2}$, proyecte $u=(1,1)$ sobre la recta compleja generada por $a=(1,i)$. Compruebe la norma del error y explique por qué usar $\ip au/\norm a^2$ como coeficiente daría una respuesta incorrecta.
\end{ejercicio}


\begin{ejercicio}[Ampliación: tres datos y dos parámetros]\label{ej:ampl-gram-3}
Ajuste una recta $y=a+bt$ a los datos $(0,1)$, $(1,2)$ y $(2,2)$ minimizando la suma de los cuadrados de los tres errores. Calcule los parámetros, los valores ajustados y el error mínimo. Explique por qué el problema tiene una solución única aunque ninguna recta pase por los tres datos.
\end{ejercicio}


\begin{ejercicio}[Ampliación: cambiar de base sin cambiar la proyección]\label{ej:ampl-gram-4}
Sea $a_1,\ldots,a_m$ una lista independiente y $C=(C_{ij})$ una matriz invertible. Defina $b_j=\sum_i a_iC_{ij}$. Pruebe que las matrices de Gram y los datos de proyección satisfacen $G_b=C^*G_aC$ y $\alpha_b=C^*\alpha_a$, y relacione las coordenadas del mismo vector proyectado. Ilustre el resultado con $a_1=(1,0)$, $a_2=(1,\eps)$, $\eps>0$, y $u=(0,1)$ en $\R^2$.
\end{ejercicio}


\begin{ejercicio}[Ampliación: listas dependientes y coeficientes no únicos]\label{ej:ampl-gram-5}
Permita ahora que $a_1,\ldots,a_m$ sean dependientes. Defina $A:\F^m\to V$ por $Ac=\sum_jc_ja_j$. Pruebe que $\ker G=\ker A$, que $\operatorname{rango}G=\dim\spanop\{a_j\}$ y que $Gc=\alpha$ sigue teniendo soluciones para $\alpha_i=\ip u{a_i}$. Muestre que todas producen el mismo mejor aproximante. Calcule el sistema para $a_1=1$, $a_2=t$, $a_3=1+t$ y $u=t^2$, con producto integral en $[0,1]$.
\end{ejercicio}


La fórmula de Gram concreta la proyección en subespacios finitos. Los ejercicios que siguen vuelven a la geometría general, donde intervienen conjuntos convexos, complementos ortogonales y familias que pueden ser infinitas.

\section{Ejercicios de proyección y complementos ortogonales}

\begin{ejercicio}[Lectura 30, ejercicio 1]\label{ej:lectura-30-1}
Sea $C$ un subconjunto convexo, completo y no vacío de un espacio con producto interno. Probar que para cada $x$ existe un único $c\in C$ tal que $\norm{x-c}=\dist(x,C)$.
\end{ejercicio}


\begin{ejercicio}[Lectura 30, ejercicio 2: Parseval sobre un subespacio]\label{ej:lectura-30-2}
Sea $M$ ortogonalmente complementado y $E\subset M$ ortonormal con $\spanop E$ denso en $M$. Para $x\in V$ y $m=P_Mx$, probar
\[
 \norm{m}^{2}=\sum_{e\in E}|\ip{x}{e}|^{2}.
\]
\end{ejercicio}


\begin{ejercicio}[Lectura 30, ejercicio 3: tres perpendiculares]\label{ej:lectura-30-3}
Sean $M\subset N$ subespacios ortogonalmente complementados. Probar $P_MP_Nx=P_Mx$.
\end{ejercicio}


\begin{ejercicio}[Lectura 30, ejercicio 4: complementos ortogonales]\label{ej:lectura-30-4}
Para subespacios $M,N$, estudiar
\begin{enumerate}[label=(\roman*)]
 \item $(M+N)^{\perp}=M^{\perp}\cap N^{\perp}$;
 \item $(M\cap N)^{\perp}=\cl{M^{\perp}+N^{\perp}}$;
 \item $M\subset(M^{\perp})^{\perp}$.
\end{enumerate}
En (ii) suponemos que $M,N$ son subespacios cerrados de un espacio de Hilbert. Explicaremos por qué no basta, en general, suponer que $M,N$ sean distinguidos en un espacio incompleto. Los incisos (i) y (iii) no requieren cerradura ni completitud.
\end{ejercicio}


\begin{ejercicio}[Lectura 30, ejercicio 5]\label{ej:lectura-30-5}
Probar que $P_M$ es lineal, continuo, idempotente y simétrico respecto al producto interno. Determinar su norma.
\end{ejercicio}


\begin{ejercicio}[Lectura 30, ejercicio 6]\label{ej:lectura-30-6}
Sea $P:V\to V$ lineal, $P^{2}=P$, y $\ip{Px}{y}=\ip{x}{Py}$ para todos $x,y$. Probar que $P$ es una proyección ortogonal.
\end{ejercicio}


\section{La descomposición como control del error}
La fórmula $x=P_Mx+(x-P_Mx)$ separa la información representable en $M$ del error perpendicular. La perpendicularidad no es solamente una imagen geométrica: demuestra unicidad, controla la norma del error y convierte la aproximación por un sistema ortonormal en una suma de cuadrados. En el capítulo~\ref{cap:nuevo15} haremos esa suma con todos los vectores de una base ortonormal.

\chapter{Bases ortonormales y desarrollos de Fourier}\label{cap:nuevo15}\label{chap:31}
\index{base!ortonormal}\index{Fourier, coeficientes de}

Una base de álgebra lineal expresa cada vector mediante una combinación \emph{finita}. En análisis necesitamos permitir límites de combinaciones finitas. La ortogonalidad hace controlable ese límite: la norma de una cola se calcula sumando cuadrados de coeficientes. Las proyecciones del capítulo~\ref{cap:nuevo14} permiten interpretar cada suma finita como la mejor aproximación disponible en las direcciones elegidas. Desarrollaremos la lectura 31 de \cite{valencia}, incluyendo el argumento de existencia, la construcción explícita de Gram--Schmidt y el ejercicio de Müntz en la norma integral que realmente aparece en el enunciado.

\section{Bases algebraicas y bases ortonormales}
\label{lectura:31}
Una \emph{base de Hamel} es un conjunto linealmente independiente cuyo espacio generado algebraicamente es todo el espacio. Por ejemplo, los vectores coordenados $(e_n)$ no constituyen una base de Hamel de $\ell^{2}$: una combinación finita tiene soporte finito, mientras $(1/n)$ tiene infinitas coordenadas no nulas.

\begin{definicion}
Un conjunto es \emph{total} si el cierre de su espacio generado es todo el espacio. Una \emph{base ortonormal} de un espacio con producto interno es un conjunto ortonormal total. En el espacio cero, la base ortonormal es el conjunto vacío.
\end{definicion}
\index{base!de Hamel}\index{conjunto total}
La independencia lineal de un conjunto ortonormal se verifica tomando productos internos: si $\sum_{j=1}^{n}\alpha_je_j=0$, el producto con $e_k$ da $\alpha_k=0$. No conviene llamar simplemente «base» a cualquier conjunto independiente total: sin ortogonalidad, estas condiciones no garantizan un desarrollo en serie con coeficientes únicos. Una base de Schauder incorpora precisamente propiedades adicionales de convergencia y unicidad. Distinguiremos siempre las expresiones \emph{base de Hamel}, \emph{base de Schauder} y \emph{base ortonormal}.

\begin{teorema}\label{ch31:existencia}
Todo espacio de Hilbert tiene una base ortonormal.
\end{teorema}
\begin{proof}
Consideremos todos los subconjuntos ortonormales de $H$, incluido el vacío, ordenados por inclusión. Si $\mathcal C$ es una cadena, su unión es ortonormal: un vector de la unión pertenece a algún miembro y tiene norma uno; dos vectores de la unión pertenecen a dos miembros comparables de la cadena y, por tanto, ambos están en uno de ellos. Si son distintos, son ortogonales. La unión es una cota superior. El lema de Zorn proporciona un conjunto ortonormal maximal $E$.

Sea $M=\cl{\spanop E}$. Es un subespacio cerrado de Hilbert, por lo que tiene proyección ortogonal. Si $M\ne H$, elijamos $x\notin M$ y sea $r=x-P_Mx$. Entonces $r\ne0$ y $r\perp M$. El vector
\[
 e=\frac{r}{\norm r}
\]
tiene norma uno y es perpendicular a cada miembro de $E$. Por tanto $E\cup\{e\}$ es ortonormal, contradiciendo la maximalidad. La normalización es necesaria: el residuo $r$ no tiene, en general, norma uno. Así $M=H$ y $E$ es una base ortonormal. Para $H=\{0\}$, el mismo argumento da $E=\varnothing$.
\end{proof}

\section{Bases de Schauder: coordenadas mediante series}\label{ampl:schauder}
Una lista de vectores puede aproximar todo el espacio mediante combinaciones finitas sin proporcionar una receta coherente de coordenadas. En una expansión en serie no se permite escoger de nuevo todos los coeficientes cada vez que aumenta la precisión: las sumas parciales deben proceder de una misma lista de coeficientes. Esta es la exigencia adicional que formaliza Schauder.

La definición y los ejemplos de sucesiones de esta sección se apoyan en Maximenko \cite{max-schauder}. Desarrollamos las pruebas y añadimos ejemplos que separan independencia, densidad, convergencia y dependencia del orden. No se afirma que todo espacio separable tenga una base de Schauder; la separabilidad es aquí una condición necesaria, no un criterio suficiente.
\index{base!de Schauder}\index{sumas parciales!operador de}\index{base!condicional}

\begin{definicion}
Una sucesión $(b_n)_{n\ge1}$ en un espacio normado $V$ es una \emph{base de Schauder} si para cada $x\in V$ existe una única sucesión de escalares $(a_n)$ tal que
\[
 x=\lim_{N\to\infty}\sum_{n=1}^Na_nb_n
\]
en la norma de $V$. Se escriben $a_n=b_n^*(x)$ y
$S_Nx=\sum_{n=1}^Nb_n^*(x)b_n$. El orden de $(b_n)$ forma parte de la definición. En dimensión finita se utiliza la versión finita, que coincide con una base algebraica.
\end{definicion}

La convergencia significa $\norm{x-S_Nx}\to0$, no sólo convergencia coordenada a coordenada en algún modelo particular. Las aplicaciones $b_n^*$ y $S_N$ son lineales: las sumas de dos expansiones convergentes dan una expansión de la suma, y su unicidad obliga a sumar los coeficientes; el argumento para multiplicación por escalares es el mismo. No usaremos aquí un teorema general sobre continuidad de estos funcionales: en los ejemplos que la requieren obtendremos cotas explícitas.

\begin{proposicion}\label{ampl:schauder-independencia}
Una base de Schauder es linealmente independiente y total.
\end{proposicion}
\begin{proof}
Si $\sum_{n=1}^Na_nb_n=0$ con algún coeficiente no nulo, al añadir coeficientes cero después de $N$ obtenemos una expansión del vector cero distinta de la expansión con todos los coeficientes nulos. Esto contradice unicidad. Además, cada vector es límite de sus sumas parciales, que pertenecen al espacio generado algebraicamente por los $b_n$. Por ello el cierre de ese espacio generado es todo $V$.
\end{proof}

\begin{proposicion}[Bases coordenadas]\label{ampl:schauder-coordenadas}
Los vectores coordenados $e_n$ forman una base de Schauder de $\ell^p$ para $1\le p<\infty$ y de $c_0$ con su norma uniforme. Sus coeficientes son las coordenadas del vector. No forman una base de Schauder de $\ell^\infty$.
\end{proposicion}
\begin{proof}
Para $x\in\ell^p$, las sumas parciales $S_Nx=\sum_{n=1}^Nx_ne_n$ satisfacen
\[
 \norm{x-S_Nx}_p^p=\sum_{n>N}|x_n|^p\longrightarrow0.
\]
En $c_0$, $\norm{x-S_Nx}_\infty=\sup_{n>N}|x_n|\to0$ por la definición de convergencia a cero. Para probar unicidad, la evaluación en una coordenada $k$ es continua porque $|x_k-y_k|\le\norm{x-y}_p$ o $\norm{x-y}_\infty$. Una serie $\sum a_ne_n$ que converge a $x$ tiene, desde su suma parcial número $k$, coordenada $k$ igual a $a_k$, de modo que $a_k=x_k$.

En $\ell^\infty$, cualquier combinación finita de $e_n$ se anula a partir de un índice. Su distancia a la sucesión constante $\mathbf1$ es al menos $1$. Así la lista ni siquiera es total en $\ell^\infty$ y no puede ser base de Schauder.
\end{proof}

\begin{proposicion}[El caso ortonormal]\label{ampl:schauder-ortonormal}
Si una base ortonormal de un espacio con producto interno puede enumerarse como $(e_n)_{n\ge1}$, esa enumeración es una base de Schauder. Los coeficientes son $\ip{x}{e_n}$. Toda permutación de la enumeración tiene la misma propiedad.
\end{proposicion}
\begin{proof}
Sea $p_N=\sum_{n=1}^N\ip{x}{e_n}e_n$. El residuo $x-p_N$ es perpendicular a $e_1,\ldots,e_N$, de modo que $p_N$ es el mejor aproximante en su espacio generado. Dado $\eps>0$, la totalidad proporciona una combinación finita $v$ de los $e_n$ con $\norm{x-v}<\eps$. Para $N$ suficientemente grande, $v$ pertenece al espacio generado por $e_1,\ldots,e_N$; entonces $\norm{x-p_N}\le\norm{x-v}<\eps$. Esto prueba convergencia sin suponer completitud adicional del ambiente.

Si $x=\lim_N\sum_{n=1}^Na_ne_n$, tomar producto interno con $e_k$ y usar su continuidad da $\ip{x}{e_k}=a_k$. La expansión es única. Después de una permutación, toda combinación finita sigue estando contenida en alguna suma inicial suficientemente larga, y el mismo argumento funciona. Para familias ortonormales no numerables se usan sumas sobre conjuntos finitos; la definición secuencial de Schauder no debe aplicarse sin esta distinción.
\end{proof}

\section{Ejercicios sobre coordenadas, densidad y orden}

\begin{ejercicio}[Ampliación: una base para las sucesiones convergentes]\label{ej:ampl-schauder-1}
En el espacio $c$ con la norma uniforme pruebe que $(\mathbf1,e_1,e_2,\ldots)$ es una base de Schauder. Determine los coeficientes, el error de las sumas parciales y cotas para los funcionales coordenados. Este ejercicio adapta la propuesta de encontrar una base de $c$ en \cite{max-schauder}; la construcción y la solución se desarrollan aquí.
\end{ejercicio}


\begin{ejercicio}[Ampliación: independencia y totalidad no bastan]\label{ej:ampl-schauder-2}
En $\ell^2$ defina $b_1=e_1$ y $b_n=e_1+n^{-1}e_n$ para $n\ge2$. Pruebe que la lista es linealmente independiente y total, pero no es una base de Schauder. Examine el vector $x=(0,1/2,1/3,\ldots)$. Ejercicio propio de esta edición.
\end{ejercicio}


\begin{ejercicio}[Ampliación: separabilidad y una obstrucción]\label{ej:ampl-schauder-3}
Pruebe que todo espacio con base de Schauder es separable. Demuestre después directamente que $\ell^\infty$ no es separable y concluya que no admite ninguna base de Schauder numerable, no sólo que su lista coordenada falla. La primera implicación y esta aplicación desarrollan la discusión de \cite{max-schauder}.
\end{ejercicio}


\begin{ejercicio}[Ampliación: transportar una base a otra]\label{ej:ampl-schauder-4}
Sea $T:V\to W$ una biyección lineal con $\norm{Tx}\le C\norm x$ y $\norm{T^{-1}y}\le D\norm y$. Demuestre que una base de Schauder $(b_n)$ de $V$ se transforma en la base $(Tb_n)$ de $W$. Aplique el resultado en $\ell^2$ a $T=I+\lambda R$, donde $R(x_1,x_2,\ldots)=(0,x_1,x_2,\ldots)$ y $|\lambda|<1$. Construya la inversa y calcule los coeficientes de la base $e_n+\lambda e_{n+1}$. Ejercicio propio de esta edición.
\end{ejercicio}


\begin{ejercicio}[Ampliación: una base cuyo orden importa]\label{ej:ampl-schauder-5}
En $c_0$ con norma uniforme, sea $s_n=e_1+\cdots+e_n$. Pruebe que $(s_n)$ es una base de Schauder, calcule sus coeficientes y la norma del operador de suma parcial. Usando $x_n=(-1)^{n-1}/n$, construya una permutación de esta base que deje de ser base de Schauder. Ejercicio propio de esta edición.
\end{ejercicio}


Hamel describe representaciones finitas; Schauder añade límites ordenados de una expansión única; una base ortonormal permite además medir el error con cuadrados de coeficientes. Desarrollaremos ahora esta última ventaja, que conduce a Parseval y a Riesz--Fischer.

\section{Coordenadas y teorema de Riesz--Fischer}
\begin{teorema}[Parseval y expansión]\label{ch31:parseval}
\index{Parseval, identidad de}
Si $E=(e_s)_{s\in S}$ es una base ortonormal de $H$, entonces
\[
 \norm{x}^{2}=\sum_{s\in S}|\ip{x}{e_s}|^{2},\qquad
 x=\sum_{s\in S}\ip{x}{e_s}e_s.
\]
La segunda suma converge en norma: para todo $\eps>0$ existe $F_0\subset S$ finito tal que cualquier suma sobre un conjunto finito $F\supset F_0$ dista menos de $\eps$ de $x$.
\end{teorema}
\begin{proof}
La identidad de Parseval del ejercicio~\ref{ej:lectura-30-2}, aplicada a $M=H$, da la primera fórmula. Para $p_F=\sum_{s\in F}\ip{x}{e_s}e_s$, la ortogonalidad da
\[
 \norm{x-p_F}^{2}=\norm{x}^{2}-\sum_{s\in F}|\ip{x}{e_s}|^{2}.
\]
Como las sumas finitas tienen supremo $\norm{x}^{2}$, existe $F_0$ para el cual la diferencia es menor que $\eps^{2}$. Al ampliar $F_0$ la suma solo aumenta, de modo que la misma estimación vale para todo $F\supset F_0$. Además el soporte de los coeficientes es numerable por Bessel; si se enumera, se obtiene una serie ordinaria con el mismo límite, independiente del orden.
\end{proof}

\begin{teorema}[Riesz--Fischer]\label{ch31:rieszfischer}
\index{Riesz--Fischer, teorema de}
La aplicación
\[
 U:H\longrightarrow\ell^{2}(S),\qquad Ux=(\ip{x}{e_s})_{s\in S},
\]
es una biyección lineal que preserva normas y productos internos.
\end{teorema}
\begin{proof}
Bessel garantiza que $Ux\in\ell^{2}(S)$. La linealidad en la primera entrada prueba que $U$ es lineal, y Parseval prueba $\norm{Ux}_2=\norm{x}$. En particular, es inyectiva.

Para la suprayectividad, sea $a\in\ell^{2}(S)$. Si su soporte es finito, $x=\sum_s a_se_s$ tiene las coordenadas requeridas. Si es infinito, enumerémoslo como $s_1,s_2,\dots$ y definamos $x_N=\sum_{j=1}^{N}a_{s_j}e_{s_j}$. Para $M>N$, Pitágoras da
\[
 \norm{x_M-x_N}^{2}=\sum_{j=N+1}^{M}|a_{s_j}|^{2}\longrightarrow0
\]
cuando $N\to\infty$, uniformemente para $M>N$. Por completitud existe $x=\lim_Nx_N\in H$. Para cada $s=s_j$, $\ip{x_N}{e_s}=a_s$ una vez que $N\geq j$; la continuidad del producto interno da $\ip{x}{e_s}=a_s$. Si $s$ está fuera del soporte, todos esos productos son cero y también lo es su límite. Luego $Ux=a$.

Finalmente, si $x_F,y_F$ son truncamientos sobre un conjunto finito común, $\ip{x_F}{y_F}=\sum_{s\in F}\ip{x}{e_s}\overline{\ip{y}{e_s}}$. Las aproximaciones en norma y la continuidad del producto interno permiten pasar al límite. La serie del lado derecho converge absolutamente por Cauchy--Schwarz en $\ell^{2}(S)$. Por ello $\ip{Ux}{Uy}_{\ell^{2}(S)}=\ip{x}{y}$.
\end{proof}
Los escalares $\ip{x}{e_s}$ se llaman \emph{coeficientes de Fourier} respecto de $E$. Esta denominación no exige funciones trigonométricas: describe cualquier sistema de coordenadas ortonormales.

\section{Ejercicios de ortonormalización, aproximación y Fourier}

\begin{ejercicio}[Lectura 31, ejercicio 1: Gram--Schmidt]\label{ej:lectura-31-1}
\index{Gram--Schmidt, procedimiento de}
Para una familia finita o numerable linealmente independiente $v_1,v_2,\dots$ en un espacio con producto interno, definir
\[
 w_1=v_1,\qquad
 w_{n+1}=v_{n+1}-\sum_{k=1}^{n}\frac{\ip{v_{n+1}}{w_k}}{\norm{w_k}^{2}}w_k.
\]
Probar la ortogonalidad, la no nulidad, la igualdad de los espacios generados y la existencia de una base ortonormal en todo espacio prehilbertiano separable.
\end{ejercicio}


\begin{ejercicio}[Lectura 31, ejercicio 2: primeros polinomios de Legendre]\label{ej:lectura-31-2}
\index{Legendre, polinomios de}
Aplicar Gram--Schmidt a $1,x,x^{2},x^{3},x^{4}$ con $\ip{f}{g}=\int_{-1}^{1}f(x)\overline{g(x)}\,dx$.
\end{ejercicio}


\begin{ejercicio}[Lectura 31, ejercicio 3]\label{ej:lectura-31-3}
Probar que $(e^{inx})_{n\in\Z}$ es ortonormal en $C([-\pi,\pi])$ con el producto interno normalizado.
\end{ejercicio}


\begin{ejercicio}[Lectura 31, ejercicio 4: Müntz en norma integral]\label{ej:lectura-31-4}
\index{Muntz@Müntz, teorema de}\index{Gram, determinante de}
Sean $1\leq n_1<n_2<\cdots$ enteros y $M_k=\spanop\{x^{n_1},\dots,x^{n_k}\}\subset C([0,1])$, con $\norm{f}_2^{2}=\int_0^1|f|^{2}$. Probar que $\bigcup_kM_k$ es denso para esta norma si y solo si $\sum_k1/n_k=\infty$. Desarrollar (i) el cociente de determinantes de Gram y la fórmula de distancia; (ii) el criterio de las series logarítmicas con sus hipótesis; (iii) el análisis del producto resultante.
\end{ejercicio}


\begin{ejercicio}[Lectura 31, ejercicio 5]\label{ej:lectura-31-5}
Probar que un espacio de Hilbert es separable si y solo si tiene una base ortonormal a lo sumo numerable.
\end{ejercicio}


\begin{ejercicio}[Lectura 31, ejercicio 6]\label{ej:lectura-31-6}
Calcular los cinco primeros coeficientes respecto de $e_0,\dots,e_4$ del ejercicio~\ref{ej:lectura-31-2} para (i) $x^2+x+1$, (ii) $\sin(\pi x)$, (iii) $\cos(\pi x)$ y (iv) $|x|$.
\end{ejercicio}


\begin{ejercicio}[Lectura 31, ejercicio 7]\label{ej:lectura-31-7}
Si $E$ es ortonormal y $x\in H$, estimar el número $n_k$ de vectores $e\in E$ con $|\ip{x}{e}|>1/k$.
\end{ejercicio}


\section{Una base como sistema de medición}
Los coeficientes de Fourier registran las componentes detectadas por cada dirección ortonormal. Bessel limita la energía registrada; Parseval afirma que una base total no pierde ninguna componente; Riesz--Fischer garantiza que toda lista de coeficientes de energía finita procede de un vector. La completitud permite pasar de las sumas finitas, puramente algebraicas, al vector que esas sumas aproximan. En el capítulo~\ref{cap:nuevo16} pasaremos de las coordenadas de un vector a la representación de los funcionales continuos y a la relación de adjunción entre operadores.

\chapter{Representación de Riesz y operadores adjuntos}\label{cap:nuevo16}\label{chap:32}
\index{Riesz!representacion en Hilbert@representación en Hilbert}\index{operador!adjunto}

Los desarrollos del capítulo~\ref{cap:nuevo15} describen un vector mediante sus componentes ortogonales. Para estudiar funcionales y operadores necesitamos relacionar esas componentes con el producto interno: una matriz y su transpuesta conjugada describen la misma relación entre vectores vista desde lados opuestos de ese producto. En Hilbert esa relación sigue existiendo sin elegir coordenadas. La herramienta que la hace posible es el teorema de representación de Riesz. La lectura 32 de \cite{valencia} lo utiliza posponiendo su prueba; aquí lo demostraremos antes de aplicarlo. Cerraremos el volumen relacionando los adjuntos con la aproximación de operadores compactos por operadores de rango finito.

\section{Cada funcional continuo es un producto interno}
\label{lectura:32}
\begin{teorema}[Representación de Riesz]\label{ch32:riesz}
Sea $H$ un espacio de Hilbert y $\varphi:H\to\F$ un funcional lineal continuo. Existe un único $v\in H$ tal que
\[
 \varphi(x)=\ip{x}{v}\quad(x\in H),\qquad \norm{\varphi}=\norm v.
\]
\end{teorema}
\begin{proof}
Si $\varphi=0$, tomamos $v=0$. Supongamos $\varphi\ne0$. Su núcleo $M$ es un subespacio cerrado: si $x_n\in M$ converge a $x$, la continuidad da $\varphi(x)=\lim\varphi(x_n)=0$. Elijamos $z$ con $\varphi(z)\ne0$ y sea $r=z-P_Mz$. Entonces $r\in M^{\perp}$ y $\varphi(r)=\varphi(z)\ne0$, por lo que $r\ne0$.

Para cada $x$, el vector $x-\varphi(x)r/\varphi(r)$ pertenece al núcleo. Tomando su producto interno con $r$,
\[
0=\ip{x}{r}-\frac{\varphi(x)}{\varphi(r)}\norm r^{2}.
\]
Así $\varphi(x)=\varphi(r)\ip{x}{r}/\norm r^{2}$. Con nuestra convención de linealidad, el vector representante es
\[
 v=\frac{\overline{\varphi(r)}}{\norm r^{2}}r,
\]
pues el escalar que multiplica a la segunda entrada sale conjugado. Si $v,w$ representan el mismo funcional, $\ip{x}{v-w}=0$ para todo $x$; con $x=v-w$ se obtiene $v=w$.

Cauchy--Schwarz da $|\varphi(x)|\leq\norm x\norm v$, de modo que $\norm{\varphi}\leq\norm v$. Si $v\ne0$, tomando $x=v/\norm v$ obtenemos $|\varphi(x)|=\norm v$, y la desigualdad opuesta. Para $v=0$ las dos normas son cero.
\end{proof}

El representante depende antilinealmente del funcional: el representante de $\alpha\varphi+\beta\psi$ es $\bar\alpha v+\bar\beta w$. Esta observación evita confundir la identificación entre $H$ y $H^*$ con una aplicación compleja lineal.

\section{Formas sesquilineales y operadores}
\begin{definicion}
Para espacios de Hilbert complejos $H_1,H_2$, una forma $\xi:H_1\times H_2\to\C$ es \emph{sesquilineal} si es lineal en la primera variable y antilineal en la segunda. Es \emph{acotada} si existe $C\geq0$ tal que
\[
 |\xi(x,y)|\leq C\norm x\norm y.
\]
\end{definicion}
\index{forma sesquilineal}
Esta condición equivale a continuidad conjunta. En efecto, si existe tal cota,
\[
 |\xi(x',y')-\xi(x,y)|
 \leq C\norm{x'-x}\norm{y'}+C\norm x\norm{y'-y},
\]
que tiende a cero cuando $(x',y')\to(x,y)$. Recíprocamente, la continuidad en $(0,0)$ proporciona $\delta>0$ tal que $|\xi(u,v)|<1$ si $\norm u,\norm v<\delta$. Para $x,y\ne0$ aplicamos esto a $u=\delta x/(2\norm x)$ y $v=\delta y/(2\norm y)$; por sesquilinealidad,
$|\xi(x,y)|\leq4\norm x\norm y/\delta^2$. Si alguno es cero, la desigualdad se cumple porque $\xi$ vale cero.

Si $T\in\Bcal(H_1,H_2)$ y $S\in\Bcal(H_2,H_1)$, las formas $\ip{Tx}{y}$ y $\ip{x}{Sy}$ son acotadas, con cotas $\norm T$ y $\norm S$ por Cauchy--Schwarz y la definición de norma de operador.

\begin{teorema}\label{ch32:formas}
Toda forma sesquilineal acotada $\xi:H_1\times H_2\to\C$ tiene una representación única
\[
 \xi(x,y)=\ip{Tx}{y},\qquad T\in\Bcal(H_1,H_2).
\]
\end{teorema}
\begin{proof}
Para cada $x$ fijo, $\varphi_x(y)=\overline{\xi(x,y)}$ es un funcional \emph{lineal} continuo sobre $H_2$ y $\norm{\varphi_x}\leq C\norm x$. Por Riesz existe un único vector $Tx$ con $\overline{\xi(x,y)}=\ip{y}{Tx}$, o equivalentemente $\xi(x,y)=\ip{Tx}{y}$.

Para $x,z\in H_1$ y $\alpha,\beta\in\C$,
\[
 \ip{T(\alpha x+\beta z)}{y}=\xi(\alpha x+\beta z,y)
 =\ip{\alpha Tx+\beta Tz}{y}\quad\text{para todo }y.
\]
La unicidad del vector representado da $T(\alpha x+\beta z)=\alpha Tx+\beta Tz$. La igualdad de normas de Riesz da $\norm{Tx}=\norm{\varphi_x}\leq C\norm x$, por lo que $T$ es acotado. Si $S$ representara la misma forma, $\ip{(T-S)x}{y}=0$ para todo $y$, y tomando $y=(T-S)x$ se obtiene $Tx=Sx$ para cada $x$.
\end{proof}

\section{Existencia, norma y ejemplos del adjunto}
\begin{teorema}\label{ch32:adjunto}
Para $T\in\Bcal(H_1,H_2)$ existe un único $T^*\in\Bcal(H_2,H_1)$ tal que
\[
 \ip{Tx}{y}_{H_2}=\ip{x}{T^*y}_{H_1}.
\]
Además $\norm{T^*}=\norm T$ y $\norm{T^*T}=\norm T^2$.
\end{teorema}
\begin{proof}
Para cada $y$, el funcional $x\mapsto\ip{Tx}{y}$ es lineal y tiene norma a lo sumo $\norm T\norm y$. Riesz proporciona el vector $T^*y$. Para probar linealidad, notemos que
\[
 \ip{x}{T^*(\alpha y+\beta z)}
 =\bar\alpha\ip{Tx}{y}+\bar\beta\ip{Tx}{z}
 =\ip{x}{\alpha T^*y+\beta T^*z}.
\]
La unicidad implica la igualdad de los segundos argumentos. Riesz también da $\norm{T^*y}\leq\norm T\norm y$, luego $\norm{T^*}\leq\norm T$. La unicidad del adjunto se obtiene de la misma identidad: si $S$ satisface la relación, $\ip{x}{(T^*-S)y}=0$ para todo $x$, y $(T^*-S)y=0$.

Tomar conjugados en la identidad del adjunto da $\ip{T^*y}{x}=\ip{y}{Tx}$; por unicidad, $(T^*)^*=T$. Aplicando la desigualdad de normas recién demostrada a $T^*$ obtenemos $\norm T\leq\norm{T^*}$ y, por tanto, igualdad.

Para cada $x$,
\[
 \norm{Tx}^{2}=\ip{x}{T^*Tx}\leq\norm x\norm{T^*Tx}
 \leq\norm{T^*T}\norm x^{2}.
\]
Tomar el supremo sobre $\norm x=1$ da $\norm T^{2}\leq\norm{T^*T}$. Por otro lado, la desigualdad de composición proporciona
$\norm{T^*T}\leq\norm{T^*}\norm T=\norm T^{2}$. Si $H_1=\{0\}$, todas las normas son cero y no hace falta tomar un supremo sobre una esfera vacía.
\end{proof}

\begin{ejemplo}[Matrices y corrimientos]
En bases ortonormales finitas, si $T$ tiene matriz $(a_{ij})$, entonces $T^*$ tiene matriz $(\bar a_{ji})$: basta expandir $\sum_{i,j}a_{ij}x_j\bar y_i$ y agrupar por $j$ para obtener $\ip{x}{T^*y}$.

En $\ell^{2}(\N)$, sean $R(x_1,x_2,\dots)=(0,x_1,x_2,\dots)$ y $L(x_1,x_2,\dots)=(x_2,x_3,\dots)$. Ambos son acotados y tienen norma uno. Como las series convergen absolutamente,
\[
 \ip{Rx}{y}=\sum_{n=2}^{\infty}x_{n-1}\bar y_n
 =\sum_{n=1}^{\infty}x_n\bar y_{n+1}=\ip{x}{Ly}.
\]
Así $R^*=L$. Se tiene $LR=I$, pero $RLx=x-x_1e_1$. El corrimiento derecho preserva el producto interno y no es suprayectivo; esta distinción será decisiva en el ejercicio~\ref{ej:lectura-32-4}, sobre operadores unitarios.
\end{ejemplo}

\section{Ejercicios de adjunción, normalidad y unitariedad}

\begin{ejercicio}[Lectura 32, ejercicio 1]\label{ej:lectura-32-1}
Probar las reglas (i) $(\alpha S+\beta T)^*=\bar\alpha S^*+\bar\beta T^*$; (ii) $(ST)^*=T^*S^*$; (iii) $T^{**}=T$; (iv) $I^*=I$; (v) $(T^{-1})^*=(T^*)^{-1}$ para $T$ invertible.
\end{ejercicio}


\begin{ejercicio}[Lectura 32, ejercicio 2]\label{ej:lectura-32-2}
\index{operador!normal}
Si $T$ es normal, es decir, $T^*T=TT^*$, probar $\norm{Tx}=\norm{T^*x}$ y que $T-\alpha I$ es normal para todo $\alpha\in\C$.
\end{ejercicio}


\begin{ejercicio}[Lectura 32, ejercicio 3]\label{ej:lectura-32-3}
\index{operador!autoadjunto}\index{operador!positivo}
Probar que $T+T^*$ y $TT^*$ son autoadjuntos, esto es, coinciden con sus adjuntos.
\end{ejercicio}


\begin{ejercicio}[Lectura 32, ejercicio 4: isometrías y unitarios]\label{ej:lectura-32-4}
\index{operador!unitario}\index{isometria@isometría!lineal}
Analizar la afirmación: $\ip{Tx}{Ty}=\ip{x}{y}$ para todo $x,y$ si y solo si $T$ es invertible y $T^{-1}=T^*$.
\end{ejercicio}


\begin{ejercicio}[Lectura 32, ejercicio 5]\label{ej:lectura-32-5}
Probar que la proyección ortogonal $P_W$ sobre un subespacio cerrado es autoadjunta.
\end{ejercicio}


\begin{ejercicio}[Lectura 32, ejercicio 6]\label{ej:lectura-32-6}
Si $T^*T=0$, probar $T=0$. Decidir si la conclusión vale cuando $TT^*=0$.
\end{ejercicio}


\begin{ejercicio}[Lectura 32, ejercicio 7]\label{ej:lectura-32-7}
Demostrar que $\norm{Tv}=\norm T\norm v$ si y solo si $T^*Tv=\norm T^{2}v$.
\end{ejercicio}


\begin{ejercicio}[Lectura 32, ejercicio 8]\label{ej:lectura-32-8}
Probar $\ker T=(T^*(H))^{\perp}$.
\end{ejercicio}


\begin{ejercicio}[Lectura 32, ejercicio 9: polarización de una forma]\label{ej:lectura-32-9}
Sea $\xi$ sesquilineal y $\widehat\xi(v)=\xi(v,v)$. Probar (i)
\[
 \xi(u,v)=\sum_{k=1}^{4}i^k\widehat\xi\left(\frac{u+i^kv}{2}\right),
\]
y (ii) que la forma cuadrática determina la forma sesquilineal.
\end{ejercicio}


\begin{ejercicio}[Lectura 32, ejercicio 10: multiplicación]\label{ej:lectura-32-10}
Sea $h\in C([a,b],\C)$, $a<b$, y $T_hf=hf$ en $C([a,b])$ con producto integral. Hallar (i) el adjunto, (ii) cuándo es normal, (iii) cuándo es autoadjunto y (iv) cuándo es unitario.
\end{ejercicio}


\section{Aproximación de operadores compactos por rango finito}
\index{operador!compacto}\index{operador!rango finito}
Los siguientes resultados enlazan este capítulo con los operadores compactos estudiados anteriormente. La geometría de Hilbert permite una aproximación en norma de operador, más fuerte que aproximar separadamente cada imagen.

\begin{teorema}\label{ch32:compactos}
Si $T:H_1\to H_2$ es compacto, existe una sucesión de operadores acotados de rango finito $T_n$ con $\norm{T_n-T}\to0$. Además, $T$ es compacto si y solo si $T^*$ es compacto.
\end{teorema}
\begin{proof}
Sea $B$ la bola unitaria de $H_1$. La compacidad del cierre de $T(B)$ implica total acotación. Para $\eps>0$, elijamos una red finita $y_1,\dots,y_m$ de radio $\eps$ para $T(B)$, y sea $F=\spanop\{y_1,\dots,y_m\}$. Al ser de dimensión finita, $F$ es cerrado y tiene proyección ortogonal $P_F$. Si $x\in B$, existe $j$ con $\norm{Tx-y_j}<\eps$. Como $P_FTx$ es el mejor aproximante en $F$,
\[
 \norm{Tx-P_FTx}\leq\norm{Tx-y_j}<\eps.
\]
Por tanto $\norm{T-P_FT}\leq\eps$. Elegir $\eps=1/n$ produce $T_n=P_{F_n}T$, de rango finito y acotado. La cota $\norm{P_{F_n}}\leq1$ asegura su acotación también cuando $F_n=\{0\}$.

Si un operador acotado $A:H_1\to H_2$ tiene rango finito, elijamos una base ortonormal $e_1,\dots,e_m$ de ese rango. Para cada $j$, el funcional $x\mapsto\ip{Ax}{e_j}$ tiene, por Riesz, un representante $v_j\in H_1$. Así
\[
 Ax=\sum_{j=1}^{m}\ip{x}{v_j}e_j,
 \qquad A^*y=\sum_{j=1}^{m}\ip{y}{e_j}v_j.
\]
La segunda fórmula se verifica sustituyéndola en la identidad del adjunto y conjugando los coeficientes en la segunda entrada. Su rango está contenido en $\spanop(v_1,\dots,v_m)$, luego $A^*$ también tiene rango finito. Todo operador acotado de rango finito es compacto: la imagen de la bola unitaria está acotada en un espacio de dimensión finita, y su cierre allí es compacto.

La continuidad de la adjunción da $\norm{T_n^*-T^*}=\norm{T_n-T}\to0$. Recordemos por qué un límite en norma de operadores compactos es compacto cuando el espacio de llegada es completo. Si $A_n\to A$ en norma, para una tolerancia $\delta>0$ elegimos $n$ con $\norm{A-A_n}<\delta/3$. Una red finita de radio $\delta/3$ para $A_n(B)$ se convierte, por desigualdad triangular, en una red de radio $2\delta/3$ para $A(B)$. Así $A(B)$ es totalmente acotado. Su cierre es completo y totalmente acotado, por lo que es compacto. Aplicado a $T_n^*$, esto demuestra que $T^*$ es compacto. Finalmente, si $T^*$ es compacto, la dirección recién probada aplicada a él muestra que $T^{**}=T$ es compacto.
\end{proof}

\section{De la geometría de Hilbert al estudio de operadores}
El recorrido desde espacios métricos hasta Hilbert muestra cómo se acumulan las estructuras. La métrica hace posible hablar de aproximación; la completitud asegura límites; la linealidad permite construir y comparar aproximantes; el producto interno identifica el mejor de ellos mediante perpendicularidad. Riesz y los adjuntos convierten esa geometría en un lenguaje para estudiar operadores. Estos resultados serán herramientas de los volúmenes de medida e integración, análisis funcional y análisis armónico.


\backmatter
\renewcommand{\chaptermark}[1]{\markboth{#1}{}}
\bookmarksetup{startatroot}
\chapter{Correspondencia con las lecturas}
Las 32 lecturas del primer curso se integran en \textbf{seis partes y dieciséis capítulos temáticos}. Se incluyen los \textbf{206 ejercicios originales numerados}, con sus incisos y soluciones. La lectura 9 contiene un ejercicio bajo un encabezado singular; la lectura 25 no tiene una lista independiente. Las verificaciones propuestas como ejercicio dentro de la teoría también se desarrollan en el texto.

Cada ejercicio identifica su lectura y número originales entre paréntesis. Su número principal sigue la numeración del capítulo del libro. La tabla permite localizar el comienzo del desarrollo de cada lectura; cuenta ejercicios numerados, no incisos ni ejemplos añadidos. Cuando un enunciado requiere corrección, el texto explica el problema y demuestra una versión válida.

Los resultados geométricos de la lectura 22 se estudian en el capítulo~\ref{cap:nuevo08}; su criterio para operadores de rango finito y los ejercicios~\ref{ej:lectura-22-1}, \ref{ej:lectura-22-2}, \ref{ej:lectura-22-3} y~\ref{ej:lectura-22-6} se sitúan en el capítulo~\ref{cap:nuevo09}, después de introducir la norma de operador. Las herramientas sobre convergencia uniforme usadas por Tietze se desarrollan en el capítulo~\ref{cap:nuevo11} y se aplican en el capítulo~\ref{cap:nuevo13}.

El criterio que caracteriza los funcionales continuos por su núcleo, procedente de la lectura 21, se demuestra en la proposición~\ref{ch21:kernel-cerrado} del capítulo~\ref{cap:nuevo09}, antes de su aplicación al rango finito. El desarrollo del dual y los ejercicios de esa lectura permanecen en el capítulo~\ref{cap:nuevo10}.

\footnotesize
\begin{longtable}{@{}r>{\raggedright\arraybackslash}p{.50\textwidth}rrr@{}}
\toprule
Lectura & Tema de origen & Cap. & Ej. & Pág.\\\midrule
\endfirsthead
\toprule Lectura & Tema de origen & Cap. & Ej. & Pág.\\\midrule
\endhead
\bottomrule\endfoot
1 & Distancias y normas & \ref{cap:nuevo01} & 7 & \pageref{lectura:01}\\
2 & Bolas, vecindades, interior y clausura & \ref{cap:nuevo01} & 7 & \pageref{lectura:02}\\
3 & Abiertos, cerrados y sucesiones & \ref{cap:nuevo01} & 9 & \pageref{lectura:03}\\
4 & La topología de un subespacio & \ref{cap:nuevo01} & 6 & \pageref{lectura:04}\\
5 & Continuidad y equivalencia topológica & \ref{cap:nuevo02} & 9 & \pageref{lectura:05}\\
6 & Densidad y numerabilidad & \ref{cap:nuevo02} & 8 & \pageref{lectura:06}\\
7 & Continuidad uniforme y control de las distancias & \ref{cap:nuevo02} & 7 & \pageref{lectura:07}\\
8 & Espacios completos y espacios de Banach & \ref{cap:nuevo03} & 7 & \pageref{lectura:08}\\
9 & La completación de un espacio métrico & \ref{cap:nuevo03} & 1 & \pageref{lectura:09}\\
10 & Categoría de Baire y consecuencias de la completitud & \ref{cap:nuevo04} & 10 & \pageref{lectura:10}\\
11 & El principio de contracción y sus aplicaciones & \ref{cap:nuevo04} & 6 & \pageref{lectura:11}\\
12 & Separación de cerrados y extensión de Tietze & \ref{cap:nuevo13} & 3 & \pageref{lectura:12}\\
13 & Compacidad: del control local al control global & \ref{cap:nuevo05} & 8 & \pageref{lectura:13}\\
14 & Compacidad y sucesiones & \ref{cap:nuevo05} & 6 & \pageref{lectura:14}\\
15 & Acotación total y compacidad local & \ref{cap:nuevo05} & 6 & \pageref{lectura:15}\\
16 & Compactificaciones y espacios universales & \ref{cap:nuevo07} & 5 & \pageref{lectura:16}\\
17 & Conexidad y continuidad & \ref{cap:nuevo06} & 7 & \pageref{lectura:17}\\
18 & Componentes conexas, caminos y arcos & \ref{cap:nuevo06} & 7 & \pageref{lectura:18}\\
19 & Geometría y series en espacios normados & \ref{cap:nuevo08} & 7 & \pageref{lectura:19}\\
20 & Operadores acotados y acotación uniforme & \ref{cap:nuevo09} & 7 & \pageref{lectura:20}\\
21 & El dual, los isomorfismos y la extensión de funcionales & \ref{cap:nuevo10} & 8 & \pageref{lectura:21}\\
22 & Qué distingue a la dimensión finita & \ref{cap:nuevo08}, \ref{cap:nuevo09} & 6 & \pageref{lectura:22}, \pageref{ch22:rango-finito-kernel}\\
23 & Operadores compactos & \ref{cap:nuevo09} & 4 & \pageref{lectura:23}\\
24 & Convexidad, separación y puntos extremos & \ref{cap:nuevo10} & 7 & \pageref{lectura:24}\\
25 & La métrica uniforme y los espacios de funciones & \ref{cap:nuevo11} & 0 & \pageref{lectura:25}\\
26 & Convergencia uniforme: límites, series, integrales y derivadas & \ref{cap:nuevo11} & 8 & \pageref{lectura:26}\\
27 & Equicontinuidad y el teorema de Arzelà--Ascoli & \ref{cap:nuevo12} & 5 & \pageref{lectura:27}\\
28 & Álgebras de funciones y aproximación de Stone--Weierstrass & \ref{cap:nuevo13} & 4 & \pageref{lectura:28}\\
29 & Producto interno y geometría de Hilbert & \ref{cap:nuevo14} & 8 & \pageref{lectura:29}\\
30 & Proyecciones ortogonales y mejor aproximación & \ref{cap:nuevo14} & 6 & \pageref{lectura:30}\\
31 & Bases ortonormales y desarrollos de Fourier & \ref{cap:nuevo15} & 7 & \pageref{lectura:31}\\
32 & Representación de Riesz y operadores adjuntos & \ref{cap:nuevo16} & 10 & \pageref{lectura:32}\\
\end{longtable}
\normalsize

\bookmarksetup{startatroot}
\chapter{Índice de notación}
La página indicada remite al comienzo del capítulo donde se introduce o se utiliza sistemáticamente el símbolo. Las letras $X,Y$ suelen designar espacios métricos; $V,W$, espacios normados; $H$, un espacio de Hilbert. Las letras locales se definen nuevamente cuando cambia el contexto.
\small
\begin{longtable}{@{}p{0.22\textwidth}p{0.66\textwidth}r@{}}
\toprule
Símbolo & Significado & Pág.\\
\midrule
\endfirsthead
\toprule Símbolo & Significado & Pág.\\\midrule
\endhead
\bottomrule\endfoot
$\N,\Z,\Q,\R,\C$ & Naturales positivos, enteros, racionales, reales y complejos & \pageref{chap:01}\\
$\F$ & Cuerpo de escalares: $\R$ o $\C$ & \pageref{chap:01}\\
$d(x,y)$ & Distancia entre $x$ e $y$ & \pageref{chap:01}\\
$\norm{x}$ & Norma del vector $x$ & \pageref{chap:01}\\
$\norm{x}_p$ & Norma de tipo $p$ en un espacio de coordenadas o sucesiones & \pageref{chap:01}\\
$\ell^p,\ \ell^\infty$ & Espacios de sucesiones sumables a la potencia $p$ o acotadas & \pageref{chap:01}\\
$B(x,r)$ & Bola abierta de centro $x$ y radio $r>0$ & \pageref{chap:02}\\
$\overline B(x,r)$ & Bola cerrada: puntos a distancia a lo sumo $r$ de $x$ & \pageref{chap:02}\\
$d(x,A),\ \dist(x,A)$ & Ínfimo de las distancias desde $x$ a puntos de $A\ne\varnothing$ & \pageref{chap:03}\\
$\diam A$ & Diámetro de $A$: supremo de las distancias entre sus puntos & \pageref{chap:02}\\
$\interior A$ & Interior de $A$ en el espacio ambiente especificado & \pageref{chap:02}\\
$\cl A$ & Clausura de $A$; si hace falta se indica el espacio ambiente & \pageref{chap:02}\\
$\partial A$ & Frontera de $A$ & \pageref{chap:05}\\
$A'$ & Conjunto de puntos de acumulación de $A$ & \pageref{chap:03}\\
$x_n\to x$ & Convergencia de la sucesión $(x_n)$ a $x$ & \pageref{chap:03}\\
$f^{-1}(A)$ & Preimagen de $A$; no presupone que $f$ sea invertible & \pageref{chap:05}\\
$\supp f$ & Soporte: clausura del conjunto donde $f$ no se anula & \pageref{chap:05}\\
$c_{00}$ & Sucesiones con un número finito de términos no nulos & \pageref{chap:06}\\
$c,\ c_0$ & Sucesiones convergentes y sucesiones que convergen a cero & \pageref{chap:08}\\
$|r|_p$ & Valor absoluto $p$-ádico de un racional, con $p$ primo & \pageref{chap:08}\\
$\Q_p$ & Cuerpo de números $p$-ádicos, completación de $\Q$ & \pageref{chap:09}\\
$\widehat X$ & Completación de un espacio métrico $X$ & \pageref{chap:09}\\
$F_\sigma,\ G_\delta$ & Unión numerable de cerrados, intersección numerable de abiertos & \pageref{chap:10}\\
$f^n$ & Iterada de una aplicación en el contexto de puntos fijos & \pageref{chap:11}\\
$\mathbf1_A$ & Función característica de $A$, igual a $1$ en $A$ y a $0$ fuera & \pageref{chap:10}\\
$C$ & Conjunto de Cantor, en el capítulo de compactificaciones & \pageref{chap:16}\\
$[0,1]^{\N}$ & Cubo de Hilbert con la topología producto & \pageref{chap:16}\\
$\mathcal N=\N^{\N}$ & Espacio de Baire & \pageref{chap:16}\\
$A+B,\ \lambda A$ & Suma de Minkowski y dilatación de conjuntos & \pageref{chap:19}\\
$\spanop A$ & Subespacio generado por combinaciones lineales finitas de $A$ & \pageref{chap:19}\\
$\Bcal(V,W)$ & Espacio de operadores lineales acotados de $V$ en $W$ & \pageref{chap:20}\\
$\norm{T}$ & Norma de operador; supremo de $\norm{Tv}$ sobre $\norm v\le1$ & \pageref{chap:20}\\
$\ker T,\ T(V)$ & Núcleo e imagen del operador $T$ & \pageref{chap:20}\\
$V^*$ & Dual continuo del espacio normado $V$ & \pageref{chap:21}\\
$V^{**}$ & Dual del dual continuo de $V$ & \pageref{chap:21}\\
$\id,\ I$ & Aplicación identidad u operador identidad, según el contexto & \pageref{chap:21}\\
$\Kcal(V,W)$ & Operadores compactos de $V$ en $W$ & \pageref{chap:23}\\
$\operatorname{co}(A)$ & Envolvente convexa de $A$ & \pageref{chap:24}\\
$\operatorname{ext}(A)$ & Conjunto de puntos extremos de un convexo $A$ & \pageref{chap:24}\\
$C(X),\ C(K)$ & Funciones continuas; $K$ designa habitualmente un compacto & \pageref{chap:25}\\
$C_b(X)$ & Funciones continuas acotadas & \pageref{chap:25}\\
$C_c(X)$ & Funciones continuas de soporte compacto & \pageref{chap:25}\\
$C_0(X)$ & Funciones continuas que se anulan en el infinito & \pageref{chap:25}\\
$B(S,Y)$ & Funciones acotadas del conjunto $S$ al espacio métrico $Y$ & \pageref{chap:25}\\
$d_\infty(f,g)$ & Distancia uniforme: $\sup_{s\in S}d(f(s),g(s))$ & \pageref{chap:25}\\
$\norm f_\infty$ & Norma uniforme: $\sup_x\norm{f(x)}$; para escalares, $\sup_x|f(x)|$ & \pageref{chap:25}\\
$f_n\rightrightarrows f$ & Convergencia uniforme & \pageref{chap:26}\\
$\ip uv$ & Producto interno, lineal en $u$ y conjugado lineal en $v$ & \pageref{chap:29}\\
$u\perp v$ & Ortogonalidad: $\ip uv=0$ & \pageref{chap:29}\\
$M^\perp$ & Complemento ortogonal de un conjunto o subespacio $M$ & \pageref{chap:29}\\
$P_M$ & Proyección ortogonal sobre $M$, cuando existe; en Hilbert basta que $M$ sea cerrado & \pageref{chap:30}\\
$\{e_s\}_{s\in S}$ & Familia ortonormal; se especifica cuándo es una base & \pageref{chap:31}\\
$\ell^2(S)$ & Familias cuadrado sumables con índices en $S$ & \pageref{chap:29}\\
$\ell^2_w$ & Sucesiones cuadrado sumables con pesos positivos $w_n$ & \pageref{chap:29}\\
$\sum_{s\in S}a_s$ & Para $a_s\ge0$, supremo de las sumas sobre subconjuntos finitos de $S$ & \pageref{chap:29}\\
$P_n$ & Polinomio de Legendre de grado $n$ en el ejercicio correspondiente & \pageref{chap:31}\\
$T^*$ & Adjunto hilbertiano de un operador acotado $T$ & \pageref{chap:32}\\
$V/M,\ [x],\ q$ & Espacio cociente, clase $x+M$ y proyección canónica & \pageref{ampl:cocientes}\\
$p,\ \bar p,\ N_p$ & Seminorma, seminorma cociente y núcleo de $p$ & \pageref{ampl:nucleo-seminorma}\\
$\norm u_\beta$ & Norma uniforme ponderada por $e^{-\beta|t-t_0|}$ & \pageref{ampl:picard-equivalencia}\\
$S,\ u_n$ & Operador de Picard y aproximaciones sucesivas & \pageref{ampl:picard-teorema}\\
$G,\ \alpha,\ c$ & Matriz de Gram, columna de productos internos y coeficientes de proyección & \pageref{ampl:gram}\\
$A^*$ & Transpuesta conjugada de una matriz compleja & \pageref{ampl:gram}\\
$b_n^*,\ S_N$ & Funcionales coordenados y sumas parciales de una base de Schauder & \pageref{ampl:schauder}\\
\end{longtable}
\normalsize

\bookmarksetup{startatroot}
\printindex
\bookmarksetup{startatroot}
\begin{thebibliography}{Val}
\addcontentsline{toc}{chapter}{Referencias}
\bibitem[Val]{valencia}
Miguel Ángel Valencia Bucio.
\emph{Análisis Matemático 1. Licenciatura en Física y Matemáticas}.
Lecturas 1--32, apuntes de curso, sin fecha indicada.
Archivos \texttt{AnM1-LFM-lect01.pdf} a \texttt{AnM1-LFM-lect32.pdf}, proporcionados por el autor de esta edición a través de su archivo de trabajo en Google Drive.
\href{https://drive.google.com/drive/folders/12qZh4h3wXB83V2dzuUmPrT-tFyURtZQJ}{Carpeta de las lecturas originales} (acceso sujeto a permisos).

\bibitem[Rud]{rudin1976}
Walter Rudin. \emph{Principles of Mathematical Analysis}.
Tercera edición. McGraw--Hill, 1976.
ISBN de la edición internacional: 978-0-07-085613-4.
\href{https://www.mheducation.co.uk/principles-of-mathematical-analysis-int-l-ed-9780070856134-emea}{Registro editorial de McGraw Hill}.

\bibitem[Mun]{munkres2000}
James R. Munkres. \emph{Topology}.
Segunda edición. Prentice Hall, 2000.
ISBN: 978-0-13-181629-9.
\href{https://books.google.com/books/about/Topology.html?id=XjoZAQAAIAAJ}{Registro bibliográfico en Google Books}.

\bibitem[Kre]{kreyszig1978}
Erwin Kreyszig. \emph{Introductory Functional Analysis with Applications}.
John Wiley \& Sons, 1978.
ISBN: 978-0-471-50731-4.
\href{https://books.google.com/books/about/Introductory_Functional_Analysis_with_Ap.html?id=Va8rAAAAYAAJ}{Registro bibliográfico en Google Books}.

\bibitem[Con]{conway1990}
John B. Conway. \emph{A Course in Functional Analysis}.
Segunda edición. Graduate Texts in Mathematics, vol.~96.
Springer--Verlag, Nueva York, 1990.
ISBN de la edición impresa de 1990: 978-0-387-97245-9.
\href{https://doi.org/10.1007/978-1-4757-4383-8}{DOI de la edición electrónica: 10.1007/978-1-4757-4383-8}.
El catálogo de Springer distingue la edición impresa de 1990 de las fechas posteriores de reimpresión y publicación electrónica.
\end{thebibliography}

\end{document}
